%%%%%%%%%%%%%%%%%%%%%%%%%%%%%%%%%%%%%%%%%%%%%%%%%%%%%%%%%%%%%%%%%%%%%%%%%%%%%
%eventueel dutch in english veranderen, eventueel twoside verwijderen       %
%%%%%%%%%%%%%%%%%%%%%%%%%%%%%%%%%%%%%%%%%%%%%%%%%%%%%%%%%%%%%%%%%%%%%%%%%%%%%
\documentclass[11pt,a4paper,twoside,english]{report}%
\usepackage[dutch]{babel}%
\selectlanguage{english}%
%%%%%%%%%%%%%%%%%%%%%%%%%%%%%%%%%%%%%%%%%%%%%%%%%%%%%%%%%%%%%%%%%%%%%%%%%%%%%
\newif\ifpublic \newif\iftwocol %nodig voor publicatiemode, niet wijzigen!  %
%%%%%%%%%%%%%%%%%%%%%%%%%%%%%%%%%%%%%%%%%%%%%%%%%%%%%%%%%%%%%%%%%%%%%%%%%%%%%
%publiceerbare versie (geen titelblad, bedanktuiging, evaluatie en bijlagen)%
%%%%%%%%%%%%%%%%%%%%%%%%%%%%%%%%%%%%%%%%%%%%%%%%%%%%%%%%%%%%%%%%%%%%%%%%%%%%%
%\publictrue        %1-kolom publicatiemode (zonder %) of verslagmode (met %)
%\twocoltrue        %2-kolommen publicatiemode
%%%%%%%%%%%%%%%%%%%%%%%%%%%%%%%%%%%%%%%%%%%%%%%%%%%%%%%%%%%%%%%%%%%%%%%%%%%%%
%nodig voor template: packages en macro's                                   %
%%%%%%%%%%%%%%%%%%%%%%%%%%%%%%%%%%%%%%%%%%%%%%%%%%%%%%%%%%%%%%%%%%%%%%%%%%%%% 
\ifpublic\iftwocol\twocolumn\fi\fi  %vlag 1 of 2 kolommen                   %
\usepackage[utf8]{inputenc}         %aanbevolen encoding, kies 1 van de 3
%\usepackage[utf-8]{inputenc}      %sommige Ubuntu en Window distributies
%\usepackage[latin1]{inputenc}     %oudere LaTeX-distributies
\usepackage{modelverslag}           %verplicht voor de template             %
\renewcommand{\redactie}[1]{\relax} %nodig voor redactionele opmerkingen    %
%%%%%%%%%%%%%%%%%%%%%%%%%%%%%%%%%%%%%%%%%%%%%%%%%%%%%%%%%%%%%%%%%%%%%%%%%%%%%
%hieronder komen alle packages die je zelf wil plaatsen of veranderen       %
%%%%%%%%%%%%%%%%%%%%%%%%%%%%%%%%%%%%%%%%%%%%%%%%%%%%%%%%%%%%%%%%%%%%%%%%%%%%%
\usepackage{algorithmic}      %aangepaste algoritmen   
\usepackage{listings}         %aangepaste listings (event. in kleuren)
%\usepackage{}
%\usepackage{}
%%%%%%%%%%%%%%%%%%%%%%%%%%%%%%%%%%%%%%%%%%%%%%%%%%%%%%%%%%%%%%%%%%%%%%%%%%%%%
%%%%%%%%%%%%%%%%% einde van de preambule, begin document %%%%%%%%%%%%%%%%%%%%
%%%%%%%%%%%%%%%%%%%%%%%%%%%%%%%%%%%%%%%%%%%%%%%%%%%%%%%%%%%%%%%%%%%%%%%%%%%%%
\begin{document}
\titelblad                   %argumenten invullen
  {Een Uppaal `model' voor TinLAb} %hier titel invullen 
  {Galvin Bartes (0799967)} %hier naam en (stamnummer) van student invullen
  {\TI}                      %hier opleiding \TI, \INF of \BI invullen
  {\today}                   %{\today} concept of {1 april 2011} definitief
  {Dhr. Oelew}           %eerste docent
  {Mevr. Elivira }                         %tweede docent
% als er meer auteurs zijn, dan wordt de tweede parameter:
% {Onbekend Talent1 (123456)\\Onbekend Talent2 (987654)\\ ...} 
\ifpublic
  {\footnotesize{\samenvatting

Het doel van een samenvatting is om potentiële lezers zo snel mogelijk
te overtuigen van de relevantie van het verslag. Als afstudeerverslagen
gepubliceerd worden, is digitaal zoeken noodzakelijk. Daarom worden in
de samenvatting (\emph{abstract}) vaak kenmerkende woorden en
uitspraken opgenomen. Een samenvatting voor een afstudeerverslag mag
niet meer dan een paar honderd woorden bevatten.
}}
\else
  \pagenumbering{roman}
  \samenvatting

Het doel van een samenvatting is om potentiële lezers zo snel mogelijk
te overtuigen van de relevantie van het verslag. Als afstudeerverslagen
gepubliceerd worden, is digitaal zoeken noodzakelijk. Daarom worden in
de samenvatting (\emph{abstract}) vaak kenmerkende woorden en
uitspraken opgenomen. Een samenvatting voor een afstudeerverslag mag
niet meer dan een paar honderd woorden bevatten.
     %(verplicht) samenvatting
  \hoofdstuk{Dankbetuiging}

Wie kan je zoal bedanken? Denk aan de begeleiders en voorbereiders van
je afstudeerproject, familieleden en andere personen die je
geadviseerd of gemotiveerd hebben.  Het is gebruikelijk om dit
voorafgaande aan het verslag te doen. Dit bedanken mag ook in de
inleiding gebeuren. Bijvoorbeeld: Bij het opstellen van dit verslag
heb ik dankbaar gebruik gemaakt van `metathesis' van \emph{Donald Craig
(donald@mun.ca)}.

  %(optioneel) dankbetuiging
  \tableofcontents         %overzicht hoofdstukken en paragrafen
  \hoofdstuk{trefwoorden}

trefwoorden
volgens de gebruikte thauserus. een thesaurus is een lijst vna goedgekeurde en geaccepteerde vaktermen, de 'controlled descriptors' met de verklaring en met de afgekeurde alternatieve vaktermen    %(optioneel) overzicht afkortingen
  \pagenumbering{arabic}
\fi
\hoofdstuk{Inleiding}
 
 

\paragraph{Algemeen}

Het ministerie van verkeer en Waterstaat wil in het kader van het klimaatakkoord en onderzoek laten uitvoeren naar de staat van het sluizenpark in Nederland. Het onderzoek moet zich richten op het ontwerpen en ontwikkelen van een geautomatiseerd sluismodel dat geschikt is voor een brede toepassing. In het onderzoek moet naar voren komen wat de huidige staat is van de sluizen met oog op veiligheid, efficiëntie, capaciteit, onderhoud, duurzaamheid en automatisering. Het onderzoek geeft aan hoe een volledig model worden opgeleverd opdat ontwerp van verschillend volledig geautomatiseerde sluizen in de toekomst geautomatiseerd kunnen worden.  

\paragraph{Recente ontwikkelingen op het gebied van sluisautomatisering}

Het ministerie van verkeer en Waterstaat wil in het kader van het klimaatakkoord en onderzoek laten uitvoeren naar de staat van het sluizenpark in Nederland. Het onderzoek moet zich richten op het ontwerpen en ontwikkelen van een geautomatiseerd sluismodel dat geschikt is voor een brede toepassing. In het onderzoek moet naar voren komen wat de huidige staat is van de sluizen met oog op veiligheid, efficiëntie, capaciteit, onderhoud, duurzaamheid en automatisering. Het onderzoek geeft aan hoe een volledig model worden opgeleverd opdat ontwerp van verschillend volledig geautomatiseerde sluizen in de toekomst geautomatiseerd kunnen worden.  
\paragraph{Wat is een sluis}

\paragraph{Wat worrdt er omschreven en wat is er geleerd}

\paragraph{Wat is uppaal}

Wat is Uppaal
Uppaal is an integrated tool environment for modeling, simulation and verification of real-time systems, developed jointly by Basic Research in Computer Science at Aalborg University in Denmark and the Department of Information Technology at Uppsala University in Sweden. It is appropriate for systems that can be modeled as a collection of non-deterministic processes with finite control structure and real-valued clocks, communicating through channels or shared variables [WPD94, LPW97b]. Typical application areas include real-time controllers and communication protocols in particular, those where timing aspects are critical.


model checking

Wat is statistical model checking?
Dit verwijst naar verschillende technieken dfie worden gebruikt voor de monitoring van een systeem. Daarbij wordt vooral gelet op een specifieke eigenschap. Met de resultaten van de statsitieken wordt de juistheid van een ontwerp beoordeeld. Statistisch model checking wordt onder andere toegepast in systeembiologie, software engineering en industriele toepassingen.
https://www-verimag.imag.fr/Statistical-Model-Checking-814.html?lang=en#:~:text=Statistical%20Model%20Checking%20(SMC)%20is,from%20state%20space%20explosion%20issues.


\cite{inriaStatsMoodCheck}
\cite{ buddeModelChecker}
\cite{AGHASuervey }


Waarom gebruiken we statistisch model checking?
To overcome the above difficulties we propose to work with Statistical Model Checking [KZHHJ09,You05,You06,SVA04,SVA05,SVA05b] an approach that has recently been proposed as an alternative to avoid an exhaustive exploration of the state-space of the model. The core idea of the approach is to conduct some simulations of the system, monitor them, and then use results from the statistic area (including sequential hypothesis testing or Monte Carlo simulation) in order to decide whether the system satisfies the property or not with some degree of confidence. By nature, SMC is a compromise between testing and classical model checking techniques. Simulation-based methods are known to be far less memory and time intensive than exhaustive ones, and are oftentimes the only option. 
https://project.inria.fr/plasma-lab/statistical-model-checking/

Alternatief
Alternatieven voor Uppaal zijn Asynchronous Events,Vesta en MRMC.
%%%%%%%%%%%%%%%%%%%%%%%%%%%%%%%%%%%%%%%%%%%%%%%%%%%%%%%%%%%%%%%%%

\paragraph{Probleemanalyse}

Na grondige analyse van het Nederlandse sluizenpark is gebleken dat renovatie van een groot aantal sluizen noodzakelijk is.  Uit een eerste verkenning is gebleken  dat het gecombineerd renoveren en automatiseren van het Nederlandsesluizenpark een aanzienlijke verbetering kan opleveren t.a.v. 
Op  het  ministerie  van  infrastructuur  enwaterstaat is helaas onvoldoende kennis van ict en systemen aanwezig om eenen ander uit te voeren 

\paragraph{Waarom nu}
In  het  kader  van  het  onlangs  afgesloten  klimaatakkoord  heeft  de  Nederlandseoverheid  daarom  besloten  over  te  gaan  tot  een  ingrijpende  renovatie  van  dediverse  sluizen  die  ons  land  rijk  is.     

\paragraph{Gewenst resultaat }


Wij vragen u een model (of een onderling samenhangend aantal modellen)aan  te  leveren,  opdat  ontwerpen  van  verschillende,  volledig  geautomatiseerdesluizen in de toekomst gerealiseerd kunnen worden. 
Zoals  gesteld  in  de  brief  is  het  de  bedoeling  dat  een  sluis  gemodelleerd  wordten  dat  bewezen  kan  worden  dat  de  te  bouwen  sluis  een  aantal  eigenschappenbezit.  

\paragraph{Scope}

He gaat om het simuleren van een geautomatiseerde sluis. Wat voor type sluis wordt niet gemeld en ook niet uit welke onderdelen. Belangrijk is dat het model werkt en dat het voldoet aan de eisen die gebaseerd zijn op basis van literatuuronderzoek, observatie, interviews, brainstorming of een andere vorm van requirements elicitation.

\paragraph{Onderzoeksvragen }

Hoe kan een geautomatiseerde sluis worden gemodeleerd met oog op ontwikkel- en onderhoudskosten,veiligheid, efficientie en capaciteit





\begin{enumerate}
	
	\item Welke requirements en kwaliteitseisen komen naar voren bij de analyse van een rampenonderzoek
	\item Welke veiligheidseisen er zijn voor sluizen in nederland. 
	\item Hoe kan in uppaal  een model worden getest dat voldoet aan de requirements/eisen volgens het rampenonderzoek?
\end{enumerate}

%%%%%%%%%%%%%%%%%%%%%%%%%%%%%%%%%%%%%%%%%%%%%%%%%%%%%%%%%%%%%%%%%


\paragraph{Design goals}
Het systeem moet minimaal aan de volgende prestatie eisen voldoen 

\begin{enumerate}
	\item   Requirements gebaseerd op rampenanalyse
	\item Model testbaar in upaal
\end{enumerate}

%%%%%%%%%%%%%%%%%%%%%%%%%%%%%%%%%%%%%%%%%%%%%%%%%%%%%%%%%%%%%%%%%
\paragraph{Welke aanpak is gekozen en welke studies liggen hieraan ten grondslag?}
https://link.springer.com/article/10.1007/s10626-020-00314-0

\paragraph{Leeswijzer}
In  de methodologie wordt de lezer uitgelegd met welke methoden de onderzoeksvragen zijn beantwoord. In het hoofdstuk Onderzoek worden alle resultaten behandeld die naar voren zijn gekomen bij het deskresearch. De analyse van de verzamelde data wordt gedaan in het hoofdstuk analyse. Hierin wordt behandeld zoekopdracht naar IoT cloud platforms, feature extractie, prijs-berekening en prijs-feature vergelijking. In het ontwerp komen de uml diagrammen en systeemschetsen naar voren. In de  de hoofdstukken Prototype, IoT cloud en Firmware wordt de implementatie behandeld van het IoT cloud platform in een bestaand project.

%%%%%%%%%%%%%%%%%%%%%%%%%%%%%%%%%%%%%%%%%%%%%%%%%%%%%%%%%%%%%%%%%



%	\input{Methoden} %(verplicht) hoofdverslag
%\hoofdstuk{methoden}
    %(verplicht) inleiding
\hoofdstuk{Hoofdverslag}

In het hoofdverslag moet een rode draad aanwezig zijn,
zodanig dat het een leesbaar artikel wordt. Deze leesbaarheid houdt in
dat het publiceerbaar moet zijn in een vakblad. Dit betekent dat voor
de onderbouwing noodzakelijke gegevens (tabellen, grafieken en
tekeningen) die de leesbaarheid kunnen verlagen, worden opgenomen in
bijlagen. Bijlagen worden niet gepubliceerd.


Plaatjes worden geplaatst als de tekst te ingewikkeld wordt.
Plaatjes moeten wel passen in de rode draad. Ook worden aan plaatjes
`leesbaarheidseisen' gesteld. Haal van te voren overbodige
informatie uit de plaatjes. Afstudeerverslagen met een grote
hoeveelheid `witruimte', `screendumps' en `clipart' zijn verdacht,
meestal duidt dit op een gebrek aan inhoudelijke informatie.


De redactie van het hoofdverslag betreft niet alleen het opmaken van
de teksten, het verwijderen van spelfouten, het controleren van
verwijzingen en het nummeren van de pagina's, maar ook het zoveel
mogelijk verwijderen van opsommingen (bullets). Hoewel opsommingen bij
het opstellen van een teksten in eerste instantie als handige
kapstokken fungeren, moeten zij uiteindelijk als gewone alinea's
of paragrafen herschreven worden.


Het verschil tussen eigen teksten en teksten van anderen moet
duidelijk worden aangegeven. Bronvermelding is verplicht, het weglaten
daarvan kan duiden op plagiaat. Een goed hoofdverslag bevat meestal
verwijzingen naar relevante en recente vakliteratuur en publicaties
van gerenommeerde onderzoeks- en onderwijsinstellingen.  Een
afstudeerverslag van een bacheloropleiding moet verplicht verwijzen
naar recente en relevante vakliteratuur. Indien een afstudeerverslag
hieraan niet voldoet, wordt het afgekeurd.


Het hoofdverslag is een verhaal. De inhoudelijke beoordeling van dit
verhaal is bij ingenieursopleidingen gebaseerd op de verantwoording
van de `keuzes' die gemaakt zijn. De argumentatie van deze gemaakte
keuzes is direct gebaseerd op de vastgestelde functionele eisen (het
pakket van eisen), de randvoorwaarden betreffende de kwaliteit en
risico's en indirect gebaseerd op de context.


Niet alleen goede maar ook de verkeerde keuzes die gemaakt zijn, zijn
belangrijk voor de beoordeling. Vaak blijken verkeerde keuzes tot
interessante conclusies te leiden. Ook zijn ze noodzakelijk om een
indruk te krijgen van de inspanning die de student heeft
geleverd. Naast de keuzes tijdens het ontwerp zijn keuzes voor
hulpmiddelen en gereedschappen relevant voor de beoordeling.  Als de
student bijvoorbeeld zelf een ontwerpmethode gekozen heeft, dan zal
die keuze verantwoord moeten worden. Als de methode voorgeschreven is
in de randvoorwaarden van de projectopdracht, dan zal aangetoond
moeten worden dat de methode op correcte manier is toegepast.


Om een goed verhaal op te stellen, moet vooraf aan enkele voorwaarden
worden voldaan. De eerste voorwaarde is de geschiktheid van het
afstudeerproject. Als een afstudeerproject niet tot keuzes leidt, kan
men zich afvragen of dat wel een echte afstudeeropdracht is. Een
afstudeerproject zonder onderzoeksaspecten is ook verdacht. Daarnaast
moet een afstudeerproject passen in het profiel van een opleiding om
beoordeelbaar te zijn. De andere voorwaarde voor goed een verhaal is
de registratie van werkzaamheden tijdens het afstudeerproject. Dit
voorkomt dat men dingen vergeet.


Laat je verslag lezen door een niet bij het afstudeerproject betrokken
persoon. Meestal ziet die al snel waarin je verslag tekort schiet.


\paragraaf{Werken met \LaTeX{}}

Het is niet verplicht om met \LaTeX{} te werken. Men mag ook gebruik
maken van andere tekstverwerkers zoals \emph{MS-Word}, Wel is het
verplicht het afstudeerverslag \LaTeX{}-geformateerd in te leveren en
van de \LaTeX{}-template \verb!modelverslag.sty! gebruik te
maken.

De \LaTeX{}-template bevat enkele macro's voor het opstellen van een
hoofdstuk (\verb!\hoofdstuk!), een paragraaf (\verb!\paragraaf!), een
afbeelding (\verb!\figuur!). De overige \LaTeX{} macro's en omgevingen
blijven bruikbaar. Bijvoorbeeld de \verb!tabular!-omgeving om tabellen
te maken:

\begin{lstlisting}{language=[LaTeX]TeX}
\begin{tabular}{formaat}
   ... 
\end{tabular}
\end{lstlisting}

\begin{center}
\begin{tabular}{|l||r|}
  \hline
  \multicolumn{2}{|c|}{Afmetingen ($1\mbox{ pt}=0,351\mbox{ mm}$)}\\
  \hline
  paper width     & \the\paperwidth\\
  text width      & \the\textwidth\\
  column width    & \the\columnwidth\\
  column seperate & \the\columnsep\\
  oddside margin  & \the\oddsidemargin\\
  evenside margin & \the\evensidemargin\\
  paper height    & \the\paperheight\\
  text height     & \the\textheight\\
  top margin      & \the\topmargin\\
  \hline
\end{tabular}
\end{center}

Een nadeel van tabellen dat ze vaak te groot zijn voor de
twocolumn-mode. Het zou mooi zijn als ze ingedrukt kunnen
worden. Bovendien is deze tabel niet-zwevend, hij wordt geplaatst
tussen de tekstdelen waar hij is ingevoerd. Dit kan bezwaarlijk zijn
bij pagina-overgangen. In dat geval kan je beter gebruikmaken van
zwevende tabellen (en figuren) die door \LaTeX{} zelf op een geschikte
plaats worden gezet. Wel moet aan een zwevende tabel een label en een
onderschrift gekoppeld worden om er naar te kunnen verwijzen. Voor een
zwevende horizontale tabel met label en onderschrift wordt in de
`template' de \verb!tabel!-omgeving aangeboden:\\

\begin{Aanpassen}
\begin{verbatim}
\begin{tabel}[afm]{formaat}{label}{onderschrift}
  ...
\end{tabel}
\end{verbatim}
\end{Aanpassen}


De \verb!tabel!-omgeving plaatst `zwevende' tabellen in verslag- en
publicatie-mode. Het eerste argument is een optioneel \verb![afm]!
argument met de defaultwaarde \verb!\normalsize! voor de afmeting van
de karakters. De mogelijke waarden voor de afmeting zijn -- van groot
tot klein -- de volgende macro's: (\verb!\huge!, \verb!\LARGE!,
\verb!\Large!, \verb!\large!, \verb!\small!, \verb!\footnotesize!,
\verb!\scriptsize! en \verb!\tiny!).

Bovendien zijn de standaard \verb!tabular! kolomformaten
\verb!r,l,c,|,||,p{lengte}! uit de tabelomgeving uitgebreid met
kolomformaten \verb!\R, \C, \L!  voor variabele celinhoud zoals het
plaatsen van meerdere regels per cel.

Een verticale tabel is mogelijk met de omgeving (\verb!TABEL!)  met
dezelfde kolomformaten mogelijkheden.  In \LaTeX{} zijn de tabellen,
vooral in de \verb!twocolumn!-mode erg lastig. Bijvoorbeeld in de
tabellen~\ref{tab:vbt} en \ref{tab:vbx} zijn twee verschillende
uitwerkingen van de tabelomgevingen:

\begin{footnotesize}
\begin{tabel}[\Large]{|r|l|}{vbt}{Vaste cellen, variabele breedte}
  \hline
  7C0 & hexadecimal \\
  3700 & octal \\ \cline{2-2}
  11111000000 & binary \\
  \hline \hline
  1984 & decimal \\
  \hline
\end{tabel}
\end{footnotesize}

%Absolute breedte in mm, cm etc. (geschikt voor één mode):
%\begin{tabel}{|>\R p{2cm}|>\L p{4cm}|}{vbx}{OpenGL libraries}
%Relatieve breedte 20% en 65% van de kolombreedte (geschikt voor alle moden).
%Procentwaarden van 0 ... 99, voor 100% neem \columnwidth zelf.
\begin{tabel}{|>\R p{\Procent{20}}|>\L
    p{\Procent{65}}|}{vbx}{Variabele cellen, variabele breedte}
\hline
 OpenGL core library & OpenGL32 voor MS-Windows en GL voor
 de meeste X-Window systemen\\
\hline
 OpenGL Utility Library & GLU\\
\hline
 Koppeling met het platform & GLX voor X-Window en WGL voor MS-Windows\\
\hline 
 OpenGL Library Utility Toolkit & GLUT, bibliotheek voor
  het openen van windows, invoer van muis en toetsenbord, menus,
  event-driven in- en uitvoer\\
\hline
\end{tabel}

Plaats afbeeldingen alleen in het hoofdverslag als ze de tekst
ondersteunen en de leesbaarheid niet verlagen.  In de tekst kan naar
afbeeldingen worden verwezen met de macro \verb!\ref{fig:label}!.

In \LaTeX{}\cite{lam1994} geschreven verslagen zijn op diverse manieren
afbeeldingen\cite{Oos1996} te plaatsen. Een van die manieren is gebruik te
maken van de macro \verb!\figuur! in de \verb!modelverslag!-package'.

`Vector graphics' figuren van het `pdf-', `eps-' en `svg-'
formaat\footnote{Een pdf-bestand kan zowel vector-graphics als
  bitmap-graphics bevatten.} met een ingewikkelde `bounding box' zijn
moeilijk op de juiste schaal te brengen. Vaak moet dat met uitproberen
bepaald worden. Het plaatsen van figuren met absolute afmetingen of
een vaste `scale' factor, kan leiden tot minder soepele oplossingen
zoals figuur~\ref{fig:PDFA}. Deze figuur heeft naast een rotatie
(\verb!angle=270!)  een vaste scale-factor (\verb!scale=0.45!) die
alleen geschikt is voor de `twocolumn-mode'.

\begin{center}
  \figuur{scale=0.45,angle=270}{plaatjes/agp.pdf}{PDFA}{Vaste breedte
    (pdf)}
\end{center}

In plaats van \verb!scale=x! kan je beter de relatieve afmeting
\verb!width=\Procent{y}! gebruiken. De waarde $y$ wordt in de
verslag-mode met uitproberen gevonden, zie figuur~\ref{fig:PDFR}.

\figuur{width=\Procent{40},angle=270}{plaatjes/agp.pdf}{PDFR}{Variabele
  breedte (pdf)}

Het afmetingsprobleem is iets gemakkelijker op te lossen met `bitmap
graphics' van het `jpg-', `gif-' en `png-' formaat omdat de figuren al
van te voren geschaald kunnen worden als de `bounding box' bij het
inlezen bekend is. De breedte (\verb!width!) kan als percentage van de
kolombreedte (\verb!width=\Procent{0 ... 99}!) worden opgegeven zoals
dat bij figuur~\ref{fig:PNGR} gedaan is. Voor een 100\% waarde neemt
men \verb!width=\columnwidth! De afmeting wordt automatisch
aangepast aan de nieuwe kolombreedte.


\begin{center}
\figuur{width=\columnwidth}{plaatjes/agp.png}{PNGR}{Variabele breedte (png)}
\end{center}


De macro \verb!\PROCENT{0...99}! is nodig voor de macro's \verb!Tabel!
en \verb!Figuur!. Deze laatste twee macro's maken het mogelijk dat
tabellen en afbeeldingen in de twocolumn-mode passen met behoud van
hun originele afmeting en detaillering (zie
figuur \ref{fig:FIXED}). De parameters van deze macro's komen overeen
met de parameters van de macro's \verb!tabel! en \verb!figuur!.

\Figuur{width=\PROCENT{40},angle=270}{plaatjes/agp.pdf}{FIXED}{Vaste breedte ook in twocolumn-mode (pdf)}

In het algemeen heeft vector-graphics een betere kwaliteit van de
weergave dan bitmap-graphics.


\paragraaf{Bijzondere tekens en afbreekproblemen}

Bijzondere tekens zoals de á, à, ä, é, è, ë, ï, ü, ç \ldots worden
probleemloos door \LaTeX{} geaccepteerd als normale utf8
karakters. Voor de uitzonderingen bestaan macro's zoals het
euro-symbool \euro{} waarvoor de macro \verb!\euro! nodig is. In
wiskundige formules kan je gebruik maken van de macro \verb!\eurom!.


In de two-columnmode zijn regels soms te lang als er gebruik gemaakt
is van \verb!verb! of \verb!verbatim! of woorden die niet goed worden
afgebroken. In dat laatste geval kan je in zo'n woord een afbreekpunt
introduceren met de twee tekens \verb!\-!. Een regel kan gecontroleerd
afgebroken door van te voren onzichtbare knikpunten te plaatsen met de
\verb!\Knak! macro. De volgende regel moet in in tegenstelling met de
twocolumnmode in de verslagmode ongeknakt worden weergegeven:

\begin{Aanpassen}
\begin{verbatim}
 ... aaaaaaa\Knak{}aaaaaaa ...
\end{verbatim}
\end{Aanpassen}


aaaaaaaaaaaaaaaaaaaaaaaaaaaaaaaaaaaaaaa\Knak{}aaaaaaaaaaaaaaaaaaaaaaaaaaaaaaaaaaaaaaa.

Voor regels waarbij de structuur niet gebroken mag worden, is de
\verb!\Knak!-methode ongeschikt, bijvoorbeeld bij scripts en
broncode. Daarentegen zorgt de \verb!Aanpassen!-omgeving ervoor dat in
de twocolumn-mode de regels met behoud van de originele structuur
worden weergegeven. Daarvoor wordt een kleinere letterafmeting
gebruikt (default de \verb!\scriptsize!). Deze omgeving werkt alleen
met niet al te lange regels. Bij zeer lange regels moet de
letterafmeting zeer klein worden waardoor de leesbaarheid in het
gedrang komt. In dat geval moet naar een andere oplossing gezocht
worden zoals het opnemen van de probleemregels (broncode en scripts)
in de bijlagen.

\begin{Aanpassen}[\tiny]
aaaaaaaaaaaaaaaaaaaaaaaaaaaaaaaaaaaaaaaaaaaaaaaaaaaaaaaaaaaaaaaaaaaaaaaaaaaaaa.
\end{Aanpassen}

Hoewel het gebruik van opsommingen (\verb!\item!), letterlijke citaten
\verb!quotation! en kaders (\verb!\fbox!) in de twocolumn-mode tot
problemen kunnen leiden, zijn ze beperkt toegestaan. Bijvoorbeeld voor
de kaders rond de teksten kan je beter gebruik maken van de
\verb!tabular!-omgeving (of de \verb!tabel!-omgeving als je geen last
wil hebben van pagina-overgangen), dan voor de standaard
\verb!\fbox!-methode. De kolom van deze omkaderde tabel moeten dan wel
een relatieve afmetingsverhouding de \verb!\columnwidth! krijgen.

\begin{Aanpassen}
\begin{verbatim}
\begin{center}
\begin{tabular}{|>\C p{\Procent{80}}|}
  \hline
  Afbreekproblemen ...
  \hline
\end{tabular}
\end{center}
\end{verbatim}
\end{Aanpassen}

\begin{center}
\begin{tabular}{|>\C p{\Procent{80}}|}
  \hline
  ~\\
  Afbreek- en andere opmaakproblemen pak je als laatste aan,
  dus bij je definitieve verslag!\\
  ~\\
  Tabellen, figuren en listingen in het hoofdverslag tot het
  noodzakelijke beperken.\\
  ~\\
  \hline
\end{tabular}
\end{center}


\paragraaf{Algoritmen en broncode\cite{wikibooks}}

Als je algoritmen met een mooie layout wilt hebben, dan zou je het
\verb!algorithmic!-pakket kunnen gebruiken. Met dit pakket kan je het
algoritme op een logische manier opbouwen met pseudotaal. Het bestand
`verslag.tex' bevat al de pakketten \verb!algorithmic! en
\verb!listings! die voor dit verslag nodig zijn. Als je zelf packages
wil toevoegen of verwijderen (afblijven van
\verb!\usepackage{moduleverslag}!)  dan moet dat in de preambule
`verslag.tex'.

\begin{Aanpassen}
\begin{verbatim}
\usepackage{algorithmic}
\end{verbatim}
\end{Aanpassen}

Een algoritme moet je maken binnen een algorithmic-omgeving, een
voorbeeld:

\begin{Aanpassen}[\small]
\begin{algorithmic}
\IF {$i\geq maxval$} 
        \STATE $i\gets 0$
\ELSE
        \IF {$i+k\leq maxval$}
                \STATE $i\gets i+k$
        \ENDIF
\ENDIF 
\end{algorithmic}
\end{Aanpassen}


Broncode kan je in een \verb!verbatim!-omgeving opnemen. De
broncoderegels zien er net zo uit zoals je ze ingetypt hebt.  Het
\verb!listings!-pakket is geavanceerder dan de
\verb!verbatim!-omgeving.

\begin{Aanpassen}
\begin{verbatim}
\usepackage{listings}
\end{verbatim}
\end{Aanpassen}

Merk even op dat alle commando's van het \verb!listings!-pakket
beginnen met \verb!lst!, dit conform de lppl-licentie.

De broncode zelf zet je in een \verb!listings!-omgeving, net zoals bij
de \verb!verbatim!-omgeving, om broncode te zetten gebruik je het
\verb!\lstinline!-commando op dezelfde manier als het
\verb!\verb!-commando. Je kunt ook broncode van een extern document laden met het commando:

\begin{Aanpassen}
\begin{verbatim}
\lstinputlisting{pathname}
\end{verbatim}
\end{Aanpassen}

Het argument `pathname' is de relatieve of absolute locatie van het
bronbestand, de map(pen) gecombineerd met de bestandsnaam. Als je
broncode van een bronbestand laadt, ben je zeker dat de broncode in je
\LaTeX{}-document altijd actueel is en hou je het \LaTeX{}-document
overzichtelijk. Als de broncode niet in dezelfde map of een submap van
het \LaTeX{}-document staat of je gebruikt absolute `pathnames', dan
is het mogelijk dat het verslag niet op andere computers gecompileerd
kan worden. Bij het inleveren van je afstudeerverslag in
\LaTeX{}-formaat zal je hiermee rekening moeten houden.


Alle opties in het \verb!listings!-pakket hebben eenzelfde structuur
\verb!sleutel=waarde!. Als je alleen 'Java' gebruikt hebt, dan kan je
deze taal voor je volledig document na de regel
\verb!\usepackage{listings}! in preambule `verslag.tex' definiëren met
\verb!\lstset{language=java}!

\lstset{language=java}

\begin{Aanpassen}
\begin{lstlisting}
public class HelloWorld {
    public static void main(String[] args) {
        System.out.println("Hello, world!");
    }
}
\end{lstlisting}
\end{Aanpassen}


De sleutel is hier dus \verb!language! en de waarde die je aan de
sleutel geeft is \verb!java!. Alles wat je als opties binnen de
\verb!\lstset!-macro zet kan je per \verb!listings!-omgeving apart
definiëren. Bijvoorbeeld html-broncode met
\verb!\begin{lstlisting}[language=html]!:

\begin{Aanpassen}
\begin{lstlisting}[language=html]
<html>
    <head>
        <title>Hello</title>
    </head>
    <body>Hello</body>
</html>
\end{lstlisting}
\end{Aanpassen}

 %(verplicht) hoofdverslag
% Modelverslag.sty Template voor de Hogeschool Rotterdan CMI 
% Versie 1.5, april 2011 - P.J. den Brok
\ProvidesPackage{modelverslag}%
%\usepackage[utf8]{inputenc}
\usepackage{hyperref}%url ondersteuning
\hypersetup{colorlinks=true,linkcolor=black,%
            citecolor=black,filecolor=black, urlcolor=black}%
\usepackage{graphicx}%
\usepackage{tabularx}%
\usepackage{calc}%
\usepackage{rotating}%
\usepackage{latexsym}%
\usepackage{amsfonts}%
\usepackage{pslatex}% 
\usepackage{eurosym}%
\newcommand{\eurom}{\mbox{\euro}}%
\newcommand{\INS}{CMI}%
\iflanguage{dutch}{\def\HR{Hogeschool Rotterdam}}{\def\HR{Rotterdam University}}%
\iflanguage{dutch}{\def\TI{Technische Informatica}}{\def\TI{Technical Informatics}}%
\iflanguage{dutch}{\def\BI{Bedrijfskundige Informatica}}{\def\BI{Business Informatics}}%
\iflanguage{dutch}{\def\INF{Informatica}}{\def\INF{Informatics}}%
\iflanguage{dutch}{\def\van{van}}{\def\van{by}}
\iflanguage{dutch}{\def\opl{Opleiding}}{\def\opl{Program}}
\iflanguage{dutch}{\def\doca{Eerste docent}}{\def\doca{First supervisor}}
\iflanguage{dutch}{\def\docb{Tweede docent}}{\def\docb{Second supervisor}}
%Speciale macro voor afbeeldingen: *.png *.gif *.jpg *.pdf. Voorbeeld:
%\figuur{scale=0.5, angle=270}{plaatjes/plot}{label}{Een eenvoudige plotter.}
\newcommand{\figuur}[4]{\begin{figure}[htpb!]\includegraphics[#1]{#2}%
 \centering\caption{#4}\label{fig:#3}\vspace{2mm}\end{figure}}%
%Speciale macro voor brede afbeeldingen: *.png *.gif *.jpg *.pdf. Voorbeeld:
%\figuur{width=\PROCENT{99}}{plaatjes/plot}{label}{Een eenvoudige
%plotter.}. In twocolumn-mode werkt alleen [tp!] plaatsing.
\newcommand{\Figuur}[4]{\begin{figure*}[htpb!]\includegraphics[#1]{#2}%
 \centering\caption{#4}\label{fig:#3}\vspace{2mm}\end{figure*}}%
% voor de redactie
\newcommand{\redactie}[1]{{\bf{} (Opm: #1) }}
% de tabelomgeving
\newcommand{\mycap}{}% noodzakelijk voor de tabelomgeving
\newcommand{\mylab}{}% noodzakelijk voor de tabelomgeving
\newcommand{\myafm}{}% noodzakelijk voor de TwoColSize omgeving
% alle formatters van de 'tabularomgeving' zijn toegestaan plus:
% Tabelomgevingen en extra formatters (zie voorbeeld hoofdverslag)
% \R p{lengte} meer regels rechts  per cel met gegeven lengte (raggedleft) 
% \C p{lengte} meer regels centrum per cel met gegeven lengte (centering) 
% \L p{lengte} meer regels links   per cel met gegeven lengte (raggedright) 
\newcommand{\BST}[1]{\let\temp=\\#1\let\\=\temp} %bescherm backslash in temp
\def\R{\BST\raggedleft\hspace{0pt}} % interne tabelformatter rechts
\def\C{\BST\centering\hspace{0pt}}  % interne tabelformatter centrum
\def\L{\BST\raggedright\hspace{0pt}}% interne tabelformatter links
\newcommand{\Knak}{\penalty0} % Regelknak tussen woorden
% Relatieve afmetingen met \Procent{0 ... 99} van de kolombreedte
\newcommand{\Procent}[1]{\columnwidth*\real{.#1}}%
% Relatieve afmetingen met \PROCENT{0 ... 99} van de tekstbreedte
\newcommand{\PROCENT}[1]{\textwidth*\real{.#1}}%
% De speciale environment 'tabel' (tabular) voor een horizontale tabel
% met een verplichte label en een verplicht onderschrift.
\newenvironment{tabel}[4][\normalsize]{\renewcommand{\mycap}{#4}%
  \renewcommand{\mylab}{#3}\begin{table}[htbp!]#1\begin{tabular}{#2}}%
 {\end{tabular}\vspace{2mm}\centering\caption{\mycap}%
  \label{tab:\mylab}\end{table}}%
% De speciale environment 'Tabel' (tabular) voor een brede horizontale
% tabel (ook in twocolumn-mode) met een verplichte label en een
% verplicht onderschrift. In twocolumn-mode werkt alleen [tp!] plaatsing.
\newenvironment{Tabel}[4][\normalsize]{\renewcommand{\mycap}{#4}%
  \renewcommand{\mylab}{#3}\begin{table*}[htpb!]#1\begin{tabular}{#2}}%
 {\end{tabular}\vspace{2mm}\centering\caption{\mycap}%
  \label{tab:\mylab}\end{table*}}%
% De speciale environment 'TABEL' macro om een grote verticale tabel
% met label en onderschrift te plaatsen op een aparte pagina
% (ongeschikt in publicatiemode)
\newenvironment{TABEL}[4][\normalsize]{\renewcommand{\mycap}{#4}%
  \renewcommand{\mylab}{#3}\begin{table}[htbp!]#1\begin{sidewaystable}{#2}}%
 {\end{sidewaystable}\vspace{2mm}\centering\caption{\mycap}%
  \label{tab:\mylab}\end{table}}%
% Een relatieve verkleining van lange regels in de two-columnmode:
\newenvironment{Aanpassen}[1][\relax]{\relax}{\relax}
% Bladspiegel afmetingen:
\textwidth      = 15cm%
\textheight     = 23cm%
\headheight     = 0cm%
\headsep        = 0cm%
\topmargin      = 0.5cm%
\oddsidemargin  = 1cm%
\evensidemargin = 0cm%
\floatsep       = 0.75cm%
\textfloatsep   = \floatsep%
\intextsep      = \floatsep%
\setlength{\parskip}{1ex plus 0.5ex minus 0.2ex}%
\sloppy%twocolummode met \verb!!'s moet sloppy, andere moden zijn ongevoeliger
\newlength{\maxwidth}
\ifpublic %publiceerbare versie: zonder bedanktuiging, evaluatie en %bijlagen
  \setlength{\parindent}{0pt}% Eur. alinea past bij listings en verbatims   
  \iftwocol\renewenvironment{Aanpassen}[1][\scriptsize]{#1}{\relax}\fi
  \newcommand{\hoofdstuk}{\subsection*}%I.v.m. pagina-overgang
  \newcommand{\paragraaf}{\subsubsection*}% I.v.m. pagina-overgang
  \newcommand{\samenvatting}[1]%
       {\begin{center}\large{\textbf{\abstractname}}\end{center}}%
  \newcommand{\titelblad}[6]%
      {\include{pagnum}% optionele start paginanummering
       \hypersetup{pdftitle={#1},
            pdfsubject={verslag},
            pdfauthor={#2},
            pdfkeywords={verslag #3},
        }%
       \begin{center}\textbf{#1}\end{center}%
       \begin{center}\footnotesize{\em{#2}}\end{center}%
       \begin{center}\footnotesize{#4}\end{center}%
       \begin{center}\footnotesize{\opl{} `\em{#3}'}\end{center}}%
  \renewenvironment{thebibliography}[1]%
     {{\normalsize{\textbf{ \\\centering\bibname \\}}}%
      \@mkboth{\MakeUppercase\bibname}{\MakeUppercase\bibname}%
      \list{\@biblabel{\@arabic\c@enumiv}}%
           {\settowidth\labelwidth{\@biblabel{#1}}%
            \leftmargin\labelwidth%
            \advance\leftmargin\labelsep%
            \@openbib@code%
            \usecounter{enumiv}%
            \let\p@enumiv\@empty%
            \renewcommand\theenumiv{\@arabic\c@enumiv}}%
      \sloppy%
      \clubpenalty4000%
      \@clubpenalty \clubpenalty%
      \widowpenalty4000%
      \sfcode`\.\@m}%
     {\def\@noitemerr%
       {\@latex@warning{Empty `thebibliography' environment}}%
      \endlist}%
\else     %normale versie verslag
  \newcommand{\hoofdstuk}[1]{\chapter*{#1}\addcontentsline{toc}{chapter}{#1}}%
  \newcommand{\paragraaf}[1]{\section*{#1}\addcontentsline{toc}{section}{#1}}%
  \newcommand{\samenvatting}[1]{\hoofdstuk{\abstractname}}%
  \newcommand{\bijlage}[1]{\chapter{#1}}%
  \newcommand{\titelblad}[6]{%
   \hypersetup{pdftitle={#1},
            pdfsubject={verslag},
            pdfauthor={#2},
            pdfkeywords={verslag #3},
           }%
   \thispagestyle{empty}%
   \begin{center}%
   ~\\~\\~\\~\\~\\~\\% truuk
   \textbf{\Large #1\\}%
   \vspace{1.0cm}%
   \van{} \\%
   \vspace{1.0cm}%
   \emph{#2} \\%
   \vspace{2.0cm}%
   \includegraphics[width=2cm]{plaatjes/logohr.png} \\%
   \vspace{1.0cm}%
   {\small \INS{}-\opl{} \emph{#3} -- \HR} \\%
   \vspace{2.0cm}%
   #4 \\%
   \vspace{\fill}%
   \end{center}%
   {\footnotesize% 
   \begin{tabular}[2]{ll}%
    \doca    & \emph{#5} \\%
    \docb    & \emph{#6} \\%
   \end{tabular}%
   }%
   \vspace{1.0cm}%
  }%
\fi%

 %(verplicht) hoofdverslag
\hoofdstuk{Requirements}

In de evaluatie reflecteer je over je eigen afstudeerproces. Daarbij
moet je vooral letten op de leereffecten. Welke competenties had je
nodig? Welke competenties kwam je tekort en moest je zelf verwerven?
Waren dit algemene of specifieke competenties?  Voldeden de
beroepscompetenties aan de standaard van het \emph{HBO-I} (analyseren,
adviseren, ontwerpen, realiseren en beheren)?  Vielen de algemene
competenties in de vijf categorieën van de \emph{Dublin
Descriptoren}\footnote{Dublin Descriptoren zijn eisen aan de
competenties voor de bachelor en master studies aan universiteiten en
hogescholen in Europa.} zoals het verkrijgen van kennis en inzicht,
het toepassen van kennis en inzicht, het maken van onderbouwde keuzen
(oordeelsvorming), het communiceren (schriftelijk en mondeling) en het
verkrijgen van leervaardigheden?

 %(verplicht) hoofdverslag
\hoofdstuk{Verificatie}

In de evaluatie reflecteer je over je eigen afstudeerproces. Daarbij
moet je vooral letten op de leereffecten. Welke competenties had je
nodig? Welke competenties kwam je tekort en moest je zelf verwerven?
Waren dit algemene of specifieke competenties?  Voldeden de
beroepscompetenties aan de standaard van het \emph{HBO-I} (analyseren,
adviseren, ontwerpen, realiseren en beheren)?  Vielen de algemene
competenties in de vijf categorieën van de \emph{Dublin
Descriptoren}\footnote{Dublin Descriptoren zijn eisen aan de
competenties voor de bachelor en master studies aan universiteiten en
hogescholen in Europa.} zoals het verkrijgen van kennis en inzicht,
het toepassen van kennis en inzicht, het maken van onderbouwde keuzen
(oordeelsvorming), het communiceren (schriftelijk en mondeling) en het
verkrijgen van leervaardigheden?

 %(verplicht) hoofdverslag
\hoofdstuk{Conclusies en aanbevelingen}

Conclusies en aanbevelingen moeten verzameld worden in een apart en
herkenbaar deel van het verslag. Hoewel in het hoofdverslag op diverse
plaatsen conclusies getrokken kunnen worden, moeten de belangrijkste
conclusies samengevoegd en samengevat worden. 


Belangrijk is dat het verschil tussen objectief controleerbare
conclusies en subjectieve aanbevelingen duidelijk wordt aangegeven.
Ook is het aan te bevelen om de belangrijkste conclusies conform de
opdrachtomschrijving te formuleren.

   %(verplicht) geaggregeerde conclusies en aanbevelingen
\ifpublic
  \iflanguage{dutch}{\def\bibname{\normalsize{Bronnen}}}
                    {\def\bibname{\normalsize{References}}}
  {\footnotesize{
\begin{thebibliography}{99}

\bibitem{lam1994} Lamport L.: \emph{\LaTeX: A Document Preparation System}, Addison-Wesley, 1994 

\bibitem{Oos1996} Oostrum van P.: \emph{Handleiding \LaTeX}, Vakgroep
  Informatica, Universiteit Utrecht, 1998,\\ 
  \url{http://people.cs.uu.nl/piet/latexhnd.pdf}

\bibitem{wikibooks} Wikibooks \LaTeX:\\
  \url{http://nl.wikibooks.org/wiki/LaTeX}


\bibitem{oid} Wikibooks \LaTeX:\\
\url{https://www.waterkant.net/suriname/2023/05/29/milieuactivist-sleur-zeer-grote-onwaarheden-staan-in-cyanide-onderzoeksrapport/}

   


%%%%%%%%%%%%%%%%%%%%%%%%%%%%%%%%%%%%%%%%%%%%%%%%%%%%%%%%%%%%%%%%%
%
%
%misc{para,
%	Author = {{IEEE Referencing}},
%	Howpublished = {\url{http://libguides.bhtafe.edu.au/content.php?pid=88814&sid=660920}},
%	Keywords = {paraphrasing, plagariasm, citation, quotation},
%	Lastchecked = {19/04/2017},
%	Title = {{How to Cite\/Quote in Your Assignment}}}
%
%misc{ieee-write,
%	Author = {{IEEE Authorship Series}},
%	Howpublished = {\url{http://ieeeauthorcenter.ieee.org/wp-content/uploads/How-to-Write-for-Technical-Periodicals-and-Conferences.pdf}},
%	Institution = {{IEEE}},
%	Keywords = {writing, paragraph, sentences, ethics},
%	Lastchecked = {18/04/2017},
%	Title = {How to Write for Technical Periodicals \& Conferences}}
%
%
%
%@article{schon2017agile,
%	title={Agile Requirements Engineering: A systematic literature review},
%	author={Sch\"on, Eva-Maria and Thomaschewski, J\"org and Escalona, Maria Jose},
%	journal={Computer Standards \& Interfaces},
%	volume={49},
%	pages={79--91},
%	year={2017},
%	publisher={Elsevier}
%}
%
%@inproceedings{royce1987managing,
%	title={Managing the development of large software systems: concepts and techniques},
%	author={Royce, Winston W},
%	booktitle={Proceedings of the 9th international conference on Software Engineering},
%	pages={328--338},
%	year=1987,
%	organization={IEEE Computer Society Press}
%}
%
%@article{leveson1993investigation,
%	title={An investigation of the Therac-25 accidents},
%	author={Leveson, Nancy G and Turner, Clark S},
%	journal={IEEE computer},
%	volume={26},
%	number={7},
%	pages={18--41},
%	year={1993}
%}
%@book{modelchecking,
%	author = {Clarke,Jr., Edmund M. and Grumberg, Orna and Peled, Doron A.},
%	title = {Model Checking},
%	year = {1999},
%	isbn = {0-262-03270-8},
%	publisher = {MIT Press},
%	address = {Cambridge, MA, USA},
%}
%@inproceedings{nuseibeh2000requirements,
%	title={Requirements engineering: a roadmap},
%	author={Nuseibeh, Bashar and Easterbrook, Steve},
%	booktitle={Proceedings of the Conference on the Future of Software Engineering},
%	pages={35--46},
%	year={2000},
%	organization={ACM}
%}
%
%
%%%%%%%%%%%%%%%%%%%%%%%%%%%%%%%%%%%%%%%%%%%%%%%%%%%%%%%%%%%%%%%%%%
%
%@online{inriaStatsMoodCheck,
%	ALTauthor = {project.inria},
%	ALTeditor = {editor},
%	title = {statistical-model-checking},
%	date = {date},
%	url = {https://project.inria.fr/plasma-lab/statistical-model-checking/},
%}
%
%
%@online{buddeModelChecker,
%	ALTauthor = {Carlos E. Budde1, Pedro R. D’Argenio2,3,4, Arnd Hartmanns1(B) , and Sean Sedwards5},
%	ALTeditor = {editor},
%	title = {A Statistical Model Checker for Nondeterminism and Rare Events},
%	date = {date},
%	url = {https://ris.utwente.nl/ws/portalfiles/portal/28200786/A_statistical_model_checker.pdf},
%}
%
%
%@online{ID,
%	ALTauthor = {GUL AGHA and  },
%	ALTeditor = {KARL PALMSKOG},
%	title = {A Survey of Statistical Model Checking
%	},
%	date = {date},
%	url = {https://dl.acm.org/doi/10.1145/3158668},
%}
%
%
%
%@online{AGHASuervey,
%	ALTauthor = {GUL AGHA and KARL PALMSKOG,},
%	ALTeditor = {editor},
%	title = {A Survey of Statistical Model Checking},
%	date = {date},
%	url = {https://dl.acm.org/doi/pdf/10.1145/3158668},
%}
%%%%%%%%%%%%%%%%%%%%%%%%%%%%%%%%%%%%%%%%%%%%%%%%%%%%%%%%%%%%%%%%%%
%
%
%
%@online{bicker21102016automatiseringsparadox,
%	ALTauthor = {author},
%	ALTeditor = {editor},
%	title = {title},
%	date = {date},
%	url = {https://www.debicker.eu/de-automatiseringsparadox/},
%}
%@online{vseautoparadox,
%	ALTauthor = {author},
%	ALTeditor = {editor},
%	title = {title},
%	date = {date},
%	url = {https://vse.nl/de-paradox-van-de-industriele-automatisering/},
%}
%
%
%
%@online{blogxot21112016slimapparaat,
%	ALTauthor = {author},
%	ALTeditor = {editor},
%	title = {title},
%	date = {date},
%	url = {https://blog.xot.nl/2016/11/21/slimme-apparaten-maken-ons-dom-en-kwetsbaar/index.html},
%}
%
%@online{ID,
%	ALTauthor = {author},
%	ALTeditor = {editor},
%	title = {title},
%	date = {date},
%	url = {https://automatie-pma.com/nieuws/industriele-automatiseringsparadox},
%}
%%%%%%%%%%%%%%%%%%%%%%%%%%%%%%%%%%%%%%%%%%%%%%%%%%%%%%%%%%%%%%%%%%
%
%
%
%
%
%
%
%
%
%@online{aviationsafety04101992airplaneCrashBijlmer,
%	ALTauthor = {author},
%	ALTeditor = {editor},
%	title = {title},
%	date = {date},
%	url = {https://aviation-safety.net/database/record.php?id=19921004-2&lang=nl  },
%}
%
%
%
%@online{boogers092002RampenRegelsRichtlijnen,
%	ALTauthor = {author},
%	ALTeditor = {editor},
%	title = {title},
%	date = {date},
%	url = { https://www.researchgate.net/publication/254815008_Rampen_regels_richtlijnen},
%}
%
%
%@online{catsr25022009Boeing737AmsterdamCrash,
%	ALTauthor = {author},
%	ALTeditor = {editor},
%	title = {title},
%	date = {date},
%	url = {https://catsr.vse.gmu.edu/SYST460/TA1951_AccidentReport.pdf  },
%}
%
%@online{INSAVienna1992Chernobyl,
%	ALTauthor = {author},
%	ALTeditor = {editor},
%	title = {title},
%	date = {date},
%	url = {https://www-pub.iaea.org/MTCD/publications/PDF/Pub913e_web.pdf  },
%}
%
%
%
%
%aviationReport
%
%@online{oVVSchietongevalOssendrecht,
%	ALTauthor = {author},
%	ALTeditor = {editor},
%	title = {title},
%	date = {date},
%	url = {https://www.youtube.com/watch?v=6jmkDClGDHo  },
%}
%
%@online{molukseTreinkaping,
%	ALTauthor = {author},
%	ALTeditor = {editor},
%	title = {title},
%	date = {date},
%	url = {https://www.youtube.com/watch?v=h99Fe9XzzHI  },
%}
%
%@online{jiang16042019TanjinExplosion,
%	ALTauthor = {author},
%	ALTeditor = {editor},
%	title = {title},
%	date = {date},
%	url = {https://www.hindawi.com/journals/joph/2019/1360805/  },
%}
%
%
%@online{hrw03082021investigateBeirutBlast,
%	ALTauthor = {author},
%	ALTeditor = {editor},
%	title = {title},
%	date = {date},
%	url = { 	https://www.hrw.org/report/2021/08/03/they-killed-us-inside/investigation-august-4-beirut-blast },
%}
%
%@online{souaibyElHussein112020Beirutstory,
%	ALTauthor = {author},
%	ALTeditor = {editor},
%	title = {title},
%	date = {date},
%	url = { https://www.researchgate.net/publication/348325979_Beirut_Explosion_the_full_story },
%}
%
%@online{ifrc2020chemicalexplosionBeirutPort,
%	ALTauthor = {author},
%	ALTeditor = {editor},
%	title = {title},
%	date = {date},
%	url = { https://reliefweb.int/sites/reliefweb.int/files/resources/CaseStudy_BeirutExplosion_TechBioHazardsweb.pdf },
%}
%
%@online{caliskan09112013747boeingkalman,
%	ALTauthor = {author},
%	ALTeditor = {editor},
%	title = {title},
%	date = {date},
%	url = {https://www.hindawi.com/journals/ijae/2014/472395/  },
%}
%
%%%%%%%%%%%%%%%%%%%%%%%%%%%%%%%%%%%%%%%%%%%%%%%%%%%%%%%%%%%%%%%%%%
%persuasive technology 
%
%@online{humanTechpersuasiveTech,
%	ALTauthor = {author},
%	ALTeditor = {editor},
%	title = {title},
%	date = {date},
%	url = {https://www.humanetech.com/youth/persuasive-technology  },
%}
%
%
%
%
%@online{rezenfeld01012018persuasiveTecgHabits,
%	ALTauthor = {author},
%	ALTeditor = {editor},
%	title = {title},
%	date = {date},
%	url = {https://spectrum.ieee.org/how-persuasive-technology-can-change-your-habits  },
%}
%
%
%@online{aldenaini28042020persuasiveTechTrends,
%	ALTauthor = {author},
%	ALTeditor = {editor},
%	title = {title},
%	date = {date},
%	url = { https://www.frontiersin.org/articles/10.3389/frai.2020.00007/full },
%}
%
%@online{larson14062017persuasivetechmanipulates,
%	ALTauthor = {author},
%	ALTeditor = {editor},
%	title = {title},
%	date = {date},
%	url = {https://psmag.com/environment/captology-fogg-invisible-manipulative-power-persuasive-technology-81301  },
%}
%
%@online{tanzem22012022persuasivetechchanginglives,
%	ALTauthor = {author},
%	ALTeditor = {editor},
%	title = {title},
%	date = {date},
%	url = {https://www.makeuseof.com/what-is-persuasive-technology/  },
%}
%
%
%
%
%@online{tikkakuddonenpersuasiveTechnology,
%	ALTauthor = {author},
%	ALTeditor = {editor},
%	title = {title},
%	date = {date},
%	url = { https://cyberpsychology.eu/article/view/12270 },
%}
%
%@online{sprongken19032018CourtProcedureDigital,
%	ALTauthor = {author},
%	ALTeditor = {editor},
%	title = {title},
%	date = {date},
%	url = { https://www.njb.nl/blogs/a-court-with-no-face-and-no-place/ },
%}
%
%@online{PROCESREGLEMENTEcourt,
%	ALTauthor = {author},
%	ALTeditor = {editor},
%	title = {title},
%	date = {date},
%	url = { http://www.e-court.nl/wp-content/uploads/2018/03/Procesreglement-e-Court-2017_20180201.pdf},
%}
%
%
%@online{ovvMortierOngevalMaliVideo,
%	ALTauthor = {author},
%	ALTeditor = {editor},
%	title = {title},
%	date = {date},
%	url = {https://www.youtube.com/watch?v=PC2ekl4SaNA  },
%}
%
%
%
%
%@online{vtmGroep29032019waterwerkencyber,
%	ALTauthor = {author},
%	ALTeditor = {editor},
%	title = {title},
%	date = {date},
%	url = { https://www.vtmgroep.nl/blog/waterwerken-in-nederland-onvoldoende-beveiligd-tegen-cyberaanvallen},
%}
%
%
%
%
%
%@online{wesemannVeiligheidLandEnWater,
%	ALTauthor = {author},
%	ALTeditor = {editor},
%	title = {title},
%	date = {date},
%	url = {
%		http://www.wesemann.nl/nl/nieuws-en-pers/274-veiligheid-op-het-water-en-op-het-land.html },
%}
%
%
%
%

%%%%%%%%%%%%%%%%%%%%%%%%%%%%%%%%%%%%%%%%%%%%%%%%%%%%%%%%%%%%%%%%%
\bibitem{ISAC_SANS_Ukraine_DUC_18Mar2016} ... \LaTeX:\\ \url{https://www.nerc.com/_layouts/15/Nerc.404/CustomFileNotFound.aspx?requestUrl=https://www.nerc.com/pa/CI/ESISAC/Documents/E-ISAC_SANS_Ukraine_DUC_18Mar2016.pdf},}
\bibitem{zetter2016GridHack,	ALTauthor = {Kim Zetter},	ALTeditor = {editor},	title = {Inside the Cunning, Unprecedented Hack of Ukraine's Power Grid},	date = {date},	url = {https://www.wired.com/2016/03/inside-cunning-unprecedented-hack-ukraines-power-grid/},}
\bibitem{2015ukrainegridattack,	ALTauthor = {wikipedia},	ALTeditor = {editor},	title = {2015 Ukraine power grid hack},	date = {date},	url = {https://en.wikipedia.org/wiki/2015_Ukraine_power_grid_hack},}
\bibitem{greenberg2017Cyberwartestlab} ... \LaTeX:\\ \url{https://www.wired.com/story/russian-hackers-attack-ukraine/},}
\bibitem{boozallen2016lightwentout,	ALTauthor = {allen},	ALTeditor = {editor},	title = {Ukrain report when the lights went out},	date = {date},	url = {https://www.boozallen.com/content/dam/boozallen/documents/2016/09/ukraine-report-when-the-lights-went-out.pdf},}
\bibitem{finklejan2016UsBlamesRussianSandworm} ... \LaTeX:\\ \url{https://www.reuters.com/article/us-ukraine-cybersecurity-sandworm-idUSKBN0UM00N20160108},}
\bibitem{wired012016ukrainegrid} ... \LaTeX:\\ \url{https://www.wired.com/2016/01/everything-we-know-about-ukraines-power-plant-hack/},}
\bibitem{icsalert20072021cyberattackukraine} ... \LaTeX:\\ \url{https://www.us-cert.gov/ics/alerts/IR-ALERT-H-16-056-01},}
\bibitem{finkle08012016russiansandwormhackers} ... \LaTeX:\\ \url{https://www.reuters.com/article/us-ukraine-cybersecurity-sandworm/u-s-firm-blames-russian-sandworm-hackers-for-ukraine-outage-idUSKBN0UM00N20160108},}
\bibitem{zinets15022017ukrainechargesrussia} ... \LaTeX:\\ \url{https://www.reuters.com/article/us-ukraine-crisis-cyber-idUSKBN15U2CN},}
\bibitem{2014russiansandworm} ... \LaTeX:\\ \url{https://www.wired.com/2014/10/russian-sandworm-hack-isight/},}
\bibitem{mcelfresh2016cyberattackhowandwhy} ... \LaTeX:\\ \url{https://theconversation.com/cyberattack-on-ukraine-grid-heres-how-it-worked-and-perhaps-why-it-was-done-52802},}
\bibitem{desarnaud2017cyberattacks} ... \LaTeX:\\ \url{https://www.ifri.org/sites/default/files/atoms/files/desarnaud_cyber_attacks_energy_infrastructures_2017_2.pdf},}
\bibitem{osti2018historycontrolsystemIncidents} ... \LaTeX:\\ \url{https://www.osti.gov/biblio/1505628},}
\bibitem{Shahzad2014ScadaProtocolsPollingScenario} ... \LaTeX:\\ \url{https://scialert.net/fulltext/?doi=tasr.2014.396.405},}
\bibitem{caseli04112016intrusiondetectioncontrolsystem} ... \LaTeX:\\ \url{https://ris.utwente.nl/ws/files/6028066/3-s2_0-B9780128015957000227.pdf},}
\bibitem{rochascadatesting} ... \LaTeX:\\ \url{https://repositorio-aberto.up.pt/bitstream/10216/119066/2/315683.pdf},}
\bibitem{levalle2020FuzzingICSProtocols} ... \LaTeX:\\ \url{https://dreamlab.net/en/blog/post/fuzzing-ics-protocols/},}
\bibitem{resch31102019IEC62351secureCommunication} ... \LaTeX:\\ \url{http://www.connectivity4ir.co.uk/article/175490/IEC-62351--Secure-communication-in-the-energy-industry.aspx},}
\bibitem{slowik2019ReassasUkraine2016Attack} ... \LaTeX:\\ \url{https://www.dragos.com/wp-content/uploads/CRASHOVERRIDE.pdf},}
\bibitem{arrizabalaga2020surveyiiotProtocols} ... \LaTeX:\\ \url{https://dl.acm.org/doi/fullHtml/10.1145/3381038},}
\bibitem{yadav2020reviewScadaArchitecture} ... \LaTeX:\\ \url{https://arxiv.org/pdf/2001.02925.pdf},}
\bibitem{cis20072021crashoverridemalware} ... \LaTeX:\\ \url{https://www.us-cert.gov/ncas/alerts/TA17-163A},}
\bibitem{holappa2017threattoElectricityNetworks} ... \LaTeX:\\ \url{https://www.nixu.com/fi/node/53},}
\bibitem{shehod2016gridadvantageus} ... \LaTeX:\\ \url{http://web.mit.edu/smadnick/www/wp/2016-22.pdf},}
\bibitem{parkwalstorm11102017russiagridattack} ... \LaTeX:\\ \url{https://jsis.washington.edu/news/cyberattack-critical-infrastructure-russia-ukrainian-power-grid-attacks/},}
\bibitem{drago2017CrashOverride} ... \LaTeX:\\ \url{https://www.dragos.com/wp-content/uploads/CrashOverride-01.pdf},}
\bibitem{wikiindustroyer} ... \LaTeX:\\ \url{https://en.wikipedia.org/wiki/Industroyer},}
\bibitem{crashoverridenetwork} ... \LaTeX:\\ \url{https://en.wikipedia.org/wiki/Crash_Override_Network},}
\bibitem{slowikvb2018crashoverride} ... \LaTeX:\\ \url{https://www.virusbulletin.com/virusbulletin/2019/03/vb2018-paper-anatomy-attack-detecting-and-defeating-crashoverride/},}
\bibitem{njccicthreat08102017crashovverrideprofile} ... \LaTeX:\\ \url{https://www.cyber.nj.gov/threat-center/threat-profiles/ics-malware-variants/crashoverride},}
\bibitem{crashoverrideindustroyermalware} ... \LaTeX:\\ \url{https://www.webopedia.com/TERM/C/crashoverride-industroyer-malware.html},}
\bibitem{incibe23062017crashoverrideback} ... \LaTeX:\\ \url{https://www.incibe-cert.es/en/blog/crashoverride-malware-ics-back-again},}
\bibitem{industroyershortfact} ... \LaTeX:\\ \url{https://rhebo.com/en/service/glossar/industroyer-25114/},}
\bibitem{vijayan2017firstmalwareCausedOutage} ... \LaTeX:\\ \url{https://www.darkreading.com/threat-intelligence/first-malware-designed-solely-for-electric-grids-caused-2016-ukraine-outage/d/d-id/1329114},}
\bibitem{ferrante2017crashoverrideredflag} ... \LaTeX:\\ \url{https://www.powermag.com/why-crashoverride-is-a-red-flag-for-u-s-power-companies/},}
\bibitem{2017cashoverrideindustroyerAlto} ... \LaTeX:\\ \url{https://blog.paloaltonetworks.com/2017/06/crashoverrideindustroyer-protections-palo-alto-networks-customers/},}
\bibitem{spinner2018crashoverrideiot} ... \LaTeX:\\ \url{https://iiot-world.com/ics-security/cybersecurity/five-cybersecurity-experts-about-crashoverride-malware-main-dangers-and-lessons-for-iiot/},}
\bibitem{abb30062017crashoverridenotification} ... \LaTeX:\\ \url{https://search.abb.com/library/Download.aspx?DocumentID=9AKK107045A1003&amp;LanguageCode=en&amp;DocumentPartId=&amp;Action=Launch},}
\bibitem{blackhatusa2017} ... \LaTeX:\\ \url{https://www.blackhat.com/us-17/briefings/schedule/#industroyercrashoverride-zero-things-cool-about-a-threat-group-targeting-the-power-grid-6159},}
\bibitem{humanitcrashalert} ... \LaTeX:\\ \url{https://humanit.asia/ta17-163a/},}
\bibitem{sroberts2017CrashoerrideCronicles} ... \LaTeX:\\ \url{https://medium.com/@sroberts/the-crash-override-chronicles-overall-8389ef178fdf},}
\bibitem{parker2020industrialsystemsattack} ... \LaTeX:\\ \url{https://www.oilandgaseng.com/articles/the-most-infamous-cyber-attacks-on-industrial-systems/},}
\bibitem{kirk2017threatIndustrialControls} ... \LaTeX:\\ \url{https://www.bankinfosecurity.com/power-grid-malware-platform-threatens-industrial-controls-a-9987},}
\bibitem{chalfant2017USElectricGrid} ... \LaTeX:\\ \url{https://thehill.com/policy/cybersecurity/337877-crash-override-malware-heightens-fears-for-us-electric-grid},}
\bibitem{fbiWarningcrashOverride} ... \LaTeX:\\ \url{https://www.inguardians.com/dhs-fbi-warn-of-attacks-against-us-energy-manufacturing-companies-and-employees/},}
\bibitem{davies2017crashoveerideUkraineAttack} ... \LaTeX:\\ \url{https://rethinkresearch.biz/articles/industroyer-crashoverride-malware-behind-ukraine-utility-attack/},}
\bibitem{slowik2018securitySessions} ... \LaTeX:\\ \url{https://electricenergyonline.com/energy/magazine/1104/article/Security-Sessions-Combating-ICS-Threats.htm},}
\bibitem{ENCS2017crashoverridemodules} ... \LaTeX:\\ \url{https://www.smart-energy.com/regional-news/europe-uk/encs-crash-override-virus/},}
\bibitem{brocklehurst2017crashoverridegridtakedown} ... \LaTeX:\\ \url{https://isssource.com/crashoverride-designed-for-grid-takedown/},}
\bibitem{20170613_crashoverride} ... \LaTeX:\\ \url{https://gigazine.net/gsc_news/en/20170613-crashoverride/},}
@onlinec{leetaru2017crashoverridehomefront} ... \LaTeX:\\ \url{https://www.forbes.com/sites/kalevleetaru/2017/06/24/crash-override-and-how-cyberwarfare-is-bringing-conflict-to-the-homefront/#42eb8984277c},}
\bibitem{fti2017redflag} ... \LaTeX:\\ \url{https://fticybersecurity.com/2017-11/crashoverride-red-flag-u-s-power-companies/},}
\bibitem{27001academy2014Whitepaper} ... \LaTeX:\\ \url{https://info.advisera.com/hubfs/27001Academy/27001Academy_FreeDownloads/NL/Checklist_of_ISO_27001_Mandatory_Documentation_NL.pdf},}
\bibitem{2017info_isoiec27019} ... \LaTeX:\\ \url{https://webstore.iec.ch/preview/info_isoiec27019%7Bed1.0%7Den.pdf},}
\bibitem{iec623512023serseries} ... \LaTeX:\\ \url{https://webstore.iec.ch/publication/6912},}
\bibitem{iec623512011Withdrawn} ... \LaTeX:\\ \url{https://webstore.iec.ch/publication/6911},}
\bibitem{IEC62351sheet} ... \LaTeX:\\ \url{https://www.ipcomm.de/protocol/IEC62351/en/sheet.html},}
\bibitem{NISTOIR2014GuidelinesCyverSec} ... \LaTeX:\\ \url{https://nvlpubs.nist.gov/nistpubs/ir/2014/NIST.IR.7628r1.pdf},}
\bibitem{ukraiinnesandwormteam} ... \LaTeX:\\ \url{https://www.fireeye.com/blog/threat-research/2016/01/ukraine-and-sandworm-team.html/cybersecurity-audit-checklist. },}
\bibitem{politicsrussagridukraine} ... \LaTeX:\\ \url{https://edition.cnn.com/2016/02/11/politics/ukraine-power-grid-attack-russia-us/index.html/ukraine-sees-russian-hand-in-cyber-attacks-on-power-grid-idUSKCN0VL18E. },}
\bibitem{Whitehead2017ukrainepoweroutage,	ALTauthor = {David E. Whitehead, Kevin Owens, Dennis Gammel, and Jess Smith},	ALTeditor = {editor},	title = {title},	date = {March 21–23, 2017},	url = {https://na.eventscloud.com/file_uploads/aed4bc20e84d2839b83c18bcba7e2876_Owens1.pdf},}
\bibitem{zetter2017moreDangerousMalware} ... \LaTeX:\\ \url{https://www.vice.com/en/article/zmeyg8/ukraine-power-grid-malware-crashoverride-industroyer},}
\bibitem{icsSecurityRussianHacking} ... \LaTeX:\\ \url{https://www.wallix.com/blog/ics-security-russian-hacking},}
\bibitem{rocha2017cybersecyrityanalysisScada} ... \LaTeX:\\ \url{https://www.sans.org/blog/confirmation-of-a-coordinated-attack-on-the-ukrainian-power-grid/},}
\bibitem{cheah2008testingiec60870} ... \LaTeX:\\ \url{http://citeseerx.ist.psu.edu/viewdoc/download;jsessionid=0513EED48102FDAD1BD940260EF12B11?doi=10.1.1.548.7490&amp;rep=rep1&amp;type=pdf},}
\bibitem{virsec2017DeepDiveIndustroyer} ... \LaTeX:\\ \url{https://virsec.com/virsec-hack-analysis-deep-dive-into-industroyer-aka-crash-override/},}
\bibitem{fauri2017EncryptionICS} ... \LaTeX:\\ \url{https://www.win.tue.nl/~setalle/2017_fauri_encryption.pdf},}
\bibitem{2017IEC61850} ... \LaTeX:\\ \url{http://blog.nettedautomation.com/2017/},}
\bibitem{2017win32industroyer} ... \LaTeX:\\ \url{https://www.welivesecurity.com/wp-content/uploads/2017/06/Win32_Industroyer.pdf},}
\bibitem{dragos2019TargetedTransStation} ... \LaTeX:\\ \url{https://www.cybersecurityintelligence.com/blog/attack-on-ukraines-power-grid-targeted-transmission-stations-4530.html},}
\bibitem{2017crashoverridenostuxnet} ... \LaTeX:\\ \url{https://arstechnica.com/information-technology/2017/06/crash-override-malware-may-sabotage-electric-grids-but-its-no-stuxnet/},}
%%%%%%%%%%%%%%%%%%%%%%%%%%%%%%%%%%%%%%%%%%%%%%%%%%%%%%%%%%%%%%%%%



%
%\chapter{Deelonderzoek naar veiligheidsrisico's en eisen voor sluizen}
%
%Gevonden weblinks in google op 07-04-2023 met zoekopdracht: "veiligheidsrisico's voor sluizen en waterwerken"
%





\bibitem{ID} ... \LaTeX:\\ \url{https://www.tweedekamer.nl/downloads/document?id=80443e97-f17e-499c-b3f2-ad608f32e1aa&title=Rapportage%20Staat%20van%20de%20infra%20RWS%20%28definitief%29.pdf}
\bibitem{ID} ... \LaTeX:\\ \url{https://www.nu.nl/internet/5814282/rekenkamer-waterwerken-niet-goed-beveiligd-tegen-cyberaanvallen.html}
\bibitem{linne21122022onderhousluis} ... \LaTeX:\\ \url{https://www.deltalimburg.nl/article/9824/Onderhoudswerkzaamheden+aan+Sluis+Linne+afgerond}
\bibitem{sluisterneuzenveiligheid} ... \LaTeX:\\ \url{https://nieuwesluisterneuzen.eu/veiligheid}
\bibitem{mrdmarinesupport} ... \LaTeX:\\ \url{https://www.mrdmarinesupport.nl/nl/maritieme-dienstverlening/ondersteuning-veiligheid/}
\bibitem{krabbendam27052021infrasitrrisicobeoordeling} ... \LaTeX:\\ \url{https://www.infrasite.nl/bouwen/2021/05/27/veiligheid-voorop-begin-project-sluis-of-brug-altijd-met-risicobeoordeling/}
\bibitem{bedieningsluisVechterweerdVilsteren} ... \LaTeX:\\ \url{https://www.wdodelta.nl/bediening-schutsluizen-vechterweerd-en-vilsteren}
\bibitem{krabbedam21052021sluisHeel} ... \LaTeX:\\ \url{https://www.infrasite.nl/waterbouw-deltas/2021/05/21/sluis-heel-onder-handen-genomen/}
\bibitem{hdsr30092022lichtprojectieswaterliniesluizen} ... \LaTeX:\\ \url{https://www.hdsr.nl/actueel/nieuws/@154100/lichtprojecties-zetten-waterliniesluizen/}
\bibitem{nos28032019rekenkamerhackwaterwerk} ... \LaTeX:\\ \url{https://nos.nl/artikel/2277937-rekenkamer-hack-aanval-op-waterwerk-niet-altijd-opgemerkt}
\bibitem{sluisLinne08012023} ... \LaTeX:\\ \url{https://varendoejesamen.nl/kenniscentrum/artikel/onderhoud-sluis-linne-afgerond}
\bibitem{gww29032021kantelendesluisdeur} ... \LaTeX:\\ \url{https://www.gww-bouw.nl/artikel/de-eerste-sluis-met-kantelende-sluisdeur/}
\bibitem{tkh03022018sluisbeheerScadaZeeland} ... \LaTeX:\\ \url{https://tkhsecurity.com/nl/waterwerken/}
\bibitem{h2o28032019dreigingsniveaurekenkamer} ... \LaTeX:\\ \url{https://www.h2owaternetwerk.nl/h2o-actueel/rekenkamer-vitale-waterwerken-nog-onvoldoende-beschermd-tegen-cyberaanvallen}
\bibitem{krammersluizencomplex} ... \LaTeX:\\ \url{https://www.magazinesrijkswaterstaat.nl/bereikbaarzeeland/2021/01/krammersluizencomplex-verleden-heden-en-toekomst}



%\chapter{Deelonderzoek wet en regelgeving voor sluizen}
%
%
%
%Gevonden weblinks in google op 07-04-2023 met zoekopdracht: "wet en regelgeving voor sluizen"



\bibitem{ID} ... \LaTeX:\\ \url{https://www.hdsr.nl/publish/pages/86927/sluizen_in_of_bij_een_waterkering_-_uitvoeringsregels.pdf}
\bibitem{ID} ... \LaTeX:\\ \url{https://api1.ibabs.eu/publicdownload.aspx?site=sluis&id=100100292}
\bibitem{ID} ... \LaTeX:\\ \url{https://services.pilz.nl/wp-content/uploads/2021/12/brochure_bruggen_2018.pdf}
\bibitem{ID} ... \LaTeX:\\ \url{https://lokaleregelgeving.overheid.nl/CVDR375606/6}
\bibitem{ID} ... \LaTeX:\\ \url{https://zoek.officielebekendmakingen.nl/stb-2019-27.html}
\bibitem{ID} ... \LaTeX:\\ \url{https://a-quin.nl/nieuws/veiligheid-van-bruggen-sluizen-waarborgen-wie-wat-hoe/}
\bibitem{ID} ... \LaTeX:\\ \url{https://www.gemeentesluis.nl/Bestuur_en_Organisatie/Wetten_Regels_Bekendmakingen}
\bibitem{ID} ... \LaTeX:\\ \url{https://www.overijssel.nl/onderwerpen/verkeer-en-vervoer/varen-in-overijssel/informatie-bedieningstijden-sluizen-en-bruggen-noordwest-overijssel/}
\bibitem{ID} ... \LaTeX:\\ \url{https://www.rijkswaterstaat.nl/water/wetten-regels-en-vergunningen}
\bibitem{ID} ... \LaTeX:\\ \url{https://www.schuttevaer.nl/nieuws/actueel/2022/11/23/binnenvaart-zit-klem-tussen-regels-en-realiteit-kapotte-steigers-en-gesperde-sluizen-dwingen-tot-doorvaren/}
\bibitem{ID} ... \LaTeX:\\ \url{https://repository.officiele-overheidspublicaties.nl/CVDR/CVDR271406/1/html/CVDR271406_1.html}
\bibitem{ID} ... \LaTeX:\\ \url{https://www.zeeland.nl/actueel/bedieningstijden-sluizen-en-bruggen}
\bibitem{ID} ... \LaTeX:\\ \url{https://www.amsterdam.nl/verkeer-vervoer/varen-amsterdam/regels-varen/}
\bibitem{ID} ... \LaTeX:\\ \url{https://www.schielandendekrimpenerwaard.nl/wat-doen-we/regels-en-afspraken-over-beheer-keur-en-leggers/}
\bibitem{ID} ... \LaTeX:\\ \url{http://www.wetboek-online.nl/wet/Wet%20tot%20samenvoeging%20van%20de%20gemeenten%20Aardenburg%20en%20Sluis.html}
\bibitem{ID} ... \LaTeX:\\ \url{https://www.rijnland.net/regels-op-een-rij/richtlijnen-en-akkoorden/alle-regelgeving-van-rijnland/}
\bibitem{ID} ... \LaTeX:\\ \url{https://www.itbb.nl/diensten/advies-ce-markering-europese-richtlijnen/}
\bibitem{ID} ... \LaTeX:\\ \url{https://www.portofamsterdam.com/nl/scheepvaart/zeevaart/regelgeving}
\bibitem{ID} ... \LaTeX:\\ \url{https://www.watersportverbond.nl/nieuws/achterstallig-onderhoud-wachtplaatsen-bruggen-en-sluizen-zuid-holland-zorgelijk/}
\bibitem{ID} ... \LaTeX:\\ \url{https://varendoejesamen.nl/nieuws}
\bibitem{ID} ... \LaTeX:\\ \url{https://www.flevoland.nl/wat-doen-we/flevowegen-vlot-en-veilig-door-flevoland/water/varen-in-flevoland/bediening-bruggen-en-sluizen}
\bibitem{ID} ... \LaTeX:\\ \url{https://eur-lex.europa.eu/legal-content/NL/TXT/PDF/?uri=CELEX:32020L0012&from=DE}
\bibitem{ID} ... \LaTeX:\\ \url{https://www.werkenvoornederland.nl/organisatie/rijkswaterstaat/ict-middelen-maken-om-bruggen-sluizen-en-tunnels-te-besturen}
\bibitem{ID} ... \LaTeX:\\ \url{https://www.lobocom.nl/infra-bruggen-sluizen}
\bibitem{ID} ... \LaTeX:\\ \url{https://waterrecreatienederland.nl/content/uploads/2018/04/richtlijnen-vaarwegen-2017.pdf}
\bibitem{ID} ... \LaTeX:\\ \url{https://www.wetterskipfryslan.nl/melden-en-regelen/vergunningen-wetten-en-regels}
\bibitem{ID} ... \LaTeX:\\ \url{https://www.onlinezeilschool.nl/sluizen/}
\bibitem{ID} ... \LaTeX:\\ \url{https://www.provincie.drenthe.nl/onderwerpen/verkeer-vervoer/vaarwegen/rondje-drenthe/bedieningstijden/}





Rampen
In dit hoofdstuk worden de resultaten van een deskresearch naar verschillende rampen behandeld
Hierbij een verslag naar de oorzaken van de rampen, de werkwijze waarop het product is ontwikkeld, de verwerking van feedback, implementatie en nazorg.
Met behulp van het 4 variabelen model wordt duidelijk gemaakt hoe het systeem is opgezet en wat daarin verkeerd is gegaan.
Het  hoofdstuk wordt afgesloten met een analyse van algemene kenmerken van de verschillende rampen die zijn onderzocht.


Het 4 variabelen model kort toegelicht
Monitored variabelen: door sensoren gekwantificeerde fenomenen uit de omgeving, bijv temperatuur

Controlled variabelen: door actuatoren \bestuurde fenomenen uit de omgeving
For example, monitored variables might be the pressure and temperature
inside a nuclear reactor while controlled variables might be visual and audible alarms, as well
as the trip signal that initiates a reactor shutdown; whenever the temperature or pressure reach
abnormal values, the alarms go off and the shutdown procedure is initiated

Input variabelen: data die de software als input gebruikt
Here, IN models the input hardware interface (sensors and analog-to-digital converters) and
relates values of monitored variables to values of input variables in the software. The input variables model the information about the environment that is available to the software. For example,
IN might model a pressure sensor that converts temperature values to analog voltages; these voltages are then converted via an A/D converter to integer values stored in a register accesible to the
software.

Output variabelen: data die de software levert als output
The output hardware interface (digital-to-analog converters and actuators) is modelled
by OUT, which relates values of the output variables of the software to values of controlled variables. An output variable might be, for instance, a boolean variable set by the software with the
understanding that the value true indicates that a reactor shutdown should occur and the value
false indicates the opposite


Bronnen:
\bibitem{ID} ... \LaTeX:\\ \url{https://www.sciencedirect.com/science/article/pii/S0167642315001033}
\bibitem{ID} ... \LaTeX:\\ \url{https://www.cas.mcmaster.ca/~lawford/papers/AVoCS2013.pdf}
\bibitem{ID} ... \LaTeX:\\ \url{https://core.ac.uk/download/pdf/38891842.pdf}
%%%%%%%%%%%%%%%%%%%%%%%%%%%%%%%%%%%%%%%%%%%%%%%%%%%%%%%%%%%%%%%%%

Therac


sheets
\bibitem{ID} ... \LaTeX:\\ \url{https://web.cs.ucdavis.edu/~rogaway/classes/188/winter04/therac-25.pdf}
\bibitem{ID} ... \LaTeX:\\ \url{https://people.physics.carleton.ca/~drogers/egs_windows_collection/tsld008.htm}
\bibitem{ID} ... \LaTeX:\\ \url{https://en.wikipedia.org/wiki/Therac-25}
\bibitem{ID} ... \LaTeX:\\ \url{https://www.youtube.com/watch?v=-7gVqBY52MY}
reproduceren van de error. IN dit stuk wordt uitgelgd hoe het product werkt en waarom bepaalde beslssingen zijn genomen in de ontwerp/productiefase
\bibitem{ID} ... \LaTeX:\\ \url{https://www.bugsnag.com/blog/bug-day-race-condition-therac-25}
kort artikel met daarin een opsomming van alle fouten in het systeem en een korte uitleg
\bibitem{ID} ... \LaTeX:\\ \url{https://www.bowdoin.edu/~allen/courses/cs260/readings/therac.pdf}
uitgebreid artikel over hoe de fout werd gereproduceerd en de resultaten daaruit voortkwamen. Alsnog werden er na de reproductie fase nog meer fouten gevonden.
\bibitem{ID} ... \LaTeX:\\ \url{https://hackaday.com/2015/10/26/killed-by-a-machine-the-therac-25/}
artikel
\bibitem{ID} ... \LaTeX:\\ \url{https://ethicsunwrapped.utexas.edu/case-study/therac-25}
onderzoeksartikel waarin de bug wordt uitgelgd: de racecondities, de bytepositie en het testen worden berkitiseerd envenals andere onderdelen van het softwareproces
\bibitem{ID} ... \LaTeX:\\ \url{https://thedailywtf.com/articles/the-therac-25-incident}
onrealistisch testplan. In dit artikel egt de auteur het belang nog eens uit van goede requirements en implementatie, niet de software is waar het probleem ligt
\bibitem{ID} ... \LaTeX:\\ \url{https://www.computer.org/csdl/magazine/co/2017/11/mco2017110008/13rRUxAStVR}
onderzoeksrapport
\bibitem{ID} ... \LaTeX:\\ \url{https://web.stanford.edu/class/cs240/old/sp2014/readings/therac-25.pdf}
geschiedenis
\bibitem{ID} ... \LaTeX:\\ \url{http://computingcases.org/case_materials/therac/case_history/Case%20History.html}
artikel
\bibitem{ID} ... \LaTeX:\\ \url{https://medium.com/swlh/software-architecture-therac-25-the-killer-radiation-machine-8a05e0705d5b}
computer error. De ongeval en de malfunction nog een keer uitgelegd
\bibitem{ID} ... \LaTeX:\\ \url{http://www.ccnr.org/fatal_dose.html}
rapport
\bibitem{ID} ... \LaTeX:\\ \url{http://sunnyday.mit.edu/papers/therac.pdf}
\bibitem{ID} ... \LaTeX:\\ \url{https://pubmed.ncbi.nlm.nih.gov/101762/}
onderzoeksartkel
\bibitem{ID} ... \LaTeX:\\ \url{http://www1.cs.columbia.edu/~junfeng/08fa-e6998/sched/readings/therac25.pdf}
\bibitem{ID} ... \LaTeX:\\ \url{https://ieeexplore.ieee.org/document/274940}
uitgebreid artikel gaat hier ook wat meer over de hardware
\bibitem{ID} ... \LaTeX:\\ \url{https://www.linkedin.com/pulse/therac-25-industrial-design-engineering-systems-wang-ph-d-cre-acb/}
artikel waarin in 3 delen de problemaiekwordt blootgesteld
\bibitem{ID} ... \LaTeX:\\ \url{http://www.cse.msu.edu/~cse470/Public/Handouts/Therac/Therac_2.html}
case study sheets
\bibitem{ID} ... \LaTeX:\\ \url{https://www.cs.jhu.edu/~cis/cista/445/Lectures/Therac.pdf}
artikel waarin vooral de fabriikant ervan langs krijgt
\bibitem{ID} ... \LaTeX:\\ \url{http://users.csc.calpoly.edu/~jdalbey/SWE/Papers/THERAC25.html}
lessons learned. Vooral de begrippen betrouwbaarheid, welgevalligheid, veilgheid en gebruiksvriendelijkheid
\bibitem{ID} ... \LaTeX:\\ \url{https://bohr.wlu.ca/cp164/therac/therac25.htm}
root-cause analysis
\bibitem{ID} ... \LaTeX:\\ \url{https://root-cause-analysis.info/2010/08/08/therac-25-radiation-overdoses/}
case study
\bibitem{ID} ... \LaTeX:\\ \url{https://dusk.geo.orst.edu/ethics/papers/Therac.Huff.pdf}
\bibitem{ID} ... \LaTeX:\\ \url{https://embeddedartistry.com/fieldatlas/case-study-therac-25/}
case study
\bibitem{ID} ... \LaTeX:\\ \url{https://www.sebokwiki.org/wiki/Medical_Radiation}
opzetten van systematische acceptaatie test met therac als voorbeeld
\bibitem{ID} ... \LaTeX:\\ \url{https://www.sciencedirect.com/science/article/pii/S1474667017448245}
artikel waarin een diagnose plaatvindt voor het bedrijf en de ingenieur/ontwerper
\bibitem{ID} ... \LaTeX:\\ \url{https://magsilva.pro.br/apps/wiki/testing/Therac_25}
rapport
\bibitem{ID} ... \LaTeX:\\ \url{http://csel.eng.ohio-state.edu/productions/pexis/readings/submod3/therac.pdf}
oorzaken aangegeven in artikel
\bibitem{ID} ... \LaTeX:\\ \url{https://www.chemeurope.com/en/encyclopedia/Therac-25.html}
het onderzoek en enkele ontwerptekeningen en oplossingen
\bibitem{ID} ... \LaTeX:\\ \url{https://pvs-studio.com/en/blog/posts/0438/}
\bibitem{ID} ... \LaTeX:\\ \url{https://www.coursera.org/lecture/software-design-threats-mitigations/therac-25-case-study-VmQPa}
\bibitem{ID} ... \LaTeX:\\ \url{https://www.semanticscholar.org/paper/The-story-of-the-Therac-25-in-LOTOS-Thomas/6c9c6024cf95aadae8b7edf1160e0e4500410eb9}
\bibitem{ID} ... \LaTeX:\\ \url{https://news.ycombinator.com/item?id=21679287}
wiki
\bibitem{ID} ... \LaTeX:\\ \url{https://en.wikibooks.org/wiki/Professionalism/Therac-25}
analyse
\bibitem{ID} ... \LaTeX:\\ \url{https://citeseerx.ist.psu.edu/viewdoc/download?doi=10.1.1.96.369&rep=rep1&type=pdf}
samenvatting
\bibitem{ID} ... \LaTeX:\\ \url{https://onlineethics.org/cases/resources-engineering-and-science-ethics/investigation-therac-25-accidents-abstract}
podcast
\bibitem{ID} ... \LaTeX:\\ \url{https://podcasts.apple.com/gb/podcast/therac-25/id1046978749?i=1000514115050}
enkele conclusies
\bibitem{ID} ... \LaTeX:\\ \url{http://www.cas.mcmaster.ca/~se4d03/therac.html}
rapport over de fouten die de verschillende partijen hebben gemaakt( overheid, ingenieurs, bedrijf, operators) en de verbeterpunten
\bibitem{ID} ... \LaTeX:\\ \url{https://www.cs.colostate.edu/~bieman/CS314/Notes/therac25.pdf}
onderzoeksrapport
\bibitem{ID} ... \LaTeX:\\ \url{https://www.cs.ucf.edu/~dcm/Teaching/COP4600-Fall2010/Literature/Therac25-Leveson.pdf}
slides online over het technisch mankement
Wat is er gebeurd, nou het volgende:
Normal radiation treatments: 6,000 rads over a 3 week period, under certain conditions Therac-25 was delivering 60,000 rads during one session.
En wat ging er mis?
Paradigm Shift

Therac-25 replaced expensive hardware safety interlocks with software controls

Real-time software
Design
Race condition caused focusing element to be incorrectly set
No indication of actual hardware settings
Error messages appeared the same regardless of how important
Error messages were difficult to understand
All errors messages could be manually overridden

\bibitem{ID} ... \LaTeX:\\ \url{https://hci.cs.siue.edu/NSF/Files/Semester/Week13-2/PPT-Text/Slide13.html}
oorzaak-gevolg diagram
\bibitem{ID} ... \LaTeX:\\ \url{https://www.thinkreliability.com/InstructorBlogs/Blog-Therac-25.pdf}
veiligheidsanalyse naar de rapportage van foutmeldingen, de beslissingsmatrix waarmee het programma wordt uitgevoerd en de software-analyse door een consultat
\bibitem{ID} ... \LaTeX:\\ \url{https://www.erenkrantz.com/Geeks/Therac-25/Side_bar_5.html}
\bibitem{ID} ... \LaTeX:\\ \url{https://bandcamp.com/therac-25}
\bibitem{ID} ... \LaTeX:\\ \url{https://sites.nd.edu/brent-marin/2017/09/21/blog-post-8-therac-25-accidents/}
\bibitem{ID} ... \LaTeX:\\ \url{https://europepmc.org/article/med/4035720}
\bibitem{ID} ... \LaTeX:\\ \url{https://itlaw.wikia.org/wiki/Therac-25}
\bibitem{ID} ... \LaTeX:\\ \url{https://www.sitepoint.com/therac-25-bad-software-kills/}
\bibitem{ID} ... \LaTeX:\\ \url{https://sqa.stackexchange.com/questions/9798/asking-for-help-with-this-therac-25-bugged-code-i-dont-understand-the-explanat}
\bibitem{ID} ... \LaTeX:\\ \url{http://www.se.rit.edu/~swen-342/activities/TheracIndividual.html}
\bibitem{ID} ... \LaTeX:\\ \url{https://www.designnews.com/automation-motion-control/engineering-disasters-deadly-zaps-therac-25}

publicaties over de therac:
\bibitem{ID} ... \LaTeX:\\ \url{https://www.cs.colostate.edu/~bieman/CS314/Notes/therac25.pdf}
\bibitem{ID} ... \LaTeX:\\ \url{https://www.cs.ucf.edu/~dcm/Teaching/COP4600-Fall2010/Literature/Therac25-Leveson.pdf}
\bibitem{ID} ... \LaTeX:\\ \url{https://onlineethics.org/cases/resources-engineering-and-science-ethics/investigation-therac-25-accidents-abstract}
\bibitem{ID} ... \LaTeX:\\ \url{https://citeseerx.ist.psu.edu/viewdoc/download?doi=10.1.1.96.369&rep=rep1&type=pdf}
\bibitem{ID} ... \LaTeX:\\ \url{http://csel.eng.ohio-state.edu/productions/pexis/readings/submod3/therac.pdf}
\bibitem{ID} ... \LaTeX:\\ \url{https://dusk.geo.orst.edu/ethics/papers/Therac.Huff.pdf}
\bibitem{ID} ... \LaTeX:\\ \url{https://www.cs.jhu.edu/~cis/cista/445/Lectures/Therac.pdf}
\bibitem{ID} ... \LaTeX:\\ \url{http://www1.cs.columbia.edu/~junfeng/08fa-e6998/sched/readings/therac25.pdf}
\bibitem{ID} ... \LaTeX:\\ \url{http://sunnyday.mit.edu/papers/therac.pdf}
\bibitem{ID} ... \LaTeX:\\ \url{https://web.stanford.edu/class/cs240/old/sp2014/readings/therac-25.pdf}




%%%%%%%%%%%%%%%%%%%%%%%%%%%%%%%%%%%%%%%%%%%%%%%%%%%%%%%%%%%%%%%%%
Krakend zorgssteem door covid-19 in suriname

vaccinatieterkort
communicatie met bevolking
communicatie met binnenland
testen van vaccinaties
besmetting vanuit eht buitenland
isolatie na vakantie en voor toeristen
tekort aan ic-personeel
tekort aan ic-bedden
tekort aan zuurtstof
tekort aan middelen}

Wat blijkt hieruit:
de impact van de crisis wereldwijd
de afhnakelijkheid van landen op goede samenwerking
Nut en noodzaak van regelgeving
Naveling van maatregels
Communicatie over beleid vanuit de overheid naar de burgers
Belang van een verzorgingstaat
Een wetenschappelijke ontwikkeling die kan inspelen op gevoelige 'trends
De impact van een lockdown op de economie
Afschaling van andere noodzakelijke no-covid zorg
De bereikbaarheid van een ziekenhuis
Waar heeft het toe geleid?
\bibitem{ID} ... \LaTeX:\\ \url{https://www.waterkant.net/suriname/2007/02/06/school-in-suriname-gesloten-om-zenuwgasvoorraad/}
\bibitem{ID} ... \LaTeX:\\ \url{https://nl.wikipedia.org/wiki/Nationaal_Co%C3%B6rdinatiecentrum_voor_Rampenbeheersing}
\bibitem{ID} ... \LaTeX:\\ \url{https://www.examenkamer.nl/index.php/27-vca-examens-in-suriname}
%%%%%%%%%%%%%%%%%%%%%%%%%%%%%%%%%%%%%%%%%%%%%%%%%%%%%%%%%%%%%%%%%

Waterramp suriname met cyanice


%%%%%%%%%%%%%%%%%%%%%%%%%%%%%%%%%%%%%%%%%%%%%%%%%%%%%%%%%%%%%%%%%
boeing 737 crashes


algemene vragen
oorzaken
\bibitem{gates18112020boeingcrisis} ... \LaTeX:\\ \url{https://www.seattletimes.com/business/boeing-aerospace/what-led-to-boeings-737-max-crisis-a-qa/}
\bibitem{boeing737maxsoftwareprobles} ... \LaTeX:\\ \url{https://www.schneier.com/blog/archives/2019/04/excellent_analy.html}
fout in de software
\bibitem{avetisov19032019boeingmalwarestate} ... \LaTeX:\\ \url{https://www.forbes.com/sites/georgeavetisov/2019/03/19/malware-at-30000-feet-what-the-737-max-says-about-the-state-of-airplane-software-security/?sh=4d26f7052a9e}
het nationaal veiligheidsbelang
\bibitem{thompson23112020nationalsecurityboeing} ... \LaTeX:\\ \url{https://www.forbes.com/sites/lorenthompson/2020/11/23/five-reasons-return-of-boeings-737-max-to-service-is-important-to-national-security/?sh=2128ea552018}
falend toezicht
\bibitem{gates21032019FAAControlSystem} ... \LaTeX:\\ \url{https://www.seattletimes.com/business/boeing-aerospace/failed-certification-faa-missed-safety-issues-in-the-737-max-system-implicated-in-the-lion-air-crash/}
onderzoeksrapport
\bibitem{faa18112020boeingreview} ... \LaTeX:\\ \url{https://www.faa.gov/foia/electronic_reading_room/boeing_reading_room/media/737_RTS_Summary.pdf}
\bibitem{wiki737maxgroundings} ... \LaTeX:\\ \url{https://en.wikipedia.org/wiki/Boeing_737_MAX_groundings}
veiligheidsrisico's
menselijke fouten
\bibitem{campbell02052019boengcrashhumanerrors} ... \LaTeX:\\ \url{https://www.theverge.com/2019/5/2/18518176/boeing-737-max-crash-problems-human-error-mcas-faa}
overzicht van crashes
\bibitem{hawkins22032019737maxairplanes} ... \LaTeX:\\ \url{https://www.theverge.com/2019/3/22/18275736/boeing-737-max-plane-crashes-grounded-problems-info-details-explained-reasons}
veiligheidsopmerking
\bibitem{thomas30082020737safest} ... \LaTeX:\\ \url{https://www.airlineratings.com/news/boeings-737-max-will-one-safest-aircraft-history/}
aanpassingen
\bibitem{boeing737maxdisplay} ... \LaTeX:\\ \url{https://www.boeing.com/commercial/737max/737-max-software-updates.page}
waarschuwingen//output signalen
\bibitem{fehrm24112020737changes} ... \LaTeX:\\ \url{https://leehamnews.com/2020/11/24/boeing-737-max-changes-beyond-mcas/}
software gerelateerde fouten
\bibitem{travis18042019737maxsoftwaredevop} ... \LaTeX:\\ \url{https://spectrum.ieee.org/aerospace/aviation/how-the-boeing-737-max-disaster-looks-to-a-software-developer}
onderzoeksrapport
de rol van de publieke opinie
\bibitem{barnett05052019737maxcrisis} ... \LaTeX:\\ \url{https://pubsonline.informs.org/do/10.1287/orms.2019.05.05/full/}
onderzoek van europese luchtvaart agentschap
\bibitem{easa27012021737maxsafereturn} ... \LaTeX:\\ \url{https://www.easa.europa.eu/newsroom-and-events/news/easa-declares-boeing-737-max-safe-return-service-europe}
veiligheidsvraagstuk
\bibitem{touitou11032019737tragedies} ... \LaTeX:\\ \url{https://phys.org/news/2019-03-boeing-max-safety-tragedies.html}
artikel over sensoren
\bibitem{hemmerdinger02022021737maxdeliveries} ... \LaTeX:\\ \url{https://www.flightglobal.com/airframers/boeing-delays-737-max-10-deliveries-two-years-to-2023/142245.article}
goedkeuring van europese luchtvaart autoriteiten
advies aan de faa
\bibitem{bielby27022021faaimprovesafety} ... \LaTeX:\\ \url{https://www.hstoday.us/subject-matter-areas/airport-aviation-security/oig-tells-faa-to-improve-safety-oversight-following-boeing-737-max-review/}
\bibitem{boyle18112020737maxupgrade} ... \LaTeX:\\ \url{https://www.geekwire.com/2020/faas-go-ahead-737-maxs-return-flight-kicks-off-massive-software-upgrade/}
\bibitem{bergstraburgess122019737maxMcasAlgorithm} ... \LaTeX:\\ \url{https://www.researchgate.net/publication/338420944_A_Promise_Theoretic_Account_of_the_Boeing_737_Max_MCAS_Algorithm_Affair}
achtergrond informatie
\bibitem{737mcas} ... \LaTeX:\\ \url{http://www.b737.org.uk/mcas.htm}
algemeen vertrouwen
\bibitem{newburger17052019boeingcrisis} ... \LaTeX:\\ \url{https://www.cnbc.com/2019/05/16/what-you-need-to-know-about-boeings-737-max-crisis.html}
toestemming europese autoriteiten
problemen
\bibitem{arstechnica22012020737problems} ... \LaTeX:\\ \url{https://arstechnica.com/information-technology/2020/01/737-max-fix-slips-to-summer-and-thats-just-one-of-boeings-problems/}
uitgebreid artikel over de onderzoeken en het vliegverbod
\bibitem{german190620217372yaftergrounded} ... \LaTeX:\\ \url{https://www.cnet.com/news/boeing-737-max-8-all-about-the-aircraft-flight-ban-and-investigations/}
computers als oorzaak
lessons learned
\bibitem{beningo02052019boeinglessons} ... \LaTeX:\\ \url{https://www.designnews.com/electronics-test/5-lessons-learn-boeing-737-max-fiasco}
\bibitem{duran05042019boeingspof} ... \LaTeX:\\ \url{https://www.eurocontrol.int/publication/effects-network-extra-standby-aircraft-and-boeing-737-max-grounding}
single point of failure
\bibitem{ID} ... \LaTeX:\\ \url{https://dmd.solutions/blog/2019/04/05/how-a-single-point-of-failure-spof-in-the-mcas-software-could-have-caused-the-boeing-737-max-crash-in-ethiopia/}
\bibitem{makichuck24012021737fearflying} ... \LaTeX:\\ \url{https://asiatimes.com/2021/01/boeings-737-max-and-the-fear-of-flying/}
lijst van tehnische aanpassingen
\bibitem{caa737modifications} ... \LaTeX:\\ \url{https://www.caa.co.uk/Consumers/Guide-to-aviation/Boeing-737-MAX/}
\bibitem{oestergaard14122020boeingdeliveries} ... \LaTeX:\\ \url{https://dsm.forecastinternational.com/wordpress/2020/12/14/airbus-and-boeing-report-november-2020-commercial-aircraft-orders-and-deliveries/}
code lek
\bibitem{reenberg787flaws} ... \LaTeX:\\ \url{https://www.wired.com/story/boeing-787-code-leak-security-flaws/}
\bibitem{fitch16092020737backlogrisks} ... \LaTeX:\\ \url{https://www.fitchratings.com/research/corporate-finance/boeing-737-max-return-backlog-risks-remain-16-09-2020}
Cultuurverandering, deregulatie, systeemwijziging of gewoon een kwestie van competentie
\bibitem{willis27082020737maxfailures} ... \LaTeX:\\ \url{https://www.aerospacetestinginternational.com/features/what-broke-the-737-max.html}
extra aanpassingen
\bibitem{ostrower11062020more737changes} ... \LaTeX:\\ \url{https://theaircurrent.com/aviation-safety/boeings-737-max-software-done-but-regulators-plot-more-changes-after-jets-return/}
wat ging er mis een analyse van een ex-iloot
De utoriteiten waren op de hoogte
\bibitem{hruska13122019faaknown737crashrate} ... \LaTeX:\\ \url{https://www.extremetech.com/extreme/303373-the-faa-knew-the-737-max-was-dangerous-and-kept-it-flying-anyway}
kwaliteiten van het alarmsysteem niet goed bekend
\bibitem{bloomberg26092019failedpred} ... \LaTeX:\\ \url{https://time.com/5687473/boeing-737-alarm-system/}
\bibitem{whiteman09072020boengcancelstock} ... \LaTeX:\\ \url{https://www.nasdaq.com/articles/boeing-gets-dealt-another-737-max-cancellation-blow.-what-it-means-for-boeing-stock-2020}
\bibitem{leopold09192019boeingreliability} ... \LaTeX:\\ \url{https://www.eetimes.com/boeing-crashes-highlight-a-worsening-reliability-crisis/}
veiligheidsvraagstuk
\bibitem{koenig11122019737crashesnofix} ... \LaTeX:\\ \url{https://www.latimes.com/business/story/2019-12-11/faa-boeing-737-max-crashes}
probleemanalyse, veiligheidsvraagstuk
\bibitem{dohertylindeman15032019737problems} ... \LaTeX:\\ \url{https://www.politico.com/story/2019/03/15/boeing-737-max-grounding-1223072}
falend toezicht
\bibitem{stodder02102019corruptoversight} ... \LaTeX:\\ \url{https://www.pogo.org/analysis/2019/10/corrupted-oversight-the-faa-boeing-and-the-737-max/}
\bibitem{afacwaLostSafeguards} ... \LaTeX:\\ \url{https://www.afacwa.org/the_inside_story_of_mcas_seattle_times}
doelstellingen en veiligheidsvraagstukken
\bibitem{swayne18032019profitssafety} ... \LaTeX:\\ \url{https://www.marxist.com/737-max-scandal-boeing-putting-profits-before-safety.htm}
\bibitem{freed26022021liftaustraliaban} ... \LaTeX:\\ \url{https://finance.yahoo.com/news/australia-lifts-ban-boeing-737-035817682.html?guccounter=1&guce_referrer=aHR0cHM6Ly93d3cuZ29vZ2xlLmNvbS8&guce_referrer_sig=AQAAAHZCJYy_0A5VS2WiPoCvH4xdrRNkmkdsv5EWJ2RLIz_AS-rxsTty6AF1_HlmJiRyWYqCXDi4p0Xs4isYkNkCq2Pfo-pQ60Xz_IfTNjm4FgoZiBMC4zpZlB6F0fwecrjE_ujAXZzG4xPJnWCd8-G3VLlPTY8h3H31eQ1i8hY9AIyy}
autoriteiten krijgen tik op de vingers
\bibitem{reed15032019softwareattention} ... \LaTeX:\\ \url{https://medium.com/@jpaulreed/the-737max-and-why-software-engineers-should-pay-attention-a041290994bd}
\bibitem{news17032019softwareexplains} ... \LaTeX:\\ \url{https://news.ycombinator.com/item?id=19414775}
\bibitem{legget21122020eu737maxsafe} ... \LaTeX:\\ \url{https://www.bbc.com/news/55366320}
\bibitem{marketscreener0103221737chinarecertification} ... \LaTeX:\\ \url{https://www.marketscreener.com/news/latest/China-studies-Boeing-737-MAX-recertification-wants-safety-concerns-fully-addressed--32569125/}
motor in brand
\bibitem{euractiv22022021737firegrounds} ... \LaTeX:\\ \url{https://www.euractiv.com/section/aviation/news/boeing-grounds-777s-after-engine-fire/}
\bibitem{benny18022019737returnUAE} ... \LaTeX:\\ \url{https://gulfnews.com/business/aviation/uae-airspace-to-see-return-of-boeing-737-max-1.1613627548923}
motor in brand gevlogen
\bibitem{biersmichel22022021777grounds} ... \LaTeX:\\ \url{https://techxplore.com/news/2021-02-boeing-urges-grounding-777s.html}
\bibitem{ID} ... \LaTeX:\\ \url{https://www.politico.eu/article/uk-temporarily-bans-some-boeing-aircraft-after-pratt-whitney-engine-incidents/}
\bibitem{reuters23022021777metalfatigue} ... \LaTeX:\\ \url{https://www.timeslive.co.za/news/world/2021-02-23-damage-to-united-boeing-777-engine-consistent-with-metal-fatigue--ntsb/}
faa was niet kritisch genoeg
\bibitem{ID} ... \LaTeX:\\ \url{https://federalnewsnetwork.com/government-news/2021/02/federal-watchdog-blasts-faa-over-certification-of-boeing-jet/}






%%%%%%%%%%%%%%%%%%%%%%%%%%%%%%%%%%%%%%%%%%%%%%%%%%%%%%%%%%%%%%%%%

china explosion 2015 tianjin

verhaal van brandweermannen
\bibitem{ID} ... \LaTeX:\\ \url{https://slate.com/human-interest/2015/08/chinese-explosion-aftermath-officials-investigate-causes-behind-warehouse-blast-and-death-of-88-firefighters.html}
artikel
\bibitem{ID} ... \LaTeX:\\ \url{https://www.chinafile.com/conversation/tianjin-explosion}
invloed van social media
\bibitem{ID} ... \LaTeX:\\ \url{https://www.economist.com/asia/2015/08/18/a-blast-in-tianjin-sets-off-an-explosion-online}
\bibitem{ID} ... \LaTeX:\\ \url{https://america.cgtn.com/2015/08/12/explosion-reported-in-tianjin-china}
\bibitem{ID} ... \LaTeX:\\ \url{https://factcheck.afp.com/no-photo-was-taken-chinese-city-tianjin-august-2015}
vergelijking van twee rampen
\bibitem{ID} ... \LaTeX:\\ \url{https://airshare.air-inc.com/how-does-the-beirut-explosion-compare-to-tianjin}
overheid en media
\bibitem{ID} ... \LaTeX:\\ \url{https://newbloommag.net/2015/08/17/tianjin-explosion/}
chemische industrie ondeer de loep
\bibitem{ID} ... \LaTeX:\\ \url{https://www.voanews.com/east-asia-pacific/tianjin-blast-puts-spotlight-chemical-industry}
\bibitem{ID} ... \LaTeX:\\ \url{https://abcnews.go.com/International/apocalyptic-aftermath-devastating-images-tianjin-china-explosions/story?id=33057017}
\bibitem{ID} ... \LaTeX:\\ \url{https://www.reachingoutacrossdurham.co.uk/osk/tianjin-explosion-2021}
\bibitem{ID} ... \LaTeX:\\ \url{https://aiche.onlinelibrary.wiley.com/doi/abs/10.1002/prs.11789}
\bibitem{ID} ... \LaTeX:\\ \url{https://www.automotivelogistics.media/thousands-of-cars-destroyed-in-tianjin-port-explosions/13570.article}
\bibitem{ID} ... \LaTeX:\\ \url{https://www.joc.com/port-news/asian-ports/port-tianjin/tianjin-port-explosions-could-be-most-expensive-maritime-disaster_20150826.html}
\bibitem{ID} ... \LaTeX:\\ \url{https://www.bloomberg.com/news/articles/2015-08-12/explosion-in-northern-china-shatters-windows-causes-injuries}
\bibitem{ID} ... \LaTeX:\\ \url{https://unece.org/fileadmin/DAM/env/documents/2016/TEIA/OECD_WGCA_24-27_OCT_2016/Session_3_Zhao_-__Introduction_of_Tianjin_Accident_-_Jinsong_Zhao.pdf}
gemaakte fouten
\bibitem{ID} ... \LaTeX:\\ \url{https://porteconomicsmanagement.org/pemp/contents/part6/port-resilience/site-2015-tianjin-port-explosions/}
\bibitem{ID} ... \LaTeX:\\ \url{https://www.alamy.com/stock-image-tianjin-china-17th-aug-2015-tianjin-explosion-aftermath-blast-site-165334778.html}
\bibitem{ID} ... \LaTeX:\\ \url{https://www.popularmechanics.com/technology/news/a16871/massive-explosions-china-city-of-tianjin/}
\bibitem{ID} ... \LaTeX:\\ \url{https://www.imago-images.com/st/0080815934}
\bibitem{ID} ... \LaTeX:\\ \url{https://www.chemistryworld.com/news/deadly-chemical-blast-at-chinese-port/8857.article}
\bibitem{ID} ... \LaTeX:\\ \url{https://www.process-worldwide.com/tianjin-explosion-from-chemical-perspective-insights-and-backgrounds-a-502381/}
vergelijking met andere explosies
\bibitem{ID} ... \LaTeX:\\ \url{https://apnews.com/article/lebanon-fires-us-news-explosions-middle-east-53f4206a7f1db0812262a15d22e1e58f}
invloed van de ramp op de industrie
\bibitem{ID} ... \LaTeX:\\ \url{https://fortune.com/2015/08/14/tianjin-port-explosion-shipping-delays/}
is er sprake van een doofpot
\bibitem{ID} ... \LaTeX:\\ \url{https://www.washingtontimes.com/news/2015/aug/20/inside-china-tianjin-explosions-cover-up-exposes-b/}
eigendomsverzekering
\bibitem{ID} ... \LaTeX:\\ \url{https://www.artemis.bm/news/tianjin-explosions-property-insurance-loss-could-reach-3-5bn-swiss-re/}
\bibitem{ID} ... \LaTeX:\\ \url{https://www.thechinastory.org/yearbooks/yearbook-2015/forum-the-abyss-%E5%9D%8E/tianjin-explosions/}
effecten op de lange termijn
\bibitem{ID} ... \LaTeX:\\ \url{https://www.flexport.com/blog/tianjin-explosion-effect-on-supply-chains/}
\bibitem{ID} ... \LaTeX:\\ \url{https://www.cicm.org.my/images/articles/CICM-Article-on-Tianjin-Blast-Oct2015.pdf}
lessons learned
\bibitem{ID} ... \LaTeX:\\ \url{https://www.genre.com/knowledge/blog/lessons-from-the-tianjin-explosion-en.html}
\bibitem{ID} ... \LaTeX:\\ \url{https://www.ft.com/content/ad62904c-44ce-11e5-b3b2-1672f710807b}
\bibitem{ID} ... \LaTeX:\\ \url{https://www.huffingtonpost.co.uk/2015/08/13/tianjin-explosion-china-shocking-footage-caught-on-camera_n_7980888.html}
\bibitem{ID} ... \LaTeX:\\ \url{https://www.thatsmags.com/china/post/19189/massive-fire-rocks-tianjin-port}
gevolgen voor de industrie
\bibitem{ID} ... \LaTeX:\\ \url{https://www.everstream.ai/risk-center/special-reports/the-jiangsu-yancheng-explosion/}
\bibitem{ID} ... \LaTeX:\\ \url{https://www.newyorker.com/news/news-desk/after-tianjin-an-outbreak-of-mistrust-in-china}
framing vanuit de chinese media
\bibitem{ID} ... \LaTeX:\\ \url{https://www.neliti.com/publications/101997/the-chinese-media-framing-of-the-2015s-tianjin-explosion}
\bibitem{ID} ... \LaTeX:\\ \url{https://www.reinsurancene.ws/chinese-insurers-settle-1-5-billion-tianjin-blast-claims/}
niewsartikel
\bibitem{ID} ... \LaTeX:\\ \url{https://www.thechemicalengineer.com/news/update-78-confirmed-dead-after-chinese-chemicals-plant-explosion/}
\bibitem{ID} ... \LaTeX:\\ \url{https://www.caixinglobal.com/2016-11-10/chinese-executive-receives-suspended-death-sentence-over-2015-tianjin-warehouse-blast-101006325.html}
toegang tot de ramplplek vanuit de okale journalistiek
\bibitem{ID} ... \LaTeX:\\ \url{https://chinadigitaltimes.net/2015/08/he-xiaoxin-how-far-can-i-go-and-how-much-can-i-do/}
artikel
\bibitem{ID} ... \LaTeX:\\ \url{https://www.wnpr.org/post/china-examines-aftermath-immense-twin-explosions-killed-dozens}
\bibitem{ID} ... \LaTeX:\\ \url{https://theconversation.com/what-is-ammonium-nitrate-the-chemical-that-exploded-in-beirut-143979}
\bibitem{ID} ... \LaTeX:\\ \url{https://chemicalwatch.com/36730/nationwide-inspections-in-china-follow-tianjin-explosion}
\bibitem{ID} ... \LaTeX:\\ \url{https://www.thehindu.com/news/international/investigation-begun-into-china-gas-explosion-as-toll-rises/article34818324.ece}
\bibitem{ID} ... \LaTeX:\\ \url{https://santiagotimes.cl/2019/03/24/64-killed-600-injured-in-china-chemical-plant-blast/}
oorzaken
\bibitem{ID} ... \LaTeX:\\ \url{https://klingecorp.com/blog/what-caused-the-tianjin-explosions/}
case study
\bibitem{ID} ... \LaTeX:\\ \url{https://www.preventionweb.net/educational/view/57235}
niewsartikel
\bibitem{ID} ... \LaTeX:\\ \url{https://www.cnbc.com/2015/08/12/explosion-in-tianjin-china.html}
chronologische uiteenzetting
\bibitem{ID} ... \LaTeX:\\ \url{https://www.aria.developpement-durable.gouv.fr/wp-content/files_mf/A46803_a46803_fiche_impel_006.pdf}
corruptie
\bibitem{ID} ... \LaTeX:\\ \url{https://www.nytimes.com/2015/08/31/world/asia/behind-tianjin-tragedy-a-company-that-flouted-regulations-and-reaped-profits.html}
mismanagement als oorzaak
\bibitem{ID} ... \LaTeX:\\ \url{https://www.nytimes.com/2016/02/06/world/asia/tianjin-explosions-were-result-of-mismanagement-china-finds.html}
\bibitem{ID} ... \LaTeX:\\ \url{https://cen.acs.org/articles/94/web/2016/02/Chinese-Investigators-Identify-Cause-Tianjin.html}
autoriteiten publiceren onderoeksrapport
\bibitem{ID} ... \LaTeX:\\ \url{https://cen.acs.org/articles/94/i7/Chinese-Investigators-Identify-Cause-Tianjin.html}
fotos van de rampplek
\bibitem{ID} ... \LaTeX:\\ \url{https://www.theatlantic.com/photo/2015/08/photos-of-the-aftermath-of-the-massive-explosions-in-tianjin-china/401228/}
\bibitem{ID} ... \LaTeX:\\ \url{https://edition.cnn.com/2015/08/13/asia/china-tianjin-explosions/index.html}
niuwesartiekel}
\bibitem{ID} ... \LaTeX:\\ \url{https://www.cbc.ca/news/world/china-explosion-tianjin-1.3189455}
verantwoordelijke
\bibitem{ID} ... \LaTeX:\\ \url{https://www.thestar.com/news/world/2016/11/09/chinese-executive-gets-death-sentence-over-tianjin-explosion-in-2015.html}
risicobeperking/controle
\bibitem{ID} ... \LaTeX:\\ \url{https://www.swissre.com/en/china/news-insights/articles/analysis-of-tianjin-port-explosion-china.html}
censuur
\bibitem{ID} ... \LaTeX:\\ \url{https://foreignpolicy.com/2015/09/10/censored-china-young-survivor-tianjin-explosion-viral-post/}
censuur
\bibitem{ID} ... \LaTeX:\\ \url{https://qz.com/756872/a-year-after-the-tianjin-blast-public-mourning-and-discussion-about-it-are-still-censored-in-china/}
verschillende artikelen
\bibitem{ID} ... \LaTeX:\\ \url{https://www.scmp.com/topics/tianjin-warehouse-explosion-2015}
\bibitem{ID} ... \LaTeX:\\ \url{https://www.wsj.com/articles/BL-CJB-27664}
\bibitem{ID} ... \LaTeX:\\ \url{https://www.nbcnews.com/news/world/tianjin-explosions-californian-witness-filmed-dramatic-china-blasts-n409701}
\bibitem{ID} ... \LaTeX:\\ \url{https://ui.adsabs.harvard.edu/abs/2016AGUFM.S13D..06P/abstract}
afwikkeling van de ramp
\bibitem{ID} ... \LaTeX:\\ \url{https://chinadialogue.net/en/pollution/9188-back-to-the-blast-zone-one-year-after-the-tianjin-explosion/}
\bibitem{ID} ... \LaTeX:\\ \url{https://www.wired.com/2015/08/chinas-huge-tianjin-explosion-looked-like-space/}
\bibitem{ID} ... \LaTeX:\\ \url{https://www.abc.net.au/news/2015-08-13/explosion-rocks-north-chinese-city-of-tianjin/6693336?nw=0}
ambtenaren onderzocht
\bibitem{ID} ... \LaTeX:\\ \url{https://thediplomat.com/2015/08/23-executives-government-officials-under-investigation-for-role-in-tianjin-explosions/}
\bibitem{ID} ... \LaTeX:\\ \url{http://america.aljazeera.com/articles/2015/8/13/at-least-50-dead-and-hundreds-injured-in-chinese-warehouse-explosion.html}
risico-inschatting
\bibitem{ID} ... \LaTeX:\\ \url{https://www.mdpi.com/2071-1050/12/3/1169/htm}
\bibitem{ID} ... \LaTeX:\\ \url{https://www.mdpi.com/2071-1050/12/3/1169/htm}
\bibitem{ID} ... \LaTeX:\\ \url{https://www.cbsnews.com/news/tianjin-port-china-massive-explosion-hundreds-injured/}
\bibitem{ID} ... \LaTeX:\\ \url{https://www.hkjcdpri.org.hk/download/casestudies/Tianjin_CASE.pdf}
\bibitem{ID} ... \LaTeX:\\ \url{https://time.com/3996168/tianjin-explosion-china-pictures/}
onderzoeksrapport
\bibitem{ID} ... \LaTeX:\\ \url{https://www.hfw.com/Tianjin-Port-explosion-August-2015}
\bibitem{ID} ... \LaTeX:\\ \url{https://news.un.org/en/story/2015/08/506912-following-tianjin-explosion-un-expert-calls-china-ensure-transparent}
\bibitem{ID} ... \LaTeX:\\ \url{https://www.france24.com/en/20150812-huge-explosions-rock-chinese-city-tianjin}
\bibitem{ID} ... \LaTeX:\\ \url{https://choice.npr.org/index.html?origin=https://www.npr.org/2015/08/14/432280627/what-caused-the-warehouse-explosions-in-tianjin-china}
123 verantwoordelijken
\bibitem{ID} ... \LaTeX:\\ \url{https://www.bbc.com/news/world-asia-china-35506311}
\bibitem{ID} ... \LaTeX:\\ \url{https://www.washingtonpost.com/gdpr-consent/?next_url=https%3a%2f%2fwww.washingtonpost.com%2fnews%2fworldviews%2fwp%2f2015%2f08%2f12%2fvideos-show-chinese-city-of-tianjin-rocked-by-enormous-explosion%2f}
lang artiekel
\bibitem{ID} ... \LaTeX:\\ \url{https://www.businessinsider.com/the-chemical-explosion-in-china-killed-more-than-100-people-and-the-devastation-is-unreal-2015-8?international=true&r=US&IR=T}
\bibitem{ID} ... \LaTeX:\\ \url{https://pubmed.ncbi.nlm.nih.gov/27311537/}
\bibitem{ID} ... \LaTeX:\\ \url{https://www.reuters.com/article/us-china-blast-insurance-idUSKCN0QM0N220150817}
\bibitem{ID} ... \LaTeX:\\ \url{https://www.sciencedirect.com/science/article/abs/pii/S0305417916300079}
\bibitem{ID} ... \LaTeX:\\ \url{https://en.wikipedia.org/wiki/2015_Tianjin_explosions}
\bibitem{ID} ... \LaTeX:\\ \url{https://www.bbc.com/news/world-asia-china-33844084}
\bibitem{ID} ... \LaTeX:\\ \url{https://www.independent.co.uk/news/world/asia/tianjin-explosion-photos-china-chemical-factory-accident-crater-revealed-a7199591.html}

veiigheidshandhaving
\bibitem{ID} ... \LaTeX:\\ \url{https://www.ilo.org/legacy/english/protection/safework/ctrl_banding/toolkit/main_guide.pdf}
\bibitem{ID} ... \LaTeX:\\ \url{https://echa.europa.eu/documents/10162/21332507/guide_chemical_safety_sme_en.pdf}
\bibitem{ID} ... \LaTeX:\\ \url{https://ec.europa.eu/taxation_customs/dds2/SAMANCTA/EN/Safety/AppendixD_EN.htm}
\bibitem{ID} ... \LaTeX:\\ \url{https://www.ilo.org/safework/info/publications/WCMS_113134/lang--en/index.htm}


%%%%%%%%%%%%%%%%%%%%%%%%%%%%%%%%%%%%%%%%%%%%%%%%%%%%%%%%%%%%%%%%%
tesla autopilot crashes


veiigheidsrisico
\bibitem{evan01042019teslaautopilotIntersection} ... \LaTeX:\\ \url{https://spectrum.ieee.org/cars-that-think/transportation/self-driving/three-small-stickers-on-road-can-steer-tesla-autopilot-into-oncoming-lane}

\bibitem{testVehicleSafetyReport} ... \LaTeX:\\ \url{https://www.tesla.com/VehicleSafetyReport}
veiligheidsrapport mbt autopilot
\bibitem{lambert31062020q2safetyreport} ... \LaTeX:\\ \url{https://electrek.co/2020/07/31/tesla-q2-2020-safety-report-strong-improvement-autopilot-accidents/}
consumentenrapport
bluetooth veiligheidsvraagstuk
\bibitem{wiredBloutoothHackTesla} ... \LaTeX:\\ \url{https://www.wired.com/story/tesla-model-x-hack-bluetooth/}
veiigheidsvraagstuk vanwege touch screen
\bibitem{preston14012021NHTSATeslaRecall} ... \LaTeX:\\ \url{https://www.consumerreports.org/car-recalls-defects/nhtsa-asks-tesla-to-recall-model-s-model-x-touch-screen-safety-issues/}
veiligheidsvraagstuk
\bibitem{cio25112020belgianTeslaHack} ... \LaTeX:\\ \url{https://cio.economictimes.indiatimes.com/news/digital-security/security-researchers-hack-steal-tesla-model-x-within-minutes/79406553}
veiligheidsvraagstuk
rapport over autopilot
\bibitem{templeton06092019HTSBReportTesla} ... \LaTeX:\\ \url{https://www.forbes.com/sites/bradtempleton/2019/09/06/ntsb-report-on-tesla-autopilot-accident-shows-whats-inside-and-its-not-pretty-for-fsd/?sh=6905e7d4dc55}
de invloed van de bestuurder bij tesla ongeluk
\bibitem{Dickey08012021SuddenTeslaAcceleration} ... \LaTeX:\\ \url{https://techcrunch.com/2021/01/08/nhtsa-tesla-sudden-unintended-acceleration-driver-error/}
veiligheidsvraagstuk
\bibitem{darkReading17112020TeslaBackup} ... \LaTeX:\\ \url{https://www.darkreading.com/threat-intelligence/security-risks-discovered-in-tesla-backup-gateway/d/d-id/1339462}
veiligheidsvraagstuk
\bibitem{leyden23032020TeslaInterfaceHack} ... \LaTeX:\\ \url{https://portswigger.net/daily-swig/web-based-attack-crashes-tesla-driver-interface}
veiigheidsvraagstuk
\bibitem{huddlestonjr03042019ChineseTeslaHack} ... \LaTeX:\\ \url{https://www.cnbc.com/2019/04/03/chinese-hackers-tricked-teslas-autopilot-into-switching-lanes.html}
veiligheidsvraagstuk
veiligheidsvraagstuk
\bibitem{heilweil26022020teslaAutopilot} ... \LaTeX:\\ \url{https://www.vox.com/recode/2020/2/26/21154502/tesla-autopilot-fatal-crashes}
rapport over ongeluk
veiligheidsvraagstuk
veiligheidsvraagstuk
\bibitem{blanco04102019NHTSATesla} ... \LaTeX:\\ \url{https://www.caranddriver.com/news/a29369387/nhtsa-tesla-safety/}
veiligheidsvraagstuk
ransomware aanval op tesla
tesla batterij is veiligheidsvraagstuk geworden
\bibitem{mitchell01072020teslabatterycooling} ... \LaTeX:\\ \url{https://www.latimes.com/business/story/2020-07-01/federal-safety-officials-probe-tesla-battery-cooling-system}
ongeluk
\bibitem{bbc26022020AutopilotCrash} ... \LaTeX:\\ \url{https://www.bbc.com/news/technology-51645566}
veiligheidsvraagstuk
veiligheidsvraagstuk
\bibitem{stumpff04052020TeslaPersonalData} ... \LaTeX:\\ \url{https://www.thedrive.com/news/33272/tesla-discarded-old-car-parts-with-customers-personal-data-passwords-report}
dodelijk ongeluk
\bibitem{levin08062018teslaautopilotsafety} ... \LaTeX:\\ \url{https://www.theguardian.com/technology/2018/jun/07/tesla-fatal-crash-silicon-valley-autopilot-mode-report}
veiligheidsvraagstuk: ransomware
veiligheidsvraagstuk: medewerker in de fout
\bibitem{cbrook06082021TeslaInsideDataThreft} ... \LaTeX:\\ \url{https://digitalguardian.com/blog/tesla-data-theft-case-illustrates-danger-insider-threat}
\bibitem{shilling25022021Tesla} ... \LaTeX:\\ \url{https://jalopnik.com/tesla-is-stopping-some-model-3-production-report-1846353323}
veiligheidsvraagstuk: hackers je systeem laten testen
verdedigen tegenover ransomware
veiligheidsrisico
prijzen omlaag
autopilot
\bibitem{randall05112019modelSurvey} ... \LaTeX:\\ \url{https://www.bloomberg.com/graphics/2019-tesla-model-3-survey/autopilot.html}
malware door een medewerker
\bibitem{ID} ... \LaTeX:\\ \url{https://www.teslarati.com/tesla-employee-fbi-thwarts-russian-cybersecurity-attack/}
dodelijk ongeluk
\bibitem{ID} ... \LaTeX:\\ \url{https://www.marketwatch.com/story/apple-engineer-killed-in-tesla-suv-crash-on-silicon-valley-freeway-was-playing-videogame-ntsb-2020-02-25}
\bibitem{fottrell03092018TeslaSecurityChecks} ... \LaTeX:\\ \url{https://www.marketwatch.com/story/nearly-100-of-teslas-stolen-in-the-us-since-2011-have-been-recovered-2018-08-10}
waarom een tesla stelen bijna onmogelijk is
\bibitem{ID} ... \LaTeX:\\ \url{https://www.welivesecurity.com/2019/03/25/white-hats-hack-tesla-keep/}
veiligheidsonderzoek
\bibitem{ID} ... \LaTeX:\\ \url{https://www.tripwire.com/state-of-security/security-data-protection/tesla-encouraging-good-faith-security-research-in-bug-bounty-program/}
softwarefout maakt diestal mogelijk
\bibitem{kirk26112020modelX} ... \LaTeX:\\ \url{https://www.bankinfosecurity.com/tesla-model-x-stolen-in-minutes-using-software-flaws-a-15462}
fouten ontdekt in onderzoek
\bibitem{ID} ... \LaTeX:\\ \url{https://www.cnet.com/roadshow/news/tesla-ev-appeal-loyalty-study/}
\bibitem{bbc24022021hyundaiBatteryFireFix} ... \LaTeX:\\ \url{https://www.bbc.com/news/technology-56156801}








tesla cloud gehacked
\bibitem{ID} ... \LaTeX:\\ \url{https://arstechnica.com/information-technology/2018/02/tesla-cloud-resources-are-hacked-to-run-cryptocurrency-mining-malware/}
\bibitem{ID} ... \LaTeX:\\ \url{https://www.motortrend.com/news/tesla-model-y-ev-safety-quality-issues-problems/}
\bibitem{ID} ... \LaTeX:\\ \url{https://securityledger.com/2019/04/hackers-remotely-steer-tesla-model-s-using-autopilot-system/}
\bibitem{ID} ... \LaTeX:\\ \url{https://www.pcmag.com/news/report-tesla-suspends-model-3-production-in-california-until-march-7}
\bibitem{ID} ... \LaTeX:\\ \url{https://www.scmp.com/business/money/article/3121173/tesla-conduct-complete-self-inspection-after-chinese-regulators}
\bibitem{ID} ... \LaTeX:\\ \url{https://www.businesswire.com/news/home/20180220005222/en/RedLock-Releases-Cloud-Security-Report-Highlighting-Focus-on-Shared-Responsibilities-Uncovers-Cloud-Related-Exposures-at-Tesla}
\bibitem{ID} ... \LaTeX:\\ \url{https://www.epa.gov/automotive-trends/highlights-automotive-trends-report}
\bibitem{ID} ... \LaTeX:\\ \url{https://www.livemint.com/Companies/o2QLbtJc9EQ7ZcpxqgFbBP/Teslas-reward-for-finding-security-bugs-Model-3.html}
\bibitem{evans11062018} ... \LaTeX:\\ \url{https://revealnews.org/blog/tesla-fired-safety-official-for-reporting-unsafe-conditions-lawsuit-says/}
\bibitem{ID} ... \LaTeX:\\ \url{https://heimdalsecurity.com/blog/security-alert-teslacrypt-4-0-unbreakable-encryption-worse-data-leakage/}
\bibitem{ID} ... \LaTeX:\\ \url{https://www.eweek.com/cloud/tesla-cloud-account-data-breach-revealed-in-redlock-security-report/}
\bibitem{hawkins22102022} ... \LaTeX:\\ \url{https://www.theverge.com/2020/10/21/21527577/tesla-full-self-driving-autopilot-beta-software-update}
\bibitem{ID} ... \LaTeX:\\ \url{file:///C:/Users/gally/Downloads/applsci-10-02749-v2.pdf}
\bibitem{gritti24062020tesladataengine} ... \LaTeX:\\ \url{https://www.braincreators.com/brainpower/insights/teslas-data-engine-and-what-we-should-all-learn-from-it}
\bibitem{ID} ... \LaTeX:\\ \url{https://bernardmarr.com/default.asp?contentID=1251}
\bibitem{ID} ... \LaTeX:\\ \url{https://arstechnica.com/cars/2019/10/how-teslas-latest-acquisition-could-accelerate-autopilot-development/}
\bibitem{bouchard07052019teslaDeepLearning} ... \LaTeX:\\ \url{https://towardsdatascience.com/teslas-deep-learning-at-scale-7eed85b235d3}
\bibitem{ID} ... \LaTeX:\\ \url{file:///C:/Users/gally/Downloads/applsci-10-02749-v2.pdf}
\bibitem{Srikanth2019teslabigdata} ... \LaTeX:\\ \url{https://www.techiexpert.com/how-tesla-is-using-artificial-intelligence-and-big-data/}
\bibitem{rangaiah25022020teslaAI} ... \LaTeX:\\ \url{https://www.analyticssteps.com/blogs/how-tesla-making-use-artificial-intelligence-its-operations}
\bibitem{marr08012018taslabigdataAI} ... \LaTeX:\\ \url{https://www.forbes.com/sites/bernardmarr/2018/01/08/the-amazing-ways-tesla-is-using-artificial-intelligence-and-big-data/?sh=5e396aa24270}
\bibitem{bdickson29072020teslalevelfive = {author},	ALTeditor = {editor},	title = {title},	date = {date},	url = {"https://bdtechtalks.com/2020/07/29/self-driving-tesla-car-deep-learning/}
\bibitem{dcruz17062022tesladesignthink} ... \LaTeX:\\ \url{https://www.mygreatlearning.com/blog/teslas-new-ai-for-self-driving-cars/}
\bibitem{ID} ... \LaTeX:\\ \url{https://www.techiexpert.com/how-tesla-is-using-artificial-intelligence-and-big-data/}





\bibitem{mcfarland22042021selfdrivingrisks} ... \LaTeX:\\ \url{https://www.cnn.com/2021/04/21/tech/tesla-full-self-driving-launch/index.html}
\bibitem{hawkins18032021fedgovinvest} ... \LaTeX:\\ \url{https://www.theverge.com/2021/3/18/22338427/tesla-autopilot-crash-michigan-nhtsa-investigation}
\bibitem{berry21042021teslacrashtexas} ... \LaTeX:\\ \url{https://www.wionews.com/technology/doctor-among-victims-of-lethal-tesla-car-crash-in-texas-378950}
\bibitem{hull23072021regulatorsaftercrash} ... \LaTeX:\\ \url{https://www.bloomberg.com/news/newsletters/2021-06-23/hyperdrive-daily-after-30-tesla-crashes-what-s-a-regulator-to-do}
\bibitem{wikiTeslaAutopilot} ... \LaTeX:\\ \url{https://en.wikipedia.org/wiki/Tesla_Autopilot}
\bibitem{nhtsaAutomatedVehiclesSafety} ... \LaTeX:\\ \url{https://www.nhtsa.gov/technology-innovation/automated-vehicles-safety}
\bibitem{dowling23042021autopilottricking} ... \LaTeX:\\ \url{https://www.caradvice.com.au/947080/elon-musk-responds-to-deadly-texas-tesla-crash-as-consumer-reports-reveals-how-autopilot-can-be-tricked/}
\bibitem{wilson19042021teslacrashregulators} ... \LaTeX:\\ \url{https://usa.streetsblog.org/2021/04/19/regulators-could-have-prevented-fatal-tesla-crash/}

\bibitem{seamans22062021aikillerap} ... \LaTeX:\\ \url{https://www.brookings.edu/research/autonomous-vehicles-as-a-killer-app-for-ai/}
\bibitem{mitchell24022020AIDataTesla} ... \LaTeX:\\ \url{https://www.latimes.com/business/story/2020-02-24/autopilot-data-secrecy}
\bibitem{denneyjdsupraFeds} ... \LaTeX:\\ \url{https://www.jdsupra.com/post/contentViewerEmbed.aspx?fid=9844cae0-aa5a-45a5-988f-7f02fa5709c1}
\bibitem{siddiqui22102020TeslaCriticism} ... \LaTeX:\\ \url{https://www.washingtonpost.com/technology/2020/10/21/tesla-self-driving/}
\bibitem{ackerman01072016TeslaImperfect} ... \LaTeX:\\ \url{https://spectrum.ieee.org/cars-that-think/transportation/self-driving/fatal-tesla-autopilot-crash-reminds-us-that-robots-arent-perfect}
\bibitem{greene04092019misuseautopilot} ... \LaTeX:\\ \url{https://thenextweb.com/news/another-tesla-owner-is-dead-because-of-autopilot}
\bibitem{michralli26112019ubserautocarcrsash} ... \LaTeX:\\ \url{https://towardsdatascience.com/another-self-driving-car-accident-another-ai-development-lesson-b2ce3dbb4444}
\bibitem{pitmann21072021wrongfullautodeath} ... \LaTeX:\\ \url{https://www.theautochannel.com/news/2021/07/21/1024631-is-it-still-wrongful-death-if-car-is-driving-itself.html}
\bibitem{stackexchange102019teslacarmistake} ... \LaTeX:\\ \url{https://ai.stackexchange.com/questions/1488/why-did-a-tesla-car-mistake-a-truck-with-a-bright-sky}
\bibitem{tasking07062017TeslaAugmentedSafety} ... \LaTeX:\\ \url{https://resources.tasking.com/p/benefits-tesla-autopilot-and-how-adas-will-save-lives}
\bibitem{griemannExaminSelfDriving} ... \LaTeX:\\ \url{https://www.jipitec.eu/issues/jipitec-9-3-2018/4806}
\bibitem{Harkey30052019SafeSystemVehicle} ... \LaTeX:\\ \url{https://static.tti.tamu.edu/conferences/traffic-safety19/presentations/lunch/harkey.pdf}
\bibitem{ID} ... \LaTeX:\\ \url{https://thepressfree.com/have-google-and-amazon-backed-the-wrong-technology/}
\bibitem{ID} ... \LaTeX:\\ \url{https://www.irishtimes.com/business/innovation/robotaxis-have-google-and-amazon-backed-the-wrong-technology-1.4626749}
\bibitem{ID} ... \LaTeX:\\ \url{https://www.afr.com/technology/how-teslas-autopilot-got-it-wrong-in-fatal-crash-20160704-gpxsje}
\bibitem{ID} ... \LaTeX:\\ \url{https://economictimes.indiatimes.com/markets/stocks/news/what-me-worry-fed-chiefs-emotional-tone-can-drive-markets-study-suggests/articleshow/84618073.cms}
\bibitem{ID} ... \LaTeX:\\ \url{https://www.ehstoday.com/safety/article/21919260/ntsb-fatal-crash-involving-tesla-autopilot-resulted-from-driver-errors-overreliance-on-automation}
\bibitem{ID} ... \LaTeX:\\ \url{https://www.vanityfair.com/news/2016/07/how-the-media-screwed-up-the-fatal-tesla-accident}


tesla crash report
\bibitem{shepardson18062021TeslaDeaths} ... \LaTeX:\\ \url{https://www.reuters.com/business/autos-transportation/us-safety-agency-says-it-has-opened-probes-into-10-tesla-crash-deaths-since-2016-2021-06-17/}
\bibitem{hawkins30062021nhtsarequiresreporting} ... \LaTeX:\\ \url{https://www.politico.com/news/2021/05/18/ntsb-tesla-owner-was-in-drivers-seat-before-april-texas-crash-489272}
\bibitem{hawkins10052021autopilotnotavailable} ... \LaTeX:\\ \url{https://www.theverge.com/2021/5/10/22429198/tesla-ntsb-texas-crash-driverless-preliminary-report}
\bibitem{szymkowski29062021nhtsaTeslaCrashReports} ... \LaTeX:\\ \url{https://www.cnet.com/roadshow/news/tesla-autopilot-nhtsa-crash-report-self-driving-car-driver-assist-system/}
\bibitem{abc1112052021AutopilotNotinTeslaCrash} ... \LaTeX:\\ \url{https://abc11.com/tesla-crash-battery-fire-national-transportation-safety-board-driverless/10619772/}
\bibitem{ankel18062021regulatorsinvestigateAutopilot} ... \LaTeX:\\ \url{https://www.businessinsider.com/tesla-autopilot-crashes-regulators-open-probes-into-30-report-2021-6?international=true&r=US&IR=T}
\bibitem{sommerfield12072021NHTSAmandateresult} ... \LaTeX:\\ \url{https://driving.ca/column/lorraine/lorraine-explains-what-the-nhtsas-self-driving-car-crash-reporting-mandate-will-find-out}
\bibitem{ID} ... \LaTeX:\\ \url{https://www.teslarati.com/tesla-model-s-crash-texas-ntsb-preliminary-report/}
\bibitem{ID} ... \LaTeX:\\ \url{https://insideevs.com/news/506498/ntsb-report-tesla-texas-crash/}
\bibitem{ID} ... \LaTeX:\\ \url{https://electrek.co/2021/06/03/tesla-tsla-crashes-report-new-orders-in-china-free-falling/}
\bibitem{ID} ... \LaTeX:\\ \url{https://www.newsy.com/stories/ntsb-releases-report-on-fatal-tesla-crash/}
\bibitem{ID} ... \LaTeX:\\ \url{https://www.ndtv.com/world-news/autopilot-not-used-in-april-tesla-crash-says-us-report-2439146}
\bibitem{ID} ... \LaTeX:\\ \url{https://www.autocar.co.nz/autocar-news-app/fatal-driverless-tesla-crash-report-shows-autopilot-not-to-blame}
\bibitem{ID} ... \LaTeX:\\ \url{https://teleperformance-waha.sabacloud.com/Saba/Web_spf/EU2PRD0152/app/dashboard}
\bibitem{ID} ... \LaTeX:\\ \url{https://www.independent.co.uk/news/world/americas/tesla-texas-crash-model-s-autopilot-b1845286.html}
\bibitem{ID} ... \LaTeX:\\ \url{https://www.wired.com/2017/01/probing-teslas-deadly-crash-feds-say-yay-self-driving/}
\bibitem{saferoardsCrashesAutonomousvehicles} ... \LaTeX:\\ \url{https://saferoads.org/wp-content/uploads/2020/03/AV-Crash-List-with-Photos-February-2020.pdf}
\bibitem{ID} ... \LaTeX:\\ \url{https://mashable.com/article/nthsa-tesla-autopilot-model-x-crash-investigation}
\bibitem{stephardson18032021revieuwingtesla} ... \LaTeX:\\ \url{https://www.usnews.com/news/top-news/articles/2021-03-18/us-safety-agency-reviewing-23-tesla-crashes-three-from-recent-weeks}
\bibitem{krishner30062021NHTSAreport} ... \LaTeX:\\ \url{https://chicago.suntimes.com/consumer-affairs/2021/6/30/22557122/nhtsa-automated-driving-crash-reports-tesla-national-highway-traffic-safety-administration}
\bibitem{gitlin11052021autopilot} ... \LaTeX:\\ \url{https://arstechnica.com/cars/2021/05/ntsb-finds-no-reason-to-suspect-autopilot-in-fatal-tesla-crash/}
\bibitem{ID} ... \LaTeX:\\ \url{https://jalopnik.com/the-ntsb-to-partially-blame-teslas-autopilot-in-fatal-c-1803136365}
\bibitem{mitchell19012017investigationstop} ... \LaTeX:\\ \url{https://www.latimes.com/business/autos/la-fi-hy-tesla-autopilot-20170119-story.html}
\bibitem{gordon10052021teslaprelimreport} ... \LaTeX:\\ \url{https://www.vice.com/en/article/z3xxaw/ntsb-releases-preliminary-report-on-tesla-crash-that-killed-two-people}
\bibitem{shaper07062018} ... \LaTeX:\\ \url{https://choice.npr.org/index.html?origin=https://www.npr.org/2018/06/07/618081406/no-driver-input-detected-in-seconds-before-deadly-tesla-crash-ntsb-finds}
\bibitem{cochran18042021nodriverTeslaCrash} ... \LaTeX:\\ \url{https://www.click2houston.com/news/local/2021/04/18/2-men-dead-after-fiery-tesla-crash-in-spring-officials-say/}
\bibitem{habib28062016NHTSATeslaReport} ... \LaTeX:\\ \url{https://static.nhtsa.gov/odi/inv/2016/INCLA-PE16007-7876.pdf}
\bibitem{firstpress11052021fatalnonautopilot} ... \LaTeX:\\ \url{https://www.firstpost.com/tech/news-analysis/tesla-model-s-involved-in-fatal-crash-in-the-us-did-not-use-autopilot-says-ntsb-report-9609911.html}
\bibitem{raynal20042021probeTeslaCrash} ... \LaTeX:\\ \url{https://www.autoweek.com/news/green-cars/a36173804/both-local-police-and-nhtsa-probe-tesla-crash/}
\bibitem{tiungteslasoftwarecrash} ... \LaTeX:\\ \url{https://www.zdnet.com/article/apple-and-tesla-under-fire-over-software-engineers-fatal-autopilot-crash/}
\bibitem{globaltimes08052021guangdongcrash,	ALTauthor = {global times},	ALTeditor = {editor},	title = {Tesla crash in Guangdong sparks new round of safety concerns},	date = {date},	url = {"https://www.globaltimes.cn/page/202105/1222902.shtml}
\bibitem{anderson30042021secondteslacrash,	ALTauthor = {Brad Anderson},	ALTeditor = {editor},	title = {NTSB To Release Report On Deadly Tesla Crash Within A Month},	date = {date},	url = {"https://www.carscoops.com/2021/04/ntsb-to-release-report-on-deadly-tesla-crash-within-a-month/}
	\bibitem{oremus21062017fatalTeslaCrash} ... \LaTeX:\\ \url{https://slate.com/technology/2017/06/a-new-report-on-what-happened-in-the-fatal-tesla-autopilot-crash.html}
	\bibitem{guardian15052021teslacrashHandsOnWheel} ... \LaTeX:\\ \url{https://www.theguardian.com/us-news/2021/may/15/tesla-fatal-california-crash-autopilot}
	\bibitem{Puzzanghera13092017TeslaSharesBlame} ... \LaTeX:\\ \url{https://www.stuff.co.nz/motoring/news/96797272/tesla-shares-some-blame-in-fatal-autopilot-crash-report}
	\bibitem{jaillet02022017teslaAutopilotLimitations} ... \LaTeX:\\ \url{https://www.overdriveonline.com/business/article/14891759/dot-report-on-fatal-2016-tesla-crash-with-tractor-trailer-blames-limitations-of-autopilot-mode}
	\bibitem{reuters03102019teslaAutoParkingFail} ... \LaTeX:\\ \url{https://venturebeat.com/2019/10/03/regulators-investigate-teslas-automated-parking-feature-following-crash-reports/}
	\bibitem{dowling23042021} ... \LaTeX:\\ \url{https://www.caradvice.com.au/947080/elon-musk-responds-to-deadly-texas-tesla-crash-as-consumer-reports-reveals-how-autopilot-can-be-tricked/}
	\bibitem{young05112021fatalTeslaReport} ... \LaTeX:\\ \url{https://www.consumeraffairs.com/news/ntsb-releases-report-on-fatal-tesla-crash-in-texas-051121.html}
	\bibitem{kierstein18032021teslaAutopilotCrashStationary} ... \LaTeX:\\ \url{https://www.motortrend.com/news/tesla-michigan-state-autopilot-crash-report/}
	\bibitem{janssen20062017teslacrashdetailflorida} ... \LaTeX:\\ \url{https://tweakers.net/nieuws/126145/onderzoeksraad-vs-publiceert-technische-details-tesla-crash-in-florida.html}
	\bibitem{ID} ... \LaTeX:\\ \url{https://gizmodo.com/new-documents-reveal-one-driver-s-agony-and-confusion-d-1841720801}
	
	
	\bibitem{ID} ... \LaTeX:\\ \url{https://www.google.com/search?q=tesla+crash+report&rlz=1C1AVUC_enNL953NL953&ei=p3kNYa6sLI_UsAeSoZrwDw&start=100&sa=N&ved=2ahUKEwjum77s_ZzyAhUPKuwKHZKQBv44WhDw0wN6BAgBEEg&biw=1920&bih=933}
	
	
	
	%%%%%%%%%%%%%%%%%%%%%%%%%%%%%%%%%%%%%%%%%%%%%%%%%%%%%%%%%%%%%%%%%
	
	vlucht 1951
	\bibitem{ID} ... \LaTeX:\\ \url{https://nl.wikipedia.org/wiki/Turkish_Airlines-vlucht_1951}
	technisch rapport
	\bibitem{ID} ... \LaTeX:\\ \url{file:///C:/Users/gally/Downloads/rapport_ta_nl_aangepast.pdf}
	beschrijving
	terugblik met overlevenden
	tijdlijn
	\bibitem{zuilen23022019Tijdlijnpoldercrash} ... \LaTeX:\\ \url{https://www.noordhollandsdagblad.nl/cnt/dmf20190221_65390940}
	artikel
	terugblik met overlevenden
	advies raad voor de veiligheid
	de overlevende, de oorzaak, regeling, herdenking, smartengeld
	verhaal van een overlevende
	herdenking
	herdenking
	bemanning deed niets met foutmelding
	parlementaire besluitenlijst
	kamervragen over de onafhankelijkheid van de raad voor veiligheid
	verhaal van een overlevende
	beschrijvend artikel van letsel en gewonden
	\bibitem{ntvg20012010letsel} ... \LaTeX:\\ \url{https://www.ntvg.nl/artikelen/vliegtuigongeval-schiphol-25-02-2009-letsels-en-verdeling-van-gewonden}
	technische fout als oorzaak
	\bibitem{wikinews04032009techfoutailines1951} ... \LaTeX:\\ \url{https://nl.wikinews.org/wiki/Technische_fout_oorzaak_vliegtuigcrash_Turkish_Airlines-vlucht_1951}
	gesprek met pieter van vollenhove voorzitter van de onderzoeksraad voor veiligheid
	onderzoeksraad voor veiligheid is onderdruk gezet
	\bibitem{luchtvaartnieuws21012020boeing737conclusies} ... \LaTeX:\\ \url{https://www.luchtvaartnieuws.nl/nieuws/categorie/72/algemeen/conclusies-crash-tk1951-na-amerikaanse-druk-afgezwakt}
	niuwesartikel
	feitenverloop
	\bibitem{adformatie280220209communicatiegebreken} ... \LaTeX:\\ \url{https://www.adformatie.nl/contentmarketing/communicatie-na-vliegramp-vertoonde-gebreken}
	zwarte doos
	\bibitem{spinnael25022009onderzoekpolderbaancrash} ... \LaTeX:\\ \url{https://flightlevel.be/244/onderzoek-polderbaan-crash-turkish-airlines-1951/}
	\bibitem{crashTurkishAirlines} ... \LaTeX:\\ \url{http://wikimapia.org/11633002/nl/Crash-Turkish-Airlines-vlucht-1951}
	\bibitem{flightradar24} ... \LaTeX:\\ \url{https://www.flightradar24.com/data/flights/tk1951}
	\bibitem{flightstatstracker} ... \LaTeX:\\ \url{https://www.flightstats.com/v2/flight-tracker/TK/1951}
	
	
	
	
	
	
	
	
	
	
	%%%%%%%%%%%%%%%%%%%%%%%%%%%%%%%%%%%%%%%%%%%%%%%%%%%%%%%%%%%%%%%%%
	de mali missie
	
	
	\bibitem{bnnvara13062018malirapport} ... \LaTeX:\\ \url{https://joop.bnnvara.nl/nieuws/rapport-haalbaarheid-en-houdbaarheid-van-mali-missie-twijfelachtig}
	\bibitem{eucal11012021malimissieverlengd} ... \LaTeX:\\ \url{https://www.consilium.europa.eu/nl/press/press-releases/2021/01/11/eucap-sahel-mali-mission-extended-until-31-january-2023-and-mandate-adjusted/}
	\bibitem{nos21052014zorgenmalimissie} ... \LaTeX:\\ \url{https://nos.nl/artikel/650637-kamer-bezorgd-over-mali-missie}
	\bibitem{meijnders} ... \LaTeX:\\ \url{https://www.bnr.nl/nieuws/10015679/koenders-positief-tegenover-verlening-mali-missie}
	\bibitem{bnrwebredactie} ... \LaTeX:\\ \url{https://www.bnr.nl/nieuws/politiek/10345553/kabinet-wil-mali-missie-stoppen-verrassend-besluit}
	\bibitem{keultjes01062016malimissiecoalitie} ... \LaTeX:\\ \url{https://www.ad.nl/nieuws/clash-om-mali-missie-dreigt-binnen-coalitie~a4151d4f/}
	\bibitem{veenhof18012019} ... \LaTeX:\\ \url{https://www.nd.nl/cultuur/boeken/536861/boek-kijkje-bij-de-mali-missie}
	\bibitem{ID} ... \LaTeX:\\ \url{https://www.youtube.com/watch?v=jmZ6uSbpCvg}
	\bibitem{isitman06012016militair} ... \LaTeX:\\ \url{https://www.ewmagazine.nl/nederland/achtergrond/2016/07/twee-nederlanse-militairen-dood-bij-oefening-mali-missie-325226/}
	\bibitem{nporadio11072016filmdemissie} ... \LaTeX:\\ \url{https://www.nporadio1.nl/nieuws/cultuur-media/9e3b076e-5401-4630-bf39-f925213c5b6b/onverwachte-openhartigheid-over-missie-in-mali}
	\bibitem{parlementairmonitor15122013mortierongeluk} ... \LaTeX:\\ \url{https://www.parlementairemonitor.nl/9353000/1/j9vvij5epmj1ey0/vjfm5p0nujzw?ctx=vj2mc67lofnr}
	
	
	sollicitatie
	de bureaucratie
	aankomst
	interview van de burgerbevolking
	steun van de bevolking minuut 15:00
	de organisatie minuut 23:00
	De militaire briefing minuut 34:00
	prioriteit minuut 39:00
	briefing minuut 40:00
	de communicatie met ministerie over inlichten minuut 44:00
	\bibitem{ID} ... \LaTeX:\\ \url{https://www.2doc.nl/documentaires/series/2doc/2016/juli/de-missie.html}
	
	%%%%%%%%%%%%%%%%%%%%%%%%%%%%%%%%%%%%%%%%%%%%%%%%%%%%%%%%%%%%%%%%%
	militair overleden door schietoefening in ossendrecht
	
	
	\bibitem{nos22032016ossendrecht} ... \LaTeX:\\ \url{https://amp.nos.nl/artikel/2094524-militair-omgekomen-bij-schietoefening-ossendrecht.html}
	
	\bibitem{ovv04042016lessenongevalossendrecht} ... \LaTeX:\\ \url{https://www.onderzoeksraad.nl/nl/page/4293/lessen-uit-schietongeval-ossendrecht}
	
	\bibitem{quekelboere10052017doodossendrecht} ... \LaTeX:\\ \url{https://www.bndestem.nl/bergen-op-zoom/dood-van-militair-sander-klap-35-in-ossendrecht-was-ongeluk-militairen-vrijuit-hij-probeerde-zijn-leven-te-redden~afe4c7a0/}
	
	
	Wat is de rol van defensie?
	Wat is er gedaan om de veligheid van de medewerkers te waarborgen?
	Waarom zijn deze regels niet nageleefd?
	Wat zijn de gevolgen?
	Zijn de acties die naderhand zijn ondernomen wel redelijk naar de slachtoffers, het nationale veiligheisbeeld en de medewerkers?
	
	
	
	
	
	
	
	%%%%%%%%%%%%%%%%%%%%%%%%%%%%%%%%%%%%%%%%%%%%%%%%%%%%%%%%%%%%%%%%%
	schipholbrand
	
	Wat is er gebeurd?
	\bibitem{wikiSchipholbrand} ... \LaTeX:\\ \url{https://nl.wikipedia.org/wiki/Schipholbrand}
	artikel
	\bibitem{schipholbrand27102005video} ... \LaTeX:\\ \url{https://www.youtube.com/watch?v=1i-hfEzxFfk}
	psychologische gevolgen
	rapport
	\bibitem{onderzoeksraad2610schipholoost} ... \LaTeX:\\ \url{https://www.onderzoeksraad.nl/nl/page/392/brand-cellencomplex-schiphol-oost-nacht-van-26-op-27-oktober}
	artikel met video
	herdenking
	impact op de persoon
	herdenking
	\bibitem{schipholbrandvideoargos} ... \LaTeX:\\ \url{https://www.vpro.nl/argos/speel~POMS_VPRO_461907~schadevergoeding-voor-ex-verdachte-schipholbrand~.html}
	chronologie
	\bibitem{nunl30052023feitenoverzicht} ... \LaTeX:\\ \url{https://www.nu.nl/binnenland/3355935/feitenoverzicht-schipholbrand-en-rechtszaken.html}
	tijdlijn
	\bibitem{ID} ... \LaTeX:\\ \url{https://www.singeluitgeverijen.nl/isbn/de-schipholbrand/}
	vervolgens van ministers
	beeldanalyse en reconstructie
	\bibitem{ID} ... \LaTeX:\\ \url{https://eenvandaag.avrotros.nl/item/schipholbrand-niet-ontstaan-in-cel-11/}
	herdenking
	korte samenvatting
	rapport
	artikel
	verwijzing naar het rapport vanuit de politieke oppositie
	beeld vanuit de gevangenisbewaarder
	nationaliteit slachtoffers schipholbrand
	verblijfsvergunning voor de slachtoffers
	gen schadevergoeding voor de verdachte
	verdachte voor de rechter
	geen schadevergoeding voor verdachte
	artikel wat ging er mis bji de schipholbrand
	brand veroorzaakt door een peuk
	smaadschrift
	bewakers worden niet vervolgd
	proces schipholbrand moet over en de brandveilgheid moet worden verbeterd
	de rol van het parlement in de evaluatie
	\bibitem{parlementairemonitorschipholbrand} ... \LaTeX:\\ \url{https://www.parlementairemonitor.nl/9353000/1/j9vvij5epmj1ey0/vi3aof7awcxg}
	onderzoeksmemo
	herdenking
	\bibitem{ID} ... \LaTeX:\\ \url{https://archief.ntr.nl/nova/page/detail/uitzendingen/3847/Den%20Haag%20Vandaag_%20herdenking%20Schipholbrand.html}
	herdenking
	invloed van de ramp op samenleving
	\bibitem{videonpoNOVA13112008} ... \LaTeX:\\ \url{https://www.npostart.nl/heropen-onderzoek-schipholbrand/13-11-2008/POMS_NTR_103332}
	opmerkelijk rapport gestolen in de nasleep
	\bibitem{schipperSchipholbrand} ... \LaTeX:\\ \url{https://www.nd.nl/nieuws/nederland/600395/schipholbrand-blijft-schrijnen}
	\bibitem{ID} ... \LaTeX:\\ \url{https://www.ed.nl/economie/om-geen-schadevergoeding-voor-verdachte-schipholbrand~a6c7c51d/63042600/?referrer=https%3A%2F%2Fwww.google.com%2F}
	\bibitem{ID} ... \LaTeX:\\ \url{https://www.groene.nl/artikel/schipholbrand-vereist-debat}
	\bibitem{rizoomes01052014schipholbrand} ... \LaTeX:\\ \url{https://www.rizoomes.nl/brandweer/brand-cellencomplex-schiphol/}
	
	
	
	publicaties
	\bibitem{heuvelkroesschipholbrandcamerabeelden} ... \LaTeX:\\ \url{http://www.msnp.nl/downloads/Onderzoeksmemo%20beeldanalyse%20Schipholbrand%20prot.pdf}
	\bibitem{ID} ... \LaTeX:\\ \url{http://www.dakweb.nl/roofs/2006-10/RH10-P30-31.pdf}
	\bibitem{ID} ... \LaTeX:\\ \url{https://www.delta.tudelft.nl/article/dood-door-zuinigheid}
	\bibitem{ID} ... \LaTeX:\\ \url{https://www.onderzoeksraad.nl/nl/page/392/brand-cellencomplex-schiphol-oost-nacht-van-26-op-27-oktober}
	Wat waren de regels destijds?
	Waren de autoriteiten in staat om op tijd in te grijpen of om erger te voorkomen?
	Wat is er gedaan om de veiligheid van illegalen en gevangenissbewaarders te verbeteren
	
	
	
	%%%%%%%%%%%%%%%%%%%%%%%%%%%%%%%%%%%%%%%%%%%%%%%%%%%%%%%%%%%%%%%%%
	vuurwerkramp enschede
	\bibitem{ID} ... \LaTeX:\\ \url{https://www.youtube.com/watch?v=OMkIsj8FsHw}
	\bibitem{ID} ... \LaTeX:\\ \url{https://depot03.archiefweb.eu/archives/archiefweb/20210703085353/http://www.vuurwerkramp.enschede.nl/publicaties/00005/#.YOAlp-gzaUk}
	
	Wat waren de afspraken omtrent vuurwerkopslag?
	Waarom werden de voorschriften neit nageleefd?
	
	
	%%%%%%%%%%%%%%%%%%%%%%%%%%%%%%%%%%%%%%%%%%%%%%%%%%%%%%%%%%%%%%%%%
	
	
	explosie in beirut
	\bibitem{landryalameddine12112020beiruthelathsystem} ... \LaTeX:\\ \url{https://bmchealthservres.biomedcentral.com/articles/10.1186/s12913-020-05906-y}
	\bibitem{ID} ... \LaTeX:\\ \url{https://news.sky.com/story/beirut-blast-cctv-captures-moment-huge-explosion-devastated-hospital-12047452}
	\bibitem{ID} ... \LaTeX:\\ \url{https://www.unodc.org/unodc/en/frontpage/2020/September/unodc-assists-lebanon-in-reestablishing-container-shipments-in-the-aftermath-of-the-port-of-beirut-explosion.html}
	\bibitem{ID} ... \LaTeX:\\ \url{https://reliefweb.int/sites/reliefweb.int/files/resources/LEB201-Lebanon-Emergency-Response.pdf}
	\bibitem{yadav07082020handlingexplosivesBeirut} ... \LaTeX:\\ \url{https://www.downtoearth.org.in/news/governance/beirut-blast-lessons-time-for-india-to-strengthen-handling-of-explosives-chemicals-72707}
	\bibitem{graham21082020rootsImpactBeirutBlast} ... \LaTeX:\\ \url{https://www.justsecurity.org/72122/the-cost-of-resilience-the-roots-and-impacts-of-the-beirut-blast/}
	\bibitem{ID} ... \LaTeX:\\ \url{https://www.fire-magazine.com/the-port-of-beirut-explosion-a-timely-reminder}
	\bibitem{neusaeter07082020beirutexplosioneval} ... \LaTeX:\\ \url{https://www.ctvnews.ca/sci-tech/mapping-the-beirut-explosion-what-the-impact-would-look-like-in-canadian-cities-1.5053932}
	
	
	secyrity:
	\bibitem{ID} ... \LaTeX:\\ \url{https://permanent.fdlp.gov/gpo45474/AN_advisory.pdf}
	
	
	secyrity:
	\bibitem{ID} ... \LaTeX:\\ \url{https://permanent.fdlp.gov/gpo45474/AN_advisory.pdf}
	
	
	
	
	
	
	
	
	%%%%%%%%%%%%%%%%%%%%%%%%%%%%%%%%%%%%%%%%%%%%%%%%%%%%%%%%%%%%%%%%%
	bijlmerramp
	
	
	
	
	%%%%%%%%%%%%%%%%%%%%%%%%%%%%%%%%%%%%%%%%%%%%%%%%%%%%%%%%%%%%%%%%%
	slmramp
	Wat is er gebeurd?
	\bibitem{ID} ... \LaTeX:\\ \url{https://www.srnieuws.com/suriname/290721/slm-ramp-herdacht/}
	\bibitem{ID} ... \LaTeX:\\ \url{https://werkgroepcaraibischeletteren.nl/documentaire-waarom-nou-jij-over-de-slm-ramp-in-89/}
	\bibitem{ID} ... \LaTeX:\\ \url{https://www.vpro.nl/speel~WO_NTR_15390142~andere-tijden-17-apr-2019-3-09-min-fouten-en-misstanden-leiden-tot-de-slm-ramp~.html}
	\bibitem{ID} ... \LaTeX:\\ \url{https://www.canonvannederland.nl/nl/kalender/06/1989-06-07}
	\bibitem{ID} ... \LaTeX:\\ \url{https://vijfeeuwenmigratie.nl/archief-herdenkingen-slm-ramp}
	\bibitem{ID} ... \LaTeX:\\ \url{https://www.hulpverleningsforum.nl/index.php?topic=84702.0}
	\bibitem{ID} ... \LaTeX:\\ \url{https://www.nporadio1.nl/fragmenten/focus/f792e720-bd85-4c18-8a71-b334d9d5de7e/2019-04-17-slm-ramp-een-paar-cowboys-hebben-achter-de-stuurknuppel-gezeten}
	\bibitem{ID} ... \LaTeX:\\ \url{https://www.waterkant.net/suriname/2017/06/07/herdenking-slm-ramp-28-jaar-geleden-suriname/}
	\bibitem{espnSLMterugblik} ... \LaTeX:\\ \url{https://www.espn.nl/video/clip?id=8744942}
	\bibitem{ID} ... \LaTeX:\\ \url{http://www.themediabrothers.nl/tag/slm-ramp/}
	\bibitem{ID} ... \LaTeX:\\ \url{https://www.rijnmond.nl/nieuws/182546/30-jaar-na-de-SLM-ramp-Ik-mis-mijn-broer-nog-elke-dag}
	\bibitem{dennisRosier01052020} ... \LaTeX:\\ \url{https://www.voetbalkrant.com/nieuws/2020-05-01/het-vergeten-verhaal-van-de-slm-ramp}
	\bibitem{hassing07062020slmramp} ... \LaTeX:\\ \url{https://www.bd.nl/sport/de-slm-ramp-en-het-hartverscheurende-verhaal-van-jerry-en-winnie-haatrecht~ae4ce105/?referrer=https%3A%2F%2Fwww.google.com%2F}
	\bibitem{amsterdamArchiefSLM} ... \LaTeX:\\ \url{https://www.amsterdam.nl/stadsarchief/nieuws/slm-ramp/}
	\bibitem{rtvOost06062019nabestaande} ... \LaTeX:\\ \url{https://www.rtvoost.nl/nieuws/313496/Nabestaande-SLM-ramp-Heb-ik-wel-mijn-broer-en-moeder-begraven}
	\bibitem{breda07062021AndroSnel} ... \LaTeX:\\ \url{https://www.bredavandaag.nl/nieuws/algemeen/337919/nac-herdenkt-andro-knel-slm-ramp-precies-32-jaar-geleden}
	\bibitem{andereTijdenSLMCrash} ... \LaTeX:\\ \url{https://www.anderetijden.nl/aflevering/792/Een-aangekondigde-vliegramp}
	\bibitem{aviationSLMCrashAccidentInvestigation} ... \LaTeX:\\ \url{https://nl.wikipedia.org/wiki/SLM-ramp}
	database
	\bibitem{aviationReport} ... \LaTeX:\\ \url{https://aviation-safety.net/database/record.php?id=19890607-2}
	rapport
	\bibitem{aviationSLMCrashAccidentInvestigation} ... \LaTeX:\\ \url{https://reports.aviation-safety.net/1989/19890607-2_DC86_N1809E.pdf}
	\bibitem{mcDonnelDouglasCommissionReportSLMCrash} ... \LaTeX:\\ \url{https://aviation-safety.net/investigation/cvr/transcripts/cvr_py764.php}
	\bibitem{wikiSRFlight764} ... \LaTeX:\\ \url{https://en.wikipedia.org/wiki/Surinam_Airways_Flight_764}
	\bibitem{ID} ... \LaTeX:\\ \url{https://web.archive.org/web/20050113010822/https://www.ntsb.gov/ntsb/brief.asp?ev_id=34510&key=0}
	\bibitem{nos07062019SLMTerugblik} ... \LaTeX:\\ \url{https://nos.nl/artikel/2287986-slm-vliegramp-van-precies-30-jaar-geleden-trof-ook-nederlands-voetbal}
	\bibitem{dagvantoenSLMCrash} ... \LaTeX:\\ \url{https://www.dagvantoen.nl/vliegtuigcrash-slm-bij-zanderij-meer-dan-170-doden/}
	\bibitem{waterkantNesty07061989} ... \LaTeX:\\ \url{https://www.waterkant.net/suriname/2006/06/07/vliegramp-suriname-op-7-juni-1989-2/}
	uitgebreid engels artikel
	\bibitem{eduNandlalSRCrash} ... \LaTeX:\\ \url{http://www.edufd.nl/planecrash/}
	ntsb investigtion
	\bibitem{oldjetsSRAirways} ... \LaTeX:\\ \url{http://www.oldjets.net/slm-dc-8-crash.html}
	uitgebreid engels artikel
	\bibitem{cloudberg02012021srflight764} ... \LaTeX:\\ \url{https://admiralcloudberg.medium.com/contract-to-kill-the-crash-of-surinam-airways-flight-764-828979c7efe2}
	persbericht
	\bibitem{apnews07061989srplanecrash} ... \LaTeX:\\ \url{https://apnews.com/article/5b240d758ee4c5422381cc7cdc98566b}
	Wat is de rol van de autoriteiten?
	Welke andere betrokkeen? Enw at is hun verantwoordelijkheid
	Hadden de negatieve gevolgen voorkomen kunnen worden?
	Hoe werd er over veiligheid gedacht?
	
	
	
	
	%%%%%%%%%%%%%%%%%%%%%%%%%%%%%%%%%%%%%%%%%%%%%%%%%%%%%%%%%%%%%%%%%
	Tsjernobyl
	\bibitem{ID} ... \LaTeX:\\ \url{https://www.youtube.com/watch?v=Xw3SFOfbR84}
	\bibitem{wikiTjernobyl} ... \LaTeX:\\ \url{https://nl.wikipedia.org/wiki/Kernramp_van_Tsjernobyl}
	\bibitem{rivmTjernobyl} ... \LaTeX:\\ \url{https://www.rivm.nl/straling-en-radioactiviteit/stralingsincidenten-en-kernongevallen/tsjernobyl}
	\bibitem{andereTijdenTjernobyl} ... \LaTeX:\\ \url{https://www.anderetijden.nl/aflevering/599/Tsjernobyl-als-Nederlandse-ramp}
	wat er is gebeurd en hoe het leven verdergaat
	\bibitem{kingskey19042022tjernobyl} ... \LaTeX:\\ \url{https://www.nationalgeographic.nl/het-leven-in-tsjernobyl-gaat-door}
	pernsioenfondsen en de tjernobyl ramp
	In 2021 worden mensen nog steeds blootgesteld blijkt ut een gezamelijk onderzoek van greenpeace en oekraiense wetenschappers
	stijging van de nucliaire activiteit gemeten in tjernobyl
	Het toerisme  aspect
	De chronologie
	\bibitem{erikbork26042023reactor4} ... \LaTeX:\\ \url{https://historianet.nl/maatschappij/rampen/tsjernobyl-atoomhel-bij-reactor-4}
	\bibitem{nosTjernobyl30jaarlater} ... \LaTeX:\\ \url{https://nos.nl/artikel/2101523-de-spookstad-van-tsjernobyl-30-jaar-later}
	Dieren in de omgeving van tjernobyl
	De chronologie
	Echtreme droogte zorgd voor gevaar
	\bibitem{knmi04052021tjernobylbosbrand} ... \LaTeX:\\ \url{https://www.knmi.nl/over-het-knmi/nieuws/35-jaar-na-tsjernobyl-liggen-branden-op-de-loer}
	\bibitem{dodonovaKVIRisicoTjernobyl} ... \LaTeX:\\ \url{https://www.kivi.nl/afdelingen/risicobeheer-en-techniek/columns/kernramp-tsjernobyl-het-dilemma-van-scherbitsky}
	Joernalistiek, entertainment en de waarheid
	\bibitem{dumarey04062020verhaalTjernobylWaarheid} ... \LaTeX:\\ \url{https://www.vrt.be/vrtnws/nl/2020/04/06/in-de-ban-van-tsjernobyl-vooruitblik/}
	Een onderzoek
	
	Huidige gevolgen van de explosie van toen
	\bibitem{sparkesNewScientistTjernoby;} ... \LaTeX:\\ \url{https://www.newscientist.nl/nieuws/steeds-meer-kernreacties-in-ontoegankelijke-ruimte-in-tsjernobyl/}
	De ramp, hoe de mensen ermee omgingen en hoe er nu geleef wordt
	
	evaluatieonderzoek en amatregeen
	\bibitem{kernenergiened26041986chronologiemaatregelen} ... \LaTeX:\\ \url{https://www.kernenergieinnederland.nl/node/308}
	\bibitem{mapszoneReactor} ... \LaTeX:\\ \url{https://www.google.com/maps/d/u/0/viewer?ie=UTF8&hl=nl&t=h&msa=0&ll=51.388923%2C30.099792&spn=0.685583%2C1.645203&z=9&source=embed&mid=1MLcOcMK_WrIJYMuTf0VVuYnMqQI}
	Invloed van de mens op de omgeving
	\bibitem{ID} ... \LaTeX:\\ \url{https://www.animalstoday.nl/mens-schadelijker-natuur-tsjernobyl/}
	Heroplevende splijtingsreacties
	docu van schooltv
	Radioactiviteit bereikt nederland
	documentaire en maatregelen
	\bibitem{kernhistoriek15062021tjernobyl} ... \LaTeX:\\ \url{https://historiek.net/kernramp-van-tsjernobyl-1986/8769/}
	Het verhaal van een overledende
	Toerisme
	toerisme
	toerisme
	Dieren in de omgevong
	Toevluchtsoord voor vluchtelingen van de oorlog met russische seperatisten
	Ouderen die terugkeerden naar hun woonplaats na de gedwongen verhuizing door de autoriteiten
	De straling neemt weer toe
	Lessen geleerd van tjernobyl
	\bibitem{nucleairforumFeitenTjernobyl} ... \LaTeX:\\ \url{https://www.nucleairforum.be/thema/veiligheid-als-prioriteit/tsjernobyl-de-feiten}
	Toerisme
	Bosbrand in tjernobyl
	invloed van de ramp op belgie
	\bibitem{kernongevalTjernobylFancGov} ... \LaTeX:\\ \url{https://fanc.fgov.be/nl/noodsituaties/zware-ongevallen-het-buitenland/1986-kernongeval-tsjernobyl}
	Boek recensie
	Fotos en berekeningen
	ontmanteling en toerisme
	Belangrijke lessen en overeenkomsten
	De journalistieke waarheid van de koude oorlog
	De lessen van
	\bibitem{arendswolters062019lessenTjernobyl} ... \LaTeX:\\ \url{https://magazines.autoriteitnvs.nl/nieuwsbrief-anvs/2019/02/de-lessen-van-tsjernobyl}
	Een toristenattractie maken van tjernobyl
	De radioactieve straling toen en nu
	de 30km zone door de ogen van toeristen
	artikel
	stedentrip
	rapport
	\bibitem{damveld08052020tjernobyl} ... \LaTeX:\\ \url{https://wisenederland.nl/wp-content/uploads/2020/06/TSJERNOBYL.pdf}
	slapend monster
	docu
	krantenartikel
	hbo serie
	docuserie
	de  nieuwe sacrofaag
	hulp aan slachtoffers
	slapende reactor
	krantenartikel
	\bibitem{deVriestjernobylHolland} ... \LaTeX:\\ \url{https://onh.nl/verhaal/besmette-melk-en-radioactieve-spinazie-tsjernobyl-in-holland}
	hbo serie
	internationale gevolgen
	toerisme
	nieuwe koepel
	media communicatie
	docu
	dieren
	\bibitem{ID} ... \LaTeX:\\ \url{https://www.amboanthos.nl/boek/nacht-in-tsjernobyl/}
	koepel
	koepel
	\bibitem{ing3enieur29042015antistralingskoepel} ... \LaTeX:\\ \url{https://www.deingenieur.nl/artikel/nieuwe-antistralingskoepel-tsjernobyl-bijna-af}
	toerisme
	toeristisch reiperspectief
	toerisme
	niwe koepel
	overschakelen naar duurzaamheid
	docu
	tjernobyl wekt nu duurazme energie
	toerisme
	overeenkomsten tjernobyl en fukushima
	drank en sla uit tjernobyl
	geen efficiente opslag is mogelijk
	
	wetenschappelijke artikelen
	
	zaterdag 26 april 1986. Er vind routineonderhoud plaats bij reactor 4, De controle wordt uitegevoerd door de dagploeg. Vnwege een test wordt jhet koelsysteem uitgeschakeld. Door omstandigheden wordt de test uitgesteld en wordt de verantwoordelijkheid overgedragen aan de avondploeg.
	De operator maakt bedieningsfouten waardoot de reactor bijna stil komt te liggen. En vervolgens probeert hij de reactor weer op gang te brengen. ondanks de snelle temperatuurstijging wordt het experiment doorgezet. Dan wordt ook het veiligheidssysteem stilgelgd. Terwijl het koelwater langzaam opwarmt, sluit hij de klep waarlangs de stoom naar de generator stroomt.
	
	De temperatuur van de reactorstaven neemt daarna snel toe. Terwijl er een oncontroleerbare kettingreactie op gang komt, laat het personeel in paniek de regelstaven zakken om de warmteontwikkeling af te remmen. Het is dan echter al te laat. Door een ontwerpfout loopt het vermogen razendsnel op tot 33.000 megawatt, ruim tien keer hoger dan normaal.
	
	In een oogwenk verandert al het koelwater in stoom. De ontploffing die daarop volgt, blaast het 2000 ton zware deksel van de reactor af.}

In de ravage vat het gloeiend hete grafiet in de reactor spontaan vlam. De uitslaande brand en een tweede explosie voeren een radioactieve rookwolk tot 8 kilometer hoogte. 
In een poging het vuur in reactor 4 te doven, storten helikopters vanuit de lucht zand, lood en boorzuur in de reactorkern. Het mag echter niet baten.

Intussen is de nucleaire brandstof zo heet geworden dat die door de bodem van het reactorvat dreigt te smelten. Als dat gebeurt, kan het bluswater onder het vat in één klap verdampen en dreigt een derde explosie die een groot deel van Europa onbewoonbaar zal maken. Om dit te voorkomen moet het water hoe dan ook worden weggepompt.

Drie brandweermannen wagen zich daarvoor in de ruimte onder de reactor, blootgesteld aan 300 sievert per uur, 300.000 keer de dosis die een Nederlander jaarlijks maximaal mag oplopen. Ze slagen daarin, maar twee van hen overlijden enkele dagen later aan acute stralingsziekte.

Hoewel geigertellers de dag na de ramp onrustbarende waarden aangeven, slaat het plaatselijk bestuur geen alarm. De bevolking is het niet gewend om vragen te stellen.

De volgende dag blijkt er wel degelijk iets ernstigs aan de hand te zijn. In een lange rij bussen worden de 135.000 inwoners op 27 april uit het besmette gebied geëvacueerd, om er nooit meer terug te keren.

De ramp is dan nog steeds geen wereldnieuws. De Sovjetautoriteiten blijken er niet eens van op de hoogte te zijn – president Gorbatsjov klaagt later dat hij via Zweden aan zijn informatie moest komen.


\bibitem{verschuur14012013tjernobylreports} ... \LaTeX:\\ \url{http://essay.utwente.nl/63353/1/Verschuur,_W._-_s0123617_(verslag).pdf}
\bibitem{paperlessarchivesTjernobyl} ... \LaTeX:\\ \url{https://www.paperlessarchives.com/chernobyl_nuclear_accident_doc.html}
\bibitem{vargos082000tjernobylconcerns} ... \LaTeX:\\ \url{https://www.pnnl.gov/main/publications/external/technical_reports/pnnl-13294.pdf}
\bibitem{mauroNuclearRiskSociety} ... \LaTeX:\\ \url{http://www.geocities.ws/scannapuerci/demauroinnovation.pdf}
\bibitem{vienna06092005LookingBack} ... \LaTeX:\\ \url{https://www-pub.iaea.org/MTCD/publications/PDF/Pub1312_web.pdf}



%%%%%%%%%%%%%%%%%%%%%%%%%%%%%%%%%%%%%%%%%%%%%%%%%%%%%%%%%%%%%%%%%

MH17




%%%%%%%%%%%%%%%%%%%%%%%%%%%%%%%%%%%%%%%%%%%%%%%%%%%%%%%%%%%%%%%%%
oekraine powergrid
\bibitem{owens21032017ukrainemitigationstrategies} ... \LaTeX:\\ \url{https://na.eventscloud.com/file_uploads/aed4bc20e84d2839b83c18bcba7e2876_Owens1.pdf}
\bibitem{ID} ... \LaTeX:\\ \url{https://www.us-cert.gov/ics/alerts/IR-ALERT-H-16-056-01}

\bibitem{cerulus2019FrontlineRussiaAttack} ... \LaTeX:\\ \url{https://www.politico.eu/article/ukraine-cyber-war-frontline-russia-malware-attacks/}
\bibitem{ID} ... \LaTeX:\\ \url{https://en.wikipedia.org/wiki/December_2015_Ukraine_power_grid_cyberattack}
\bibitem{grammatikis2019AttackIEC6087505104} ... \LaTeX:\\ \url{https://www.researchgate.net/publication/333671061_Attacking_IEC-60870-5-104_SCADA_Systems}
\bibitem{ID} ... \LaTeX:\\ \url{https://ris.utwente.nl/ws/files/6028066/3-s2_0-B9780128015957000227.pdf}
\bibitem{hidajat2016ScadaSimulator} ... \LaTeX:\\ \url{https://www.diva-portal.org/smash/get/diva2:1046339/FULLTEXT01.pdf}
\bibitem{ID} ... \LaTeX:\\ \url{https://www.semanticscholar.org/paper/Cybersecurity-analysis-of-a-SCADA-system-under-and-Rocha/dfa7c12551ebe7b24da8d806e87e946051a57cb9}
\bibitem{ID} ... \LaTeX:\\ \url{https://tutcris.tut.fi/portal/files/16294332/jafary_1534.pdf}
\bibitem{ID} ... \LaTeX:\\ \url{http://blog.nettedautomation.com/2017/}
\bibitem{uscert20072021crashmalware} ... \LaTeX:\\ \url{https://www.us-cert.gov/ncas/alerts/TA17-163A}
\bibitem{zetter12062017malwareanalysis} ... \LaTeX:\\ \url{https://www.vice.com/en_us/article/zmeyg8/ukraine-power-grid-malware-crashoverride-industroyer}
\bibitem{icsRussianHackingCyberWeapon} ... \LaTeX:\\ \url{http://blog.wallix.com/ics-security-russian-hacking}
\bibitem{usgovC2M2} ... \LaTeX:\\ \url{https://www.energy.gov/ceser/activities/cybersecurity-critical-energy-infrastructure/energy-sector-cybersecurity-0}

bediening werking schutsluizen pdf
\bibitem{varendoejesamenVeiligSluisvaren} ... \LaTeX:\\ \url{https://www.varendoejesamen.nl/storage/app/media/downloads/vlot-en-veilig-door-brug-en-sluis-.pdf}
\bibitem{vlaamsewaterwegen012014} ... \LaTeX:\\ \url{http://www.scarphout.be/assets/bedieningstijden2014.pdf}
\bibitem{bardetsluizenAmsterdam} ... \LaTeX:\\ \url{https://www.theobakker.net/pdf/sluizen.pdf}
\bibitem{dvsbedieningsluizenenbruggen} ... \LaTeX:\\ \url{http://www.watersportalmanak.nl/files/File/Brugbediengstijden_watersport.pdf}
\bibitem{crowstappenplanmachinerichtlijnen} ... \LaTeX:\\ \url{https://www.crow.nl/downloads/pdf/verkeer-en-vervoer/verkeersmanagement/verkeersregelinstallaties/stappenplan-machinerichtlijnen_web.aspx}
\bibitem{rijksoverheidrwsonderzoeksrapporten} ... \LaTeX:\\ \url{https://puc.overheid.nl/rijkswaterstaat/doc/PUC_95170_31/}
\bibitem{bedieningstijdensluizenenbruggen} ... \LaTeX:\\ \url{http://wsv.wsvdegors.nl/wp-content/uploads/2017/05/Bedieningstijden_201701.pdf}
\bibitem{sluiscomplexkornwerderzand} ... \LaTeX:\\ \url{https://www.commissiemer.nl/projectdocumenten/00004717.pdf}
\bibitem{varengroningendrenthe} ... \LaTeX:\\ \url{https://tasmanroutes.nl/wp-content/uploads/docs/1900-bedieningstijden-groningen-drenthe.pdf}
\bibitem{wallemuldergetijhoogteverschillen} ... \LaTeX:\\ \url{http://www.vliz.be/docs/groterede/GR21_Zeesluis.pdf}
\bibitem{biemandssluizenLith} ... \LaTeX:\\ \url{https://www.bhic.nl/media/document/file/rien-biemans-sluis-en-stuw-bij-lith.pdf}
\bibitem{weilerburgers11062018zoutindringingschutsluizen} ... \LaTeX:\\ \url{https://www.nattekunstwerkenvandetoekomst.nl/upload/documents/tinymce/KpNK-2017-SKW-01c001-v1-Zoutindringing-door-schutsluizen-overzicht-projecten-en-aanzet-formulering-tbv-netwerkmodellen.pdf}
\bibitem{rwsrichtlijnvaarwegen2011} ... \LaTeX:\\ \url{https://www.arnhemspeil.nl/nap/dok/2011-12-00-rijkswaterstaat-richtlijnen-vaarwegen.pdf}
\bibitem{slotterwisscha02062016swimwaywaddenzee} ... \LaTeX:\\ \url{https://rijkewaddenzee.nl/wp-content/uploads/2016/08/Inventarisatie-toestand-vispasseerbaarheid-zoet-zout-overgangen-Waddenzee-2-6-2016-PRW-rapportage-Definitief.pdf}
\bibitem{visvriendelijksluisbeheer06012014} ... \LaTeX:\\ \url{http://www.nevepaling.nl/files/Image//nederlands/informatiecentrum/2014-definitieve-voorkeursvariantennotitie-visvriendelijk-sluisbeheer-afsluitdijk-en-houtribdijk//2014_definitieve_voorkeursvariantennotitie_visvriendelijk_sluisbeheer_afsluitdijk_en_houtribdijk.pdf}
\bibitem{ID} ... \LaTeX:\\ \url{https://www.ifv.nl/kennisplein/Documents/20120614-BwNL-Handboek-brandbeveiligingsinstallaties.pdf}
\bibitem{aubel112016sluiscomplexHeumen} ... \LaTeX:\\ \url{https://ienc-kennisportaal.nl/wp-content/uploads/2017/01/Objectbeschrijving-Heumen.pdf}
\bibitem{arends052005schutsluisstolwijkschevliet} ... \LaTeX:\\ \url{https://library.wur.nl/edepot/websites/stolwijkersluis/presentatie-data/data/pdf/TUDelft-bouwhistorisch-onderzoek.pdf}
\bibitem{multidomeinbediening} ... \LaTeX:\\ \url{https://www.icentrale.nl/wp-content/uploads/bsk-pdf-manager/2019/01/20170929_Project-2.02-Deliverable-Gehele-werkpakket-2.02.pdf}
\bibitem{historischesluizen} ... \LaTeX:\\ \url{https://www.stowa.nl/sites/default/files/assets/PUBLICATIES/Publicaties%202000-2010/Publicaties%202000-2004/STOWA%202004-XX%20boekenreeks%2020.pdf}
\bibitem{nmmag2017centralebediening} ... \LaTeX:\\ \url{https://www.nm-magazine.nl/pdf/NM_Magazine_2017-3.pdf}
\bibitem{knooppuntvaarwegenFryslanGroningenDrenthe} ... \LaTeX:\\ \url{https://www.varendoejesamen.nl/storage/app/media/knooppunten/knooppuntenboekje_03_Friesland_Groningen_Drenthe.pdf}
\bibitem{afsluitdijk052015} ... \LaTeX:\\ \url{https://deafsluitdijk.nl/wp-content/uploads/2014/05/Plan-project-MER-Afsluitdijk.pdf}
\bibitem{ID} ... \LaTeX:\\ \url{file:///C:/Users/gally/Downloads/vaarroutekaart_provincie_drenthe.pdf}
\bibitem{ID} ... \LaTeX:\\ \url{file:///C:/Users/gally/Downloads/6227_watermanagement_nl_dv.pdf}
\bibitem{ID} ... \LaTeX:\\ \url{file:///C:/Users/gally/Downloads/applsci-11-00092-v3.pdf}
\bibitem{ID} ... \LaTeX:\\ \url{file:///C:/Users/gally/Downloads/Bedieningstijden_sluizen_en_bruggen_2004.pdf.pdf}
\bibitem{ID} ... \LaTeX:\\ \url{file:///C:/Users/gally/Downloads/Bedieningstijden.pdf}
\bibitem{ID} ... \LaTeX:\\ \url{file:///C:/Users/gally/Downloads/bijlagerapport_c_-_analyse_geavanceerd_definitief_v1_0.pdf}
\bibitem{ID} ... \LaTeX:\\ \url{file:///C:/Users/gally/Downloads/BIT-advies+Bediening+op+Afstand,+sluizen+en+bruggen+in+Friesland.pdf}
\bibitem{ID} ... \LaTeX:\\ \url{file:///C:/Users/gally/Downloads/conceptverordeningnautischbeheerzuid-holland.pdf}
\bibitem{ID} ... \LaTeX:\\ \url{file:///C:/Users/gally/Downloads/De_Deltawerken_Cultuurhistorie_ontwerpgeschiedenis_web-A.pdf}
\bibitem{ID} ... \LaTeX:\\ \url{file:///C:/Users/gally/Downloads/duurzaamheid_bij_de_ontwikkeling_van_reevesluis.pdf}
\bibitem{ID} ... \LaTeX:\\ \url{file:///C:/Users/gally/Downloads/gebruikershandleiding-databank-vismigratie.pdf}
\bibitem{ID} ... \LaTeX:\\ \url{file:///C:/Users/gally/Downloads/onderzoek_vispasseerbaarheid_sluizen_zuid_holland_2016_definitief_16-5-20171.pdf}
\bibitem{ID} ... \LaTeX:\\ \url{file:///C:/Users/gally/Downloads/rapport-veiligheid-van-op-afstand-bediende-bruggen.pdf}
\bibitem{ID} ... \LaTeX:\\ \url{file:///C:/Users/gally/Downloads/richtlijnen-vaarwegen-2017_tcm21-127359%20(1).pdf}
\bibitem{ID} ... \LaTeX:\\ \url{file:///C:/Users/gally/Downloads/richtlijnen-vaarwegen-2017_tcm21-127359.pdf}
\bibitem{ID} ... \LaTeX:\\ \url{file:///C:/Users/gally/Downloads/TPE_2342144_20210511223747_PEU_52457169.pdf}


Wat hebben alle bovenstaande rampen/ongelukken gemeen? Veiligheid.
Bij de therac waren er diverse problemen: communicatie, doorontwikkeling, controle en toetsing
Was het makkelijk te onderzoeken? Waarom?
Bij de boeing 737 crashes was het probleem van controle en communicatie naar medewerkers
Was het makkelijk te onderzoeken? Waarom?

Uit de evaluatie van de china explosion 2015 tianjin komt naar voren dat communicatie, transparantie en veiligheid niet altijd prioriteit hadden bij de lokale autoriteiten
Was het makkelijk te onderzoeken? Waarom?

Bij de tesla autopilot crashes komen soms onvoldoende onderbouwde ontwerpkeuzes naar voren die niet goed zij  afgewogen tegenover het gedrag van de bestuurder
vlucht 1951
Was het makkelijk te onderzoeken? Waarom?

De ramp in Tsjernobyl toont aan hoe autoriteiten een ramp in de doofpot proberen te stoppen
Was het makkelijk te onderzoeken? Waarom?



Wat heb ik geleerd
Ik heb erg veel geleerd van het veilig opzetten van VPN’s. Een VPN opzettenhad ik namelijk nog nooit gedaan. Het opzetten van SSH en het aanmaken vanVM’s was al bekend. Ook had ik nog nooit met UDP sockets geprogrammeerd.Verder heb ik geleerd hoe ik in de praktijk een VM in een VLAN kan zetten enhoe VLAN’s netwerken van elkaar kunnen scheiden.Het leukste onderdeel van het project, was dat wonderbaarlijk mijn gekozenoplossing elegant werkte. UDP Servers en clients zijn gerealiseerd met minderdan enkele regels logisch scipt. Ik had aan genomen dat het werken met socketsin shell absoluut rampzalig zou uitpakken. Ik ben blij dat het opdracht zo vrijwas, zodat ik experimenteel kon zijn met mijn implementatie.




\bibitem{ID} ... \LaTeX:\\ \url{https://www.uni-saarland.de/fileadmin/user_upload/Professoren/FreyG/DS_KT_GF_INCOM_May_2012.pdf}
vanaf 2.1 tot en met 5

\bibitem{ID} ... \LaTeX:\\ \url{http://www.lasid.ufba.br/publicacoes/artigos/Integrating+UML+and+UPPAAL+for+Designing,+Specifying+and+Verifying+Component-Based+Real-Time+Systems.pdf}

hf7
Reachability: i.e. some condition an posssibly be satisfied
Safety: i.e. some condition will never occur
Liveness: i.e. some condition wille eventually become true [] eventually or leadsto
hf 8
Het systeem is deadlockvrij
De wachttijd is altijd gelijk aan de invaarttijd _2x de nivlleertijd en de invaartijd van de overkant

\bibitem{ID} ... \LaTeX:\\ \url{https://www.diva-portal.org/smash/get/diva2:495691/FULLTEXT01.pdf}

blz 6 tot en met 10
\bibitem{ID} ... \LaTeX:\\ \url{https://www.cister-labs.pt/docs/formal_verification_of_aadl_models_using_uppaal/1331/view.pdf}

hf 3 geeft een voorbeeld van een template met guard en acies
De volgende automata worden gebruikt met hun lokale variabelen

De volgende globale variabelen

Een lijst met relevante einschappen van een schutsluis:

\bibitem{ID} ... \LaTeX:\\ \url{https://iopscience.iop.org/article/10.1088/1742-6596/1821/1/012031/pdf}

hf 5
deadlock

\bibitem{ID} ... \LaTeX:\\ \url{http://www.es.mdh.se/pdf_publications/2934.pdf}

hf 3 tool support
Modelling in UML
Code generation
Domain Model
Behaviour model
State Hierarchy
Transitions
Trigger methods
Time events
Effects
Requirements
Environment model

hf 4
\bibitem{ID} ... \LaTeX:\\ \url{https://files.ifi.uzh.ch/stiller/CLOSER%202014/WEBIST/WEBIST/Internet%20Technology/Full%20Papers/WEBIST_2014_130_CR.pdf}


\bibitem{ID} ... \LaTeX:\\ \url{https://files.ifi.uzh.ch/stiller/CLOSER%202014/WEBIST/WEBIST/Internet%20Technology/Full%20Papers/WEBIST_2014_130_CR.pdf}

4.2 5 en 6
Het Sluisbeheeerder model wordt getoond in fuguur[]. Het model is een uitbreiding van een schutsluis met alle condities en effecten. De kleuren in de automation verwijizen naar de kleuren in de staat van de automata . De template begint met een initiele lokatie start. De sluisbeheerder initieert het proces door een aangekomen schip te registreren metbehulp van een sychronizate met het channel... over de edge richtng de lokatie "aanmelden." Dit symboliseert een opstartprocedure, ook wordt een functie enqueeu_aanmeldLijst() gebruikt om de juiste waarden te geven aan lokale en globale avariabelen. De lokatie aanmelden regisseert het opstellen van schepen boven of beneden van de sluiskolk. De template Schip synchronizeerd met de template Sluisbeheerder met het channel move_down[id] of move_up[id] en bereikt daarmee de volgende lokatie afhankelijk af de sluis boven of beneden is worden de schepen die in de opstellijst voorkomen, max 2, klaargemaakt voor invaren.. De templates Stoplicht en sluisdeur synchroniseren met de channels ... call_Deur en call_stoplicht.
Het Sluisbeheerder model gebruikt de variabelen clock x, wachttijd_beneden, wachttijd_boven als invariant tussen de lokaties. Om op de hoogete te zijn van de invaar-/uitvaart van de verschillende schepen worden lijsten bijgehouden: list_wachtrij_beneden, list_pos_invaren_beneden, list_schepenInSluis, list_wachtrij_boven en list_pos_invaren_boven.

Het model voltooit de volgende transitie op basis van de waarde van de boolean sluis_bove en sluis_beneden. en de lokale klok variabele x.
Vanaf de locatie invaarverbod_gecontroleerd  wordt gecontroleerd of er nog invarende schepen zijn die in de sluiskolk passen.
Op de lokatie sluiskolk gereed zijn er 1 of meer schepen in de sluis. Als er nog plek is in de sluiskolk n er is nog een schip klaar om in te varen dan wordt dit gecontroleerd, de functie enqueu() voegt het schip toe aan de queue van de sluiskolk. De functie deque() verwijdert de schip van de lijst met invarende schepen. De variabele sluis_boven of sluis_beneden is waar, bij de switch voor het sluiten van deuren en het aanroepe van het stoplicht nr gelang de positie van  de laate binnenvarende schip (boven of beneden). Hierna bereikt de automation sluiskolk_afgesloten.



De lokatie start_nivelleren kiest op basis van de variabelen sluis_boven en de variabelen sluis_beneden het nivellereingsprogramma.
Heet nivellereingsprogramma is Aof B. De keuze voor het programma wordt bepaald door de variabelen van het schip dat in de sluis zit.

De lokatie klaarmaken_voor_openen wordt bereikt als de   hoogte van de sluis  door het nivellereingsprogramma is bereikt.
De positie van de kluis is bepaald door de schepen in de sluis. Vanuit deze lokatie wordt gekeken off de stoplichten gereed moeten worden gemaakt en of de sluisdeuren open mogen.
Hierna volgt een transiie waarin de stoplichte op groen worden gezet en de sluisdeuren worden geopend voor de uitvaart van de schepen in de sluis.
Als alle schepen zijn uitgevaren die uit moeten varen, worden de stoplichten op groen gezet en de deuren gesloten.


De lokatie uitvaren_toegestaan heeft een verbinding(edge) met de lokatie sluis_afsluiten.
Er is een select statement, e:id_t gebruikt als onderdeel van het prototocol om alle uitvarende schepen uit de queue van de sluiskolk te halen, en wordt dan ook gebruikt door de synchronisatie met de channel leave om de schepen uit de sluiskolk te begeleiden. De edge hieraan gekoppeld bevat de functie deque() om de variabelen  van de sluiskolk te resetten.

Vanuit de positie van de sluis worden de schepen gesignaleerd op een invaarverbod en worden de deuren van de sluis gesloten.
De lokatie sluiskolk_afgesloten is bereikt.

Ship [guards, invariants, assignents, synchronizations, properties,aannames]
De template Schip begint bij de Init lokatie. De lokatie is verbonden met de lokatie aangekomen met een edge waarbij een synchronizatie wordt aangeroepen met de template sluisbeheerder. De clock wordt op nul gezet. De lokatie aangekomen is verbonden met de lokatie aangemeld. De edge bevat een synchronizatie waarmee de edge een synchronizatie uitvoert met de template Sluisbheheerder.
De volgende lokatie is  controleren. De edge waarmee de lokatie aangemeld in verbinding staat met de lokatie cnotroleren heeft een synchronisatie voor de template Sluisbeheerder. De lokatie controleren heeft ook een edge met de lokatie wachten. Een schip max maximaal 30 seconden wachten op de lokatie wachten voordat er een mogelijkheid is om opniew in aanmerking te komen voor een controle. Als een schip langer dan 30 tijdseenheden moet wachten de is er een mogelijkheid voor het schip te vertrekken. Hierbij eindigt het schip het invaarproces. Een schip kan dus na aanvaren maximaal 20 seconden wachten om toestemming te krijgen voor een positie invaren anders wordt deze verwezen naar een wachtrij.
Hierna volgdde lokate invarene. De lokatie invarene implieert dat een schip in een invaarproces is dat eindigt in de lokatie gestopt.
Hierop volgd de lokatie nivelleer_start. Hierop wordt een nivelleer_proces gestart. Daarbij is ee synchronisatie met de template Sluisbeheerder.
De lokatie nivelleer_stop is een lokate waarin het nivelleerproces al is gestopt. Van hieruit is er een edge met de lokatie klaar voor vertrek. De edge synchroniseert hiermee met de template Sluisbeheerder.
De lokatie klaar_voor_vertrek is verbonden met de lokatie Init. Met een guard x>=3 tijdseenheden mag een schip vertrekken.


Deur
De deur bevat de volgende lokaties: dicht, openend, open en sluitende.
Een deur sluit niet in een enkele actie. Het proces die een deur dooploopt zijn de processen openend en sluitende. De finale lokaties zijn open en dicht.

Nivelleermachine
De nivelleermachine begint bij de lokatie uit. Met een synchronisatie wordt een nivelleermachine aangezet. De automatie kiest een programma en werkt deze uit in de lokatie bezig. Als ht programma is afgerond volgt de lokatie klaar. Na elk nivelleerproces wordt de machine uitgezet

Stoplicht
Een stoplicht heeft twee lokaties: rood en groen.



%%%%%%%%%%%%%%%%%%%%%%%%%%%%%%%%%%%%%%%%%%%%%%%%%%%%%%%%%%%%%%%%%

Bijlage A performance
\bibitem{kumarUppaalDMAMACProtocol} ... \LaTeX:\\ \url{https://home.hvl.no/ansatte/aaks/articles/2015IKT617.pdf}

test specification
\bibitem{larsenRealtimeUppaalTesting} ... \LaTeX:\\ \url{https://d-nb.info/987511998/34}

sheet 24 tot 65
\bibitem{proenza102008UppaalModelChecker} ... \LaTeX:\\ \url{http://ppedreiras.av.it.pt/resources/empse0809/slides/TheUppaalModelChecker-Julian.pdf}


2.3.4.2
4.7

coffie apparaat

\bibitem{uppaalCoffeeMachine} ... \LaTeX:\\ \url{https://www.comp.nus.edu.sg/~cs5270/Notes/chapt6a.pdf}



what is a good software specification
\bibitem{fvaandrager2322010Goodmodel} ... \LaTeX:\\ \url{http://www.cs.ru.nl/~fvaan/PV/what_is_a_good_model.html#:~:text=A%20good%20model%20has%20a%20clearly%20specified%20purpose%20and%20(ideally,code%20generation%2C%20and%20test%20generation.}

\bibitem{onix01102022devopmodel} ... \LaTeX:\\ \url{https://onix-systems.com/blog/7-basic-software-development-models-which-one-to-choose}
\bibitem{sulemani04012021softwareprocesmodel} ... \LaTeX:\\ \url{https://www.educative.io/blog/software-process-model-types}

\bibitem{globalluxsoft18102017softdev} ... \LaTeX:\\ \url{https://medium.com/globalluxsoft/5-popular-software-development-models-with-their-pros-and-cons-12a486b569dc}
\bibitem{wiegers30052022SRS} ... \LaTeX:\\ \url{https://www.jamasoftware.com/blog/characteristics-of-excellent-requirements/}
\bibitem{muller06092020goodspecification} ... \LaTeX:\\ \url{https://www.gaudisite.nl/ValidationOfRequirementsSlides.pdf}
\bibitem{informit30062008reqmanagement} ... \LaTeX:\\ \url{https://www.informit.com/articles/article.aspx?p=1152528&seqNum=4}
\bibitem{altexsoft15092020writingSRS} ... \LaTeX:\\ \url{https://www.altexsoft.com/blog/software-requirements-specification/}
\bibitem{ID} ... \LaTeX:\\ \url{E:\Backup Mijn Documenten\Hogeschool vakken\TINLab advnced algorithms\tinlab_advanced_algoriths\achtergrondinfo research}
sheet 28 transitorische relaties vertalen van ctl naar ltl
\bibitem{ID} ... \LaTeX:\\ \url{file:///E:/Backup%20Mijn%20Documenten/Hogeschool%20vakken/TINLab%20advnced%20algorithms/tinlab_advanced_algoriths/achtergrondinfo%20research/buchi/lec16_Buchi+LTL.pdf}
\bibitem{ID} ... \LaTeX:\\ \url{file:///E:/Backup%20Mijn%20Documenten/Hogeschool%20vakken/TINLab%20advnced%20algorithms/tinlab_advanced_algoriths/achtergrondinfo%20research/buchi/lect4.pdf}
\bibitem{ID} ... \LaTeX:\\ \url{file:///E:/Backup%20Mijn%20Documenten/Hogeschool%20vakken/TINLab%20advnced%20algorithms/tinlab_advanced_algoriths/achtergrondinfo%20research/buchi/lecture8.pdf}
transitie relaties in LTL sheet 8
\bibitem{ID} ... \LaTeX:\\ \url{file:///E:/Backup%20Mijn%20Documenten/Hogeschool%20vakken/TINLab%20advnced%20algorithms/tinlab_advanced_algoriths/achtergrondinfo%20research/CS%20267%20Automated%20Verification/l2.pdf}
\bibitem{ID} ... \LaTeX:\\ \url{file:///E:/Backup%20Mijn%20Documenten/Hogeschool%20vakken/TINLab%20advnced%20algorithms/tinlab_advanced_algoriths/achtergrondinfo%20research/FORMAL%20METHODS/slide3.pdf}
\bibitem{ID} ... \LaTeX:\\ \url{file:///E:/Backup%20Mijn%20Documenten/Hogeschool%20vakken/TINLab%20advnced%20algorithms/tinlab_advanced_algoriths/achtergrondinfo%20research/FORMAL%20METHODS/slide4.pdf}
hf 4.2
\bibitem{ID} ... \LaTeX:\\ \url{file:///E:/Backup%20Mijn%20Documenten/Hogeschool%20vakken/TINLab%20advnced%20algorithms/tinlab_advanced_algoriths/achtergrondinfo%20research/properties%20ctl/Chapter-4-Formal-Methods-LTL-CTL-TRAFFIC-LIGHT-EXAMPLE-pages-18-24.pdf}


\bibitem{ID} ... \LaTeX:\\ \url{E:\Backup Mijn Documenten\Hogeschool vakken\TINLab advnced algorithms\tinlab_advanced_algoriths\lesmateriaal\modelchecking.pdf}





\bibitem{ID} ... \LaTeX:\\ \url{E:\Backup Mijn Documenten\Hogeschool vakken\TINLab advnced algorithms\tinlab_advanced_algoriths\lesmateriaal}

parallelle compositie
\bibitem{ID} ... \LaTeX:\\ \url{file:///E:/Backup%20Mijn%20Documenten/Hogeschool%20vakken/TINLab%20advnced%20algorithms/tinlab_advanced_algoriths/achtergrondinfo%20research/properties%20ctl/model.pdf}

Urgent locations
Is hetzelfde als het toevoegen van een clock x, met een invariant x<=o op de locatie. Zolang een systeem in een urgente locatie zit mag er geen tijd verstrijken
Bjivoorbeeld als een sluis klaar is engeen schpeen in de sluis. Dan moet er een urgentie zijn dat alle schepen waar mogelijk worden opgesteld voor invaren. Als er geen schepen in de wachtrij en er staan geenschepen klaar om in te varen dn is er misschien urgentie om aan de andere kant schepen op te halen.
Commited locations
Als een of meerdere locaties ingesteld zijn als committed. Een committed state kan niet vertragen  en de volgende transitie moet een transitie zijn waarin de uitgaande edge komt van een committed edge


zeno gedrag: de mogelijkheid dat in een eindige hoeveelheid tijd een oneindig antal handelingen kan worden verricht.
Bijvoorbeeld tijdens het nivelleren
Bij het opstellen van schepen
Bij het laten wachten van schepen
Bij het invaren van schepen
\bibitem{ID} ... \LaTeX:\\ \url{file:///E:/Backup%20Mijn%20Documenten/Hogeschool%20vakken/TINLab%20advnced%20algorithms/tinlab_advanced_algoriths/achtergrondinfo%20research/properties%20ctl/std.pdf}



\bibitem{ID} ... \LaTeX:\\ \url{https://wayback.archive-it.org/9650/20200409062940/http:/p3-raw.greenpeace.org/international/Global/international/publications/nuclear/2016/Nuclear_Scars.pdf}
\bibitem{ID} ... \LaTeX:\\ \url{https://bdtechtalks.com/2020/07/29/self-driving-tesla-car-deep-learning/}



%%%%%%%%%%%%%%%%%%%%%%%%%%%%%%%%%%%%%%%%%%%%%%%%%%%%%%%%%%%%%%%%%
critical safety systems chemicals
\bibitem{esc16022021scsdeveloping} ... \LaTeX:\\ \url{https://esc.uk.net/safety-critical-systems}
\bibitem{oecd2008chemsafeperfindct} ... \LaTeX:\\ \url{https://www.oecd.org/chemicalsafety/chemical-accidents/41269710.pdf}
\bibitem{issa2003chemicalsID} ... \LaTeX:\\ \url{https://safety-work.org/fileadmin/safety-work/articles/Verwechslung_von_Chemikalien/Stoffverwechslung_e.pdf}
\bibitem{sommerville2008CriticalSystems} ... \LaTeX:\\ \url{https://ifs.host.cs.st-andrews.ac.uk/Books/SE9/Web/Dependability/CritSys.html}
\bibitem{identifyinglaboratoryHazards} ... \LaTeX:\\ \url{https://www.acs.org/content/dam/acsorg/about/governance/committees/chemicalsafety/publications/identifying-and-evaluating-hazards-in-research-laboratories.pdf}
\bibitem{ID} ... \LaTeX:\\ \url{https://www.computer.org/csdl/magazine/so/2017/04/mso2017040049/13rRUxCitHw}
\bibitem{ID} ... \LaTeX:\\ \url{https://msquair.files.wordpress.com/2012/06/assca-guiding-philosophic-principles-on-the-design-and-acquisition-of-safety-critical-systems-v1-6.pdf}
\bibitem{ID} ... \LaTeX:\\ \url{https://epsc.be/Documents/PS+Fundamentals/_/EPSC_Process%20Safety%20Fundamentals%20-%20Booklet_March2021.pdf}
\bibitem{winceckCriticalToSafety} ... \LaTeX:\\ \url{https://www.icheme.org/media/8976/xxiv-poster-11.pdf}
\bibitem{chambersHazardAnalysisSCS} ... \LaTeX:\\ \url{https://crpit.scem.westernsydney.edu.au/confpapers/CRPITV55Chambers.pdf}
\bibitem{rslater1998SCSAnalysis} ... \LaTeX:\\ \url{https://users.ece.cmu.edu/~koopman/des_s99/safety_critical/}




%%%%%%%%%%%%%%%%%%%%%%%%%%%%%%%%%%%%%%%%%%%%%%%%%%%%%%%%%%%%%%%%%
critical safety systems airplanes
\bibitem{ID} ... \LaTeX:\\ \url{file:///C:/Users/gally/Downloads/AGARDAG300.pdf}
\bibitem{brat2015verifysafetyflightcritical} ... \LaTeX:\\ \url{https://arxiv.org/abs/1502.02605}
\bibitem{knightchallengessafetyCritical} ... \LaTeX:\\ \url{https://users.encs.concordia.ca/~ymzhang/courses/reliability/ICSE02Knight.pdf}
\bibitem{ID} ... \LaTeX:\\ \url{https://www.jstor.org/stable/44682826}
\bibitem{johnson2006devsafetycritical} ... \LaTeX:\\ \url{http://www.dcs.gla.ac.uk/~johnson/teaching/safety/slides/pt2.pdf}
\bibitem{yeagerSafetyCritical} ... \LaTeX:\\ \url{https://sites.google.com/site/cis115textbook/safety-critical-systems}
\bibitem{ID} ... \LaTeX:\\ \url{https://www.dau.edu/tools/se-brainbook/Pages/Design%20Considerations/Critical-Safety-Item.aspx}
\bibitem{ID} ... \LaTeX:\\ \url{https://mcdpinc.com/safety-critical-systems}
\bibitem{fallsafedesign} ... \LaTeX:\\ \url{https://faculty.up.edu/lulay/MEStudentPage/failsafe.pdf}
\bibitem{2008manualflightsafetyParts} ... \LaTeX:\\ \url{https://www.enidine.com/CorporateSite/media/itt/Resources/Distributors/EndUserDocuments/Suppliers_Documents/QAM03_Rev_E.pdf}
\bibitem{arForce2015VerificationExpectations} ... \LaTeX:\\ \url{https://daytonaero.com/wp-content/uploads/AC-17-01.pdf}
\bibitem{harvardRiskResearchaircraft} ... \LaTeX:\\ \url{https://rmas.fad.harvard.edu/pages/chartered-private-aircraft-0}
\bibitem{ID} ... \LaTeX:\\ \url{https://pubmed.ncbi.nlm.nih.gov/7966484/}
\bibitem{ID} ... \LaTeX:\\ \url{https://nebula.esa.int/content/assessment-methodology-certification-safety-gnc-critical-space-systems}
\bibitem{ID} ... \LaTeX:\\ \url{https://www.aopa.org/training-and-safety/online-learning/safety-spotlights/aircraft-systems}
\bibitem{lalaArchitecturalPrinciples} ... \LaTeX:\\ \url{https://www.cs.unc.edu/~anderson/teach/comp790/papers/safety_critical_arch.pdf}
\bibitem{andersenromanski2011verificationsafetycritical} ... \LaTeX:\\ \url{https://queue.acm.org/detail.cfm?id=2024356}
\bibitem{aviatioLawGeneralDefinitions} ... \LaTeX:\\ \url{https://www.law.cornell.edu/cfr/text/14/1.1}
\bibitem{ntsb2006safetyreportTransportAirplanes} ... \LaTeX:\\ \url{http://libraryonline.erau.edu/online-full-text/ntsb/safety-reports/SR06-02.pdf}
\bibitem{mitNotesSafetyCritical} ... \LaTeX:\\ \url{https://www.cs.uct.ac.za/mit_notes/human_computer_interaction/htmls/ch02s10.html}
\bibitem{ID} ... \LaTeX:\\ \url{https://flightsafety.org/}
\bibitem{kochenfender2020aisafetycritical} ... \LaTeX:\\ \url{https://engineering.stanford.edu/magazine/article/mykel-kochenderfer-ai-and-safety-critical-systems}
\bibitem{humanfactorsAviationSafety} ... \LaTeX:\\ \url{https://www.faasafety.gov/files/gslac/courses/content/258/1097/AMT_Handbook_Addendum_Human_Factors.pdf}
\bibitem{eulegislator2000safetyregulation} ... \LaTeX:\\ \url{https://www.eurocontrol.int/sites/default/files/2019-06/src-doc-1-e1.0.pdf}
\bibitem{ID} ... \LaTeX:\\ \url{http://aerossurance.com/safety-management/critical-maintenance-tasks/}
\bibitem{gao2020aviationcybersecurity} ... \LaTeX:\\ \url{https://www.gao.gov/assets/gao-21-86.pdf}
\bibitem{roleofcodeinsoftware} ... \LaTeX:\\ \url{https://criticalsoftware.com/en/news/coding-the-skies}
\bibitem{airlie2018designparameters} ... \LaTeX:\\ \url{https://aviation.stackexchange.com/questions/46677/what-are-the-design-parameters-for-airliner-safety}
\bibitem{ID} ... \LaTeX:\\ \url{https://www.cantwell.senate.gov/news/press-releases/cantwells-comprehensive-bipartisan-bicameral-aircraft-safety-and-certification-reforms-signed-into-law}
\bibitem{ID} ... \LaTeX:\\ \url{https://www.forbes.com/advisor/travel-rewards/737-max-what-is-safety-anyway/}
\bibitem{humanfactors2010Aviation} ... \LaTeX:\\ \url{https://www.tandfonline.com/doi/full/10.1080/00140130903521587}
\bibitem{ID} ... \LaTeX:\\ \url{https://www.doi.gov/aviation/safety}
\bibitem{courtin2018safetyconsiderations} ... \LaTeX:\\ \url{https://dspace.mit.edu/bitstream/handle/1721.1/118438/ICAT_2018_07_Christoper%20Courtin_Report.pdf?sequence=1&isAllowed=y}
\bibitem{ID} ... \LaTeX:\\ \url{https://www.defence.gov.au/dasp/Docs/Manuals/7001053/eTAMMweb/1049.htm}
\bibitem{prentice2014Failedaviationprogram} ... \LaTeX:\\ \url{https://www.aviationpros.com/aircraft/commercial-airline/article/10239806/staying-legal-another-failed-faa-safety-program}
\bibitem{civilAviationConsulting} ... \LaTeX:\\ \url{https://www.iata.org/en/services/consulting/safety-operations/}
\bibitem{olivercalvardpotocnik2017Aviationautomation} ... \LaTeX:\\ \url{https://hbr.org/2017/09/the-tragic-crash-of-flight-af447-shows-the-unlikely-but-catastrophic-consequences-of-automation}
\bibitem{ID} ... \LaTeX:\\ \url{https://www.infosys.com/industries/communication-services/documents/landing-gear-design-and-development.pdf}
\bibitem{airforce2000SystemSafety} ... \LaTeX:\\ \url{https://www.acqnotes.com/Attachments/AF_System-Safety-HNDBK.pdf}
\bibitem{fed2019SafeSecureSUAS} ... \LaTeX:\\ \url{https://www.transportation.gov/testimony/state-airline-safety-federal-oversight-commercial-aviation}
\bibitem{AirportSafety} ... \LaTeX:\\ \url{https://www.federalregister.gov/documents/2019/02/13/2019-00758/safe-and-secure-operations-of-small-unmanned-aircraft-systems}
\bibitem{ID} ... \LaTeX:\\ \url{https://archive.etsc.eu/documents/safety%20in%20airports.pdf}
\bibitem{passengerSafety} ... \LaTeX:\\ \url{https://journals.sagepub.com/doi/pdf/10.1177/002029400403700202}
\bibitem{ID} ... \LaTeX:\\ \url{https://www.unmannedsystems.ca/wp-content/uploads/2019/01/DRAFT-AC-922-001-RPAS-SAFETY-ASSURANCE.pdf}
\bibitem{ID} ... \LaTeX:\\ \url{https://www.ccsdualsnap.com/pressure-switches-in-aerospace-applications/}
\bibitem{britishColumbia2020GuideSafetyCritical} ... \LaTeX:\\ \url{https://www.egbc.ca/getmedia/78073fda-5a83-4f0f-b12f-0a40dcbbc29d/EGBC-Safety-Critical-Software-V1-0.pdf.aspx}
\bibitem{uscongres2019aircraftresearch} ... \LaTeX:\\ \url{https://readwrite.com/2018/12/21/air-travel-is-far-safer-than-you-think-heres-why/}
\bibitem{ID} ... \LaTeX:\\ \url{https://fas.org/sgp/crs/misc/R45939.pdf}
\bibitem{ariAssociation2018} ... \LaTeX:\\ \url{https://cdn.ymaws.com/www.astna.org/resource/collection/4392B20B-D0DB-4E76-959C-6989214920E9/ASTNA_Safety_Position_Paper_2018_FINAL.pdf}
\bibitem{ID} ... \LaTeX:\\ \url{https://transportation.house.gov/imo/media/doc/2020.09.15%20FINAL%20737%20MAX%20Report%20for%20Public%20Release.pdf}
\bibitem{ID} ... \LaTeX:\\ \url{https://www.transportstyrelsen.se/globalassets/global/luftfart/seminarier_och_information/seminarier-2016/luftvardighet-camo-och-145-verkstader/11b-critical-task-fpl.pdf}
\bibitem{fox2005HelicopterSafety} ... \LaTeX:\\ \url{https://www.h-a-c.ca/IHSS_Helicopter_Safety_History_05.pdf}
\bibitem{fulvio1993safetycriticalsystems} ... \LaTeX:\\ \url{https://assembly.coe.int/nw/xml/XRef/X2H-Xref-ViewHTML.asp?FileID=7144&lang=EN}
\bibitem{rakas2018criticalsystemsFailures} ... \LaTeX:\\ \url{https://www.skybrary.aero/index.php/Cockpit_Automation_-_Advantages_and_Safety_Challenges}
\bibitem{ID} ... \LaTeX:\\ \url{https://ntrs.nasa.gov/citations/20120014507}
\bibitem{cockpitAutomation} ... \LaTeX:\\ \url{https://www.sciencedirect.com/science/article/abs/pii/S092575351730601X}
\bibitem{ID} ... \LaTeX:\\ \url{https://www.semanticscholar.org/paper/Safety-critical-avionics-for-the-777-primary-flight-Yeh/8facf90f4a9051c3ab8ce11e39d0893118268d90}
\bibitem{ID} ... \LaTeX:\\ \url{https://www.easa.europa.eu/faq/19013}
\bibitem{belcaatrovalidateSafetyCritical} ... \LaTeX:\\ \url{https://ntrs.nasa.gov/api/citations/20120014507/downloads/20120014507.pdf}
\bibitem{ID,	civilAviationsAuthority = {author},	ALTeditor = {editor},	title = {title},	date = {date},	url = {"https://publicapps.caa.co.uk/modalapplication.aspx?catid=1&pagetype=65&appid=11&mode=list&type=subcat&id=32}
	\bibitem{ID} ... \LaTeX:\\ \url{https://definitions.uslegal.com/f/flight-safety-critical-aircraft-part/}
	\bibitem{ID} ... \LaTeX:\\ \url{https://nbaa.org/nbaa-aviation-groups-ask-congress-to-prevent-5g-interference-to-critical-safety-systems/}
	\bibitem{ID} ... \LaTeX:\\ \url{https://www.dlr.de/ft/en/desktopdefault.aspx/tabid-1360/1856_read-36215/}
	\bibitem{ID} ... \LaTeX:\\ \url{https://www.fsd.lrg.tum.de/research/safety-critical/}
	\bibitem{knight2010SafetyCritical} ... \LaTeX:\\ \url{https://ieeexplore.ieee.org/document/1007998}
	\bibitem{AAA052005IdentifySCHitems} ... \LaTeX:\\ \url{https://www.faa.gov/about/office_org/headquarters_offices/ast/licenses_permits/media/RLVGuide_01-05_05.pdf}
	\bibitem{devTopics01032020} ... \LaTeX:\\ \url{https://smallbusinessprogramming.com/safety-critical-software-15-things-every-developer-should-know/}
	\bibitem{ID} ... \LaTeX:\\ \url{https://coreavi.com/the-future-of-safety-critical-systems-in-the-emerging-autonomous-world/}
	\bibitem{valdes2018SafetybyAutomation} ... \LaTeX:\\ \url{https://www.intechopen.com/chapters/59838}
	\bibitem{2015whensafetymanagementsystemsfail} ... \LaTeX:\\ \url{https://verticalmag.com/features/whensafetymanagementsystemsfail/}
	
	
	%%%%%%%%%%%%%%%%%%%%%%%%%%%%%%%%%%%%%%%%%%%%%%%%%%%%%%%%%%%%%%%%%
	critical safety systems fireworks
	\bibitem{ID} ... \LaTeX:\\ \url{https://www.hsdl.org/c/firework-safety/}
	\bibitem{ID} ... \LaTeX:\\ \url{https://www.cpsc.gov/Safety-Education/Safety-Education-Centers/Fireworks}
	\bibitem{ID} ... \LaTeX:\\ \url{https://www.seattletimes.com/subscribe/signup-offers/?pw=redirect&subsource=paywall&return=https://www.seattletimes.com/opinion/editorials/firework-safety-even-more-critical-after-heat-wave/}
	\bibitem{ID} ... \LaTeX:\\ \url{https://www.nrcan.gc.ca/sites/www.nrcan.gc.ca/files/mineralsmetals/pdf/mms-smm/expl-expl/20170828-G05-09E_ACC.pdf}
	\bibitem{fireworksinjuries2021duringcovid} ... \LaTeX:\\ \url{https://www.prnewswire.com/news-releases/fireworks-related-injuries-and-deaths-spiked-during-the-covid-19-pandemic-301322243.html}
	\bibitem{osha2017safetymanagementexplosives} ... \LaTeX:\\ \url{https://www.osha.gov/sites/default/files/publications/OSHA3912.pdf}
	\bibitem{ID} ... \LaTeX:\\ \url{https://www.firelinx.com/wp-content/uploads/2021/02/FLX-Issues-in-Firing-System-Safety.pdf}
	\bibitem{ID} ... \LaTeX:\\ \url{http://www.eig2.org.uk/wp-content/uploads/WTOFD-Blue-Guide.pdf}
	\bibitem{hse2014explosiveregulations} ... \LaTeX:\\ \url{https://www.hse.gov.uk/explosives/er2014-fireworks-retail-prem.pdf}
	\bibitem{ID} ... \LaTeX:\\ \url{https://www.firerescue1.com/firefighter-safety/articles/11-fireworks-safety-videos-from-the-serious-to-the-humorous-fHy0M4pT2gjcQ8jA/}
	\bibitem{pirone2016lessonslearnedfireworks} ... \LaTeX:\\ \url{https://www.aidic.it/cet/16/53/044.pdf}
	\bibitem{ID} ... \LaTeX:\\ \url{http://www.alarmascasas.com.mx/sites/default/files/85006-0061%20--%20FireWorks%20Brochure.pdf}
	\bibitem{ID} ... \LaTeX:\\ \url{https://www.ehs.ufl.edu/programs/fire/fireworks/}
	\bibitem{ID} ... \LaTeX:\\ \url{https://www.interlogix.com.au/documents/FireWorks%20Features%20and%20Operation%20(fire%20only).pdf}
	\bibitem{ID} ... \LaTeX:\\ \url{https://townhall.virginia.gov/l/GetFile.cfm?File=C:%5CTownHall%5Cdocroot%5CGuidanceDocs%5C960%5CGDoc_DFP_4448_v1.pdf}
	\bibitem{ID} ... \LaTeX:\\ \url{https://ec.europa.eu/growth/sectors/chemicals/specific-chemicals_en}
	\bibitem{safehandlepyrotechnics} ... \LaTeX:\\ \url{http://www.iiakm.org/ojakm/articles/2015/volume3_3/OJAKM_Volume3_3pp27-36.pdf}
	\bibitem{ID} ... \LaTeX:\\ \url{https://www.bristol.gov.uk/documents/20182/1175006/Fireworks+in+retail+premises/6aa6ee24-5b74-43b4-a1d9-747689b1dbc9}
	\bibitem{ID} ... \LaTeX:\\ \url{https://www.eversys.com.br/imagens/uploads/arqs/bra_arquivos/04-software-gerenciador-fireworks-brochura.pdf}
	\bibitem{ID} ... \LaTeX:\\ \url{http://www.doiserbia.nb.rs/img/doi/0354-9836/2016/0354-98361500050G.pdf}
	\bibitem{costinFireworksEmbedded} ... \LaTeX:\\ \url{http://s3.eurecom.fr/docs/wisec14_Costin.pdf}
	\bibitem{ID} ... \LaTeX:\\ \url{https://www.firetechsystems.com/assets/uploads/2018/09/FireWorks-Brochure.pdf}
	\bibitem{ritz2020firesafety} ... \LaTeX:\\ \url{https://blog.ritzsafety.com/fireworks-safety-tips}
	\bibitem{lawson2018criticalsystems} ... \LaTeX:\\ \url{https://www.engineerlive.com/content/fire-detection-and-protection-through-safety-critical-systems}
	%%%%%%%%%%%%%%%%%%%%%%%%%%%%%%%%%%%%%%%%%%%%%%%%%%%%%%%%%%%%%%%%%
%    
%    algemene vragen
%    oorzaken
%    @online{gates18112020boeingcrisis,	ALTauthor = {author},	ALTeditor = {editor},	title = {title},	date = {date},	url = {"https://www.seattletimes.com/business/boeing-aerospace/what-led-to-boeings-737-max-crisis-a-qa/"},}
%    @online{boeing737maxsoftwareprobles,	ALTauthor = {author},	ALTeditor = {editor},	title = {title},	date = {date},	url = {"https://www.schneier.com/blog/archives/2019/04/excellent_analy.html"},}
%    fout in de software
%    @online{avetisov19032019boeingmalwarestate,	ALTauthor = {author},	ALTeditor = {editor},	title = {title},	date = {date},	url = {"https://www.forbes.com/sites/georgeavetisov/2019/03/19/malware-at-30000-feet-what-the-737-max-says-about-the-state-of-airplane-software-security/?sh=4d26f7052a9e"},}
%    het nationaal veiligheidsbelang
%    @online{thompson23112020nationalsecurityboeing,	ALTauthor = {author},	ALTeditor = {editor},	title = {title},	date = {date},	url = {"https://www.forbes.com/sites/lorenthompson/2020/11/23/five-reasons-return-of-boeings-737-max-to-service-is-important-to-national-security/?sh=2128ea552018"},}
%    falend toezicht
%    @online{gates21032019FAAControlSystem,	ALTauthor = {author},	ALTeditor = {editor},	title = {title},	date = {date},	url = {"https://www.seattletimes.com/business/boeing-aerospace/failed-certification-faa-missed-safety-issues-in-the-737-max-system-implicated-in-the-lion-air-crash/"},}
%    onderzoeksrapport
%    @online{faa18112020boeingreview,	ALTauthor = {author},	ALTeditor = {editor},	title = {title},	date = {date},	url = {"https://www.faa.gov/foia/electronic_reading_room/boeing_reading_room/media/737_RTS_Summary.pdf"},}
%    @online{wiki737maxgroundings,	ALTauthor = {author},	ALTeditor = {editor},	title = {title},	date = {date},	url = {"https://en.wikipedia.org/wiki/Boeing_737_MAX_groundings"},}
%    veiligheidsrisico's
%    menselijke fouten
%    @online{campbell02052019boengcrashhumanerrors,	ALTauthor = {author},	ALTeditor = {editor},	title = {title},	date = {date},	url = {"https://www.theverge.com/2019/5/2/18518176/boeing-737-max-crash-problems-human-error-mcas-faa"},}
%    overzicht van crashes
%    @online{hawkins22032019737maxairplanes,	ALTauthor = {author},	ALTeditor = {editor},	title = {title},	date = {date},	url = {"https://www.theverge.com/2019/3/22/18275736/boeing-737-max-plane-crashes-grounded-problems-info-details-explained-reasons"},}
%    veiligheidsopmerking
%    @online{thomas30082020737safest,	ALTauthor = {author},	ALTeditor = {editor},	title = {title},	date = {date},	url = {"https://www.airlineratings.com/news/boeings-737-max-will-one-safest-aircraft-history/"},}
%    aanpassingen
%    @online{boeing737maxdisplay,	ALTauthor = {author},	ALTeditor = {editor},	title = {title},	date = {date},	url = {"https://www.boeing.com/commercial/737max/737-max-software-updates.page"},}
%    waarschuwingen//output signalen
%    @online{fehrm24112020737changes,	ALTauthor = {author},	ALTeditor = {editor},	title = {title},	date = {date},	url = {"https://leehamnews.com/2020/11/24/boeing-737-max-changes-beyond-mcas/"},}
%    software gerelateerde fouten
%    @online{travis18042019737maxsoftwaredevop,	ALTauthor = {author},	ALTeditor = {editor},	title = {title},	date = {date},	url = {"https://spectrum.ieee.org/aerospace/aviation/how-the-boeing-737-max-disaster-looks-to-a-software-developer"},}
%    onderzoeksrapport
%    de rol van de publieke opinie
%    @online{barnett05052019737maxcrisis,	ALTauthor = {author},	ALTeditor = {editor},	title = {title},	date = {date},	url = {"https://pubsonline.informs.org/do/10.1287/orms.2019.05.05/full/"},}
%    onderzoek van europese luchtvaart agentschap
%    @online{easa27012021737maxsafereturn,	ALTauthor = {author},	ALTeditor = {editor},	title = {title},	date = {date},	url = {"https://www.easa.europa.eu/newsroom-and-events/news/easa-declares-boeing-737-max-safe-return-service-europe"},}
%    veiligheidsvraagstuk
%    @online{touitou11032019737tragedies,	ALTauthor = {author},	ALTeditor = {editor},	title = {title},	date = {date},	url = {"https://phys.org/news/2019-03-boeing-max-safety-tragedies.html"},}
%    artikel over sensoren
%    @online{hemmerdinger02022021737maxdeliveries,	ALTauthor = {author},	ALTeditor = {editor},	title = {title},	date = {date},	url = {"https://www.flightglobal.com/airframers/boeing-delays-737-max-10-deliveries-two-years-to-2023/142245.article"},}
%    goedkeuring van europese luchtvaart autoriteiten
%    advies aan de faa
%    @online{bielby27022021faaimprovesafety,	ALTauthor = {author},	ALTeditor = {editor},	title = {title},	date = {date},	url = {"https://www.hstoday.us/subject-matter-areas/airport-aviation-security/oig-tells-faa-to-improve-safety-oversight-following-boeing-737-max-review/"},}
%    @online{boyle18112020737maxupgrade,	ALTauthor = {author},	ALTeditor = {editor},	title = {title},	date = {date},	url = {"https://www.geekwire.com/2020/faas-go-ahead-737-maxs-return-flight-kicks-off-massive-software-upgrade/"},}
%    @online{bergstraburgess122019737maxMcasAlgorithm,	ALTauthor = {author},	ALTeditor = {editor},	title = {title},	date = {date},	url = {"https://www.researchgate.net/publication/338420944_A_Promise_Theoretic_Account_of_the_Boeing_737_Max_MCAS_Algorithm_Affair"},}
%    achtergrond informatie
%    @online{737mcas,	ALTauthor = {author},	ALTeditor = {editor},	title = {title},	date = {date},	url = {"http://www.b737.org.uk/mcas.htm"},}
%    algemeen vertrouwen
%    @online{newburger17052019boeingcrisis,	ALTauthor = {author},	ALTeditor = {editor},	title = {title},	date = {date},	url = {"https://www.cnbc.com/2019/05/16/what-you-need-to-know-about-boeings-737-max-crisis.html"},}
%    toestemming europese autoriteiten
%    problemen
%    @online{arstechnica22012020737problems,	ALTauthor = {author},	ALTeditor = {editor},	title = {title},	date = {date},	url = {"https://arstechnica.com/information-technology/2020/01/737-max-fix-slips-to-summer-and-thats-just-one-of-boeings-problems/"},}
%    uitgebreid artikel over de onderzoeken en het vliegverbod
%    @online{german190620217372yaftergrounded,	ALTauthor = {author},	ALTeditor = {editor},	title = {title},	date = {date},	url = {"https://www.cnet.com/news/boeing-737-max-8-all-about-the-aircraft-flight-ban-and-investigations/"},}
%    computers als oorzaak
%    lessons learned
%    @online{beningo02052019boeinglessons,	ALTauthor = {author},	ALTeditor = {editor},	title = {title},	date = {date},	url = {"https://www.designnews.com/electronics-test/5-lessons-learn-boeing-737-max-fiasco"},}
%    @online{duran05042019boeingspof,	ALTauthor = {author},	ALTeditor = {editor},	title = {title},	date = {date},	url = {"https://www.eurocontrol.int/publication/effects-network-extra-standby-aircraft-and-boeing-737-max-grounding"},}
%    single point of failure
%    @online{ID,	ALTauthor = {author},	ALTeditor = {editor},	title = {title},	date = {date},	url = {"https://dmd.solutions/blog/2019/04/05/how-a-single-point-of-failure-spof-in-the-mcas-software-could-have-caused-the-boeing-737-max-crash-in-ethiopia/"},}
%    @online{makichuck24012021737fearflying,	ALTauthor = {author},	ALTeditor = {editor},	title = {title},	date = {date},	url = {"https://asiatimes.com/2021/01/boeings-737-max-and-the-fear-of-flying/"},}
%    lijst van tehnische aanpassingen
%    @online{caa737modifications,	ALTauthor = {author},	ALTeditor = {editor},	title = {title},	date = {date},	url = {"https://www.caa.co.uk/Consumers/Guide-to-aviation/Boeing-737-MAX/"},}
%    @online{oestergaard14122020boeingdeliveries,	ALTauthor = {author},	ALTeditor = {editor},	title = {title},	date = {date},	url = {"https://dsm.forecastinternational.com/wordpress/2020/12/14/airbus-and-boeing-report-november-2020-commercial-aircraft-orders-and-deliveries/"},}
%    code lek
%    @online{reenberg787flaws,	ALTauthor = {author},	ALTeditor = {editor},	title = {title},	date = {date},	url = {"https://www.wired.com/story/boeing-787-code-leak-security-flaws/"},}
%    @online{fitch16092020737backlogrisks,	ALTauthor = {author},	ALTeditor = {editor},	title = {title},	date = {date},	url = {"https://www.fitchratings.com/research/corporate-finance/boeing-737-max-return-backlog-risks-remain-16-09-2020"},}
%    Cultuurverandering, deregulatie, systeemwijziging of gewoon een kwestie van competentie
%    @online{willis27082020737maxfailures,	ALTauthor = {author},	ALTeditor = {editor},	title = {title},	date = {date},	url = {"https://www.aerospacetestinginternational.com/features/what-broke-the-737-max.html"},}
%    extra aanpassingen
%    @online{ostrower11062020more737changes,	ALTauthor = {author},	ALTeditor = {editor},	title = {title},	date = {date},	url = {"https://theaircurrent.com/aviation-safety/boeings-737-max-software-done-but-regulators-plot-more-changes-after-jets-return/"},}
%    wat ging er mis een analyse van een ex-iloot
%    De utoriteiten waren op de hoogte
%    @online{hruska13122019faaknown737crashrate,	ALTauthor = {author},	ALTeditor = {editor},	title = {title},	date = {date},	url = {"https://www.extremetech.com/extreme/303373-the-faa-knew-the-737-max-was-dangerous-and-kept-it-flying-anyway"},}
%    kwaliteiten van het alarmsysteem niet goed bekend
%    @online{bloomberg26092019failedpred,	ALTauthor = {author},	ALTeditor = {editor},	title = {title},	date = {date},	url = {"https://time.com/5687473/boeing-737-alarm-system/"},}
%    @online{whiteman09072020boengcancelstock,	ALTauthor = {author},	ALTeditor = {editor},	title = {title},	date = {date},	url = {"https://www.nasdaq.com/articles/boeing-gets-dealt-another-737-max-cancellation-blow.-what-it-means-for-boeing-stock-2020"},}
%    @online{leopold09192019boeingreliability,	ALTauthor = {author},	ALTeditor = {editor},	title = {title},	date = {date},	url = {"https://www.eetimes.com/boeing-crashes-highlight-a-worsening-reliability-crisis/"},}
%    veiligheidsvraagstuk
%    @online{koenig11122019737crashesnofix,	ALTauthor = {author},	ALTeditor = {editor},	title = {title},	date = {date},	url = {"https://www.latimes.com/business/story/2019-12-11/faa-boeing-737-max-crashes"},}
%    probleemanalyse, veiligheidsvraagstuk
%    @online{dohertylindeman15032019737problems,	ALTauthor = {author},	ALTeditor = {editor},	title = {title},	date = {date},	url = {"https://www.politico.com/story/2019/03/15/boeing-737-max-grounding-1223072"},}
%    falend toezicht
%    @online{stodder02102019corruptoversight,	ALTauthor = {author},	ALTeditor = {editor},	title = {title},	date = {date},	url = {"https://www.pogo.org/analysis/2019/10/corrupted-oversight-the-faa-boeing-and-the-737-max/"},}
%    @online{afacwaLostSafeguards,	ALTauthor = {author},	ALTeditor = {editor},	title = {title},	date = {date},	url = {"https://www.afacwa.org/the_inside_story_of_mcas_seattle_times"},}
%    doelstellingen en veiligheidsvraagstukken
%    @online{swayne18032019profitssafety,	ALTauthor = {author},	ALTeditor = {editor},	title = {title},	date = {date},	url = {"https://www.marxist.com/737-max-scandal-boeing-putting-profits-before-safety.htm"},}
%    @online{freed26022021liftaustraliaban,	ALTauthor = {author},	ALTeditor = {editor},	title = {title},	date = {date},	url = {"https://finance.yahoo.com/news/australia-lifts-ban-boeing-737-035817682.html?guccounter=1&guce_referrer=aHR0cHM6Ly93d3cuZ29vZ2xlLmNvbS8&guce_referrer_sig=AQAAAHZCJYy_0A5VS2WiPoCvH4xdrRNkmkdsv5EWJ2RLIz_AS-rxsTty6AF1_HlmJiRyWYqCXDi4p0Xs4isYkNkCq2Pfo-pQ60Xz_IfTNjm4FgoZiBMC4zpZlB6F0fwecrjE_ujAXZzG4xPJnWCd8-G3VLlPTY8h3H31eQ1i8hY9AIyy"},}
%    autoriteiten krijgen tik op de vingers
%    @online{reed15032019softwareattention,	ALTauthor = {author},	ALTeditor = {editor},	title = {title},	date = {date},	url = {"https://medium.com/@jpaulreed/the-737max-and-why-software-engineers-should-pay-attention-a041290994bd"},}
%    @online{news17032019softwareexplains,	ALTauthor = {author},	ALTeditor = {editor},	title = {title},	date = {date},	url = {"https://news.ycombinator.com/item?id=19414775"},}
%    @online{legget21122020eu737maxsafe,	ALTauthor = {author},	ALTeditor = {editor},	title = {title},	date = {date},	url = {"https://www.bbc.com/news/55366320"},}
%    @online{marketscreener0103221737chinarecertification,	ALTauthor = {author},	ALTeditor = {editor},	title = {title},	date = {date},	url = {"https://www.marketscreener.com/news/latest/China-studies-Boeing-737-MAX-recertification-wants-safety-concerns-fully-addressed--32569125/"},}
%    motor in brand
%    @online{euractiv22022021737firegrounds,	ALTauthor = {author},	ALTeditor = {editor},	title = {title},	date = {date},	url = {"https://www.euractiv.com/section/aviation/news/boeing-grounds-777s-after-engine-fire/"},}
%    @online{benny18022019737returnUAE,	ALTauthor = {author},	ALTeditor = {editor},	title = {title},	date = {date},	url = {"https://gulfnews.com/business/aviation/uae-airspace-to-see-return-of-boeing-737-max-1.1613627548923"},}
%    motor in brand gevlogen
%    @online{biersmichel22022021777grounds,	ALTauthor = {author},	ALTeditor = {editor},	title = {title},	date = {date},	url = {"https://techxplore.com/news/2021-02-boeing-urges-grounding-777s.html"},}
%    @online{ID,	ALTauthor = {author},	ALTeditor = {editor},	title = {title},	date = {date},	url = {"https://www.politico.eu/article/uk-temporarily-bans-some-boeing-aircraft-after-pratt-whitney-engine-incidents/"},}
%    @online{reuters23022021777metalfatigue,	ALTauthor = {author},	ALTeditor = {editor},	title = {title},	date = {date},	url = {"https://www.timeslive.co.za/news/world/2021-02-23-damage-to-united-boeing-777-engine-consistent-with-metal-fatigue--ntsb/"},}
%    faa was niet kritisch genoeg
%    @online{ID,	ALTauthor = {author},	ALTeditor = {editor},	title = {title},	date = {date},	url = {"https://federalnewsnetwork.com/government-news/2021/02/federal-watchdog-blasts-faa-over-certification-of-boeing-jet/"},}
%    
%    
%    
%    
%    
%    
%    
%    \cite{bnnvara13062018malirapport}
%    \cite{eucal11012021malimissieverlengd}
%    \cite{nos21052014zorgenmalimissie}
%    \cite{meijnders}
%    \cite{bnrwebredactie}
%    \cite{keultjes01062016malimissiecoalitie}
%    \cite{veenhof18012019}
%    
%    \cite{isitman06012016militair}
%    \cite{nporadio11072016filmdemissie}
%    \cite{parlementairmonitor15122013mortierongeluk}
%    
%    sollicitatie
%    de bureaucratie
%    aankomst
%    interview van de burgerbevolking
%    steun van de bevolking minuut 15:00
%    de organisatie minuut 23:00
%    De militaire briefing minuut 34:00
%    prioriteit minuut 39:00
%    briefing minuut 40:00
%    de communicatie met ministerie over inlichten minuut 44:00
%    @online{DemissieFilm,	ALTauthor = {author},	ALTeditor = {editor},	title = {title},	date = {date},	url = {"https://www.2doc.nl/documentaires/series/2doc/2016/juli/de-missie.html"},}
%    \cite{DemissieFilm}
%    
%    
%    
%    
%    
%    
%    verhaal van brandweermannen
%    
%    \cite{staff31082015tanjinblastunrevealed}
%    artikel
%    
%    \cite{chinafile18082015tanjinexplosion}
%    invloed van social media
%    \bibitem{ID} ... \LaTeX:\\ \url{https://www.economist.com/asia/2015/08/18/a-blast-in-tianjin-sets-off-an-explosion-online}
%    \cite{}
%    \bibitem{ID} ... \LaTeX:\\ \url{https://america.cgtn.com/2015/08/12/explosion-reported-in-tianjin-china}
%    \cite{}
%    \bibitem{ID} ... \LaTeX:\\ \url{https://factcheck.afp.com/no-photo-was-taken-chinese-city-tianjin-august-2015}
%    \cite{}
%    vergelijking van twee rampen
%    \bibitem{ID} ... \LaTeX:\\ \url{https://airshare.air-inc.com/how-does-the-beirut-explosion-compare-to-tianjin}
%    \cite{}
%    overheid en media
%    \bibitem{ID} ... \LaTeX:\\ \url{https://newbloommag.net/2015/08/17/tianjin-explosion/}
%    \cite{}
%    chemische industrie ondeer de loep
%    \bibitem{ID} ... \LaTeX:\\ \url{https://www.voanews.com/east-asia-pacific/tianjin-blast-puts-spotlight-chemical-industry}
%    \cite{}
%    \bibitem{ID} ... \LaTeX:\\ \url{https://abcnews.go.com/International/apocalyptic-aftermath-devastating-images-tianjin-china-explosions/story?id=33057017}
%    \cite{}
%    \bibitem{ID} ... \LaTeX:\\ \url{https://www.reachingoutacrossdurham.co.uk/osk/tianjin-explosion-2021}
%    \cite{}
%    \bibitem{pinghuang2410201TanjinFactreport} ... \LaTeX:\\ \url{https://aiche.onlinelibrary.wiley.com/doi/abs/10.1002/prs.11789}
%    \cite{pinghuang2410201TanjinFactreport}
%    \bibitem{ID} ... \LaTeX:\\ \url{https://www.automotivelogistics.media/thousands-of-cars-destroyed-in-tianjin-port-explosions/13570.article}
%    \cite{}
%    \bibitem{ID} ... \LaTeX:\\ \url{https://www.joc.com/port-news/asian-ports/port-tianjin/tianjin-port-explosions-could-be-most-expensive-maritime-disaster_20150826.html}
%    \cite{}
%    \bibitem{ID} ... \LaTeX:\\ \url{https://www.bloomberg.com/news/articles/2015-08-12/explosion-in-northern-china-shatters-windows-causes-injuries}
%    \cite{}
%    \bibitem{ID} ... \LaTeX:\\ \url{https://unece.org/fileadmin/DAM/env/documents/2016/TEIA/OECD_WGCA_24-27_OCT_2016/Session_3_Zhao_-__Introduction_of_Tianjin_Accident_-_Jinsong_Zhao.pdf}
%    \cite{}
%    gemaakte fouten
%    \bibitem{portoTanjinExplosionSight} ... \LaTeX:\\ \url{https://porteconomicsmanagement.org/pemp/contents/part6/port-resilience/site-2015-tianjin-port-explosions/}
%    \cite{portoTanjinExplosionSight}
%    \bibitem{ID} ... \LaTeX:\\ \url{https://www.alamy.com/stock-image-tianjin-china-17th-aug-2015-tianjin-explosion-aftermath-blast-site-165334778.html}
%    \cite{}
%    \bibitem{ID} ... \LaTeX:\\ \url{https://www.popularmechanics.com/technology/news/a16871/massive-explosions-china-city-of-tianjin/}
%    \cite{}
%    \bibitem{imago17082015TanjinApartmentImages} ... \LaTeX:\\ \url{https://www.imago-images.com/st/0080815934}
%    \cite{imago17082015TanjinApartmentImages}
%    \bibitem{trager14082015Chemicalblast} ... \LaTeX:\\ \url{https://www.chemistryworld.com/news/deadly-chemical-blast-at-chinese-port/8857.article}
%    \cite{trager14082015Chemicalblast}
%    \bibitem{pangeramo27082015TanjinExplosion} ... \LaTeX:\\ \url{https://www.process-worldwide.com/tianjin-explosion-from-chemical-perspective-insights-and-backgrounds-a-502381/}
%    \cite{pangeramo27082015TanjinExplosion}
%    vergelijking met andere explosies
%    \bibitem{ap06082020ammaniumnitrate} ... \LaTeX:\\ \url{https://apnews.com/article/lebanon-fires-us-news-explosions-middle-east-53f4206a7f1db0812262a15d22e1e58f}
%    \cite{ap06082020ammaniumnitrate}
%    invloed van de ramp op de industrie
%    \bibitem{morris14082015TanjinIndustryImpact} ... \LaTeX:\\ \url{https://fortune.com/2015/08/14/tianjin-port-explosion-shipping-delays/}
%    \cite{morris14082015TanjinIndustryImpact}
%    is er sprake van een doofpot
%    \bibitem{milesyu20082015exposingtoxicgovlines} ... \LaTeX:\\ \url{https://www.washingtontimes.com/news/2015/aug/20/inside-china-tianjin-explosions-cover-up-exposes-b/}
%    \cite{milesyu20082015exposingtoxicgovlines}
%    eigendomsverzekering
%    \bibitem{artemis30032016tanjininsurance} ... \LaTeX:\\ \url{https://www.artemis.bm/news/tianjin-explosions-property-insurance-loss-could-reach-3-5bn-swiss-re/}
%    \cite{artemis30032016tanjininsurance}
%    \bibitem{aidenxiatanjinblast} ... \LaTeX:\\ \url{https://www.thechinastory.org/yearbooks/yearbook-2015/forum-the-abyss-%E5%9D%8E/tianjin-explosions/}
%    \cite{aidenxiatanjinblast}
%    effecten op de lange termijn
%    \bibitem{danwangTanjinflexreport} ... \LaTeX:\\ \url{https://www.flexport.com/blog/tianjin-explosion-effect-on-supply-chains/}
%    \cite{danwangTanjinflexreport}
%    \bibitem{keyHighlightsTanjin} ... \LaTeX:\\ \url{https://www.cicm.org.my/images/articles/CICM-Article-on-Tianjin-Blast-Oct2015.pdf}
%    \cite{keyHighlightsTanjin}
%    lessons learned
%    \bibitem{ID} ... \LaTeX:\\ \url{https://www.genre.com/knowledge/blog/lessons-from-the-tianjin-explosion-en.html}
%    \cite{}
%    \bibitem{ID} ... \LaTeX:\\ \url{https://www.ft.com/content/ad62904c-44ce-11e5-b3b2-1672f710807b}
%    \cite{}
%    \bibitem{hartley13082015videofootage} ... \LaTeX:\\ \url{https://www.huffingtonpost.co.uk/2015/08/13/tianjin-explosion-china-shocking-footage-caught-on-camera_n_7980888.html}
%    \cite{hartley13082015videofootage}
%    \bibitem{odonnel01062017firetanjinblast2015} ... \LaTeX:\\ \url{https://www.thatsmags.com/china/post/19189/massive-fire-rocks-tianjin-port}
%    \cite{odonnel01062017firetanjinblast2015}
%    gevolgen voor de industrie
%    \bibitem{ID} ... \LaTeX:\\ \url{https://www.everstream.ai/risk-center/special-reports/the-jiangsu-yancheng-explosion/}
%    \cite{}
%    \bibitem{fan15082015newyorkermistrustchina} ... \LaTeX:\\ \url{https://www.newyorker.com/news/news-desk/after-tianjin-an-outbreak-of-mistrust-in-china}
%    \cite{fan15082015newyorkermistrustchina}
%    framing vanuit de chinese media
%    \bibitem{yanlidongchinamediaframingTanjin} ... \LaTeX:\\ \url{https://www.neliti.com/publications/101997/the-chinese-media-framing-of-the-2015s-tianjin-explosion}
%    \cite{yanlidongchinamediaframingTanjin}
%    \bibitem{evans27092017TnjinInsurance} ... \LaTeX:\\ \url{https://www.reinsurancene.ws/chinese-insurers-settle-1-5-billion-tianjin-blast-claims/}
%    \cite{evans27092017TnjinInsurance}
%    niewsartikel
%    \bibitem{jasi26032019chineschemplant} ... \LaTeX:\\ \url{https://www.thechemicalengineer.com/news/update-78-confirmed-dead-after-chinese-chemicals-plant-explosion/}
%    \cite{jasi26032019chineschemplant}
%    \bibitem{shiqingTanjinExecutiveSentence} ... \LaTeX:\\ \url{https://www.caixinglobal.com/2016-11-10/chinese-executive-receives-suspended-death-sentence-over-2015-tianjin-warehouse-blast-101006325.html}
%    \cite{shiqingTanjinExecutiveSentence}
%    toegang tot de ramplplek vanuit de okale journalistiek
%    \bibitem{sophiebeach15082015} ... \LaTeX:\\ \url{https://chinadigitaltimes.net/2015/08/he-xiaoxin-how-far-can-i-go-and-how-much-can-i-do/}
%    \cite{sophiebeach15082015}
%    artikel
%    \bibitem{ID} ... \LaTeX:\\ \url{https://www.wnpr.org/post/china-examines-aftermath-immense-twin-explosions-killed-dozens}
%    \cite{}
%    \bibitem{hamzeh05082020BeirutBlast} ... \LaTeX:\\ \url{https://theconversation.com/what-is-ammonium-nitrate-the-chemical-that-exploded-in-beirut-143979}
%    \cite{hamzeh05082020BeirutBlast}
%    \bibitem{chemwatch18082015TanjiinExplosion} ... \LaTeX:\\ \url{https://chemicalwatch.com/36730/nationwide-inspections-in-china-follow-tianjin-explosion}
%    \cite{chemwatch18082015TanjiinExplosion}
%    \bibitem{thehindu15062019chinaExplosion} ... \LaTeX:\\ \url{https://www.thehindu.com/news/international/investigation-begun-into-china-gas-explosion-as-toll-rises/article34818324.ece}
%    \cite{thehindu15062019chinaExplosion}
%    \bibitem{santagotimes24032019chinablast} ... \LaTeX:\\ \url{https://santiagotimes.cl/2019/03/24/64-killed-600-injured-in-china-chemical-plant-blast/}
%    \cite{santagotimes24032019chinablast}
%    oorzaken
%    \bibitem{klingecorp28042020causedTanjin} ... \LaTeX:\\ \url{https://klingecorp.com/blog/what-caused-the-tianjin-explosions/}
%    \cite{klingecorp28042020causedTanjin}
%    case study
%    \bibitem{mcgarryExplosions2017 ... \LaTeX:\\ \url{https://www.preventionweb.net/educational/view/57235}
%    	\cite{mcgarryExplosions2017}
%    	niewsartikel
%    	\bibitem{roswnfeld13082015TanjinReports} ... \LaTeX:\\ \url{https://www.cnbc.com/2015/08/12/explosion-in-tianjin-china.html}
%    	\cite{roswnfeld13082015TanjinReports}
%    	chronologische uiteenzetting
%    	\bibitem{aria12082015explosionaTanjin} ... \LaTeX:\\ \url{https://www.aria.developpement-durable.gouv.fr/wp-content/files_mf/A46803_a46803_fiche_impel_006.pdf}
%    	\cite{aria12082015explosionaTanjin}
%    	corruptie
%    	\bibitem{ID} ... \LaTeX:\\ \url{https://www.nytimes.com/2015/08/31/world/asia/behind-tianjin-tragedy-a-company-that-flouted-regulations-and-reaped-profits.html}
%    	\cite{}
%    	mismanagement als oorzaak
%    	\bibitem{ID} ... \LaTeX:\\ \url{https://www.nytimes.com/2016/02/06/world/asia/tianjin-explosions-were-result-of-mismanagement-china-finds.html}
%    	\cite{}
%    	\bibitem{ID} ... \LaTeX:\\ \url{https://cen.acs.org/articles/94/web/2016/02/Chinese-Investigators-Identify-Cause-Tianjin.html}
%    	\cite{}
%    	autoriteiten publiceren onderoeksrapport
%    	\bibitem{tremblay11022016chineseInvestigatorsTanjin} ... \LaTeX:\\ \url{https://cen.acs.org/articles/94/i7/Chinese-Investigators-Identify-Cause-Tianjin.html}
%    	\cite{tremblay11022016chineseInvestigatorsTanjin}
%    	fotos van de rampplek
%    	\bibitem{taylor13082015TanjinExplosianAftermath} ... \LaTeX:\\ \url{https://www.theatlantic.com/photo/2015/08/photos-of-the-aftermath-of-the-massive-explosions-in-tianjin-china/401228/}
%    	\cite{taylor13082015TanjinExplosianAftermath}
%    	\bibitem{ID} ... \LaTeX:\\ \url{https://edition.cnn.com/2015/08/13/asia/china-tianjin-explosions/index.html}
%    	\cite{}
%    	niuwesartiekel}
%    \bibitem{associatedPresss13082013} ... \LaTeX:\\ \url{https://www.cbc.ca/news/world/china-explosion-tianjin-1.3189455}
%    \cite{associatedPresss13082013}
%    verantwoordelijke
%    \bibitem{ID} ... \LaTeX:\\ \url{https://www.thestar.com/news/world/2016/11/09/chinese-executive-gets-death-sentence-over-tianjin-explosion-in-2015.html}
%    \cite{}
%    risicobeperking/controle
%    \bibitem{ID} ... \LaTeX:\\ \url{https://www.swissre.com/en/china/news-insights/articles/analysis-of-tianjin-port-explosion-china.html}
%    \cite{}
%    censuur
%    \bibitem{ID} ... \LaTeX:\\ \url{https://foreignpolicy.com/2015/09/10/censored-china-young-survivor-tianjin-explosion-viral-post/}
%    \cite{}
%    censuur
%    \bibitem{ID} ... \LaTeX:\\ \url{https://qz.com/756872/a-year-after-the-tianjin-blast-public-mourning-and-discussion-about-it-are-still-censored-in-china/}
%    \cite{}
%    verschillende artikelen
%    \bibitem{ID} ... \LaTeX:\\ \url{https://www.scmp.com/topics/tianjin-warehouse-explosion-2015}
%    \cite{}
%    \bibitem{ID} ... \LaTeX:\\ \url{https://www.wsj.com/articles/BL-CJB-27664}
%    \cite{}
%    \bibitem{ID} ... \LaTeX:\\ \url{https://www.nbcnews.com/news/world/tianjin-explosions-californian-witness-filmed-dramatic-china-blasts-n409701}
%    \cite{}
%    \bibitem{ID} ... \LaTeX:\\ \url{https://ui.adsabs.harvard.edu/abs/2016AGUFM.S13D..06P/abstract}
%    afwikkeling van de ramp
%    \cite{}
%    \bibitem{ID} ... \LaTeX:\\ \url{https://chinadialogue.net/en/pollution/9188-back-to-the-blast-zone-one-year-after-the-tianjin-explosion/}
%    \cite{}
%    \bibitem{ID} ... \LaTeX:\\ \url{https://www.wired.com/2015/08/chinas-huge-tianjin-explosion-looked-like-space/}
%    \cite{}
%    \bibitem{ID} ... \LaTeX:\\ \url{https://www.abc.net.au/news/2015-08-13/explosion-rocks-north-chinese-city-of-tianjin/6693336?nw=0}
%    \cite{}
%    ambtenaren onderzocht
%    
%    risico-inschatting
%    \bibitem{ID} ... \LaTeX:\\ \url{https://www.mdpi.com/2071-1050/12/3/1169/htm}
%    \cite{}
%    \bibitem{ID} ... \LaTeX:\\ \url{https://www.mdpi.com/2071-1050/12/3/1169/htm}
%    \cite{}
%    \bibitem{ID} ... \LaTeX:\\ \url{https://www.cbsnews.com/news/tianjin-port-china-massive-explosion-hundreds-injured/}
%    \cite{}
%    \bibitem{ID} ... \LaTeX:\\ \url{https://www.hkjcdpri.org.hk/download/casestudies/Tianjin_CASE.pdf}
%    \cite{}
%    \bibitem{ID} ... \LaTeX:\\ \url{https://time.com/3996168/tianjin-explosion-china-pictures/}
%    \cite{}
%    onderzoeksrapport
%    \bibitem{ID} ... \LaTeX:\\ \url{https://www.hfw.com/Tianjin-Port-explosion-August-2015}
%    \cite{}
%    \bibitem{un20082015InvestigationTanjin} ... \LaTeX:\\ \url{https://news.un.org/en/story/2015/08/506912-following-tianjin-explosion-un-expert-calls-china-ensure-transparent}
%    \cite{un20082015InvestigationTanjin}
%    \bibitem{france2412082015TnjinExplosion} ... \LaTeX:\\ \url{https://www.france24.com/en/20150812-huge-explosions-rock-chinese-city-tianjin}
%    \cite{france2412082015TnjinExplosion}
%    \bibitem{npr14082015TanjinCause} ... \LaTeX:\\ \url{https://choice.npr.org/index.html?origin=https://www.npr.org/2015/08/14/432280627/what-caused-the-warehouse-explosions-in-tianjin-china}
%    \cite{npr14082015TanjinCause}
%    123 verantwoordelijken
%    \bibitem{bbc05022016TanjinResponsibles} ... \LaTeX:\\ \url{https://www.bbc.com/news/world-asia-china-35506311}
%    \cite{bbc05022016TanjinResponsibles}
%    \bibitem{ID} ... \LaTeX:\\ \url{https://www.washingtonpost.com/gdpr-consent/?next_url=https%3a%2f%2fwww.washingtonpost.com%2fnews%2fworldviews%2fwp%2f2015%2f08%2f12%2fvideos-show-chinese-city-of-tianjin-rocked-by-enormous-explosion%2f}
%    \cite{}
%    lang artiekel
%    \bibitem{CBodeen15082015TanjinExplosion} ... \LaTeX:\\ \url{https://www.businessinsider.com/the-chemical-explosion-in-china-killed-more-than-100-people-and-the-devastation-is-unreal-2015-8?international=true&r=US&IR=T}
%    \cite{CBodeen15082015TanjinExplosion}
%    \bibitem{ID} ... \LaTeX:\\ \url{https://pubmed.ncbi.nlm.nih.gov/27311537/}
%    \cite{}
%    \bibitem{reutersTanjinInsurance} ... \LaTeX:\\ \url{https://www.reuters.com/article/us-china-blast-insurance-idUSKCN0QM0N220150817}
%    \cite{reutersTanjinInsurance}
%    \bibitem{yu082016evaluationTanjin2015} ... \LaTeX:\\ \url{https://www.sciencedirect.com/science/article/abs/pii/S0305417916300079}
%    \cite{yu082016evaluationTanjin2015}
%    \bibitem{wiki2015TanjinExplosions} ... \LaTeX:\\ \url{https://en.wikipedia.org/wiki/2015_Tianjin_explosions}
%    \cite{wiki2015TanjinExplosions}
%    \bibitem{bbc17082015whathappenedTanjin} ... \LaTeX:\\ \url{https://www.bbc.com/news/world-asia-china-33844084}
%    \cite{bbc17082015whathappenedTanjin}
%    \bibitem{mortimer19082016taijinexplosioncrater} ... \LaTeX:\\ \url{https://www.independent.co.uk/news/world/asia/tianjin-explosion-photos-china-chemical-factory-accident-crater-revealed-a7199591.html}
%    \cite{mortimer19082016taijinexplosioncrater}
%    veiigheidshandhaving
%    \bibitem{internationallabourofficeChmControlTooliit} ... \LaTeX:\\ \url{https://www.ilo.org/legacy/english/protection/safework/ctrl_banding/toolkit/main_guide.pdf}
%    \cite{internationallabourofficeChmControlTooliit}
%    \bibitem{ID} ... \LaTeX:\\ \url{https://echa.europa.eu/documents/10162/21332507/guide_chemical_safety_sme_en.pdf}
%    \cite{}
%    \bibitem{euTaxationCustomsICSC} ... \LaTeX:\\ \url{https://ec.europa.eu/taxation_customs/dds2/SAMANCTA/EN/Safety/AppendixD_EN.htm}
%    \cite{euTaxationCustomsICSC}
%    \bibitem{iloWHOChemSafetyCards} ... \LaTeX:\\ \url{https://www.ilo.org/safework/info/publications/WCMS_113134/lang--en/index.htm}
%    \cite{iloWHOChemSafetyCards}
%    
%    
%    
%    Wat is er gebeurd?
%    @online{wikiSchipholbrand,	ALTauthor = {author},	ALTeditor = {editor},	title = {title},	date = {date},	url = {"https://nl.wikipedia.org/wiki/Schipholbrand"},}
%    artikel
%    @online{schipholbrand27102005video,	ALTauthor = {author},	ALTeditor = {editor},	title = {title},	date = {date},	url = {"https://www.youtube.com/watch?v=1i-hfEzxFfk"},}
%    psychologische gevolgen
%    rapport
%    @online{onderzoeksraad2610schipholoost,	ALTauthor = {author},	ALTeditor = {editor},	title = {title},	date = {date},	url = {"https://www.onderzoeksraad.nl/nl/page/392/brand-cellencomplex-schiphol-oost-nacht-van-26-op-27-oktober"},}
%    artikel met video
%    herdenking
%    impact op de persoon
%    herdenking
%    @online{schipholbrandvideoargos,	ALTauthor = {author},	ALTeditor = {editor},	title = {title},	date = {date},	url = {"https://www.vpro.nl/argos/speel~POMS_VPRO_461907~schadevergoeding-voor-ex-verdachte-schipholbrand~.html"},}
%    chronologie
%    @online{nunl30052023feitenoverzicht,	ALTauthor = {author},	ALTeditor = {editor},	title = {title},	date = {date},	url = {"https://www.nu.nl/binnenland/3355935/feitenoverzicht-schipholbrand-en-rechtszaken.html"},}
%    tijdlijn
%    @online{ID,	ALTauthor = {author},	ALTeditor = {editor},	title = {title},	date = {date},	url = {"https://www.singeluitgeverijen.nl/isbn/de-schipholbrand/"},}
%    vervolgens van ministers
%    beeldanalyse en reconstructie
%    @online{ID,	ALTauthor = {author},	ALTeditor = {editor},	title = {title},	date = {date},	url = {"https://eenvandaag.avrotros.nl/item/schipholbrand-niet-ontstaan-in-cel-11/"},}
%    herdenking
%    korte samenvatting
%    rapport
%    artikel
%    verwijzing naar het rapport vanuit de politieke oppositie
%    beeld vanuit de gevangenisbewaarder
%    nationaliteit slachtoffers schipholbrand
%    verblijfsvergunning voor de slachtoffers
%    gen schadevergoeding voor de verdachte
%    verdachte voor de rechter
%    geen schadevergoeding voor verdachte
%    artikel wat ging er mis bji de schipholbrand
%    brand veroorzaakt door een peuk
%    smaadschrift
%    bewakers worden niet vervolgd
%    proces schipholbrand moet over en de brandveilgheid moet worden verbeterd
%    de rol van het parlement in de evaluatie
%    @online{parlementairemonitorschipholbrand,	ALTauthor = {author},	ALTeditor = {editor},	title = {title},	date = {date},	url = {"https://www.parlementairemonitor.nl/9353000/1/j9vvij5epmj1ey0/vi3aof7awcxg"},}
%    onderzoeksmemo
%    herdenking
%    @online{ID,	ALTauthor = {author},	ALTeditor = {editor},	title = {title},	date = {date},	url = {"https://archief.ntr.nl/nova/page/detail/uitzendingen/3847/Den%20Haag%20Vandaag_%20herdenking%20Schipholbrand.html"},}
%    herdenking
%    invloed van de ramp op samenleving
%    @online{videonpoNOVA13112008,	ALTauthor = {author},	ALTeditor = {editor},	title = {title},	date = {date},	url = {"https://www.npostart.nl/heropen-onderzoek-schipholbrand/13-11-2008/POMS_NTR_103332"},}
%    opmerkelijk rapport gestolen in de nasleep
%    @online{rizoomes01052014schipholbrand,	ALTauthor = {author},	ALTeditor = {editor},	title = {title},	date = {date},	url = {"https://www.rizoomes.nl/brandweer/brand-cellencomplex-schiphol/"},}
%    
%    
%    
%    publicaties
%    @online{heuvelkroesschipholbrandcamerabeelden,	ALTauthor = {author},	ALTeditor = {editor},	title = {title},	date = {date},	url = {"http://www.msnp.nl/downloads/Onderzoeksmemo%20beeldanalyse%20Schipholbrand%20prot.pdf"},}
%    Wat waren de regels destijds?
%    Waren de autoriteiten in staat om op tijd in te grijpen of om erger te voorkomen?
%    Wat is er gedaan om de veiligheid van illegalen en gevangenissbewaarders te verbeteren
%    
%    
%    
%    algemene vragen
%    oorzaken
%    \bibitem{gates18112020boeingcrisis} ... \LaTeX:\\ \url{https://www.seattletimes.com/business/boeing-aerospace/what-led-to-boeings-737-max-crisis-a-qa/}
%    \cite{gates18112020boeingcrisis}
%    \bibitem{boeing737maxsoftwareprobles} ... \LaTeX:\\ \url{https://www.schneier.com/blog/archives/2019/04/excellent_analy.html}
%    \cite{boeing737maxsoftwareprobles}
%    fout in de software
%    \bibitem{avetisov19032019boeingmalwarestate} ... \LaTeX:\\ \url{https://www.forbes.com/sites/georgeavetisov/2019/03/19/malware-at-30000-feet-what-the-737-max-says-about-the-state-of-airplane-software-security/?sh=4d26f7052a9e}
%    \cite{avetisov19032019boeingmalwarestate}
%    het nationaal veiligheidsbelang
%    \bibitem{thompson23112020nationalsecurityboeing} ... \LaTeX:\\ \url{https://www.forbes.com/sites/lorenthompson/2020/11/23/five-reasons-return-of-boeings-737-max-to-service-is-important-to-national-security/?sh=2128ea552018}
%    \cite{thompson23112020nationalsecurityboeing}
%    falend toezicht
%    \bibitem{gates21032019FAAControlSystem} ... \LaTeX:\\ \url{https://www.seattletimes.com/business/boeing-aerospace/failed-certification-faa-missed-safety-issues-in-the-737-max-system-implicated-in-the-lion-air-crash/}
%    \cite{gates21032019FAAControlSystem}
%    onderzoeksrapport
%    \bibitem{faa18112020boeingreview} ... \LaTeX:\\ \url{https://www.faa.gov/foia/electronic_reading_room/boeing_reading_room/media/737_RTS_Summary.pdf}
%    \cite{faa18112020boeingreview}
%    \bibitem{wiki737maxgroundings} ... \LaTeX:\\ \url{https://en.wikipedia.org/wiki/Boeing_737_MAX_groundings}
%    \cite{wiki737maxgroundings}
%    veiligheidsrisico's
%    menselijke fouten
%    \bibitem{campbell02052019boengcrashhumanerrors} ... \LaTeX:\\ \url{https://www.theverge.com/2019/5/2/18518176/boeing-737-max-crash-problems-human-error-mcas-faa}
%    \cite{campbell02052019boengcrashhumanerrors}
%    overzicht van crashes
%    \bibitem{hawkins22032019737maxairplanes} ... \LaTeX:\\ \url{https://www.theverge.com/2019/3/22/18275736/boeing-737-max-plane-crashes-grounded-problems-info-details-explained-reasons}
%    \cite{hawkins22032019737maxairplanes}
%    veiligheidsopmerking
%    \bibitem{thomas30082020737safest} ... \LaTeX:\\ \url{https://www.airlineratings.com/news/boeings-737-max-will-one-safest-aircraft-history/}
%    \cite{thomas30082020737safest}
%    aanpassingen
%    \bibitem{boeing737maxdisplay} ... \LaTeX:\\ \url{https://www.boeing.com/commercial/737max/737-max-software-updates.page}
%    \cite{boeing737maxdisplay}
%    waarschuwingen//output signalen
%    \bibitem{fehrm24112020737changes} ... \LaTeX:\\ \url{https://leehamnews.com/2020/11/24/boeing-737-max-changes-beyond-mcas/}
%    \cite{fehrm24112020737changes}
%    software gerelateerde fouten
%    \bibitem{travis18042019737maxsoftwaredevop} ... \LaTeX:\\ \url{https://spectrum.ieee.org/aerospace/aviation/how-the-boeing-737-max-disaster-looks-to-a-software-developer}
%    \cite{travis18042019737maxsoftwaredevop}
%    onderzoeksrapport
%    de rol van de publieke opinie
%    \bibitem{barnett05052019737maxcrisis} ... \LaTeX:\\ \url{https://pubsonline.informs.org/do/10.1287/orms.2019.05.05/full/}
%    \cite{barnett05052019737maxcrisis}
%    onderzoek van europese luchtvaart agentschap
%    \bibitem{easa27012021737maxsafereturn} ... \LaTeX:\\ \url{https://www.easa.europa.eu/newsroom-and-events/news/easa-declares-boeing-737-max-safe-return-service-europe}
%    \cite{easa27012021737maxsafereturn}
%    \veiligheidsvraagstuk
%    \bibitem{touitou11032019737tragedies} ... \LaTeX:\\ \url{https://phys.org/news/2019-03-boeing-max-safety-tragedies.html}
%    \cite{touitou11032019737tragedies}
%    artikel over sensoren
%    \bibitem{hemmerdinger02022021737maxdeliveries} ... \LaTeX:\\ \url{https://www.flightglobal.com/airframers/boeing-delays-737-max-10-deliveries-two-years-to-2023/142245.article}
%    \cite{hemmerdinger02022021737maxdeliveries}
%    goedkeuring van europese luchtvaart autoriteiten
%    advies aan de faa
%    \bibitem{bielby27022021faaimprovesafety} ... \LaTeX:\\ \url{https://www.hstoday.us/subject-matter-areas/airport-aviation-security/oig-tells-faa-to-improve-safety-oversight-following-boeing-737-max-review/}
%    \cite{bielby27022021faaimprovesafety}
%    \bibitem{boyle18112020737maxupgrade} ... \LaTeX:\\ \url{https://www.geekwire.com/2020/faas-go-ahead-737-maxs-return-flight-kicks-off-massive-software-upgrade/}
%    \cite{boyle18112020737maxupgrade}
%    \bibitem{bergstraburgess122019737maxMcasAlgorithm} ... \LaTeX:\\ \url{https://www.researchgate.net/publication/338420944_A_Promise_Theoretic_Account_of_the_Boeing_737_Max_MCAS_Algorithm_Affair}
%    \cite{bergstraburgess122019737maxMcasAlgorithm}
%    achtergrond informatie
%    \bibitem{737mcas} ... \LaTeX:\\ \url{http://www.b737.org.uk/mcas.htm}
%    \cite{737mcas}
%    algemeen vertrouwen
%    \bibitem{newburger17052019boeingcrisis} ... \LaTeX:\\ \url{https://www.cnbc.com/2019/05/16/what-you-need-to-know-about-boeings-737-max-crisis.html}
%    \cite{newburger17052019boeingcrisis}
%    toestemming europese autoriteiten
%    problemen
%    \bibitem{arstechnica22012020737problems} ... \LaTeX:\\ \url{https://arstechnica.com/information-technology/2020/01/737-max-fix-slips-to-summer-and-thats-just-one-of-boeings-problems/}
%    \cite{arstechnica22012020737problems}
%    uitgebreid artikel over de onderzoeken en het vliegverbod
%    \bibitem{german190620217372yaftergrounded} ... \LaTeX:\\ \url{https://www.cnet.com/news/boeing-737-max-8-all-about-the-aircraft-flight-ban-and-investigations/}
%    \cite{german190620217372yaftergrounded}
%    computers als oorzaak
%    lessons learned
%    \bibitem{beningo02052019boeinglessons} ... \LaTeX:\\ \url{https://www.designnews.com/electronics-test/5-lessons-learn-boeing-737-max-fiasco}
%    \cite{beningo02052019boeinglessons}
%    \bibitem{duran05042019boeingspof} ... \LaTeX:\\ \url{https://www.eurocontrol.int/publication/effects-network-extra-standby-aircraft-and-boeing-737-max-grounding}
%    \cite{duran05042019boeingspof}
%    single point of failure
%    \bibitem{ID} ... \LaTeX:\\ \url{https://dmd.solutions/blog/2019/04/05/how-a-single-point-of-failure-spof-in-the-mcas-software-could-have-caused-the-boeing-737-max-crash-in-ethiopia/}
%    \cite{}
%    \bibitem{makichuck24012021737fearflying} ... \LaTeX:\\ \url{https://asiatimes.com/2021/01/boeings-737-max-and-the-fear-of-flying/}
%    \cite{makichuck24012021737fearflying}
%    lijst van tehnische aanpassingen
%    \bibitem{caa737modifications} ... \LaTeX:\\ \url{https://www.caa.co.uk/Consumers/Guide-to-aviation/Boeing-737-MAX/}
%    \cite{caa737modifications}
%    \bibitem{oestergaard14122020boeingdeliveries} ... \LaTeX:\\ \url{https://dsm.forecastinternational.com/wordpress/2020/12/14/airbus-and-boeing-report-november-2020-commercial-aircraft-orders-and-deliveries/}
%    \cite{oestergaard14122020boeingdeliveries}
%    code lek
%    \bibitem{reenberg787flaws} ... \LaTeX:\\ \url{https://www.wired.com/story/boeing-787-code-leak-security-flaws/}
%    \cite{reenberg787flaws}
%    \bibitem{fitch16092020737backlogrisks} ... \LaTeX:\\ \url{https://www.fitchratings.com/research/corporate-finance/boeing-737-max-return-backlog-risks-remain-16-09-2020}
%    \cite{fitch16092020737backlogrisks}
%    Cultuurverandering, deregulatie, systeemwijziging of gewoon een kwestie van competentie
%    \bibitem{willis27082020737maxfailures} ... \LaTeX:\\ \url{https://www.aerospacetestinginternational.com/features/what-broke-the-737-max.html}
%    \cite{willis27082020737maxfailures}
%    extra aanpassingen
%    \bibitem{ostrower11062020more737changes} ... \LaTeX:\\ \url{https://theaircurrent.com/aviation-safety/boeings-737-max-software-done-but-regulators-plot-more-changes-after-jets-return/}
%    \cite{ostrower11062020more737changes}
%    wat ging er mis een analyse van een ex-iloot
%    De utoriteiten waren op de hoogte
%    \bibitem{hruska13122019faaknown737crashrate} ... \LaTeX:\\ \url{https://www.extremetech.com/extreme/303373-the-faa-knew-the-737-max-was-dangerous-and-kept-it-flying-anyway}
%    \cite{hruska13122019faaknown737crashrate}
%    kwaliteiten van het alarmsysteem niet goed bekend
%    \bibitem{bloomberg26092019failedpred} ... \LaTeX:\\ \url{https://time.com/5687473/boeing-737-alarm-system/}
%    \cite{bloomberg26092019failedpred}
%    \bibitem{whiteman09072020boengcancelstock} ... \LaTeX:\\ \url{https://www.nasdaq.com/articles/boeing-gets-dealt-another-737-max-cancellation-blow.-what-it-means-for-boeing-stock-2020}
%    \cite{whiteman09072020boengcancelstock}
%    \bibitem{leopold09192019boeingreliability} ... \LaTeX:\\ \url{https://www.eetimes.com/boeing-crashes-highlight-a-worsening-reliability-crisis/}
%    \cite{leopold09192019boeingreliability}
%    veiligheidsvraagstuk
%    \bibitem{koenig11122019737crashesnofix} ... \LaTeX:\\ \url{https://www.latimes.com/business/story/2019-12-11/faa-boeing-737-max-crashes}
%    \cite{koenig11122019737crashesnofix}
%    probleemanalyse, veiligheidsvraagstuk
%    \bibitem{dohertylindeman15032019737problems} ... \LaTeX:\\ \url{https://www.politico.com/story/2019/03/15/boeing-737-max-grounding-1223072}
%    \cite{dohertylindeman15032019737problems}
%    falend toezicht
%    \bibitem{stodder02102019corruptoversight} ... \LaTeX:\\ \url{https://www.pogo.org/analysis/2019/10/corrupted-oversight-the-faa-boeing-and-the-737-max/}
%    \cite{stodder02102019corruptoversight}
%    \bibitem{afacwaLostSafeguards} ... \LaTeX:\\ \url{https://www.afacwa.org/the_inside_story_of_mcas_seattle_times}
%    \cite{afacwaLostSafeguards}
%    doelstellingen en veiligheidsvraagstukken
%    \bibitem{swayne18032019profitssafety} ... \LaTeX:\\ \url{https://www.marxist.com/737-max-scandal-boeing-putting-profits-before-safety.htm}
%    \cite{swayne18032019profitssafety}
%    \bibitem{freed26022021liftaustraliaban} ... \LaTeX:\\ \url{https://finance.yahoo.com/news/australia-lifts-ban-boeing-737-035817682.html?guccounter=1&guce_referrer=aHR0cHM6Ly93d3cuZ29vZ2xlLmNvbS8&guce_referrer_sig=AQAAAHZCJYy_0A5VS2WiPoCvH4xdrRNkmkdsv5EWJ2RLIz_AS-rxsTty6AF1_HlmJiRyWYqCXDi4p0Xs4isYkNkCq2Pfo-pQ60Xz_IfTNjm4FgoZiBMC4zpZlB6F0fwecrjE_ujAXZzG4xPJnWCd8-G3VLlPTY8h3H31eQ1i8hY9AIyy}
%    \cite{freed26022021liftaustraliaban}
%    autoriteiten krijgen tik op de vingers
%    \bibitem{reed15032019softwareattention} ... \LaTeX:\\ \url{https://medium.com/@jpaulreed/the-737max-and-why-software-engineers-should-pay-attention-a041290994bd}
%    \cite{reed15032019softwareattention}
%    \bibitem{news17032019softwareexplains} ... \LaTeX:\\ \url{https://news.ycombinator.com/item?id=19414775}
%    \cite{news17032019softwareexplains}
%    \bibitem{legget21122020eu737maxsafe} ... \LaTeX:\\ \url{https://www.bbc.com/news/55366320}
%    \cite{legget21122020eu737maxsafe}
%    \bibitem{marketscreener0103221737chinarecertification} ... \LaTeX:\\ \url{https://www.marketscreener.com/news/latest/China-studies-Boeing-737-MAX-recertification-wants-safety-concerns-fully-addressed--32569125/}
%    \cite{marketscreener0103221737chinarecertification}
%    motor in brand
%    \bibitem{euractiv22022021737firegrounds} ... \LaTeX:\\ \url{https://www.euractiv.com/section/aviation/news/boeing-grounds-777s-after-engine-fire/}
%    \cite{euractiv22022021737firegrounds}
%    \bibitem{benny18022019737returnUAE} ... \LaTeX:\\ \url{https://gulfnews.com/business/aviation/uae-airspace-to-see-return-of-boeing-737-max-1.1613627548923}
%    \cite{benny18022019737returnUAE}
%    motor in brand gevlogen
%    \bibitem{biersmichel22022021777grounds} ... \LaTeX:\\ \url{https://techxplore.com/news/2021-02-boeing-urges-grounding-777s.html}
%    \cite{biersmichel22022021777grounds}
%    \bibitem{ID} ... \LaTeX:\\ \url{https://www.politico.eu/article/uk-temporarily-bans-some-boeing-aircraft-after-pratt-whitney-engine-incidents/}
%    \cite{}
%    \bibitem{reuters23022021777metalfatigue} ... \LaTeX:\\ \url{https://www.timeslive.co.za/news/world/2021-02-23-damage-to-united-boeing-777-engine-consistent-with-metal-fatigue--ntsb/}
%    \cite{reuters23022021777metalfatigue}
%    faa was niet kritisch genoeg
%    \bibitem{ID} ... \LaTeX:\\ \url{https://federalnewsnetwork.com/government-news/2021/02/federal-watchdog-blasts-faa-over-certification-of-boeing-jet/}
%    \cite{}
%    
%    
%    
%    algemene vragen
%    oorzaken
%    \bibitem{{gates18112020boeingcrisis} ... \LaTeX:\\ \url{https://www.seattletimes.com/business/boeing-aerospace/what-led-to-boeings-737-max-crisis-a-qa/}
%    	
%    	\bibitem{{boeing737maxsoftwareprobles} ... \LaTeX:\\ \url{https://www.schneier.com/blog/archives/2019/04/excellent_analy.html}
%    		fout in de software
%    		\bibitem{{avetisov19032019boeingmalwarestate} ... \LaTeX:\\ \url{https://www.forbes.com/sites/georgeavetisov/2019/03/19/malware-at-30000-feet-what-the-737-max-says-about-the-state-of-airplane-software-security/?sh=4d26f7052a9e}
%    			het nationaal veiligheidsbelang
%    			\bibitem{{thompson23112020nationalsecurityboeing} ... \LaTeX:\\ \url{https://www.forbes.com/sites/lorenthompson/2020/11/23/five-reasons-return-of-boeings-737-max-to-service-is-important-to-national-security/?sh=2128ea552018}
%    				falend toezicht
%    				\bibitem{{gates21032019FAAControlSystem} ... \LaTeX:\\ \url{https://www.seattletimes.com/business/boeing-aerospace/failed-certification-faa-missed-safety-issues-in-the-737-max-system-implicated-in-the-lion-air-crash/}
%    					onderzoeksrapport
%    					\bibitem{{faa18112020boeingreview} ... \LaTeX:\\ \url{https://www.faa.gov/foia/electronic_reading_room/boeing_reading_room/media/737_RTS_Summary.pdf}
%    						\bibitem{{wiki737maxgroundings} ... \LaTeX:\\ \url{https://en.wikipedia.org/wiki/Boeing_737_MAX_groundings}
%    							veiligheidsrisico's
%    							menselijke fouten
%    							\bibitem{{campbell02052019boengcrashhumanerrors} ... \LaTeX:\\ \url{https://www.theverge.com/2019/5/2/18518176/boeing-737-max-crash-problems-human-error-mcas-faa}
%    								overzicht van crashes
%    								\bibitem{{hawkins22032019737maxairplanes} ... \LaTeX:\\ \url{https://www.theverge.com/2019/3/22/18275736/boeing-737-max-plane-crashes-grounded-problems-info-details-explained-reasons}
%    									veiligheidsopmerking
%    									\bibitem{{thomas30082020737safest} ... \LaTeX:\\ \url{https://www.airlineratings.com/news/boeings-737-max-will-one-safest-aircraft-history/}
%    										aanpassingen
%    										\bibitem{{boeing737maxdisplay} ... \LaTeX:\\ \url{https://www.boeing.com/commercial/737max/737-max-software-updates.page}
%    											waarschuwingen//output signalen
%    											\bibitem{{fehrm24112020737changes} ... \LaTeX:\\ \url{https://leehamnews.com/2020/11/24/boeing-737-max-changes-beyond-mcas/}
%    												software gerelateerde fouten
%    												\bibitem{{travis18042019737maxsoftwaredevop} ... \LaTeX:\\ \url{https://spectrum.ieee.org/aerospace/aviation/how-the-boeing-737-max-disaster-looks-to-a-software-developer}
%    													onderzoeksrapport
%    													de rol van de publieke opinie
%    													\bibitem{{barnett05052019737maxcrisis} ... \LaTeX:\\ \url{https://pubsonline.informs.org/do/10.1287/orms.2019.05.05/full/}
%    														onderzoek van europese luchtvaart agentschap
%    														\bibitem{{easa27012021737maxsafereturn} ... \LaTeX:\\ \url{https://www.easa.europa.eu/newsroom-and-events/news/easa-declares-boeing-737-max-safe-return-service-europe}
%    															veiligheidsvraagstuk
%    															\bibitem{{touitou11032019737tragedies} ... \LaTeX:\\ \url{https://phys.org/news/2019-03-boeing-max-safety-tragedies.html}
%    																artikel over sensoren
%    																\bibitem{{hemmerdinger02022021737maxdeliveries} ... \LaTeX:\\ \url{https://www.flightglobal.com/airframers/boeing-delays-737-max-10-deliveries-two-years-to-2023/142245.article}
%    																	goedkeuring van europese luchtvaart autoriteiten
%    																	advies aan de faa
%    																	\bibitem{{bielby27022021faaimprovesafety} ... \LaTeX:\\ \url{https://www.hstoday.us/subject-matter-areas/airport-aviation-security/oig-tells-faa-to-improve-safety-oversight-following-boeing-737-max-review/}
%    																		\bibitem{{boyle18112020737maxupgrade} ... \LaTeX:\\ \url{https://www.geekwire.com/2020/faas-go-ahead-737-maxs-return-flight-kicks-off-massive-software-upgrade/}
%    																			\bibitem{{bergstraburgess122019737maxMcasAlgorithm} ... \LaTeX:\\ \url{https://www.researchgate.net/publication/338420944_A_Promise_Theoretic_Account_of_the_Boeing_737_Max_MCAS_Algorithm_Affair}
%    																				achtergrond informatie
%    																				\bibitem{{737mcas} ... \LaTeX:\\ \url{http://www.b737.org.uk/mcas.htm}
%    																					algemeen vertrouwen
%    																					\bibitem{{newburger17052019boeingcrisis} ... \LaTeX:\\ \url{https://www.cnbc.com/2019/05/16/what-you-need-to-know-about-boeings-737-max-crisis.html}
%    																						toestemming europese autoriteiten
%    																						problemen
%    																						\bibitem{{arstechnica22012020737problems} ... \LaTeX:\\ \url{https://arstechnica.com/information-technology/2020/01/737-max-fix-slips-to-summer-and-thats-just-one-of-boeings-problems/}
%    																							uitgebreid artikel over de onderzoeken en het vliegverbod
%    																							\bibitem{{german190620217372yaftergrounded} ... \LaTeX:\\ \url{https://www.cnet.com/news/boeing-737-max-8-all-about-the-aircraft-flight-ban-and-investigations/}
%    																								computers als oorzaak
%    																								lessons learned
%    																								\bibitem{{beningo02052019boeinglessons} ... \LaTeX:\\ \url{https://www.designnews.com/electronics-test/5-lessons-learn-boeing-737-max-fiasco}
%    																									\bibitem{{duran05042019boeingspof} ... \LaTeX:\\ \url{https://www.eurocontrol.int/publication/effects-network-extra-standby-aircraft-and-boeing-737-max-grounding}
%    																										single point of failure
%    																										\bibitem{{ID} ... \LaTeX:\\ \url{https://dmd.solutions/blog/2019/04/05/how-a-single-point-of-failure-spof-in-the-mcas-software-could-have-caused-the-boeing-737-max-crash-in-ethiopia/}
%    																											\bibitem{{makichuck24012021737fearflying} ... \LaTeX:\\ \url{https://asiatimes.com/2021/01/boeings-737-max-and-the-fear-of-flying/}
%    																												lijst van tehnische aanpassingen
%    																												\bibitem{{caa737modifications} ... \LaTeX:\\ \url{https://www.caa.co.uk/Consumers/Guide-to-aviation/Boeing-737-MAX/}
%    																													\bibitem{{oestergaard14122020boeingdeliveries} ... \LaTeX:\\ \url{https://dsm.forecastinternational.com/wordpress/2020/12/14/airbus-and-boeing-report-november-2020-commercial-aircraft-orders-and-deliveries/}
%    																														code lek
%    																														\bibitem{{reenberg787flaws} ... \LaTeX:\\ \url{https://www.wired.com/story/boeing-787-code-leak-security-flaws/}
%    																															\bibitem{{fitch16092020737backlogrisks} ... \LaTeX:\\ \url{https://www.fitchratings.com/research/corporate-finance/boeing-737-max-return-backlog-risks-remain-16-09-2020}
%    																																Cultuurverandering, deregulatie, systeemwijziging of gewoon een kwestie van competentie
%    																																\bibitem{{willis27082020737maxfailures} ... \LaTeX:\\ \url{https://www.aerospacetestinginternational.com/features/what-broke-the-737-max.html}
%    																																	extra aanpassingen
%    																																	\bibitem{{ostrower11062020more737changes} ... \LaTeX:\\ \url{https://theaircurrent.com/aviation-safety/boeings-737-max-software-done-but-regulators-plot-more-changes-after-jets-return/}
%    																																		wat ging er mis een analyse van een ex-iloot
%    																																		De utoriteiten waren op de hoogte
%    																																		\bibitem{{hruska13122019faaknown737crashrate} ... \LaTeX:\\ \url{https://www.extremetech.com/extreme/303373-the-faa-knew-the-737-max-was-dangerous-and-kept-it-flying-anyway}
%    																																			kwaliteiten van het alarmsysteem niet goed bekend
%    																																			\bibitem{{bloomberg26092019failedpred} ... \LaTeX:\\ \url{https://time.com/5687473/boeing-737-alarm-system/}
%    																																				\bibitem{{whiteman09072020boengcancelstock} ... \LaTeX:\\ \url{https://www.nasdaq.com/articles/boeing-gets-dealt-another-737-max-cancellation-blow.-what-it-means-for-boeing-stock-2020}
%    																																					\bibitem{{leopold09192019boeingreliability} ... \LaTeX:\\ \url{https://www.eetimes.com/boeing-crashes-highlight-a-worsening-reliability-crisis/}
%    																																						veiligheidsvraagstuk
%    																																						\bibitem{{koenig11122019737crashesnofix} ... \LaTeX:\\ \url{https://www.latimes.com/business/story/2019-12-11/faa-boeing-737-max-crashes}
%    																																							probleemanalyse, veiligheidsvraagstuk
%    																																							\bibitem{{dohertylindeman15032019737problems} ... \LaTeX:\\ \url{https://www.politico.com/story/2019/03/15/boeing-737-max-grounding-1223072}
%    																																								falend toezicht
%    																																								\bibitem{{stodder02102019corruptoversight} ... \LaTeX:\\ \url{https://www.pogo.org/analysis/2019/10/corrupted-oversight-the-faa-boeing-and-the-737-max/}
%    																																									\bibitem{{afacwaLostSafeguards} ... \LaTeX:\\ \url{https://www.afacwa.org/the_inside_story_of_mcas_seattle_times}
%    																																										doelstellingen en veiligheidsvraagstukken
%    																																										\bibitem{{swayne18032019profitssafety} ... \LaTeX:\\ \url{https://www.marxist.com/737-max-scandal-boeing-putting-profits-before-safety.htm}
%    																																											\bibitem{{freed26022021liftaustraliaban} ... \LaTeX:\\ \url{https://finance.yahoo.com/news/australia-lifts-ban-boeing-737-035817682.html?guccounter=1&guce_referrer=aHR0cHM6Ly93d3cuZ29vZ2xlLmNvbS8&guce_referrer_sig=AQAAAHZCJYy_0A5VS2WiPoCvH4xdrRNkmkdsv5EWJ2RLIz_AS-rxsTty6AF1_HlmJiRyWYqCXDi4p0Xs4isYkNkCq2Pfo-pQ60Xz_IfTNjm4FgoZiBMC4zpZlB6F0fwecrjE_ujAXZzG4xPJnWCd8-G3VLlPTY8h3H31eQ1i8hY9AIyy}
%    																																												autoriteiten krijgen tik op de vingers
%    																																												\bibitem{{reed15032019softwareattention} ... \LaTeX:\\ \url{https://medium.com/@jpaulreed/the-737max-and-why-software-engineers-should-pay-attention-a041290994bd}
%    																																													\bibitem{{news17032019softwareexplains} ... \LaTeX:\\ \url{https://news.ycombinator.com/item?id=19414775}
%    																																														\bibitem{{legget21122020eu737maxsafe} ... \LaTeX:\\ \url{https://www.bbc.com/news/55366320}
%    																																															\bibitem{{marketscreener0103221737chinarecertification} ... \LaTeX:\\ \url{https://www.marketscreener.com/news/latest/China-studies-Boeing-737-MAX-recertification-wants-safety-concerns-fully-addressed--32569125/}
%    																																																motor in brand
%    																																																\bibitem{{euractiv22022021737firegrounds} ... \LaTeX:\\ \url{https://www.euractiv.com/section/aviation/news/boeing-grounds-777s-after-engine-fire/}
%    																																																	\bibitem{{benny18022019737returnUAE} ... \LaTeX:\\ \url{https://gulfnews.com/business/aviation/uae-airspace-to-see-return-of-boeing-737-max-1.1613627548923}
%    																																																		motor in brand gevlogen
%    																																																		\bibitem{{biersmichel22022021777grounds} ... \LaTeX:\\ \url{https://techxplore.com/news/2021-02-boeing-urges-grounding-777s.html}
%    																																																			\bibitem{{ID} ... \LaTeX:\\ \url{https://www.politico.eu/article/uk-temporarily-bans-some-boeing-aircraft-after-pratt-whitney-engine-incidents/}
%    																																																				\bibitem{{reuters23022021777metalfatigue} ... \LaTeX:\\ \url{https://www.timeslive.co.za/news/world/2021-02-23-damage-to-united-boeing-777-engine-consistent-with-metal-fatigue--ntsb/}
%    																																																					faa was niet kritisch genoeg
%    																																																					\bibitem{{ID} ... \LaTeX:\\ \url{https://federalnewsnetwork.com/government-news/2021/02/federal-watchdog-blasts-faa-over-certification-of-boeing-jet/}
%    																																																						\cite{}
    																																																						
  
\end{thebibliography}

%
%\newpage
%\bibliography{references}
%

\vfill
}}   %(verplicht) bronvermeldingen
\else
  \iflanguage{dutch}{\def\bibname{Bronnen}}{\def\bibname{References}}
  
\begin{thebibliography}{99}

\bibitem{lam1994} Lamport L.: \emph{\LaTeX: A Document Preparation System}, Addison-Wesley, 1994 

\bibitem{Oos1996} Oostrum van P.: \emph{Handleiding \LaTeX}, Vakgroep
  Informatica, Universiteit Utrecht, 1998,\\ 
  \url{http://people.cs.uu.nl/piet/latexhnd.pdf}

\bibitem{wikibooks} Wikibooks \LaTeX:\\
  \url{http://nl.wikibooks.org/wiki/LaTeX}


\bibitem{oid} Wikibooks \LaTeX:\\
\url{https://www.waterkant.net/suriname/2023/05/29/milieuactivist-sleur-zeer-grote-onwaarheden-staan-in-cyanide-onderzoeksrapport/}

   


%%%%%%%%%%%%%%%%%%%%%%%%%%%%%%%%%%%%%%%%%%%%%%%%%%%%%%%%%%%%%%%%%
%
%
%misc{para,
%	Author = {{IEEE Referencing}},
%	Howpublished = {\url{http://libguides.bhtafe.edu.au/content.php?pid=88814&sid=660920}},
%	Keywords = {paraphrasing, plagariasm, citation, quotation},
%	Lastchecked = {19/04/2017},
%	Title = {{How to Cite\/Quote in Your Assignment}}}
%
%misc{ieee-write,
%	Author = {{IEEE Authorship Series}},
%	Howpublished = {\url{http://ieeeauthorcenter.ieee.org/wp-content/uploads/How-to-Write-for-Technical-Periodicals-and-Conferences.pdf}},
%	Institution = {{IEEE}},
%	Keywords = {writing, paragraph, sentences, ethics},
%	Lastchecked = {18/04/2017},
%	Title = {How to Write for Technical Periodicals \& Conferences}}
%
%
%
%@article{schon2017agile,
%	title={Agile Requirements Engineering: A systematic literature review},
%	author={Sch\"on, Eva-Maria and Thomaschewski, J\"org and Escalona, Maria Jose},
%	journal={Computer Standards \& Interfaces},
%	volume={49},
%	pages={79--91},
%	year={2017},
%	publisher={Elsevier}
%}
%
%@inproceedings{royce1987managing,
%	title={Managing the development of large software systems: concepts and techniques},
%	author={Royce, Winston W},
%	booktitle={Proceedings of the 9th international conference on Software Engineering},
%	pages={328--338},
%	year=1987,
%	organization={IEEE Computer Society Press}
%}
%
%@article{leveson1993investigation,
%	title={An investigation of the Therac-25 accidents},
%	author={Leveson, Nancy G and Turner, Clark S},
%	journal={IEEE computer},
%	volume={26},
%	number={7},
%	pages={18--41},
%	year={1993}
%}
%@book{modelchecking,
%	author = {Clarke,Jr., Edmund M. and Grumberg, Orna and Peled, Doron A.},
%	title = {Model Checking},
%	year = {1999},
%	isbn = {0-262-03270-8},
%	publisher = {MIT Press},
%	address = {Cambridge, MA, USA},
%}
%@inproceedings{nuseibeh2000requirements,
%	title={Requirements engineering: a roadmap},
%	author={Nuseibeh, Bashar and Easterbrook, Steve},
%	booktitle={Proceedings of the Conference on the Future of Software Engineering},
%	pages={35--46},
%	year={2000},
%	organization={ACM}
%}
%
%
%%%%%%%%%%%%%%%%%%%%%%%%%%%%%%%%%%%%%%%%%%%%%%%%%%%%%%%%%%%%%%%%%%
%
%@online{inriaStatsMoodCheck,
%	ALTauthor = {project.inria},
%	ALTeditor = {editor},
%	title = {statistical-model-checking},
%	date = {date},
%	url = {https://project.inria.fr/plasma-lab/statistical-model-checking/},
%}
%
%
%@online{buddeModelChecker,
%	ALTauthor = {Carlos E. Budde1, Pedro R. D’Argenio2,3,4, Arnd Hartmanns1(B) , and Sean Sedwards5},
%	ALTeditor = {editor},
%	title = {A Statistical Model Checker for Nondeterminism and Rare Events},
%	date = {date},
%	url = {https://ris.utwente.nl/ws/portalfiles/portal/28200786/A_statistical_model_checker.pdf},
%}
%
%
%@online{ID,
%	ALTauthor = {GUL AGHA and  },
%	ALTeditor = {KARL PALMSKOG},
%	title = {A Survey of Statistical Model Checking
%	},
%	date = {date},
%	url = {https://dl.acm.org/doi/10.1145/3158668},
%}
%
%
%
%@online{AGHASuervey,
%	ALTauthor = {GUL AGHA and KARL PALMSKOG,},
%	ALTeditor = {editor},
%	title = {A Survey of Statistical Model Checking},
%	date = {date},
%	url = {https://dl.acm.org/doi/pdf/10.1145/3158668},
%}
%%%%%%%%%%%%%%%%%%%%%%%%%%%%%%%%%%%%%%%%%%%%%%%%%%%%%%%%%%%%%%%%%%
%
%
%
%@online{bicker21102016automatiseringsparadox,
%	ALTauthor = {author},
%	ALTeditor = {editor},
%	title = {title},
%	date = {date},
%	url = {https://www.debicker.eu/de-automatiseringsparadox/},
%}
%@online{vseautoparadox,
%	ALTauthor = {author},
%	ALTeditor = {editor},
%	title = {title},
%	date = {date},
%	url = {https://vse.nl/de-paradox-van-de-industriele-automatisering/},
%}
%
%
%
%@online{blogxot21112016slimapparaat,
%	ALTauthor = {author},
%	ALTeditor = {editor},
%	title = {title},
%	date = {date},
%	url = {https://blog.xot.nl/2016/11/21/slimme-apparaten-maken-ons-dom-en-kwetsbaar/index.html},
%}
%
%@online{ID,
%	ALTauthor = {author},
%	ALTeditor = {editor},
%	title = {title},
%	date = {date},
%	url = {https://automatie-pma.com/nieuws/industriele-automatiseringsparadox},
%}
%%%%%%%%%%%%%%%%%%%%%%%%%%%%%%%%%%%%%%%%%%%%%%%%%%%%%%%%%%%%%%%%%%
%
%
%
%
%
%
%
%
%
%@online{aviationsafety04101992airplaneCrashBijlmer,
%	ALTauthor = {author},
%	ALTeditor = {editor},
%	title = {title},
%	date = {date},
%	url = {https://aviation-safety.net/database/record.php?id=19921004-2&lang=nl  },
%}
%
%
%
%@online{boogers092002RampenRegelsRichtlijnen,
%	ALTauthor = {author},
%	ALTeditor = {editor},
%	title = {title},
%	date = {date},
%	url = { https://www.researchgate.net/publication/254815008_Rampen_regels_richtlijnen},
%}
%
%
%@online{catsr25022009Boeing737AmsterdamCrash,
%	ALTauthor = {author},
%	ALTeditor = {editor},
%	title = {title},
%	date = {date},
%	url = {https://catsr.vse.gmu.edu/SYST460/TA1951_AccidentReport.pdf  },
%}
%
%@online{INSAVienna1992Chernobyl,
%	ALTauthor = {author},
%	ALTeditor = {editor},
%	title = {title},
%	date = {date},
%	url = {https://www-pub.iaea.org/MTCD/publications/PDF/Pub913e_web.pdf  },
%}
%
%
%
%
%aviationReport
%
%@online{oVVSchietongevalOssendrecht,
%	ALTauthor = {author},
%	ALTeditor = {editor},
%	title = {title},
%	date = {date},
%	url = {https://www.youtube.com/watch?v=6jmkDClGDHo  },
%}
%
%@online{molukseTreinkaping,
%	ALTauthor = {author},
%	ALTeditor = {editor},
%	title = {title},
%	date = {date},
%	url = {https://www.youtube.com/watch?v=h99Fe9XzzHI  },
%}
%
%@online{jiang16042019TanjinExplosion,
%	ALTauthor = {author},
%	ALTeditor = {editor},
%	title = {title},
%	date = {date},
%	url = {https://www.hindawi.com/journals/joph/2019/1360805/  },
%}
%
%
%@online{hrw03082021investigateBeirutBlast,
%	ALTauthor = {author},
%	ALTeditor = {editor},
%	title = {title},
%	date = {date},
%	url = { 	https://www.hrw.org/report/2021/08/03/they-killed-us-inside/investigation-august-4-beirut-blast },
%}
%
%@online{souaibyElHussein112020Beirutstory,
%	ALTauthor = {author},
%	ALTeditor = {editor},
%	title = {title},
%	date = {date},
%	url = { https://www.researchgate.net/publication/348325979_Beirut_Explosion_the_full_story },
%}
%
%@online{ifrc2020chemicalexplosionBeirutPort,
%	ALTauthor = {author},
%	ALTeditor = {editor},
%	title = {title},
%	date = {date},
%	url = { https://reliefweb.int/sites/reliefweb.int/files/resources/CaseStudy_BeirutExplosion_TechBioHazardsweb.pdf },
%}
%
%@online{caliskan09112013747boeingkalman,
%	ALTauthor = {author},
%	ALTeditor = {editor},
%	title = {title},
%	date = {date},
%	url = {https://www.hindawi.com/journals/ijae/2014/472395/  },
%}
%
%%%%%%%%%%%%%%%%%%%%%%%%%%%%%%%%%%%%%%%%%%%%%%%%%%%%%%%%%%%%%%%%%%
%persuasive technology 
%
%@online{humanTechpersuasiveTech,
%	ALTauthor = {author},
%	ALTeditor = {editor},
%	title = {title},
%	date = {date},
%	url = {https://www.humanetech.com/youth/persuasive-technology  },
%}
%
%
%
%
%@online{rezenfeld01012018persuasiveTecgHabits,
%	ALTauthor = {author},
%	ALTeditor = {editor},
%	title = {title},
%	date = {date},
%	url = {https://spectrum.ieee.org/how-persuasive-technology-can-change-your-habits  },
%}
%
%
%@online{aldenaini28042020persuasiveTechTrends,
%	ALTauthor = {author},
%	ALTeditor = {editor},
%	title = {title},
%	date = {date},
%	url = { https://www.frontiersin.org/articles/10.3389/frai.2020.00007/full },
%}
%
%@online{larson14062017persuasivetechmanipulates,
%	ALTauthor = {author},
%	ALTeditor = {editor},
%	title = {title},
%	date = {date},
%	url = {https://psmag.com/environment/captology-fogg-invisible-manipulative-power-persuasive-technology-81301  },
%}
%
%@online{tanzem22012022persuasivetechchanginglives,
%	ALTauthor = {author},
%	ALTeditor = {editor},
%	title = {title},
%	date = {date},
%	url = {https://www.makeuseof.com/what-is-persuasive-technology/  },
%}
%
%
%
%
%@online{tikkakuddonenpersuasiveTechnology,
%	ALTauthor = {author},
%	ALTeditor = {editor},
%	title = {title},
%	date = {date},
%	url = { https://cyberpsychology.eu/article/view/12270 },
%}
%
%@online{sprongken19032018CourtProcedureDigital,
%	ALTauthor = {author},
%	ALTeditor = {editor},
%	title = {title},
%	date = {date},
%	url = { https://www.njb.nl/blogs/a-court-with-no-face-and-no-place/ },
%}
%
%@online{PROCESREGLEMENTEcourt,
%	ALTauthor = {author},
%	ALTeditor = {editor},
%	title = {title},
%	date = {date},
%	url = { http://www.e-court.nl/wp-content/uploads/2018/03/Procesreglement-e-Court-2017_20180201.pdf},
%}
%
%
%@online{ovvMortierOngevalMaliVideo,
%	ALTauthor = {author},
%	ALTeditor = {editor},
%	title = {title},
%	date = {date},
%	url = {https://www.youtube.com/watch?v=PC2ekl4SaNA  },
%}
%
%
%
%
%@online{vtmGroep29032019waterwerkencyber,
%	ALTauthor = {author},
%	ALTeditor = {editor},
%	title = {title},
%	date = {date},
%	url = { https://www.vtmgroep.nl/blog/waterwerken-in-nederland-onvoldoende-beveiligd-tegen-cyberaanvallen},
%}
%
%
%
%
%
%@online{wesemannVeiligheidLandEnWater,
%	ALTauthor = {author},
%	ALTeditor = {editor},
%	title = {title},
%	date = {date},
%	url = {
%		http://www.wesemann.nl/nl/nieuws-en-pers/274-veiligheid-op-het-water-en-op-het-land.html },
%}
%
%
%
%

%%%%%%%%%%%%%%%%%%%%%%%%%%%%%%%%%%%%%%%%%%%%%%%%%%%%%%%%%%%%%%%%%
\bibitem{ISAC_SANS_Ukraine_DUC_18Mar2016} ... \LaTeX:\\ \url{https://www.nerc.com/_layouts/15/Nerc.404/CustomFileNotFound.aspx?requestUrl=https://www.nerc.com/pa/CI/ESISAC/Documents/E-ISAC_SANS_Ukraine_DUC_18Mar2016.pdf},}
\bibitem{zetter2016GridHack,	ALTauthor = {Kim Zetter},	ALTeditor = {editor},	title = {Inside the Cunning, Unprecedented Hack of Ukraine's Power Grid},	date = {date},	url = {https://www.wired.com/2016/03/inside-cunning-unprecedented-hack-ukraines-power-grid/},}
\bibitem{2015ukrainegridattack,	ALTauthor = {wikipedia},	ALTeditor = {editor},	title = {2015 Ukraine power grid hack},	date = {date},	url = {https://en.wikipedia.org/wiki/2015_Ukraine_power_grid_hack},}
\bibitem{greenberg2017Cyberwartestlab} ... \LaTeX:\\ \url{https://www.wired.com/story/russian-hackers-attack-ukraine/},}
\bibitem{boozallen2016lightwentout,	ALTauthor = {allen},	ALTeditor = {editor},	title = {Ukrain report when the lights went out},	date = {date},	url = {https://www.boozallen.com/content/dam/boozallen/documents/2016/09/ukraine-report-when-the-lights-went-out.pdf},}
\bibitem{finklejan2016UsBlamesRussianSandworm} ... \LaTeX:\\ \url{https://www.reuters.com/article/us-ukraine-cybersecurity-sandworm-idUSKBN0UM00N20160108},}
\bibitem{wired012016ukrainegrid} ... \LaTeX:\\ \url{https://www.wired.com/2016/01/everything-we-know-about-ukraines-power-plant-hack/},}
\bibitem{icsalert20072021cyberattackukraine} ... \LaTeX:\\ \url{https://www.us-cert.gov/ics/alerts/IR-ALERT-H-16-056-01},}
\bibitem{finkle08012016russiansandwormhackers} ... \LaTeX:\\ \url{https://www.reuters.com/article/us-ukraine-cybersecurity-sandworm/u-s-firm-blames-russian-sandworm-hackers-for-ukraine-outage-idUSKBN0UM00N20160108},}
\bibitem{zinets15022017ukrainechargesrussia} ... \LaTeX:\\ \url{https://www.reuters.com/article/us-ukraine-crisis-cyber-idUSKBN15U2CN},}
\bibitem{2014russiansandworm} ... \LaTeX:\\ \url{https://www.wired.com/2014/10/russian-sandworm-hack-isight/},}
\bibitem{mcelfresh2016cyberattackhowandwhy} ... \LaTeX:\\ \url{https://theconversation.com/cyberattack-on-ukraine-grid-heres-how-it-worked-and-perhaps-why-it-was-done-52802},}
\bibitem{desarnaud2017cyberattacks} ... \LaTeX:\\ \url{https://www.ifri.org/sites/default/files/atoms/files/desarnaud_cyber_attacks_energy_infrastructures_2017_2.pdf},}
\bibitem{osti2018historycontrolsystemIncidents} ... \LaTeX:\\ \url{https://www.osti.gov/biblio/1505628},}
\bibitem{Shahzad2014ScadaProtocolsPollingScenario} ... \LaTeX:\\ \url{https://scialert.net/fulltext/?doi=tasr.2014.396.405},}
\bibitem{caseli04112016intrusiondetectioncontrolsystem} ... \LaTeX:\\ \url{https://ris.utwente.nl/ws/files/6028066/3-s2_0-B9780128015957000227.pdf},}
\bibitem{rochascadatesting} ... \LaTeX:\\ \url{https://repositorio-aberto.up.pt/bitstream/10216/119066/2/315683.pdf},}
\bibitem{levalle2020FuzzingICSProtocols} ... \LaTeX:\\ \url{https://dreamlab.net/en/blog/post/fuzzing-ics-protocols/},}
\bibitem{resch31102019IEC62351secureCommunication} ... \LaTeX:\\ \url{http://www.connectivity4ir.co.uk/article/175490/IEC-62351--Secure-communication-in-the-energy-industry.aspx},}
\bibitem{slowik2019ReassasUkraine2016Attack} ... \LaTeX:\\ \url{https://www.dragos.com/wp-content/uploads/CRASHOVERRIDE.pdf},}
\bibitem{arrizabalaga2020surveyiiotProtocols} ... \LaTeX:\\ \url{https://dl.acm.org/doi/fullHtml/10.1145/3381038},}
\bibitem{yadav2020reviewScadaArchitecture} ... \LaTeX:\\ \url{https://arxiv.org/pdf/2001.02925.pdf},}
\bibitem{cis20072021crashoverridemalware} ... \LaTeX:\\ \url{https://www.us-cert.gov/ncas/alerts/TA17-163A},}
\bibitem{holappa2017threattoElectricityNetworks} ... \LaTeX:\\ \url{https://www.nixu.com/fi/node/53},}
\bibitem{shehod2016gridadvantageus} ... \LaTeX:\\ \url{http://web.mit.edu/smadnick/www/wp/2016-22.pdf},}
\bibitem{parkwalstorm11102017russiagridattack} ... \LaTeX:\\ \url{https://jsis.washington.edu/news/cyberattack-critical-infrastructure-russia-ukrainian-power-grid-attacks/},}
\bibitem{drago2017CrashOverride} ... \LaTeX:\\ \url{https://www.dragos.com/wp-content/uploads/CrashOverride-01.pdf},}
\bibitem{wikiindustroyer} ... \LaTeX:\\ \url{https://en.wikipedia.org/wiki/Industroyer},}
\bibitem{crashoverridenetwork} ... \LaTeX:\\ \url{https://en.wikipedia.org/wiki/Crash_Override_Network},}
\bibitem{slowikvb2018crashoverride} ... \LaTeX:\\ \url{https://www.virusbulletin.com/virusbulletin/2019/03/vb2018-paper-anatomy-attack-detecting-and-defeating-crashoverride/},}
\bibitem{njccicthreat08102017crashovverrideprofile} ... \LaTeX:\\ \url{https://www.cyber.nj.gov/threat-center/threat-profiles/ics-malware-variants/crashoverride},}
\bibitem{crashoverrideindustroyermalware} ... \LaTeX:\\ \url{https://www.webopedia.com/TERM/C/crashoverride-industroyer-malware.html},}
\bibitem{incibe23062017crashoverrideback} ... \LaTeX:\\ \url{https://www.incibe-cert.es/en/blog/crashoverride-malware-ics-back-again},}
\bibitem{industroyershortfact} ... \LaTeX:\\ \url{https://rhebo.com/en/service/glossar/industroyer-25114/},}
\bibitem{vijayan2017firstmalwareCausedOutage} ... \LaTeX:\\ \url{https://www.darkreading.com/threat-intelligence/first-malware-designed-solely-for-electric-grids-caused-2016-ukraine-outage/d/d-id/1329114},}
\bibitem{ferrante2017crashoverrideredflag} ... \LaTeX:\\ \url{https://www.powermag.com/why-crashoverride-is-a-red-flag-for-u-s-power-companies/},}
\bibitem{2017cashoverrideindustroyerAlto} ... \LaTeX:\\ \url{https://blog.paloaltonetworks.com/2017/06/crashoverrideindustroyer-protections-palo-alto-networks-customers/},}
\bibitem{spinner2018crashoverrideiot} ... \LaTeX:\\ \url{https://iiot-world.com/ics-security/cybersecurity/five-cybersecurity-experts-about-crashoverride-malware-main-dangers-and-lessons-for-iiot/},}
\bibitem{abb30062017crashoverridenotification} ... \LaTeX:\\ \url{https://search.abb.com/library/Download.aspx?DocumentID=9AKK107045A1003&amp;LanguageCode=en&amp;DocumentPartId=&amp;Action=Launch},}
\bibitem{blackhatusa2017} ... \LaTeX:\\ \url{https://www.blackhat.com/us-17/briefings/schedule/#industroyercrashoverride-zero-things-cool-about-a-threat-group-targeting-the-power-grid-6159},}
\bibitem{humanitcrashalert} ... \LaTeX:\\ \url{https://humanit.asia/ta17-163a/},}
\bibitem{sroberts2017CrashoerrideCronicles} ... \LaTeX:\\ \url{https://medium.com/@sroberts/the-crash-override-chronicles-overall-8389ef178fdf},}
\bibitem{parker2020industrialsystemsattack} ... \LaTeX:\\ \url{https://www.oilandgaseng.com/articles/the-most-infamous-cyber-attacks-on-industrial-systems/},}
\bibitem{kirk2017threatIndustrialControls} ... \LaTeX:\\ \url{https://www.bankinfosecurity.com/power-grid-malware-platform-threatens-industrial-controls-a-9987},}
\bibitem{chalfant2017USElectricGrid} ... \LaTeX:\\ \url{https://thehill.com/policy/cybersecurity/337877-crash-override-malware-heightens-fears-for-us-electric-grid},}
\bibitem{fbiWarningcrashOverride} ... \LaTeX:\\ \url{https://www.inguardians.com/dhs-fbi-warn-of-attacks-against-us-energy-manufacturing-companies-and-employees/},}
\bibitem{davies2017crashoveerideUkraineAttack} ... \LaTeX:\\ \url{https://rethinkresearch.biz/articles/industroyer-crashoverride-malware-behind-ukraine-utility-attack/},}
\bibitem{slowik2018securitySessions} ... \LaTeX:\\ \url{https://electricenergyonline.com/energy/magazine/1104/article/Security-Sessions-Combating-ICS-Threats.htm},}
\bibitem{ENCS2017crashoverridemodules} ... \LaTeX:\\ \url{https://www.smart-energy.com/regional-news/europe-uk/encs-crash-override-virus/},}
\bibitem{brocklehurst2017crashoverridegridtakedown} ... \LaTeX:\\ \url{https://isssource.com/crashoverride-designed-for-grid-takedown/},}
\bibitem{20170613_crashoverride} ... \LaTeX:\\ \url{https://gigazine.net/gsc_news/en/20170613-crashoverride/},}
@onlinec{leetaru2017crashoverridehomefront} ... \LaTeX:\\ \url{https://www.forbes.com/sites/kalevleetaru/2017/06/24/crash-override-and-how-cyberwarfare-is-bringing-conflict-to-the-homefront/#42eb8984277c},}
\bibitem{fti2017redflag} ... \LaTeX:\\ \url{https://fticybersecurity.com/2017-11/crashoverride-red-flag-u-s-power-companies/},}
\bibitem{27001academy2014Whitepaper} ... \LaTeX:\\ \url{https://info.advisera.com/hubfs/27001Academy/27001Academy_FreeDownloads/NL/Checklist_of_ISO_27001_Mandatory_Documentation_NL.pdf},}
\bibitem{2017info_isoiec27019} ... \LaTeX:\\ \url{https://webstore.iec.ch/preview/info_isoiec27019%7Bed1.0%7Den.pdf},}
\bibitem{iec623512023serseries} ... \LaTeX:\\ \url{https://webstore.iec.ch/publication/6912},}
\bibitem{iec623512011Withdrawn} ... \LaTeX:\\ \url{https://webstore.iec.ch/publication/6911},}
\bibitem{IEC62351sheet} ... \LaTeX:\\ \url{https://www.ipcomm.de/protocol/IEC62351/en/sheet.html},}
\bibitem{NISTOIR2014GuidelinesCyverSec} ... \LaTeX:\\ \url{https://nvlpubs.nist.gov/nistpubs/ir/2014/NIST.IR.7628r1.pdf},}
\bibitem{ukraiinnesandwormteam} ... \LaTeX:\\ \url{https://www.fireeye.com/blog/threat-research/2016/01/ukraine-and-sandworm-team.html/cybersecurity-audit-checklist. },}
\bibitem{politicsrussagridukraine} ... \LaTeX:\\ \url{https://edition.cnn.com/2016/02/11/politics/ukraine-power-grid-attack-russia-us/index.html/ukraine-sees-russian-hand-in-cyber-attacks-on-power-grid-idUSKCN0VL18E. },}
\bibitem{Whitehead2017ukrainepoweroutage,	ALTauthor = {David E. Whitehead, Kevin Owens, Dennis Gammel, and Jess Smith},	ALTeditor = {editor},	title = {title},	date = {March 21–23, 2017},	url = {https://na.eventscloud.com/file_uploads/aed4bc20e84d2839b83c18bcba7e2876_Owens1.pdf},}
\bibitem{zetter2017moreDangerousMalware} ... \LaTeX:\\ \url{https://www.vice.com/en/article/zmeyg8/ukraine-power-grid-malware-crashoverride-industroyer},}
\bibitem{icsSecurityRussianHacking} ... \LaTeX:\\ \url{https://www.wallix.com/blog/ics-security-russian-hacking},}
\bibitem{rocha2017cybersecyrityanalysisScada} ... \LaTeX:\\ \url{https://www.sans.org/blog/confirmation-of-a-coordinated-attack-on-the-ukrainian-power-grid/},}
\bibitem{cheah2008testingiec60870} ... \LaTeX:\\ \url{http://citeseerx.ist.psu.edu/viewdoc/download;jsessionid=0513EED48102FDAD1BD940260EF12B11?doi=10.1.1.548.7490&amp;rep=rep1&amp;type=pdf},}
\bibitem{virsec2017DeepDiveIndustroyer} ... \LaTeX:\\ \url{https://virsec.com/virsec-hack-analysis-deep-dive-into-industroyer-aka-crash-override/},}
\bibitem{fauri2017EncryptionICS} ... \LaTeX:\\ \url{https://www.win.tue.nl/~setalle/2017_fauri_encryption.pdf},}
\bibitem{2017IEC61850} ... \LaTeX:\\ \url{http://blog.nettedautomation.com/2017/},}
\bibitem{2017win32industroyer} ... \LaTeX:\\ \url{https://www.welivesecurity.com/wp-content/uploads/2017/06/Win32_Industroyer.pdf},}
\bibitem{dragos2019TargetedTransStation} ... \LaTeX:\\ \url{https://www.cybersecurityintelligence.com/blog/attack-on-ukraines-power-grid-targeted-transmission-stations-4530.html},}
\bibitem{2017crashoverridenostuxnet} ... \LaTeX:\\ \url{https://arstechnica.com/information-technology/2017/06/crash-override-malware-may-sabotage-electric-grids-but-its-no-stuxnet/},}
%%%%%%%%%%%%%%%%%%%%%%%%%%%%%%%%%%%%%%%%%%%%%%%%%%%%%%%%%%%%%%%%%



%
%\chapter{Deelonderzoek naar veiligheidsrisico's en eisen voor sluizen}
%
%Gevonden weblinks in google op 07-04-2023 met zoekopdracht: "veiligheidsrisico's voor sluizen en waterwerken"
%





\bibitem{ID} ... \LaTeX:\\ \url{https://www.tweedekamer.nl/downloads/document?id=80443e97-f17e-499c-b3f2-ad608f32e1aa&title=Rapportage%20Staat%20van%20de%20infra%20RWS%20%28definitief%29.pdf}
\bibitem{ID} ... \LaTeX:\\ \url{https://www.nu.nl/internet/5814282/rekenkamer-waterwerken-niet-goed-beveiligd-tegen-cyberaanvallen.html}
\bibitem{linne21122022onderhousluis} ... \LaTeX:\\ \url{https://www.deltalimburg.nl/article/9824/Onderhoudswerkzaamheden+aan+Sluis+Linne+afgerond}
\bibitem{sluisterneuzenveiligheid} ... \LaTeX:\\ \url{https://nieuwesluisterneuzen.eu/veiligheid}
\bibitem{mrdmarinesupport} ... \LaTeX:\\ \url{https://www.mrdmarinesupport.nl/nl/maritieme-dienstverlening/ondersteuning-veiligheid/}
\bibitem{krabbendam27052021infrasitrrisicobeoordeling} ... \LaTeX:\\ \url{https://www.infrasite.nl/bouwen/2021/05/27/veiligheid-voorop-begin-project-sluis-of-brug-altijd-met-risicobeoordeling/}
\bibitem{bedieningsluisVechterweerdVilsteren} ... \LaTeX:\\ \url{https://www.wdodelta.nl/bediening-schutsluizen-vechterweerd-en-vilsteren}
\bibitem{krabbedam21052021sluisHeel} ... \LaTeX:\\ \url{https://www.infrasite.nl/waterbouw-deltas/2021/05/21/sluis-heel-onder-handen-genomen/}
\bibitem{hdsr30092022lichtprojectieswaterliniesluizen} ... \LaTeX:\\ \url{https://www.hdsr.nl/actueel/nieuws/@154100/lichtprojecties-zetten-waterliniesluizen/}
\bibitem{nos28032019rekenkamerhackwaterwerk} ... \LaTeX:\\ \url{https://nos.nl/artikel/2277937-rekenkamer-hack-aanval-op-waterwerk-niet-altijd-opgemerkt}
\bibitem{sluisLinne08012023} ... \LaTeX:\\ \url{https://varendoejesamen.nl/kenniscentrum/artikel/onderhoud-sluis-linne-afgerond}
\bibitem{gww29032021kantelendesluisdeur} ... \LaTeX:\\ \url{https://www.gww-bouw.nl/artikel/de-eerste-sluis-met-kantelende-sluisdeur/}
\bibitem{tkh03022018sluisbeheerScadaZeeland} ... \LaTeX:\\ \url{https://tkhsecurity.com/nl/waterwerken/}
\bibitem{h2o28032019dreigingsniveaurekenkamer} ... \LaTeX:\\ \url{https://www.h2owaternetwerk.nl/h2o-actueel/rekenkamer-vitale-waterwerken-nog-onvoldoende-beschermd-tegen-cyberaanvallen}
\bibitem{krammersluizencomplex} ... \LaTeX:\\ \url{https://www.magazinesrijkswaterstaat.nl/bereikbaarzeeland/2021/01/krammersluizencomplex-verleden-heden-en-toekomst}



%\chapter{Deelonderzoek wet en regelgeving voor sluizen}
%
%
%
%Gevonden weblinks in google op 07-04-2023 met zoekopdracht: "wet en regelgeving voor sluizen"



\bibitem{ID} ... \LaTeX:\\ \url{https://www.hdsr.nl/publish/pages/86927/sluizen_in_of_bij_een_waterkering_-_uitvoeringsregels.pdf}
\bibitem{ID} ... \LaTeX:\\ \url{https://api1.ibabs.eu/publicdownload.aspx?site=sluis&id=100100292}
\bibitem{ID} ... \LaTeX:\\ \url{https://services.pilz.nl/wp-content/uploads/2021/12/brochure_bruggen_2018.pdf}
\bibitem{ID} ... \LaTeX:\\ \url{https://lokaleregelgeving.overheid.nl/CVDR375606/6}
\bibitem{ID} ... \LaTeX:\\ \url{https://zoek.officielebekendmakingen.nl/stb-2019-27.html}
\bibitem{ID} ... \LaTeX:\\ \url{https://a-quin.nl/nieuws/veiligheid-van-bruggen-sluizen-waarborgen-wie-wat-hoe/}
\bibitem{ID} ... \LaTeX:\\ \url{https://www.gemeentesluis.nl/Bestuur_en_Organisatie/Wetten_Regels_Bekendmakingen}
\bibitem{ID} ... \LaTeX:\\ \url{https://www.overijssel.nl/onderwerpen/verkeer-en-vervoer/varen-in-overijssel/informatie-bedieningstijden-sluizen-en-bruggen-noordwest-overijssel/}
\bibitem{ID} ... \LaTeX:\\ \url{https://www.rijkswaterstaat.nl/water/wetten-regels-en-vergunningen}
\bibitem{ID} ... \LaTeX:\\ \url{https://www.schuttevaer.nl/nieuws/actueel/2022/11/23/binnenvaart-zit-klem-tussen-regels-en-realiteit-kapotte-steigers-en-gesperde-sluizen-dwingen-tot-doorvaren/}
\bibitem{ID} ... \LaTeX:\\ \url{https://repository.officiele-overheidspublicaties.nl/CVDR/CVDR271406/1/html/CVDR271406_1.html}
\bibitem{ID} ... \LaTeX:\\ \url{https://www.zeeland.nl/actueel/bedieningstijden-sluizen-en-bruggen}
\bibitem{ID} ... \LaTeX:\\ \url{https://www.amsterdam.nl/verkeer-vervoer/varen-amsterdam/regels-varen/}
\bibitem{ID} ... \LaTeX:\\ \url{https://www.schielandendekrimpenerwaard.nl/wat-doen-we/regels-en-afspraken-over-beheer-keur-en-leggers/}
\bibitem{ID} ... \LaTeX:\\ \url{http://www.wetboek-online.nl/wet/Wet%20tot%20samenvoeging%20van%20de%20gemeenten%20Aardenburg%20en%20Sluis.html}
\bibitem{ID} ... \LaTeX:\\ \url{https://www.rijnland.net/regels-op-een-rij/richtlijnen-en-akkoorden/alle-regelgeving-van-rijnland/}
\bibitem{ID} ... \LaTeX:\\ \url{https://www.itbb.nl/diensten/advies-ce-markering-europese-richtlijnen/}
\bibitem{ID} ... \LaTeX:\\ \url{https://www.portofamsterdam.com/nl/scheepvaart/zeevaart/regelgeving}
\bibitem{ID} ... \LaTeX:\\ \url{https://www.watersportverbond.nl/nieuws/achterstallig-onderhoud-wachtplaatsen-bruggen-en-sluizen-zuid-holland-zorgelijk/}
\bibitem{ID} ... \LaTeX:\\ \url{https://varendoejesamen.nl/nieuws}
\bibitem{ID} ... \LaTeX:\\ \url{https://www.flevoland.nl/wat-doen-we/flevowegen-vlot-en-veilig-door-flevoland/water/varen-in-flevoland/bediening-bruggen-en-sluizen}
\bibitem{ID} ... \LaTeX:\\ \url{https://eur-lex.europa.eu/legal-content/NL/TXT/PDF/?uri=CELEX:32020L0012&from=DE}
\bibitem{ID} ... \LaTeX:\\ \url{https://www.werkenvoornederland.nl/organisatie/rijkswaterstaat/ict-middelen-maken-om-bruggen-sluizen-en-tunnels-te-besturen}
\bibitem{ID} ... \LaTeX:\\ \url{https://www.lobocom.nl/infra-bruggen-sluizen}
\bibitem{ID} ... \LaTeX:\\ \url{https://waterrecreatienederland.nl/content/uploads/2018/04/richtlijnen-vaarwegen-2017.pdf}
\bibitem{ID} ... \LaTeX:\\ \url{https://www.wetterskipfryslan.nl/melden-en-regelen/vergunningen-wetten-en-regels}
\bibitem{ID} ... \LaTeX:\\ \url{https://www.onlinezeilschool.nl/sluizen/}
\bibitem{ID} ... \LaTeX:\\ \url{https://www.provincie.drenthe.nl/onderwerpen/verkeer-vervoer/vaarwegen/rondje-drenthe/bedieningstijden/}





Rampen
In dit hoofdstuk worden de resultaten van een deskresearch naar verschillende rampen behandeld
Hierbij een verslag naar de oorzaken van de rampen, de werkwijze waarop het product is ontwikkeld, de verwerking van feedback, implementatie en nazorg.
Met behulp van het 4 variabelen model wordt duidelijk gemaakt hoe het systeem is opgezet en wat daarin verkeerd is gegaan.
Het  hoofdstuk wordt afgesloten met een analyse van algemene kenmerken van de verschillende rampen die zijn onderzocht.


Het 4 variabelen model kort toegelicht
Monitored variabelen: door sensoren gekwantificeerde fenomenen uit de omgeving, bijv temperatuur

Controlled variabelen: door actuatoren \bestuurde fenomenen uit de omgeving
For example, monitored variables might be the pressure and temperature
inside a nuclear reactor while controlled variables might be visual and audible alarms, as well
as the trip signal that initiates a reactor shutdown; whenever the temperature or pressure reach
abnormal values, the alarms go off and the shutdown procedure is initiated

Input variabelen: data die de software als input gebruikt
Here, IN models the input hardware interface (sensors and analog-to-digital converters) and
relates values of monitored variables to values of input variables in the software. The input variables model the information about the environment that is available to the software. For example,
IN might model a pressure sensor that converts temperature values to analog voltages; these voltages are then converted via an A/D converter to integer values stored in a register accesible to the
software.

Output variabelen: data die de software levert als output
The output hardware interface (digital-to-analog converters and actuators) is modelled
by OUT, which relates values of the output variables of the software to values of controlled variables. An output variable might be, for instance, a boolean variable set by the software with the
understanding that the value true indicates that a reactor shutdown should occur and the value
false indicates the opposite


Bronnen:
\bibitem{ID} ... \LaTeX:\\ \url{https://www.sciencedirect.com/science/article/pii/S0167642315001033}
\bibitem{ID} ... \LaTeX:\\ \url{https://www.cas.mcmaster.ca/~lawford/papers/AVoCS2013.pdf}
\bibitem{ID} ... \LaTeX:\\ \url{https://core.ac.uk/download/pdf/38891842.pdf}
%%%%%%%%%%%%%%%%%%%%%%%%%%%%%%%%%%%%%%%%%%%%%%%%%%%%%%%%%%%%%%%%%

Therac


sheets
\bibitem{ID} ... \LaTeX:\\ \url{https://web.cs.ucdavis.edu/~rogaway/classes/188/winter04/therac-25.pdf}
\bibitem{ID} ... \LaTeX:\\ \url{https://people.physics.carleton.ca/~drogers/egs_windows_collection/tsld008.htm}
\bibitem{ID} ... \LaTeX:\\ \url{https://en.wikipedia.org/wiki/Therac-25}
\bibitem{ID} ... \LaTeX:\\ \url{https://www.youtube.com/watch?v=-7gVqBY52MY}
reproduceren van de error. IN dit stuk wordt uitgelgd hoe het product werkt en waarom bepaalde beslssingen zijn genomen in de ontwerp/productiefase
\bibitem{ID} ... \LaTeX:\\ \url{https://www.bugsnag.com/blog/bug-day-race-condition-therac-25}
kort artikel met daarin een opsomming van alle fouten in het systeem en een korte uitleg
\bibitem{ID} ... \LaTeX:\\ \url{https://www.bowdoin.edu/~allen/courses/cs260/readings/therac.pdf}
uitgebreid artikel over hoe de fout werd gereproduceerd en de resultaten daaruit voortkwamen. Alsnog werden er na de reproductie fase nog meer fouten gevonden.
\bibitem{ID} ... \LaTeX:\\ \url{https://hackaday.com/2015/10/26/killed-by-a-machine-the-therac-25/}
artikel
\bibitem{ID} ... \LaTeX:\\ \url{https://ethicsunwrapped.utexas.edu/case-study/therac-25}
onderzoeksartikel waarin de bug wordt uitgelgd: de racecondities, de bytepositie en het testen worden berkitiseerd envenals andere onderdelen van het softwareproces
\bibitem{ID} ... \LaTeX:\\ \url{https://thedailywtf.com/articles/the-therac-25-incident}
onrealistisch testplan. In dit artikel egt de auteur het belang nog eens uit van goede requirements en implementatie, niet de software is waar het probleem ligt
\bibitem{ID} ... \LaTeX:\\ \url{https://www.computer.org/csdl/magazine/co/2017/11/mco2017110008/13rRUxAStVR}
onderzoeksrapport
\bibitem{ID} ... \LaTeX:\\ \url{https://web.stanford.edu/class/cs240/old/sp2014/readings/therac-25.pdf}
geschiedenis
\bibitem{ID} ... \LaTeX:\\ \url{http://computingcases.org/case_materials/therac/case_history/Case%20History.html}
artikel
\bibitem{ID} ... \LaTeX:\\ \url{https://medium.com/swlh/software-architecture-therac-25-the-killer-radiation-machine-8a05e0705d5b}
computer error. De ongeval en de malfunction nog een keer uitgelegd
\bibitem{ID} ... \LaTeX:\\ \url{http://www.ccnr.org/fatal_dose.html}
rapport
\bibitem{ID} ... \LaTeX:\\ \url{http://sunnyday.mit.edu/papers/therac.pdf}
\bibitem{ID} ... \LaTeX:\\ \url{https://pubmed.ncbi.nlm.nih.gov/101762/}
onderzoeksartkel
\bibitem{ID} ... \LaTeX:\\ \url{http://www1.cs.columbia.edu/~junfeng/08fa-e6998/sched/readings/therac25.pdf}
\bibitem{ID} ... \LaTeX:\\ \url{https://ieeexplore.ieee.org/document/274940}
uitgebreid artikel gaat hier ook wat meer over de hardware
\bibitem{ID} ... \LaTeX:\\ \url{https://www.linkedin.com/pulse/therac-25-industrial-design-engineering-systems-wang-ph-d-cre-acb/}
artikel waarin in 3 delen de problemaiekwordt blootgesteld
\bibitem{ID} ... \LaTeX:\\ \url{http://www.cse.msu.edu/~cse470/Public/Handouts/Therac/Therac_2.html}
case study sheets
\bibitem{ID} ... \LaTeX:\\ \url{https://www.cs.jhu.edu/~cis/cista/445/Lectures/Therac.pdf}
artikel waarin vooral de fabriikant ervan langs krijgt
\bibitem{ID} ... \LaTeX:\\ \url{http://users.csc.calpoly.edu/~jdalbey/SWE/Papers/THERAC25.html}
lessons learned. Vooral de begrippen betrouwbaarheid, welgevalligheid, veilgheid en gebruiksvriendelijkheid
\bibitem{ID} ... \LaTeX:\\ \url{https://bohr.wlu.ca/cp164/therac/therac25.htm}
root-cause analysis
\bibitem{ID} ... \LaTeX:\\ \url{https://root-cause-analysis.info/2010/08/08/therac-25-radiation-overdoses/}
case study
\bibitem{ID} ... \LaTeX:\\ \url{https://dusk.geo.orst.edu/ethics/papers/Therac.Huff.pdf}
\bibitem{ID} ... \LaTeX:\\ \url{https://embeddedartistry.com/fieldatlas/case-study-therac-25/}
case study
\bibitem{ID} ... \LaTeX:\\ \url{https://www.sebokwiki.org/wiki/Medical_Radiation}
opzetten van systematische acceptaatie test met therac als voorbeeld
\bibitem{ID} ... \LaTeX:\\ \url{https://www.sciencedirect.com/science/article/pii/S1474667017448245}
artikel waarin een diagnose plaatvindt voor het bedrijf en de ingenieur/ontwerper
\bibitem{ID} ... \LaTeX:\\ \url{https://magsilva.pro.br/apps/wiki/testing/Therac_25}
rapport
\bibitem{ID} ... \LaTeX:\\ \url{http://csel.eng.ohio-state.edu/productions/pexis/readings/submod3/therac.pdf}
oorzaken aangegeven in artikel
\bibitem{ID} ... \LaTeX:\\ \url{https://www.chemeurope.com/en/encyclopedia/Therac-25.html}
het onderzoek en enkele ontwerptekeningen en oplossingen
\bibitem{ID} ... \LaTeX:\\ \url{https://pvs-studio.com/en/blog/posts/0438/}
\bibitem{ID} ... \LaTeX:\\ \url{https://www.coursera.org/lecture/software-design-threats-mitigations/therac-25-case-study-VmQPa}
\bibitem{ID} ... \LaTeX:\\ \url{https://www.semanticscholar.org/paper/The-story-of-the-Therac-25-in-LOTOS-Thomas/6c9c6024cf95aadae8b7edf1160e0e4500410eb9}
\bibitem{ID} ... \LaTeX:\\ \url{https://news.ycombinator.com/item?id=21679287}
wiki
\bibitem{ID} ... \LaTeX:\\ \url{https://en.wikibooks.org/wiki/Professionalism/Therac-25}
analyse
\bibitem{ID} ... \LaTeX:\\ \url{https://citeseerx.ist.psu.edu/viewdoc/download?doi=10.1.1.96.369&rep=rep1&type=pdf}
samenvatting
\bibitem{ID} ... \LaTeX:\\ \url{https://onlineethics.org/cases/resources-engineering-and-science-ethics/investigation-therac-25-accidents-abstract}
podcast
\bibitem{ID} ... \LaTeX:\\ \url{https://podcasts.apple.com/gb/podcast/therac-25/id1046978749?i=1000514115050}
enkele conclusies
\bibitem{ID} ... \LaTeX:\\ \url{http://www.cas.mcmaster.ca/~se4d03/therac.html}
rapport over de fouten die de verschillende partijen hebben gemaakt( overheid, ingenieurs, bedrijf, operators) en de verbeterpunten
\bibitem{ID} ... \LaTeX:\\ \url{https://www.cs.colostate.edu/~bieman/CS314/Notes/therac25.pdf}
onderzoeksrapport
\bibitem{ID} ... \LaTeX:\\ \url{https://www.cs.ucf.edu/~dcm/Teaching/COP4600-Fall2010/Literature/Therac25-Leveson.pdf}
slides online over het technisch mankement
Wat is er gebeurd, nou het volgende:
Normal radiation treatments: 6,000 rads over a 3 week period, under certain conditions Therac-25 was delivering 60,000 rads during one session.
En wat ging er mis?
Paradigm Shift

Therac-25 replaced expensive hardware safety interlocks with software controls

Real-time software
Design
Race condition caused focusing element to be incorrectly set
No indication of actual hardware settings
Error messages appeared the same regardless of how important
Error messages were difficult to understand
All errors messages could be manually overridden

\bibitem{ID} ... \LaTeX:\\ \url{https://hci.cs.siue.edu/NSF/Files/Semester/Week13-2/PPT-Text/Slide13.html}
oorzaak-gevolg diagram
\bibitem{ID} ... \LaTeX:\\ \url{https://www.thinkreliability.com/InstructorBlogs/Blog-Therac-25.pdf}
veiligheidsanalyse naar de rapportage van foutmeldingen, de beslissingsmatrix waarmee het programma wordt uitgevoerd en de software-analyse door een consultat
\bibitem{ID} ... \LaTeX:\\ \url{https://www.erenkrantz.com/Geeks/Therac-25/Side_bar_5.html}
\bibitem{ID} ... \LaTeX:\\ \url{https://bandcamp.com/therac-25}
\bibitem{ID} ... \LaTeX:\\ \url{https://sites.nd.edu/brent-marin/2017/09/21/blog-post-8-therac-25-accidents/}
\bibitem{ID} ... \LaTeX:\\ \url{https://europepmc.org/article/med/4035720}
\bibitem{ID} ... \LaTeX:\\ \url{https://itlaw.wikia.org/wiki/Therac-25}
\bibitem{ID} ... \LaTeX:\\ \url{https://www.sitepoint.com/therac-25-bad-software-kills/}
\bibitem{ID} ... \LaTeX:\\ \url{https://sqa.stackexchange.com/questions/9798/asking-for-help-with-this-therac-25-bugged-code-i-dont-understand-the-explanat}
\bibitem{ID} ... \LaTeX:\\ \url{http://www.se.rit.edu/~swen-342/activities/TheracIndividual.html}
\bibitem{ID} ... \LaTeX:\\ \url{https://www.designnews.com/automation-motion-control/engineering-disasters-deadly-zaps-therac-25}

publicaties over de therac:
\bibitem{ID} ... \LaTeX:\\ \url{https://www.cs.colostate.edu/~bieman/CS314/Notes/therac25.pdf}
\bibitem{ID} ... \LaTeX:\\ \url{https://www.cs.ucf.edu/~dcm/Teaching/COP4600-Fall2010/Literature/Therac25-Leveson.pdf}
\bibitem{ID} ... \LaTeX:\\ \url{https://onlineethics.org/cases/resources-engineering-and-science-ethics/investigation-therac-25-accidents-abstract}
\bibitem{ID} ... \LaTeX:\\ \url{https://citeseerx.ist.psu.edu/viewdoc/download?doi=10.1.1.96.369&rep=rep1&type=pdf}
\bibitem{ID} ... \LaTeX:\\ \url{http://csel.eng.ohio-state.edu/productions/pexis/readings/submod3/therac.pdf}
\bibitem{ID} ... \LaTeX:\\ \url{https://dusk.geo.orst.edu/ethics/papers/Therac.Huff.pdf}
\bibitem{ID} ... \LaTeX:\\ \url{https://www.cs.jhu.edu/~cis/cista/445/Lectures/Therac.pdf}
\bibitem{ID} ... \LaTeX:\\ \url{http://www1.cs.columbia.edu/~junfeng/08fa-e6998/sched/readings/therac25.pdf}
\bibitem{ID} ... \LaTeX:\\ \url{http://sunnyday.mit.edu/papers/therac.pdf}
\bibitem{ID} ... \LaTeX:\\ \url{https://web.stanford.edu/class/cs240/old/sp2014/readings/therac-25.pdf}




%%%%%%%%%%%%%%%%%%%%%%%%%%%%%%%%%%%%%%%%%%%%%%%%%%%%%%%%%%%%%%%%%
Krakend zorgssteem door covid-19 in suriname

vaccinatieterkort
communicatie met bevolking
communicatie met binnenland
testen van vaccinaties
besmetting vanuit eht buitenland
isolatie na vakantie en voor toeristen
tekort aan ic-personeel
tekort aan ic-bedden
tekort aan zuurtstof
tekort aan middelen}

Wat blijkt hieruit:
de impact van de crisis wereldwijd
de afhnakelijkheid van landen op goede samenwerking
Nut en noodzaak van regelgeving
Naveling van maatregels
Communicatie over beleid vanuit de overheid naar de burgers
Belang van een verzorgingstaat
Een wetenschappelijke ontwikkeling die kan inspelen op gevoelige 'trends
De impact van een lockdown op de economie
Afschaling van andere noodzakelijke no-covid zorg
De bereikbaarheid van een ziekenhuis
Waar heeft het toe geleid?
\bibitem{ID} ... \LaTeX:\\ \url{https://www.waterkant.net/suriname/2007/02/06/school-in-suriname-gesloten-om-zenuwgasvoorraad/}
\bibitem{ID} ... \LaTeX:\\ \url{https://nl.wikipedia.org/wiki/Nationaal_Co%C3%B6rdinatiecentrum_voor_Rampenbeheersing}
\bibitem{ID} ... \LaTeX:\\ \url{https://www.examenkamer.nl/index.php/27-vca-examens-in-suriname}
%%%%%%%%%%%%%%%%%%%%%%%%%%%%%%%%%%%%%%%%%%%%%%%%%%%%%%%%%%%%%%%%%

Waterramp suriname met cyanice


%%%%%%%%%%%%%%%%%%%%%%%%%%%%%%%%%%%%%%%%%%%%%%%%%%%%%%%%%%%%%%%%%
boeing 737 crashes


algemene vragen
oorzaken
\bibitem{gates18112020boeingcrisis} ... \LaTeX:\\ \url{https://www.seattletimes.com/business/boeing-aerospace/what-led-to-boeings-737-max-crisis-a-qa/}
\bibitem{boeing737maxsoftwareprobles} ... \LaTeX:\\ \url{https://www.schneier.com/blog/archives/2019/04/excellent_analy.html}
fout in de software
\bibitem{avetisov19032019boeingmalwarestate} ... \LaTeX:\\ \url{https://www.forbes.com/sites/georgeavetisov/2019/03/19/malware-at-30000-feet-what-the-737-max-says-about-the-state-of-airplane-software-security/?sh=4d26f7052a9e}
het nationaal veiligheidsbelang
\bibitem{thompson23112020nationalsecurityboeing} ... \LaTeX:\\ \url{https://www.forbes.com/sites/lorenthompson/2020/11/23/five-reasons-return-of-boeings-737-max-to-service-is-important-to-national-security/?sh=2128ea552018}
falend toezicht
\bibitem{gates21032019FAAControlSystem} ... \LaTeX:\\ \url{https://www.seattletimes.com/business/boeing-aerospace/failed-certification-faa-missed-safety-issues-in-the-737-max-system-implicated-in-the-lion-air-crash/}
onderzoeksrapport
\bibitem{faa18112020boeingreview} ... \LaTeX:\\ \url{https://www.faa.gov/foia/electronic_reading_room/boeing_reading_room/media/737_RTS_Summary.pdf}
\bibitem{wiki737maxgroundings} ... \LaTeX:\\ \url{https://en.wikipedia.org/wiki/Boeing_737_MAX_groundings}
veiligheidsrisico's
menselijke fouten
\bibitem{campbell02052019boengcrashhumanerrors} ... \LaTeX:\\ \url{https://www.theverge.com/2019/5/2/18518176/boeing-737-max-crash-problems-human-error-mcas-faa}
overzicht van crashes
\bibitem{hawkins22032019737maxairplanes} ... \LaTeX:\\ \url{https://www.theverge.com/2019/3/22/18275736/boeing-737-max-plane-crashes-grounded-problems-info-details-explained-reasons}
veiligheidsopmerking
\bibitem{thomas30082020737safest} ... \LaTeX:\\ \url{https://www.airlineratings.com/news/boeings-737-max-will-one-safest-aircraft-history/}
aanpassingen
\bibitem{boeing737maxdisplay} ... \LaTeX:\\ \url{https://www.boeing.com/commercial/737max/737-max-software-updates.page}
waarschuwingen//output signalen
\bibitem{fehrm24112020737changes} ... \LaTeX:\\ \url{https://leehamnews.com/2020/11/24/boeing-737-max-changes-beyond-mcas/}
software gerelateerde fouten
\bibitem{travis18042019737maxsoftwaredevop} ... \LaTeX:\\ \url{https://spectrum.ieee.org/aerospace/aviation/how-the-boeing-737-max-disaster-looks-to-a-software-developer}
onderzoeksrapport
de rol van de publieke opinie
\bibitem{barnett05052019737maxcrisis} ... \LaTeX:\\ \url{https://pubsonline.informs.org/do/10.1287/orms.2019.05.05/full/}
onderzoek van europese luchtvaart agentschap
\bibitem{easa27012021737maxsafereturn} ... \LaTeX:\\ \url{https://www.easa.europa.eu/newsroom-and-events/news/easa-declares-boeing-737-max-safe-return-service-europe}
veiligheidsvraagstuk
\bibitem{touitou11032019737tragedies} ... \LaTeX:\\ \url{https://phys.org/news/2019-03-boeing-max-safety-tragedies.html}
artikel over sensoren
\bibitem{hemmerdinger02022021737maxdeliveries} ... \LaTeX:\\ \url{https://www.flightglobal.com/airframers/boeing-delays-737-max-10-deliveries-two-years-to-2023/142245.article}
goedkeuring van europese luchtvaart autoriteiten
advies aan de faa
\bibitem{bielby27022021faaimprovesafety} ... \LaTeX:\\ \url{https://www.hstoday.us/subject-matter-areas/airport-aviation-security/oig-tells-faa-to-improve-safety-oversight-following-boeing-737-max-review/}
\bibitem{boyle18112020737maxupgrade} ... \LaTeX:\\ \url{https://www.geekwire.com/2020/faas-go-ahead-737-maxs-return-flight-kicks-off-massive-software-upgrade/}
\bibitem{bergstraburgess122019737maxMcasAlgorithm} ... \LaTeX:\\ \url{https://www.researchgate.net/publication/338420944_A_Promise_Theoretic_Account_of_the_Boeing_737_Max_MCAS_Algorithm_Affair}
achtergrond informatie
\bibitem{737mcas} ... \LaTeX:\\ \url{http://www.b737.org.uk/mcas.htm}
algemeen vertrouwen
\bibitem{newburger17052019boeingcrisis} ... \LaTeX:\\ \url{https://www.cnbc.com/2019/05/16/what-you-need-to-know-about-boeings-737-max-crisis.html}
toestemming europese autoriteiten
problemen
\bibitem{arstechnica22012020737problems} ... \LaTeX:\\ \url{https://arstechnica.com/information-technology/2020/01/737-max-fix-slips-to-summer-and-thats-just-one-of-boeings-problems/}
uitgebreid artikel over de onderzoeken en het vliegverbod
\bibitem{german190620217372yaftergrounded} ... \LaTeX:\\ \url{https://www.cnet.com/news/boeing-737-max-8-all-about-the-aircraft-flight-ban-and-investigations/}
computers als oorzaak
lessons learned
\bibitem{beningo02052019boeinglessons} ... \LaTeX:\\ \url{https://www.designnews.com/electronics-test/5-lessons-learn-boeing-737-max-fiasco}
\bibitem{duran05042019boeingspof} ... \LaTeX:\\ \url{https://www.eurocontrol.int/publication/effects-network-extra-standby-aircraft-and-boeing-737-max-grounding}
single point of failure
\bibitem{ID} ... \LaTeX:\\ \url{https://dmd.solutions/blog/2019/04/05/how-a-single-point-of-failure-spof-in-the-mcas-software-could-have-caused-the-boeing-737-max-crash-in-ethiopia/}
\bibitem{makichuck24012021737fearflying} ... \LaTeX:\\ \url{https://asiatimes.com/2021/01/boeings-737-max-and-the-fear-of-flying/}
lijst van tehnische aanpassingen
\bibitem{caa737modifications} ... \LaTeX:\\ \url{https://www.caa.co.uk/Consumers/Guide-to-aviation/Boeing-737-MAX/}
\bibitem{oestergaard14122020boeingdeliveries} ... \LaTeX:\\ \url{https://dsm.forecastinternational.com/wordpress/2020/12/14/airbus-and-boeing-report-november-2020-commercial-aircraft-orders-and-deliveries/}
code lek
\bibitem{reenberg787flaws} ... \LaTeX:\\ \url{https://www.wired.com/story/boeing-787-code-leak-security-flaws/}
\bibitem{fitch16092020737backlogrisks} ... \LaTeX:\\ \url{https://www.fitchratings.com/research/corporate-finance/boeing-737-max-return-backlog-risks-remain-16-09-2020}
Cultuurverandering, deregulatie, systeemwijziging of gewoon een kwestie van competentie
\bibitem{willis27082020737maxfailures} ... \LaTeX:\\ \url{https://www.aerospacetestinginternational.com/features/what-broke-the-737-max.html}
extra aanpassingen
\bibitem{ostrower11062020more737changes} ... \LaTeX:\\ \url{https://theaircurrent.com/aviation-safety/boeings-737-max-software-done-but-regulators-plot-more-changes-after-jets-return/}
wat ging er mis een analyse van een ex-iloot
De utoriteiten waren op de hoogte
\bibitem{hruska13122019faaknown737crashrate} ... \LaTeX:\\ \url{https://www.extremetech.com/extreme/303373-the-faa-knew-the-737-max-was-dangerous-and-kept-it-flying-anyway}
kwaliteiten van het alarmsysteem niet goed bekend
\bibitem{bloomberg26092019failedpred} ... \LaTeX:\\ \url{https://time.com/5687473/boeing-737-alarm-system/}
\bibitem{whiteman09072020boengcancelstock} ... \LaTeX:\\ \url{https://www.nasdaq.com/articles/boeing-gets-dealt-another-737-max-cancellation-blow.-what-it-means-for-boeing-stock-2020}
\bibitem{leopold09192019boeingreliability} ... \LaTeX:\\ \url{https://www.eetimes.com/boeing-crashes-highlight-a-worsening-reliability-crisis/}
veiligheidsvraagstuk
\bibitem{koenig11122019737crashesnofix} ... \LaTeX:\\ \url{https://www.latimes.com/business/story/2019-12-11/faa-boeing-737-max-crashes}
probleemanalyse, veiligheidsvraagstuk
\bibitem{dohertylindeman15032019737problems} ... \LaTeX:\\ \url{https://www.politico.com/story/2019/03/15/boeing-737-max-grounding-1223072}
falend toezicht
\bibitem{stodder02102019corruptoversight} ... \LaTeX:\\ \url{https://www.pogo.org/analysis/2019/10/corrupted-oversight-the-faa-boeing-and-the-737-max/}
\bibitem{afacwaLostSafeguards} ... \LaTeX:\\ \url{https://www.afacwa.org/the_inside_story_of_mcas_seattle_times}
doelstellingen en veiligheidsvraagstukken
\bibitem{swayne18032019profitssafety} ... \LaTeX:\\ \url{https://www.marxist.com/737-max-scandal-boeing-putting-profits-before-safety.htm}
\bibitem{freed26022021liftaustraliaban} ... \LaTeX:\\ \url{https://finance.yahoo.com/news/australia-lifts-ban-boeing-737-035817682.html?guccounter=1&guce_referrer=aHR0cHM6Ly93d3cuZ29vZ2xlLmNvbS8&guce_referrer_sig=AQAAAHZCJYy_0A5VS2WiPoCvH4xdrRNkmkdsv5EWJ2RLIz_AS-rxsTty6AF1_HlmJiRyWYqCXDi4p0Xs4isYkNkCq2Pfo-pQ60Xz_IfTNjm4FgoZiBMC4zpZlB6F0fwecrjE_ujAXZzG4xPJnWCd8-G3VLlPTY8h3H31eQ1i8hY9AIyy}
autoriteiten krijgen tik op de vingers
\bibitem{reed15032019softwareattention} ... \LaTeX:\\ \url{https://medium.com/@jpaulreed/the-737max-and-why-software-engineers-should-pay-attention-a041290994bd}
\bibitem{news17032019softwareexplains} ... \LaTeX:\\ \url{https://news.ycombinator.com/item?id=19414775}
\bibitem{legget21122020eu737maxsafe} ... \LaTeX:\\ \url{https://www.bbc.com/news/55366320}
\bibitem{marketscreener0103221737chinarecertification} ... \LaTeX:\\ \url{https://www.marketscreener.com/news/latest/China-studies-Boeing-737-MAX-recertification-wants-safety-concerns-fully-addressed--32569125/}
motor in brand
\bibitem{euractiv22022021737firegrounds} ... \LaTeX:\\ \url{https://www.euractiv.com/section/aviation/news/boeing-grounds-777s-after-engine-fire/}
\bibitem{benny18022019737returnUAE} ... \LaTeX:\\ \url{https://gulfnews.com/business/aviation/uae-airspace-to-see-return-of-boeing-737-max-1.1613627548923}
motor in brand gevlogen
\bibitem{biersmichel22022021777grounds} ... \LaTeX:\\ \url{https://techxplore.com/news/2021-02-boeing-urges-grounding-777s.html}
\bibitem{ID} ... \LaTeX:\\ \url{https://www.politico.eu/article/uk-temporarily-bans-some-boeing-aircraft-after-pratt-whitney-engine-incidents/}
\bibitem{reuters23022021777metalfatigue} ... \LaTeX:\\ \url{https://www.timeslive.co.za/news/world/2021-02-23-damage-to-united-boeing-777-engine-consistent-with-metal-fatigue--ntsb/}
faa was niet kritisch genoeg
\bibitem{ID} ... \LaTeX:\\ \url{https://federalnewsnetwork.com/government-news/2021/02/federal-watchdog-blasts-faa-over-certification-of-boeing-jet/}






%%%%%%%%%%%%%%%%%%%%%%%%%%%%%%%%%%%%%%%%%%%%%%%%%%%%%%%%%%%%%%%%%

china explosion 2015 tianjin

verhaal van brandweermannen
\bibitem{ID} ... \LaTeX:\\ \url{https://slate.com/human-interest/2015/08/chinese-explosion-aftermath-officials-investigate-causes-behind-warehouse-blast-and-death-of-88-firefighters.html}
artikel
\bibitem{ID} ... \LaTeX:\\ \url{https://www.chinafile.com/conversation/tianjin-explosion}
invloed van social media
\bibitem{ID} ... \LaTeX:\\ \url{https://www.economist.com/asia/2015/08/18/a-blast-in-tianjin-sets-off-an-explosion-online}
\bibitem{ID} ... \LaTeX:\\ \url{https://america.cgtn.com/2015/08/12/explosion-reported-in-tianjin-china}
\bibitem{ID} ... \LaTeX:\\ \url{https://factcheck.afp.com/no-photo-was-taken-chinese-city-tianjin-august-2015}
vergelijking van twee rampen
\bibitem{ID} ... \LaTeX:\\ \url{https://airshare.air-inc.com/how-does-the-beirut-explosion-compare-to-tianjin}
overheid en media
\bibitem{ID} ... \LaTeX:\\ \url{https://newbloommag.net/2015/08/17/tianjin-explosion/}
chemische industrie ondeer de loep
\bibitem{ID} ... \LaTeX:\\ \url{https://www.voanews.com/east-asia-pacific/tianjin-blast-puts-spotlight-chemical-industry}
\bibitem{ID} ... \LaTeX:\\ \url{https://abcnews.go.com/International/apocalyptic-aftermath-devastating-images-tianjin-china-explosions/story?id=33057017}
\bibitem{ID} ... \LaTeX:\\ \url{https://www.reachingoutacrossdurham.co.uk/osk/tianjin-explosion-2021}
\bibitem{ID} ... \LaTeX:\\ \url{https://aiche.onlinelibrary.wiley.com/doi/abs/10.1002/prs.11789}
\bibitem{ID} ... \LaTeX:\\ \url{https://www.automotivelogistics.media/thousands-of-cars-destroyed-in-tianjin-port-explosions/13570.article}
\bibitem{ID} ... \LaTeX:\\ \url{https://www.joc.com/port-news/asian-ports/port-tianjin/tianjin-port-explosions-could-be-most-expensive-maritime-disaster_20150826.html}
\bibitem{ID} ... \LaTeX:\\ \url{https://www.bloomberg.com/news/articles/2015-08-12/explosion-in-northern-china-shatters-windows-causes-injuries}
\bibitem{ID} ... \LaTeX:\\ \url{https://unece.org/fileadmin/DAM/env/documents/2016/TEIA/OECD_WGCA_24-27_OCT_2016/Session_3_Zhao_-__Introduction_of_Tianjin_Accident_-_Jinsong_Zhao.pdf}
gemaakte fouten
\bibitem{ID} ... \LaTeX:\\ \url{https://porteconomicsmanagement.org/pemp/contents/part6/port-resilience/site-2015-tianjin-port-explosions/}
\bibitem{ID} ... \LaTeX:\\ \url{https://www.alamy.com/stock-image-tianjin-china-17th-aug-2015-tianjin-explosion-aftermath-blast-site-165334778.html}
\bibitem{ID} ... \LaTeX:\\ \url{https://www.popularmechanics.com/technology/news/a16871/massive-explosions-china-city-of-tianjin/}
\bibitem{ID} ... \LaTeX:\\ \url{https://www.imago-images.com/st/0080815934}
\bibitem{ID} ... \LaTeX:\\ \url{https://www.chemistryworld.com/news/deadly-chemical-blast-at-chinese-port/8857.article}
\bibitem{ID} ... \LaTeX:\\ \url{https://www.process-worldwide.com/tianjin-explosion-from-chemical-perspective-insights-and-backgrounds-a-502381/}
vergelijking met andere explosies
\bibitem{ID} ... \LaTeX:\\ \url{https://apnews.com/article/lebanon-fires-us-news-explosions-middle-east-53f4206a7f1db0812262a15d22e1e58f}
invloed van de ramp op de industrie
\bibitem{ID} ... \LaTeX:\\ \url{https://fortune.com/2015/08/14/tianjin-port-explosion-shipping-delays/}
is er sprake van een doofpot
\bibitem{ID} ... \LaTeX:\\ \url{https://www.washingtontimes.com/news/2015/aug/20/inside-china-tianjin-explosions-cover-up-exposes-b/}
eigendomsverzekering
\bibitem{ID} ... \LaTeX:\\ \url{https://www.artemis.bm/news/tianjin-explosions-property-insurance-loss-could-reach-3-5bn-swiss-re/}
\bibitem{ID} ... \LaTeX:\\ \url{https://www.thechinastory.org/yearbooks/yearbook-2015/forum-the-abyss-%E5%9D%8E/tianjin-explosions/}
effecten op de lange termijn
\bibitem{ID} ... \LaTeX:\\ \url{https://www.flexport.com/blog/tianjin-explosion-effect-on-supply-chains/}
\bibitem{ID} ... \LaTeX:\\ \url{https://www.cicm.org.my/images/articles/CICM-Article-on-Tianjin-Blast-Oct2015.pdf}
lessons learned
\bibitem{ID} ... \LaTeX:\\ \url{https://www.genre.com/knowledge/blog/lessons-from-the-tianjin-explosion-en.html}
\bibitem{ID} ... \LaTeX:\\ \url{https://www.ft.com/content/ad62904c-44ce-11e5-b3b2-1672f710807b}
\bibitem{ID} ... \LaTeX:\\ \url{https://www.huffingtonpost.co.uk/2015/08/13/tianjin-explosion-china-shocking-footage-caught-on-camera_n_7980888.html}
\bibitem{ID} ... \LaTeX:\\ \url{https://www.thatsmags.com/china/post/19189/massive-fire-rocks-tianjin-port}
gevolgen voor de industrie
\bibitem{ID} ... \LaTeX:\\ \url{https://www.everstream.ai/risk-center/special-reports/the-jiangsu-yancheng-explosion/}
\bibitem{ID} ... \LaTeX:\\ \url{https://www.newyorker.com/news/news-desk/after-tianjin-an-outbreak-of-mistrust-in-china}
framing vanuit de chinese media
\bibitem{ID} ... \LaTeX:\\ \url{https://www.neliti.com/publications/101997/the-chinese-media-framing-of-the-2015s-tianjin-explosion}
\bibitem{ID} ... \LaTeX:\\ \url{https://www.reinsurancene.ws/chinese-insurers-settle-1-5-billion-tianjin-blast-claims/}
niewsartikel
\bibitem{ID} ... \LaTeX:\\ \url{https://www.thechemicalengineer.com/news/update-78-confirmed-dead-after-chinese-chemicals-plant-explosion/}
\bibitem{ID} ... \LaTeX:\\ \url{https://www.caixinglobal.com/2016-11-10/chinese-executive-receives-suspended-death-sentence-over-2015-tianjin-warehouse-blast-101006325.html}
toegang tot de ramplplek vanuit de okale journalistiek
\bibitem{ID} ... \LaTeX:\\ \url{https://chinadigitaltimes.net/2015/08/he-xiaoxin-how-far-can-i-go-and-how-much-can-i-do/}
artikel
\bibitem{ID} ... \LaTeX:\\ \url{https://www.wnpr.org/post/china-examines-aftermath-immense-twin-explosions-killed-dozens}
\bibitem{ID} ... \LaTeX:\\ \url{https://theconversation.com/what-is-ammonium-nitrate-the-chemical-that-exploded-in-beirut-143979}
\bibitem{ID} ... \LaTeX:\\ \url{https://chemicalwatch.com/36730/nationwide-inspections-in-china-follow-tianjin-explosion}
\bibitem{ID} ... \LaTeX:\\ \url{https://www.thehindu.com/news/international/investigation-begun-into-china-gas-explosion-as-toll-rises/article34818324.ece}
\bibitem{ID} ... \LaTeX:\\ \url{https://santiagotimes.cl/2019/03/24/64-killed-600-injured-in-china-chemical-plant-blast/}
oorzaken
\bibitem{ID} ... \LaTeX:\\ \url{https://klingecorp.com/blog/what-caused-the-tianjin-explosions/}
case study
\bibitem{ID} ... \LaTeX:\\ \url{https://www.preventionweb.net/educational/view/57235}
niewsartikel
\bibitem{ID} ... \LaTeX:\\ \url{https://www.cnbc.com/2015/08/12/explosion-in-tianjin-china.html}
chronologische uiteenzetting
\bibitem{ID} ... \LaTeX:\\ \url{https://www.aria.developpement-durable.gouv.fr/wp-content/files_mf/A46803_a46803_fiche_impel_006.pdf}
corruptie
\bibitem{ID} ... \LaTeX:\\ \url{https://www.nytimes.com/2015/08/31/world/asia/behind-tianjin-tragedy-a-company-that-flouted-regulations-and-reaped-profits.html}
mismanagement als oorzaak
\bibitem{ID} ... \LaTeX:\\ \url{https://www.nytimes.com/2016/02/06/world/asia/tianjin-explosions-were-result-of-mismanagement-china-finds.html}
\bibitem{ID} ... \LaTeX:\\ \url{https://cen.acs.org/articles/94/web/2016/02/Chinese-Investigators-Identify-Cause-Tianjin.html}
autoriteiten publiceren onderoeksrapport
\bibitem{ID} ... \LaTeX:\\ \url{https://cen.acs.org/articles/94/i7/Chinese-Investigators-Identify-Cause-Tianjin.html}
fotos van de rampplek
\bibitem{ID} ... \LaTeX:\\ \url{https://www.theatlantic.com/photo/2015/08/photos-of-the-aftermath-of-the-massive-explosions-in-tianjin-china/401228/}
\bibitem{ID} ... \LaTeX:\\ \url{https://edition.cnn.com/2015/08/13/asia/china-tianjin-explosions/index.html}
niuwesartiekel}
\bibitem{ID} ... \LaTeX:\\ \url{https://www.cbc.ca/news/world/china-explosion-tianjin-1.3189455}
verantwoordelijke
\bibitem{ID} ... \LaTeX:\\ \url{https://www.thestar.com/news/world/2016/11/09/chinese-executive-gets-death-sentence-over-tianjin-explosion-in-2015.html}
risicobeperking/controle
\bibitem{ID} ... \LaTeX:\\ \url{https://www.swissre.com/en/china/news-insights/articles/analysis-of-tianjin-port-explosion-china.html}
censuur
\bibitem{ID} ... \LaTeX:\\ \url{https://foreignpolicy.com/2015/09/10/censored-china-young-survivor-tianjin-explosion-viral-post/}
censuur
\bibitem{ID} ... \LaTeX:\\ \url{https://qz.com/756872/a-year-after-the-tianjin-blast-public-mourning-and-discussion-about-it-are-still-censored-in-china/}
verschillende artikelen
\bibitem{ID} ... \LaTeX:\\ \url{https://www.scmp.com/topics/tianjin-warehouse-explosion-2015}
\bibitem{ID} ... \LaTeX:\\ \url{https://www.wsj.com/articles/BL-CJB-27664}
\bibitem{ID} ... \LaTeX:\\ \url{https://www.nbcnews.com/news/world/tianjin-explosions-californian-witness-filmed-dramatic-china-blasts-n409701}
\bibitem{ID} ... \LaTeX:\\ \url{https://ui.adsabs.harvard.edu/abs/2016AGUFM.S13D..06P/abstract}
afwikkeling van de ramp
\bibitem{ID} ... \LaTeX:\\ \url{https://chinadialogue.net/en/pollution/9188-back-to-the-blast-zone-one-year-after-the-tianjin-explosion/}
\bibitem{ID} ... \LaTeX:\\ \url{https://www.wired.com/2015/08/chinas-huge-tianjin-explosion-looked-like-space/}
\bibitem{ID} ... \LaTeX:\\ \url{https://www.abc.net.au/news/2015-08-13/explosion-rocks-north-chinese-city-of-tianjin/6693336?nw=0}
ambtenaren onderzocht
\bibitem{ID} ... \LaTeX:\\ \url{https://thediplomat.com/2015/08/23-executives-government-officials-under-investigation-for-role-in-tianjin-explosions/}
\bibitem{ID} ... \LaTeX:\\ \url{http://america.aljazeera.com/articles/2015/8/13/at-least-50-dead-and-hundreds-injured-in-chinese-warehouse-explosion.html}
risico-inschatting
\bibitem{ID} ... \LaTeX:\\ \url{https://www.mdpi.com/2071-1050/12/3/1169/htm}
\bibitem{ID} ... \LaTeX:\\ \url{https://www.mdpi.com/2071-1050/12/3/1169/htm}
\bibitem{ID} ... \LaTeX:\\ \url{https://www.cbsnews.com/news/tianjin-port-china-massive-explosion-hundreds-injured/}
\bibitem{ID} ... \LaTeX:\\ \url{https://www.hkjcdpri.org.hk/download/casestudies/Tianjin_CASE.pdf}
\bibitem{ID} ... \LaTeX:\\ \url{https://time.com/3996168/tianjin-explosion-china-pictures/}
onderzoeksrapport
\bibitem{ID} ... \LaTeX:\\ \url{https://www.hfw.com/Tianjin-Port-explosion-August-2015}
\bibitem{ID} ... \LaTeX:\\ \url{https://news.un.org/en/story/2015/08/506912-following-tianjin-explosion-un-expert-calls-china-ensure-transparent}
\bibitem{ID} ... \LaTeX:\\ \url{https://www.france24.com/en/20150812-huge-explosions-rock-chinese-city-tianjin}
\bibitem{ID} ... \LaTeX:\\ \url{https://choice.npr.org/index.html?origin=https://www.npr.org/2015/08/14/432280627/what-caused-the-warehouse-explosions-in-tianjin-china}
123 verantwoordelijken
\bibitem{ID} ... \LaTeX:\\ \url{https://www.bbc.com/news/world-asia-china-35506311}
\bibitem{ID} ... \LaTeX:\\ \url{https://www.washingtonpost.com/gdpr-consent/?next_url=https%3a%2f%2fwww.washingtonpost.com%2fnews%2fworldviews%2fwp%2f2015%2f08%2f12%2fvideos-show-chinese-city-of-tianjin-rocked-by-enormous-explosion%2f}
lang artiekel
\bibitem{ID} ... \LaTeX:\\ \url{https://www.businessinsider.com/the-chemical-explosion-in-china-killed-more-than-100-people-and-the-devastation-is-unreal-2015-8?international=true&r=US&IR=T}
\bibitem{ID} ... \LaTeX:\\ \url{https://pubmed.ncbi.nlm.nih.gov/27311537/}
\bibitem{ID} ... \LaTeX:\\ \url{https://www.reuters.com/article/us-china-blast-insurance-idUSKCN0QM0N220150817}
\bibitem{ID} ... \LaTeX:\\ \url{https://www.sciencedirect.com/science/article/abs/pii/S0305417916300079}
\bibitem{ID} ... \LaTeX:\\ \url{https://en.wikipedia.org/wiki/2015_Tianjin_explosions}
\bibitem{ID} ... \LaTeX:\\ \url{https://www.bbc.com/news/world-asia-china-33844084}
\bibitem{ID} ... \LaTeX:\\ \url{https://www.independent.co.uk/news/world/asia/tianjin-explosion-photos-china-chemical-factory-accident-crater-revealed-a7199591.html}

veiigheidshandhaving
\bibitem{ID} ... \LaTeX:\\ \url{https://www.ilo.org/legacy/english/protection/safework/ctrl_banding/toolkit/main_guide.pdf}
\bibitem{ID} ... \LaTeX:\\ \url{https://echa.europa.eu/documents/10162/21332507/guide_chemical_safety_sme_en.pdf}
\bibitem{ID} ... \LaTeX:\\ \url{https://ec.europa.eu/taxation_customs/dds2/SAMANCTA/EN/Safety/AppendixD_EN.htm}
\bibitem{ID} ... \LaTeX:\\ \url{https://www.ilo.org/safework/info/publications/WCMS_113134/lang--en/index.htm}


%%%%%%%%%%%%%%%%%%%%%%%%%%%%%%%%%%%%%%%%%%%%%%%%%%%%%%%%%%%%%%%%%
tesla autopilot crashes


veiigheidsrisico
\bibitem{evan01042019teslaautopilotIntersection} ... \LaTeX:\\ \url{https://spectrum.ieee.org/cars-that-think/transportation/self-driving/three-small-stickers-on-road-can-steer-tesla-autopilot-into-oncoming-lane}

\bibitem{testVehicleSafetyReport} ... \LaTeX:\\ \url{https://www.tesla.com/VehicleSafetyReport}
veiligheidsrapport mbt autopilot
\bibitem{lambert31062020q2safetyreport} ... \LaTeX:\\ \url{https://electrek.co/2020/07/31/tesla-q2-2020-safety-report-strong-improvement-autopilot-accidents/}
consumentenrapport
bluetooth veiligheidsvraagstuk
\bibitem{wiredBloutoothHackTesla} ... \LaTeX:\\ \url{https://www.wired.com/story/tesla-model-x-hack-bluetooth/}
veiigheidsvraagstuk vanwege touch screen
\bibitem{preston14012021NHTSATeslaRecall} ... \LaTeX:\\ \url{https://www.consumerreports.org/car-recalls-defects/nhtsa-asks-tesla-to-recall-model-s-model-x-touch-screen-safety-issues/}
veiligheidsvraagstuk
\bibitem{cio25112020belgianTeslaHack} ... \LaTeX:\\ \url{https://cio.economictimes.indiatimes.com/news/digital-security/security-researchers-hack-steal-tesla-model-x-within-minutes/79406553}
veiligheidsvraagstuk
rapport over autopilot
\bibitem{templeton06092019HTSBReportTesla} ... \LaTeX:\\ \url{https://www.forbes.com/sites/bradtempleton/2019/09/06/ntsb-report-on-tesla-autopilot-accident-shows-whats-inside-and-its-not-pretty-for-fsd/?sh=6905e7d4dc55}
de invloed van de bestuurder bij tesla ongeluk
\bibitem{Dickey08012021SuddenTeslaAcceleration} ... \LaTeX:\\ \url{https://techcrunch.com/2021/01/08/nhtsa-tesla-sudden-unintended-acceleration-driver-error/}
veiligheidsvraagstuk
\bibitem{darkReading17112020TeslaBackup} ... \LaTeX:\\ \url{https://www.darkreading.com/threat-intelligence/security-risks-discovered-in-tesla-backup-gateway/d/d-id/1339462}
veiligheidsvraagstuk
\bibitem{leyden23032020TeslaInterfaceHack} ... \LaTeX:\\ \url{https://portswigger.net/daily-swig/web-based-attack-crashes-tesla-driver-interface}
veiigheidsvraagstuk
\bibitem{huddlestonjr03042019ChineseTeslaHack} ... \LaTeX:\\ \url{https://www.cnbc.com/2019/04/03/chinese-hackers-tricked-teslas-autopilot-into-switching-lanes.html}
veiligheidsvraagstuk
veiligheidsvraagstuk
\bibitem{heilweil26022020teslaAutopilot} ... \LaTeX:\\ \url{https://www.vox.com/recode/2020/2/26/21154502/tesla-autopilot-fatal-crashes}
rapport over ongeluk
veiligheidsvraagstuk
veiligheidsvraagstuk
\bibitem{blanco04102019NHTSATesla} ... \LaTeX:\\ \url{https://www.caranddriver.com/news/a29369387/nhtsa-tesla-safety/}
veiligheidsvraagstuk
ransomware aanval op tesla
tesla batterij is veiligheidsvraagstuk geworden
\bibitem{mitchell01072020teslabatterycooling} ... \LaTeX:\\ \url{https://www.latimes.com/business/story/2020-07-01/federal-safety-officials-probe-tesla-battery-cooling-system}
ongeluk
\bibitem{bbc26022020AutopilotCrash} ... \LaTeX:\\ \url{https://www.bbc.com/news/technology-51645566}
veiligheidsvraagstuk
veiligheidsvraagstuk
\bibitem{stumpff04052020TeslaPersonalData} ... \LaTeX:\\ \url{https://www.thedrive.com/news/33272/tesla-discarded-old-car-parts-with-customers-personal-data-passwords-report}
dodelijk ongeluk
\bibitem{levin08062018teslaautopilotsafety} ... \LaTeX:\\ \url{https://www.theguardian.com/technology/2018/jun/07/tesla-fatal-crash-silicon-valley-autopilot-mode-report}
veiligheidsvraagstuk: ransomware
veiligheidsvraagstuk: medewerker in de fout
\bibitem{cbrook06082021TeslaInsideDataThreft} ... \LaTeX:\\ \url{https://digitalguardian.com/blog/tesla-data-theft-case-illustrates-danger-insider-threat}
\bibitem{shilling25022021Tesla} ... \LaTeX:\\ \url{https://jalopnik.com/tesla-is-stopping-some-model-3-production-report-1846353323}
veiligheidsvraagstuk: hackers je systeem laten testen
verdedigen tegenover ransomware
veiligheidsrisico
prijzen omlaag
autopilot
\bibitem{randall05112019modelSurvey} ... \LaTeX:\\ \url{https://www.bloomberg.com/graphics/2019-tesla-model-3-survey/autopilot.html}
malware door een medewerker
\bibitem{ID} ... \LaTeX:\\ \url{https://www.teslarati.com/tesla-employee-fbi-thwarts-russian-cybersecurity-attack/}
dodelijk ongeluk
\bibitem{ID} ... \LaTeX:\\ \url{https://www.marketwatch.com/story/apple-engineer-killed-in-tesla-suv-crash-on-silicon-valley-freeway-was-playing-videogame-ntsb-2020-02-25}
\bibitem{fottrell03092018TeslaSecurityChecks} ... \LaTeX:\\ \url{https://www.marketwatch.com/story/nearly-100-of-teslas-stolen-in-the-us-since-2011-have-been-recovered-2018-08-10}
waarom een tesla stelen bijna onmogelijk is
\bibitem{ID} ... \LaTeX:\\ \url{https://www.welivesecurity.com/2019/03/25/white-hats-hack-tesla-keep/}
veiligheidsonderzoek
\bibitem{ID} ... \LaTeX:\\ \url{https://www.tripwire.com/state-of-security/security-data-protection/tesla-encouraging-good-faith-security-research-in-bug-bounty-program/}
softwarefout maakt diestal mogelijk
\bibitem{kirk26112020modelX} ... \LaTeX:\\ \url{https://www.bankinfosecurity.com/tesla-model-x-stolen-in-minutes-using-software-flaws-a-15462}
fouten ontdekt in onderzoek
\bibitem{ID} ... \LaTeX:\\ \url{https://www.cnet.com/roadshow/news/tesla-ev-appeal-loyalty-study/}
\bibitem{bbc24022021hyundaiBatteryFireFix} ... \LaTeX:\\ \url{https://www.bbc.com/news/technology-56156801}








tesla cloud gehacked
\bibitem{ID} ... \LaTeX:\\ \url{https://arstechnica.com/information-technology/2018/02/tesla-cloud-resources-are-hacked-to-run-cryptocurrency-mining-malware/}
\bibitem{ID} ... \LaTeX:\\ \url{https://www.motortrend.com/news/tesla-model-y-ev-safety-quality-issues-problems/}
\bibitem{ID} ... \LaTeX:\\ \url{https://securityledger.com/2019/04/hackers-remotely-steer-tesla-model-s-using-autopilot-system/}
\bibitem{ID} ... \LaTeX:\\ \url{https://www.pcmag.com/news/report-tesla-suspends-model-3-production-in-california-until-march-7}
\bibitem{ID} ... \LaTeX:\\ \url{https://www.scmp.com/business/money/article/3121173/tesla-conduct-complete-self-inspection-after-chinese-regulators}
\bibitem{ID} ... \LaTeX:\\ \url{https://www.businesswire.com/news/home/20180220005222/en/RedLock-Releases-Cloud-Security-Report-Highlighting-Focus-on-Shared-Responsibilities-Uncovers-Cloud-Related-Exposures-at-Tesla}
\bibitem{ID} ... \LaTeX:\\ \url{https://www.epa.gov/automotive-trends/highlights-automotive-trends-report}
\bibitem{ID} ... \LaTeX:\\ \url{https://www.livemint.com/Companies/o2QLbtJc9EQ7ZcpxqgFbBP/Teslas-reward-for-finding-security-bugs-Model-3.html}
\bibitem{evans11062018} ... \LaTeX:\\ \url{https://revealnews.org/blog/tesla-fired-safety-official-for-reporting-unsafe-conditions-lawsuit-says/}
\bibitem{ID} ... \LaTeX:\\ \url{https://heimdalsecurity.com/blog/security-alert-teslacrypt-4-0-unbreakable-encryption-worse-data-leakage/}
\bibitem{ID} ... \LaTeX:\\ \url{https://www.eweek.com/cloud/tesla-cloud-account-data-breach-revealed-in-redlock-security-report/}
\bibitem{hawkins22102022} ... \LaTeX:\\ \url{https://www.theverge.com/2020/10/21/21527577/tesla-full-self-driving-autopilot-beta-software-update}
\bibitem{ID} ... \LaTeX:\\ \url{file:///C:/Users/gally/Downloads/applsci-10-02749-v2.pdf}
\bibitem{gritti24062020tesladataengine} ... \LaTeX:\\ \url{https://www.braincreators.com/brainpower/insights/teslas-data-engine-and-what-we-should-all-learn-from-it}
\bibitem{ID} ... \LaTeX:\\ \url{https://bernardmarr.com/default.asp?contentID=1251}
\bibitem{ID} ... \LaTeX:\\ \url{https://arstechnica.com/cars/2019/10/how-teslas-latest-acquisition-could-accelerate-autopilot-development/}
\bibitem{bouchard07052019teslaDeepLearning} ... \LaTeX:\\ \url{https://towardsdatascience.com/teslas-deep-learning-at-scale-7eed85b235d3}
\bibitem{ID} ... \LaTeX:\\ \url{file:///C:/Users/gally/Downloads/applsci-10-02749-v2.pdf}
\bibitem{Srikanth2019teslabigdata} ... \LaTeX:\\ \url{https://www.techiexpert.com/how-tesla-is-using-artificial-intelligence-and-big-data/}
\bibitem{rangaiah25022020teslaAI} ... \LaTeX:\\ \url{https://www.analyticssteps.com/blogs/how-tesla-making-use-artificial-intelligence-its-operations}
\bibitem{marr08012018taslabigdataAI} ... \LaTeX:\\ \url{https://www.forbes.com/sites/bernardmarr/2018/01/08/the-amazing-ways-tesla-is-using-artificial-intelligence-and-big-data/?sh=5e396aa24270}
\bibitem{bdickson29072020teslalevelfive = {author},	ALTeditor = {editor},	title = {title},	date = {date},	url = {"https://bdtechtalks.com/2020/07/29/self-driving-tesla-car-deep-learning/}
\bibitem{dcruz17062022tesladesignthink} ... \LaTeX:\\ \url{https://www.mygreatlearning.com/blog/teslas-new-ai-for-self-driving-cars/}
\bibitem{ID} ... \LaTeX:\\ \url{https://www.techiexpert.com/how-tesla-is-using-artificial-intelligence-and-big-data/}





\bibitem{mcfarland22042021selfdrivingrisks} ... \LaTeX:\\ \url{https://www.cnn.com/2021/04/21/tech/tesla-full-self-driving-launch/index.html}
\bibitem{hawkins18032021fedgovinvest} ... \LaTeX:\\ \url{https://www.theverge.com/2021/3/18/22338427/tesla-autopilot-crash-michigan-nhtsa-investigation}
\bibitem{berry21042021teslacrashtexas} ... \LaTeX:\\ \url{https://www.wionews.com/technology/doctor-among-victims-of-lethal-tesla-car-crash-in-texas-378950}
\bibitem{hull23072021regulatorsaftercrash} ... \LaTeX:\\ \url{https://www.bloomberg.com/news/newsletters/2021-06-23/hyperdrive-daily-after-30-tesla-crashes-what-s-a-regulator-to-do}
\bibitem{wikiTeslaAutopilot} ... \LaTeX:\\ \url{https://en.wikipedia.org/wiki/Tesla_Autopilot}
\bibitem{nhtsaAutomatedVehiclesSafety} ... \LaTeX:\\ \url{https://www.nhtsa.gov/technology-innovation/automated-vehicles-safety}
\bibitem{dowling23042021autopilottricking} ... \LaTeX:\\ \url{https://www.caradvice.com.au/947080/elon-musk-responds-to-deadly-texas-tesla-crash-as-consumer-reports-reveals-how-autopilot-can-be-tricked/}
\bibitem{wilson19042021teslacrashregulators} ... \LaTeX:\\ \url{https://usa.streetsblog.org/2021/04/19/regulators-could-have-prevented-fatal-tesla-crash/}

\bibitem{seamans22062021aikillerap} ... \LaTeX:\\ \url{https://www.brookings.edu/research/autonomous-vehicles-as-a-killer-app-for-ai/}
\bibitem{mitchell24022020AIDataTesla} ... \LaTeX:\\ \url{https://www.latimes.com/business/story/2020-02-24/autopilot-data-secrecy}
\bibitem{denneyjdsupraFeds} ... \LaTeX:\\ \url{https://www.jdsupra.com/post/contentViewerEmbed.aspx?fid=9844cae0-aa5a-45a5-988f-7f02fa5709c1}
\bibitem{siddiqui22102020TeslaCriticism} ... \LaTeX:\\ \url{https://www.washingtonpost.com/technology/2020/10/21/tesla-self-driving/}
\bibitem{ackerman01072016TeslaImperfect} ... \LaTeX:\\ \url{https://spectrum.ieee.org/cars-that-think/transportation/self-driving/fatal-tesla-autopilot-crash-reminds-us-that-robots-arent-perfect}
\bibitem{greene04092019misuseautopilot} ... \LaTeX:\\ \url{https://thenextweb.com/news/another-tesla-owner-is-dead-because-of-autopilot}
\bibitem{michralli26112019ubserautocarcrsash} ... \LaTeX:\\ \url{https://towardsdatascience.com/another-self-driving-car-accident-another-ai-development-lesson-b2ce3dbb4444}
\bibitem{pitmann21072021wrongfullautodeath} ... \LaTeX:\\ \url{https://www.theautochannel.com/news/2021/07/21/1024631-is-it-still-wrongful-death-if-car-is-driving-itself.html}
\bibitem{stackexchange102019teslacarmistake} ... \LaTeX:\\ \url{https://ai.stackexchange.com/questions/1488/why-did-a-tesla-car-mistake-a-truck-with-a-bright-sky}
\bibitem{tasking07062017TeslaAugmentedSafety} ... \LaTeX:\\ \url{https://resources.tasking.com/p/benefits-tesla-autopilot-and-how-adas-will-save-lives}
\bibitem{griemannExaminSelfDriving} ... \LaTeX:\\ \url{https://www.jipitec.eu/issues/jipitec-9-3-2018/4806}
\bibitem{Harkey30052019SafeSystemVehicle} ... \LaTeX:\\ \url{https://static.tti.tamu.edu/conferences/traffic-safety19/presentations/lunch/harkey.pdf}
\bibitem{ID} ... \LaTeX:\\ \url{https://thepressfree.com/have-google-and-amazon-backed-the-wrong-technology/}
\bibitem{ID} ... \LaTeX:\\ \url{https://www.irishtimes.com/business/innovation/robotaxis-have-google-and-amazon-backed-the-wrong-technology-1.4626749}
\bibitem{ID} ... \LaTeX:\\ \url{https://www.afr.com/technology/how-teslas-autopilot-got-it-wrong-in-fatal-crash-20160704-gpxsje}
\bibitem{ID} ... \LaTeX:\\ \url{https://economictimes.indiatimes.com/markets/stocks/news/what-me-worry-fed-chiefs-emotional-tone-can-drive-markets-study-suggests/articleshow/84618073.cms}
\bibitem{ID} ... \LaTeX:\\ \url{https://www.ehstoday.com/safety/article/21919260/ntsb-fatal-crash-involving-tesla-autopilot-resulted-from-driver-errors-overreliance-on-automation}
\bibitem{ID} ... \LaTeX:\\ \url{https://www.vanityfair.com/news/2016/07/how-the-media-screwed-up-the-fatal-tesla-accident}


tesla crash report
\bibitem{shepardson18062021TeslaDeaths} ... \LaTeX:\\ \url{https://www.reuters.com/business/autos-transportation/us-safety-agency-says-it-has-opened-probes-into-10-tesla-crash-deaths-since-2016-2021-06-17/}
\bibitem{hawkins30062021nhtsarequiresreporting} ... \LaTeX:\\ \url{https://www.politico.com/news/2021/05/18/ntsb-tesla-owner-was-in-drivers-seat-before-april-texas-crash-489272}
\bibitem{hawkins10052021autopilotnotavailable} ... \LaTeX:\\ \url{https://www.theverge.com/2021/5/10/22429198/tesla-ntsb-texas-crash-driverless-preliminary-report}
\bibitem{szymkowski29062021nhtsaTeslaCrashReports} ... \LaTeX:\\ \url{https://www.cnet.com/roadshow/news/tesla-autopilot-nhtsa-crash-report-self-driving-car-driver-assist-system/}
\bibitem{abc1112052021AutopilotNotinTeslaCrash} ... \LaTeX:\\ \url{https://abc11.com/tesla-crash-battery-fire-national-transportation-safety-board-driverless/10619772/}
\bibitem{ankel18062021regulatorsinvestigateAutopilot} ... \LaTeX:\\ \url{https://www.businessinsider.com/tesla-autopilot-crashes-regulators-open-probes-into-30-report-2021-6?international=true&r=US&IR=T}
\bibitem{sommerfield12072021NHTSAmandateresult} ... \LaTeX:\\ \url{https://driving.ca/column/lorraine/lorraine-explains-what-the-nhtsas-self-driving-car-crash-reporting-mandate-will-find-out}
\bibitem{ID} ... \LaTeX:\\ \url{https://www.teslarati.com/tesla-model-s-crash-texas-ntsb-preliminary-report/}
\bibitem{ID} ... \LaTeX:\\ \url{https://insideevs.com/news/506498/ntsb-report-tesla-texas-crash/}
\bibitem{ID} ... \LaTeX:\\ \url{https://electrek.co/2021/06/03/tesla-tsla-crashes-report-new-orders-in-china-free-falling/}
\bibitem{ID} ... \LaTeX:\\ \url{https://www.newsy.com/stories/ntsb-releases-report-on-fatal-tesla-crash/}
\bibitem{ID} ... \LaTeX:\\ \url{https://www.ndtv.com/world-news/autopilot-not-used-in-april-tesla-crash-says-us-report-2439146}
\bibitem{ID} ... \LaTeX:\\ \url{https://www.autocar.co.nz/autocar-news-app/fatal-driverless-tesla-crash-report-shows-autopilot-not-to-blame}
\bibitem{ID} ... \LaTeX:\\ \url{https://teleperformance-waha.sabacloud.com/Saba/Web_spf/EU2PRD0152/app/dashboard}
\bibitem{ID} ... \LaTeX:\\ \url{https://www.independent.co.uk/news/world/americas/tesla-texas-crash-model-s-autopilot-b1845286.html}
\bibitem{ID} ... \LaTeX:\\ \url{https://www.wired.com/2017/01/probing-teslas-deadly-crash-feds-say-yay-self-driving/}
\bibitem{saferoardsCrashesAutonomousvehicles} ... \LaTeX:\\ \url{https://saferoads.org/wp-content/uploads/2020/03/AV-Crash-List-with-Photos-February-2020.pdf}
\bibitem{ID} ... \LaTeX:\\ \url{https://mashable.com/article/nthsa-tesla-autopilot-model-x-crash-investigation}
\bibitem{stephardson18032021revieuwingtesla} ... \LaTeX:\\ \url{https://www.usnews.com/news/top-news/articles/2021-03-18/us-safety-agency-reviewing-23-tesla-crashes-three-from-recent-weeks}
\bibitem{krishner30062021NHTSAreport} ... \LaTeX:\\ \url{https://chicago.suntimes.com/consumer-affairs/2021/6/30/22557122/nhtsa-automated-driving-crash-reports-tesla-national-highway-traffic-safety-administration}
\bibitem{gitlin11052021autopilot} ... \LaTeX:\\ \url{https://arstechnica.com/cars/2021/05/ntsb-finds-no-reason-to-suspect-autopilot-in-fatal-tesla-crash/}
\bibitem{ID} ... \LaTeX:\\ \url{https://jalopnik.com/the-ntsb-to-partially-blame-teslas-autopilot-in-fatal-c-1803136365}
\bibitem{mitchell19012017investigationstop} ... \LaTeX:\\ \url{https://www.latimes.com/business/autos/la-fi-hy-tesla-autopilot-20170119-story.html}
\bibitem{gordon10052021teslaprelimreport} ... \LaTeX:\\ \url{https://www.vice.com/en/article/z3xxaw/ntsb-releases-preliminary-report-on-tesla-crash-that-killed-two-people}
\bibitem{shaper07062018} ... \LaTeX:\\ \url{https://choice.npr.org/index.html?origin=https://www.npr.org/2018/06/07/618081406/no-driver-input-detected-in-seconds-before-deadly-tesla-crash-ntsb-finds}
\bibitem{cochran18042021nodriverTeslaCrash} ... \LaTeX:\\ \url{https://www.click2houston.com/news/local/2021/04/18/2-men-dead-after-fiery-tesla-crash-in-spring-officials-say/}
\bibitem{habib28062016NHTSATeslaReport} ... \LaTeX:\\ \url{https://static.nhtsa.gov/odi/inv/2016/INCLA-PE16007-7876.pdf}
\bibitem{firstpress11052021fatalnonautopilot} ... \LaTeX:\\ \url{https://www.firstpost.com/tech/news-analysis/tesla-model-s-involved-in-fatal-crash-in-the-us-did-not-use-autopilot-says-ntsb-report-9609911.html}
\bibitem{raynal20042021probeTeslaCrash} ... \LaTeX:\\ \url{https://www.autoweek.com/news/green-cars/a36173804/both-local-police-and-nhtsa-probe-tesla-crash/}
\bibitem{tiungteslasoftwarecrash} ... \LaTeX:\\ \url{https://www.zdnet.com/article/apple-and-tesla-under-fire-over-software-engineers-fatal-autopilot-crash/}
\bibitem{globaltimes08052021guangdongcrash,	ALTauthor = {global times},	ALTeditor = {editor},	title = {Tesla crash in Guangdong sparks new round of safety concerns},	date = {date},	url = {"https://www.globaltimes.cn/page/202105/1222902.shtml}
\bibitem{anderson30042021secondteslacrash,	ALTauthor = {Brad Anderson},	ALTeditor = {editor},	title = {NTSB To Release Report On Deadly Tesla Crash Within A Month},	date = {date},	url = {"https://www.carscoops.com/2021/04/ntsb-to-release-report-on-deadly-tesla-crash-within-a-month/}
	\bibitem{oremus21062017fatalTeslaCrash} ... \LaTeX:\\ \url{https://slate.com/technology/2017/06/a-new-report-on-what-happened-in-the-fatal-tesla-autopilot-crash.html}
	\bibitem{guardian15052021teslacrashHandsOnWheel} ... \LaTeX:\\ \url{https://www.theguardian.com/us-news/2021/may/15/tesla-fatal-california-crash-autopilot}
	\bibitem{Puzzanghera13092017TeslaSharesBlame} ... \LaTeX:\\ \url{https://www.stuff.co.nz/motoring/news/96797272/tesla-shares-some-blame-in-fatal-autopilot-crash-report}
	\bibitem{jaillet02022017teslaAutopilotLimitations} ... \LaTeX:\\ \url{https://www.overdriveonline.com/business/article/14891759/dot-report-on-fatal-2016-tesla-crash-with-tractor-trailer-blames-limitations-of-autopilot-mode}
	\bibitem{reuters03102019teslaAutoParkingFail} ... \LaTeX:\\ \url{https://venturebeat.com/2019/10/03/regulators-investigate-teslas-automated-parking-feature-following-crash-reports/}
	\bibitem{dowling23042021} ... \LaTeX:\\ \url{https://www.caradvice.com.au/947080/elon-musk-responds-to-deadly-texas-tesla-crash-as-consumer-reports-reveals-how-autopilot-can-be-tricked/}
	\bibitem{young05112021fatalTeslaReport} ... \LaTeX:\\ \url{https://www.consumeraffairs.com/news/ntsb-releases-report-on-fatal-tesla-crash-in-texas-051121.html}
	\bibitem{kierstein18032021teslaAutopilotCrashStationary} ... \LaTeX:\\ \url{https://www.motortrend.com/news/tesla-michigan-state-autopilot-crash-report/}
	\bibitem{janssen20062017teslacrashdetailflorida} ... \LaTeX:\\ \url{https://tweakers.net/nieuws/126145/onderzoeksraad-vs-publiceert-technische-details-tesla-crash-in-florida.html}
	\bibitem{ID} ... \LaTeX:\\ \url{https://gizmodo.com/new-documents-reveal-one-driver-s-agony-and-confusion-d-1841720801}
	
	
	\bibitem{ID} ... \LaTeX:\\ \url{https://www.google.com/search?q=tesla+crash+report&rlz=1C1AVUC_enNL953NL953&ei=p3kNYa6sLI_UsAeSoZrwDw&start=100&sa=N&ved=2ahUKEwjum77s_ZzyAhUPKuwKHZKQBv44WhDw0wN6BAgBEEg&biw=1920&bih=933}
	
	
	
	%%%%%%%%%%%%%%%%%%%%%%%%%%%%%%%%%%%%%%%%%%%%%%%%%%%%%%%%%%%%%%%%%
	
	vlucht 1951
	\bibitem{ID} ... \LaTeX:\\ \url{https://nl.wikipedia.org/wiki/Turkish_Airlines-vlucht_1951}
	technisch rapport
	\bibitem{ID} ... \LaTeX:\\ \url{file:///C:/Users/gally/Downloads/rapport_ta_nl_aangepast.pdf}
	beschrijving
	terugblik met overlevenden
	tijdlijn
	\bibitem{zuilen23022019Tijdlijnpoldercrash} ... \LaTeX:\\ \url{https://www.noordhollandsdagblad.nl/cnt/dmf20190221_65390940}
	artikel
	terugblik met overlevenden
	advies raad voor de veiligheid
	de overlevende, de oorzaak, regeling, herdenking, smartengeld
	verhaal van een overlevende
	herdenking
	herdenking
	bemanning deed niets met foutmelding
	parlementaire besluitenlijst
	kamervragen over de onafhankelijkheid van de raad voor veiligheid
	verhaal van een overlevende
	beschrijvend artikel van letsel en gewonden
	\bibitem{ntvg20012010letsel} ... \LaTeX:\\ \url{https://www.ntvg.nl/artikelen/vliegtuigongeval-schiphol-25-02-2009-letsels-en-verdeling-van-gewonden}
	technische fout als oorzaak
	\bibitem{wikinews04032009techfoutailines1951} ... \LaTeX:\\ \url{https://nl.wikinews.org/wiki/Technische_fout_oorzaak_vliegtuigcrash_Turkish_Airlines-vlucht_1951}
	gesprek met pieter van vollenhove voorzitter van de onderzoeksraad voor veiligheid
	onderzoeksraad voor veiligheid is onderdruk gezet
	\bibitem{luchtvaartnieuws21012020boeing737conclusies} ... \LaTeX:\\ \url{https://www.luchtvaartnieuws.nl/nieuws/categorie/72/algemeen/conclusies-crash-tk1951-na-amerikaanse-druk-afgezwakt}
	niuwesartikel
	feitenverloop
	\bibitem{adformatie280220209communicatiegebreken} ... \LaTeX:\\ \url{https://www.adformatie.nl/contentmarketing/communicatie-na-vliegramp-vertoonde-gebreken}
	zwarte doos
	\bibitem{spinnael25022009onderzoekpolderbaancrash} ... \LaTeX:\\ \url{https://flightlevel.be/244/onderzoek-polderbaan-crash-turkish-airlines-1951/}
	\bibitem{crashTurkishAirlines} ... \LaTeX:\\ \url{http://wikimapia.org/11633002/nl/Crash-Turkish-Airlines-vlucht-1951}
	\bibitem{flightradar24} ... \LaTeX:\\ \url{https://www.flightradar24.com/data/flights/tk1951}
	\bibitem{flightstatstracker} ... \LaTeX:\\ \url{https://www.flightstats.com/v2/flight-tracker/TK/1951}
	
	
	
	
	
	
	
	
	
	
	%%%%%%%%%%%%%%%%%%%%%%%%%%%%%%%%%%%%%%%%%%%%%%%%%%%%%%%%%%%%%%%%%
	de mali missie
	
	
	\bibitem{bnnvara13062018malirapport} ... \LaTeX:\\ \url{https://joop.bnnvara.nl/nieuws/rapport-haalbaarheid-en-houdbaarheid-van-mali-missie-twijfelachtig}
	\bibitem{eucal11012021malimissieverlengd} ... \LaTeX:\\ \url{https://www.consilium.europa.eu/nl/press/press-releases/2021/01/11/eucap-sahel-mali-mission-extended-until-31-january-2023-and-mandate-adjusted/}
	\bibitem{nos21052014zorgenmalimissie} ... \LaTeX:\\ \url{https://nos.nl/artikel/650637-kamer-bezorgd-over-mali-missie}
	\bibitem{meijnders} ... \LaTeX:\\ \url{https://www.bnr.nl/nieuws/10015679/koenders-positief-tegenover-verlening-mali-missie}
	\bibitem{bnrwebredactie} ... \LaTeX:\\ \url{https://www.bnr.nl/nieuws/politiek/10345553/kabinet-wil-mali-missie-stoppen-verrassend-besluit}
	\bibitem{keultjes01062016malimissiecoalitie} ... \LaTeX:\\ \url{https://www.ad.nl/nieuws/clash-om-mali-missie-dreigt-binnen-coalitie~a4151d4f/}
	\bibitem{veenhof18012019} ... \LaTeX:\\ \url{https://www.nd.nl/cultuur/boeken/536861/boek-kijkje-bij-de-mali-missie}
	\bibitem{ID} ... \LaTeX:\\ \url{https://www.youtube.com/watch?v=jmZ6uSbpCvg}
	\bibitem{isitman06012016militair} ... \LaTeX:\\ \url{https://www.ewmagazine.nl/nederland/achtergrond/2016/07/twee-nederlanse-militairen-dood-bij-oefening-mali-missie-325226/}
	\bibitem{nporadio11072016filmdemissie} ... \LaTeX:\\ \url{https://www.nporadio1.nl/nieuws/cultuur-media/9e3b076e-5401-4630-bf39-f925213c5b6b/onverwachte-openhartigheid-over-missie-in-mali}
	\bibitem{parlementairmonitor15122013mortierongeluk} ... \LaTeX:\\ \url{https://www.parlementairemonitor.nl/9353000/1/j9vvij5epmj1ey0/vjfm5p0nujzw?ctx=vj2mc67lofnr}
	
	
	sollicitatie
	de bureaucratie
	aankomst
	interview van de burgerbevolking
	steun van de bevolking minuut 15:00
	de organisatie minuut 23:00
	De militaire briefing minuut 34:00
	prioriteit minuut 39:00
	briefing minuut 40:00
	de communicatie met ministerie over inlichten minuut 44:00
	\bibitem{ID} ... \LaTeX:\\ \url{https://www.2doc.nl/documentaires/series/2doc/2016/juli/de-missie.html}
	
	%%%%%%%%%%%%%%%%%%%%%%%%%%%%%%%%%%%%%%%%%%%%%%%%%%%%%%%%%%%%%%%%%
	militair overleden door schietoefening in ossendrecht
	
	
	\bibitem{nos22032016ossendrecht} ... \LaTeX:\\ \url{https://amp.nos.nl/artikel/2094524-militair-omgekomen-bij-schietoefening-ossendrecht.html}
	
	\bibitem{ovv04042016lessenongevalossendrecht} ... \LaTeX:\\ \url{https://www.onderzoeksraad.nl/nl/page/4293/lessen-uit-schietongeval-ossendrecht}
	
	\bibitem{quekelboere10052017doodossendrecht} ... \LaTeX:\\ \url{https://www.bndestem.nl/bergen-op-zoom/dood-van-militair-sander-klap-35-in-ossendrecht-was-ongeluk-militairen-vrijuit-hij-probeerde-zijn-leven-te-redden~afe4c7a0/}
	
	
	Wat is de rol van defensie?
	Wat is er gedaan om de veligheid van de medewerkers te waarborgen?
	Waarom zijn deze regels niet nageleefd?
	Wat zijn de gevolgen?
	Zijn de acties die naderhand zijn ondernomen wel redelijk naar de slachtoffers, het nationale veiligheisbeeld en de medewerkers?
	
	
	
	
	
	
	
	%%%%%%%%%%%%%%%%%%%%%%%%%%%%%%%%%%%%%%%%%%%%%%%%%%%%%%%%%%%%%%%%%
	schipholbrand
	
	Wat is er gebeurd?
	\bibitem{wikiSchipholbrand} ... \LaTeX:\\ \url{https://nl.wikipedia.org/wiki/Schipholbrand}
	artikel
	\bibitem{schipholbrand27102005video} ... \LaTeX:\\ \url{https://www.youtube.com/watch?v=1i-hfEzxFfk}
	psychologische gevolgen
	rapport
	\bibitem{onderzoeksraad2610schipholoost} ... \LaTeX:\\ \url{https://www.onderzoeksraad.nl/nl/page/392/brand-cellencomplex-schiphol-oost-nacht-van-26-op-27-oktober}
	artikel met video
	herdenking
	impact op de persoon
	herdenking
	\bibitem{schipholbrandvideoargos} ... \LaTeX:\\ \url{https://www.vpro.nl/argos/speel~POMS_VPRO_461907~schadevergoeding-voor-ex-verdachte-schipholbrand~.html}
	chronologie
	\bibitem{nunl30052023feitenoverzicht} ... \LaTeX:\\ \url{https://www.nu.nl/binnenland/3355935/feitenoverzicht-schipholbrand-en-rechtszaken.html}
	tijdlijn
	\bibitem{ID} ... \LaTeX:\\ \url{https://www.singeluitgeverijen.nl/isbn/de-schipholbrand/}
	vervolgens van ministers
	beeldanalyse en reconstructie
	\bibitem{ID} ... \LaTeX:\\ \url{https://eenvandaag.avrotros.nl/item/schipholbrand-niet-ontstaan-in-cel-11/}
	herdenking
	korte samenvatting
	rapport
	artikel
	verwijzing naar het rapport vanuit de politieke oppositie
	beeld vanuit de gevangenisbewaarder
	nationaliteit slachtoffers schipholbrand
	verblijfsvergunning voor de slachtoffers
	gen schadevergoeding voor de verdachte
	verdachte voor de rechter
	geen schadevergoeding voor verdachte
	artikel wat ging er mis bji de schipholbrand
	brand veroorzaakt door een peuk
	smaadschrift
	bewakers worden niet vervolgd
	proces schipholbrand moet over en de brandveilgheid moet worden verbeterd
	de rol van het parlement in de evaluatie
	\bibitem{parlementairemonitorschipholbrand} ... \LaTeX:\\ \url{https://www.parlementairemonitor.nl/9353000/1/j9vvij5epmj1ey0/vi3aof7awcxg}
	onderzoeksmemo
	herdenking
	\bibitem{ID} ... \LaTeX:\\ \url{https://archief.ntr.nl/nova/page/detail/uitzendingen/3847/Den%20Haag%20Vandaag_%20herdenking%20Schipholbrand.html}
	herdenking
	invloed van de ramp op samenleving
	\bibitem{videonpoNOVA13112008} ... \LaTeX:\\ \url{https://www.npostart.nl/heropen-onderzoek-schipholbrand/13-11-2008/POMS_NTR_103332}
	opmerkelijk rapport gestolen in de nasleep
	\bibitem{schipperSchipholbrand} ... \LaTeX:\\ \url{https://www.nd.nl/nieuws/nederland/600395/schipholbrand-blijft-schrijnen}
	\bibitem{ID} ... \LaTeX:\\ \url{https://www.ed.nl/economie/om-geen-schadevergoeding-voor-verdachte-schipholbrand~a6c7c51d/63042600/?referrer=https%3A%2F%2Fwww.google.com%2F}
	\bibitem{ID} ... \LaTeX:\\ \url{https://www.groene.nl/artikel/schipholbrand-vereist-debat}
	\bibitem{rizoomes01052014schipholbrand} ... \LaTeX:\\ \url{https://www.rizoomes.nl/brandweer/brand-cellencomplex-schiphol/}
	
	
	
	publicaties
	\bibitem{heuvelkroesschipholbrandcamerabeelden} ... \LaTeX:\\ \url{http://www.msnp.nl/downloads/Onderzoeksmemo%20beeldanalyse%20Schipholbrand%20prot.pdf}
	\bibitem{ID} ... \LaTeX:\\ \url{http://www.dakweb.nl/roofs/2006-10/RH10-P30-31.pdf}
	\bibitem{ID} ... \LaTeX:\\ \url{https://www.delta.tudelft.nl/article/dood-door-zuinigheid}
	\bibitem{ID} ... \LaTeX:\\ \url{https://www.onderzoeksraad.nl/nl/page/392/brand-cellencomplex-schiphol-oost-nacht-van-26-op-27-oktober}
	Wat waren de regels destijds?
	Waren de autoriteiten in staat om op tijd in te grijpen of om erger te voorkomen?
	Wat is er gedaan om de veiligheid van illegalen en gevangenissbewaarders te verbeteren
	
	
	
	%%%%%%%%%%%%%%%%%%%%%%%%%%%%%%%%%%%%%%%%%%%%%%%%%%%%%%%%%%%%%%%%%
	vuurwerkramp enschede
	\bibitem{ID} ... \LaTeX:\\ \url{https://www.youtube.com/watch?v=OMkIsj8FsHw}
	\bibitem{ID} ... \LaTeX:\\ \url{https://depot03.archiefweb.eu/archives/archiefweb/20210703085353/http://www.vuurwerkramp.enschede.nl/publicaties/00005/#.YOAlp-gzaUk}
	
	Wat waren de afspraken omtrent vuurwerkopslag?
	Waarom werden de voorschriften neit nageleefd?
	
	
	%%%%%%%%%%%%%%%%%%%%%%%%%%%%%%%%%%%%%%%%%%%%%%%%%%%%%%%%%%%%%%%%%
	
	
	explosie in beirut
	\bibitem{landryalameddine12112020beiruthelathsystem} ... \LaTeX:\\ \url{https://bmchealthservres.biomedcentral.com/articles/10.1186/s12913-020-05906-y}
	\bibitem{ID} ... \LaTeX:\\ \url{https://news.sky.com/story/beirut-blast-cctv-captures-moment-huge-explosion-devastated-hospital-12047452}
	\bibitem{ID} ... \LaTeX:\\ \url{https://www.unodc.org/unodc/en/frontpage/2020/September/unodc-assists-lebanon-in-reestablishing-container-shipments-in-the-aftermath-of-the-port-of-beirut-explosion.html}
	\bibitem{ID} ... \LaTeX:\\ \url{https://reliefweb.int/sites/reliefweb.int/files/resources/LEB201-Lebanon-Emergency-Response.pdf}
	\bibitem{yadav07082020handlingexplosivesBeirut} ... \LaTeX:\\ \url{https://www.downtoearth.org.in/news/governance/beirut-blast-lessons-time-for-india-to-strengthen-handling-of-explosives-chemicals-72707}
	\bibitem{graham21082020rootsImpactBeirutBlast} ... \LaTeX:\\ \url{https://www.justsecurity.org/72122/the-cost-of-resilience-the-roots-and-impacts-of-the-beirut-blast/}
	\bibitem{ID} ... \LaTeX:\\ \url{https://www.fire-magazine.com/the-port-of-beirut-explosion-a-timely-reminder}
	\bibitem{neusaeter07082020beirutexplosioneval} ... \LaTeX:\\ \url{https://www.ctvnews.ca/sci-tech/mapping-the-beirut-explosion-what-the-impact-would-look-like-in-canadian-cities-1.5053932}
	
	
	secyrity:
	\bibitem{ID} ... \LaTeX:\\ \url{https://permanent.fdlp.gov/gpo45474/AN_advisory.pdf}
	
	
	secyrity:
	\bibitem{ID} ... \LaTeX:\\ \url{https://permanent.fdlp.gov/gpo45474/AN_advisory.pdf}
	
	
	
	
	
	
	
	
	%%%%%%%%%%%%%%%%%%%%%%%%%%%%%%%%%%%%%%%%%%%%%%%%%%%%%%%%%%%%%%%%%
	bijlmerramp
	
	
	
	
	%%%%%%%%%%%%%%%%%%%%%%%%%%%%%%%%%%%%%%%%%%%%%%%%%%%%%%%%%%%%%%%%%
	slmramp
	Wat is er gebeurd?
	\bibitem{ID} ... \LaTeX:\\ \url{https://www.srnieuws.com/suriname/290721/slm-ramp-herdacht/}
	\bibitem{ID} ... \LaTeX:\\ \url{https://werkgroepcaraibischeletteren.nl/documentaire-waarom-nou-jij-over-de-slm-ramp-in-89/}
	\bibitem{ID} ... \LaTeX:\\ \url{https://www.vpro.nl/speel~WO_NTR_15390142~andere-tijden-17-apr-2019-3-09-min-fouten-en-misstanden-leiden-tot-de-slm-ramp~.html}
	\bibitem{ID} ... \LaTeX:\\ \url{https://www.canonvannederland.nl/nl/kalender/06/1989-06-07}
	\bibitem{ID} ... \LaTeX:\\ \url{https://vijfeeuwenmigratie.nl/archief-herdenkingen-slm-ramp}
	\bibitem{ID} ... \LaTeX:\\ \url{https://www.hulpverleningsforum.nl/index.php?topic=84702.0}
	\bibitem{ID} ... \LaTeX:\\ \url{https://www.nporadio1.nl/fragmenten/focus/f792e720-bd85-4c18-8a71-b334d9d5de7e/2019-04-17-slm-ramp-een-paar-cowboys-hebben-achter-de-stuurknuppel-gezeten}
	\bibitem{ID} ... \LaTeX:\\ \url{https://www.waterkant.net/suriname/2017/06/07/herdenking-slm-ramp-28-jaar-geleden-suriname/}
	\bibitem{espnSLMterugblik} ... \LaTeX:\\ \url{https://www.espn.nl/video/clip?id=8744942}
	\bibitem{ID} ... \LaTeX:\\ \url{http://www.themediabrothers.nl/tag/slm-ramp/}
	\bibitem{ID} ... \LaTeX:\\ \url{https://www.rijnmond.nl/nieuws/182546/30-jaar-na-de-SLM-ramp-Ik-mis-mijn-broer-nog-elke-dag}
	\bibitem{dennisRosier01052020} ... \LaTeX:\\ \url{https://www.voetbalkrant.com/nieuws/2020-05-01/het-vergeten-verhaal-van-de-slm-ramp}
	\bibitem{hassing07062020slmramp} ... \LaTeX:\\ \url{https://www.bd.nl/sport/de-slm-ramp-en-het-hartverscheurende-verhaal-van-jerry-en-winnie-haatrecht~ae4ce105/?referrer=https%3A%2F%2Fwww.google.com%2F}
	\bibitem{amsterdamArchiefSLM} ... \LaTeX:\\ \url{https://www.amsterdam.nl/stadsarchief/nieuws/slm-ramp/}
	\bibitem{rtvOost06062019nabestaande} ... \LaTeX:\\ \url{https://www.rtvoost.nl/nieuws/313496/Nabestaande-SLM-ramp-Heb-ik-wel-mijn-broer-en-moeder-begraven}
	\bibitem{breda07062021AndroSnel} ... \LaTeX:\\ \url{https://www.bredavandaag.nl/nieuws/algemeen/337919/nac-herdenkt-andro-knel-slm-ramp-precies-32-jaar-geleden}
	\bibitem{andereTijdenSLMCrash} ... \LaTeX:\\ \url{https://www.anderetijden.nl/aflevering/792/Een-aangekondigde-vliegramp}
	\bibitem{aviationSLMCrashAccidentInvestigation} ... \LaTeX:\\ \url{https://nl.wikipedia.org/wiki/SLM-ramp}
	database
	\bibitem{aviationReport} ... \LaTeX:\\ \url{https://aviation-safety.net/database/record.php?id=19890607-2}
	rapport
	\bibitem{aviationSLMCrashAccidentInvestigation} ... \LaTeX:\\ \url{https://reports.aviation-safety.net/1989/19890607-2_DC86_N1809E.pdf}
	\bibitem{mcDonnelDouglasCommissionReportSLMCrash} ... \LaTeX:\\ \url{https://aviation-safety.net/investigation/cvr/transcripts/cvr_py764.php}
	\bibitem{wikiSRFlight764} ... \LaTeX:\\ \url{https://en.wikipedia.org/wiki/Surinam_Airways_Flight_764}
	\bibitem{ID} ... \LaTeX:\\ \url{https://web.archive.org/web/20050113010822/https://www.ntsb.gov/ntsb/brief.asp?ev_id=34510&key=0}
	\bibitem{nos07062019SLMTerugblik} ... \LaTeX:\\ \url{https://nos.nl/artikel/2287986-slm-vliegramp-van-precies-30-jaar-geleden-trof-ook-nederlands-voetbal}
	\bibitem{dagvantoenSLMCrash} ... \LaTeX:\\ \url{https://www.dagvantoen.nl/vliegtuigcrash-slm-bij-zanderij-meer-dan-170-doden/}
	\bibitem{waterkantNesty07061989} ... \LaTeX:\\ \url{https://www.waterkant.net/suriname/2006/06/07/vliegramp-suriname-op-7-juni-1989-2/}
	uitgebreid engels artikel
	\bibitem{eduNandlalSRCrash} ... \LaTeX:\\ \url{http://www.edufd.nl/planecrash/}
	ntsb investigtion
	\bibitem{oldjetsSRAirways} ... \LaTeX:\\ \url{http://www.oldjets.net/slm-dc-8-crash.html}
	uitgebreid engels artikel
	\bibitem{cloudberg02012021srflight764} ... \LaTeX:\\ \url{https://admiralcloudberg.medium.com/contract-to-kill-the-crash-of-surinam-airways-flight-764-828979c7efe2}
	persbericht
	\bibitem{apnews07061989srplanecrash} ... \LaTeX:\\ \url{https://apnews.com/article/5b240d758ee4c5422381cc7cdc98566b}
	Wat is de rol van de autoriteiten?
	Welke andere betrokkeen? Enw at is hun verantwoordelijkheid
	Hadden de negatieve gevolgen voorkomen kunnen worden?
	Hoe werd er over veiligheid gedacht?
	
	
	
	
	%%%%%%%%%%%%%%%%%%%%%%%%%%%%%%%%%%%%%%%%%%%%%%%%%%%%%%%%%%%%%%%%%
	Tsjernobyl
	\bibitem{ID} ... \LaTeX:\\ \url{https://www.youtube.com/watch?v=Xw3SFOfbR84}
	\bibitem{wikiTjernobyl} ... \LaTeX:\\ \url{https://nl.wikipedia.org/wiki/Kernramp_van_Tsjernobyl}
	\bibitem{rivmTjernobyl} ... \LaTeX:\\ \url{https://www.rivm.nl/straling-en-radioactiviteit/stralingsincidenten-en-kernongevallen/tsjernobyl}
	\bibitem{andereTijdenTjernobyl} ... \LaTeX:\\ \url{https://www.anderetijden.nl/aflevering/599/Tsjernobyl-als-Nederlandse-ramp}
	wat er is gebeurd en hoe het leven verdergaat
	\bibitem{kingskey19042022tjernobyl} ... \LaTeX:\\ \url{https://www.nationalgeographic.nl/het-leven-in-tsjernobyl-gaat-door}
	pernsioenfondsen en de tjernobyl ramp
	In 2021 worden mensen nog steeds blootgesteld blijkt ut een gezamelijk onderzoek van greenpeace en oekraiense wetenschappers
	stijging van de nucliaire activiteit gemeten in tjernobyl
	Het toerisme  aspect
	De chronologie
	\bibitem{erikbork26042023reactor4} ... \LaTeX:\\ \url{https://historianet.nl/maatschappij/rampen/tsjernobyl-atoomhel-bij-reactor-4}
	\bibitem{nosTjernobyl30jaarlater} ... \LaTeX:\\ \url{https://nos.nl/artikel/2101523-de-spookstad-van-tsjernobyl-30-jaar-later}
	Dieren in de omgeving van tjernobyl
	De chronologie
	Echtreme droogte zorgd voor gevaar
	\bibitem{knmi04052021tjernobylbosbrand} ... \LaTeX:\\ \url{https://www.knmi.nl/over-het-knmi/nieuws/35-jaar-na-tsjernobyl-liggen-branden-op-de-loer}
	\bibitem{dodonovaKVIRisicoTjernobyl} ... \LaTeX:\\ \url{https://www.kivi.nl/afdelingen/risicobeheer-en-techniek/columns/kernramp-tsjernobyl-het-dilemma-van-scherbitsky}
	Joernalistiek, entertainment en de waarheid
	\bibitem{dumarey04062020verhaalTjernobylWaarheid} ... \LaTeX:\\ \url{https://www.vrt.be/vrtnws/nl/2020/04/06/in-de-ban-van-tsjernobyl-vooruitblik/}
	Een onderzoek
	
	Huidige gevolgen van de explosie van toen
	\bibitem{sparkesNewScientistTjernoby;} ... \LaTeX:\\ \url{https://www.newscientist.nl/nieuws/steeds-meer-kernreacties-in-ontoegankelijke-ruimte-in-tsjernobyl/}
	De ramp, hoe de mensen ermee omgingen en hoe er nu geleef wordt
	
	evaluatieonderzoek en amatregeen
	\bibitem{kernenergiened26041986chronologiemaatregelen} ... \LaTeX:\\ \url{https://www.kernenergieinnederland.nl/node/308}
	\bibitem{mapszoneReactor} ... \LaTeX:\\ \url{https://www.google.com/maps/d/u/0/viewer?ie=UTF8&hl=nl&t=h&msa=0&ll=51.388923%2C30.099792&spn=0.685583%2C1.645203&z=9&source=embed&mid=1MLcOcMK_WrIJYMuTf0VVuYnMqQI}
	Invloed van de mens op de omgeving
	\bibitem{ID} ... \LaTeX:\\ \url{https://www.animalstoday.nl/mens-schadelijker-natuur-tsjernobyl/}
	Heroplevende splijtingsreacties
	docu van schooltv
	Radioactiviteit bereikt nederland
	documentaire en maatregelen
	\bibitem{kernhistoriek15062021tjernobyl} ... \LaTeX:\\ \url{https://historiek.net/kernramp-van-tsjernobyl-1986/8769/}
	Het verhaal van een overledende
	Toerisme
	toerisme
	toerisme
	Dieren in de omgevong
	Toevluchtsoord voor vluchtelingen van de oorlog met russische seperatisten
	Ouderen die terugkeerden naar hun woonplaats na de gedwongen verhuizing door de autoriteiten
	De straling neemt weer toe
	Lessen geleerd van tjernobyl
	\bibitem{nucleairforumFeitenTjernobyl} ... \LaTeX:\\ \url{https://www.nucleairforum.be/thema/veiligheid-als-prioriteit/tsjernobyl-de-feiten}
	Toerisme
	Bosbrand in tjernobyl
	invloed van de ramp op belgie
	\bibitem{kernongevalTjernobylFancGov} ... \LaTeX:\\ \url{https://fanc.fgov.be/nl/noodsituaties/zware-ongevallen-het-buitenland/1986-kernongeval-tsjernobyl}
	Boek recensie
	Fotos en berekeningen
	ontmanteling en toerisme
	Belangrijke lessen en overeenkomsten
	De journalistieke waarheid van de koude oorlog
	De lessen van
	\bibitem{arendswolters062019lessenTjernobyl} ... \LaTeX:\\ \url{https://magazines.autoriteitnvs.nl/nieuwsbrief-anvs/2019/02/de-lessen-van-tsjernobyl}
	Een toristenattractie maken van tjernobyl
	De radioactieve straling toen en nu
	de 30km zone door de ogen van toeristen
	artikel
	stedentrip
	rapport
	\bibitem{damveld08052020tjernobyl} ... \LaTeX:\\ \url{https://wisenederland.nl/wp-content/uploads/2020/06/TSJERNOBYL.pdf}
	slapend monster
	docu
	krantenartikel
	hbo serie
	docuserie
	de  nieuwe sacrofaag
	hulp aan slachtoffers
	slapende reactor
	krantenartikel
	\bibitem{deVriestjernobylHolland} ... \LaTeX:\\ \url{https://onh.nl/verhaal/besmette-melk-en-radioactieve-spinazie-tsjernobyl-in-holland}
	hbo serie
	internationale gevolgen
	toerisme
	nieuwe koepel
	media communicatie
	docu
	dieren
	\bibitem{ID} ... \LaTeX:\\ \url{https://www.amboanthos.nl/boek/nacht-in-tsjernobyl/}
	koepel
	koepel
	\bibitem{ing3enieur29042015antistralingskoepel} ... \LaTeX:\\ \url{https://www.deingenieur.nl/artikel/nieuwe-antistralingskoepel-tsjernobyl-bijna-af}
	toerisme
	toeristisch reiperspectief
	toerisme
	niwe koepel
	overschakelen naar duurzaamheid
	docu
	tjernobyl wekt nu duurazme energie
	toerisme
	overeenkomsten tjernobyl en fukushima
	drank en sla uit tjernobyl
	geen efficiente opslag is mogelijk
	
	wetenschappelijke artikelen
	
	zaterdag 26 april 1986. Er vind routineonderhoud plaats bij reactor 4, De controle wordt uitegevoerd door de dagploeg. Vnwege een test wordt jhet koelsysteem uitgeschakeld. Door omstandigheden wordt de test uitgesteld en wordt de verantwoordelijkheid overgedragen aan de avondploeg.
	De operator maakt bedieningsfouten waardoot de reactor bijna stil komt te liggen. En vervolgens probeert hij de reactor weer op gang te brengen. ondanks de snelle temperatuurstijging wordt het experiment doorgezet. Dan wordt ook het veiligheidssysteem stilgelgd. Terwijl het koelwater langzaam opwarmt, sluit hij de klep waarlangs de stoom naar de generator stroomt.
	
	De temperatuur van de reactorstaven neemt daarna snel toe. Terwijl er een oncontroleerbare kettingreactie op gang komt, laat het personeel in paniek de regelstaven zakken om de warmteontwikkeling af te remmen. Het is dan echter al te laat. Door een ontwerpfout loopt het vermogen razendsnel op tot 33.000 megawatt, ruim tien keer hoger dan normaal.
	
	In een oogwenk verandert al het koelwater in stoom. De ontploffing die daarop volgt, blaast het 2000 ton zware deksel van de reactor af.}

In de ravage vat het gloeiend hete grafiet in de reactor spontaan vlam. De uitslaande brand en een tweede explosie voeren een radioactieve rookwolk tot 8 kilometer hoogte. 
In een poging het vuur in reactor 4 te doven, storten helikopters vanuit de lucht zand, lood en boorzuur in de reactorkern. Het mag echter niet baten.

Intussen is de nucleaire brandstof zo heet geworden dat die door de bodem van het reactorvat dreigt te smelten. Als dat gebeurt, kan het bluswater onder het vat in één klap verdampen en dreigt een derde explosie die een groot deel van Europa onbewoonbaar zal maken. Om dit te voorkomen moet het water hoe dan ook worden weggepompt.

Drie brandweermannen wagen zich daarvoor in de ruimte onder de reactor, blootgesteld aan 300 sievert per uur, 300.000 keer de dosis die een Nederlander jaarlijks maximaal mag oplopen. Ze slagen daarin, maar twee van hen overlijden enkele dagen later aan acute stralingsziekte.

Hoewel geigertellers de dag na de ramp onrustbarende waarden aangeven, slaat het plaatselijk bestuur geen alarm. De bevolking is het niet gewend om vragen te stellen.

De volgende dag blijkt er wel degelijk iets ernstigs aan de hand te zijn. In een lange rij bussen worden de 135.000 inwoners op 27 april uit het besmette gebied geëvacueerd, om er nooit meer terug te keren.

De ramp is dan nog steeds geen wereldnieuws. De Sovjetautoriteiten blijken er niet eens van op de hoogte te zijn – president Gorbatsjov klaagt later dat hij via Zweden aan zijn informatie moest komen.


\bibitem{verschuur14012013tjernobylreports} ... \LaTeX:\\ \url{http://essay.utwente.nl/63353/1/Verschuur,_W._-_s0123617_(verslag).pdf}
\bibitem{paperlessarchivesTjernobyl} ... \LaTeX:\\ \url{https://www.paperlessarchives.com/chernobyl_nuclear_accident_doc.html}
\bibitem{vargos082000tjernobylconcerns} ... \LaTeX:\\ \url{https://www.pnnl.gov/main/publications/external/technical_reports/pnnl-13294.pdf}
\bibitem{mauroNuclearRiskSociety} ... \LaTeX:\\ \url{http://www.geocities.ws/scannapuerci/demauroinnovation.pdf}
\bibitem{vienna06092005LookingBack} ... \LaTeX:\\ \url{https://www-pub.iaea.org/MTCD/publications/PDF/Pub1312_web.pdf}



%%%%%%%%%%%%%%%%%%%%%%%%%%%%%%%%%%%%%%%%%%%%%%%%%%%%%%%%%%%%%%%%%

MH17




%%%%%%%%%%%%%%%%%%%%%%%%%%%%%%%%%%%%%%%%%%%%%%%%%%%%%%%%%%%%%%%%%
oekraine powergrid
\bibitem{owens21032017ukrainemitigationstrategies} ... \LaTeX:\\ \url{https://na.eventscloud.com/file_uploads/aed4bc20e84d2839b83c18bcba7e2876_Owens1.pdf}
\bibitem{ID} ... \LaTeX:\\ \url{https://www.us-cert.gov/ics/alerts/IR-ALERT-H-16-056-01}

\bibitem{cerulus2019FrontlineRussiaAttack} ... \LaTeX:\\ \url{https://www.politico.eu/article/ukraine-cyber-war-frontline-russia-malware-attacks/}
\bibitem{ID} ... \LaTeX:\\ \url{https://en.wikipedia.org/wiki/December_2015_Ukraine_power_grid_cyberattack}
\bibitem{grammatikis2019AttackIEC6087505104} ... \LaTeX:\\ \url{https://www.researchgate.net/publication/333671061_Attacking_IEC-60870-5-104_SCADA_Systems}
\bibitem{ID} ... \LaTeX:\\ \url{https://ris.utwente.nl/ws/files/6028066/3-s2_0-B9780128015957000227.pdf}
\bibitem{hidajat2016ScadaSimulator} ... \LaTeX:\\ \url{https://www.diva-portal.org/smash/get/diva2:1046339/FULLTEXT01.pdf}
\bibitem{ID} ... \LaTeX:\\ \url{https://www.semanticscholar.org/paper/Cybersecurity-analysis-of-a-SCADA-system-under-and-Rocha/dfa7c12551ebe7b24da8d806e87e946051a57cb9}
\bibitem{ID} ... \LaTeX:\\ \url{https://tutcris.tut.fi/portal/files/16294332/jafary_1534.pdf}
\bibitem{ID} ... \LaTeX:\\ \url{http://blog.nettedautomation.com/2017/}
\bibitem{uscert20072021crashmalware} ... \LaTeX:\\ \url{https://www.us-cert.gov/ncas/alerts/TA17-163A}
\bibitem{zetter12062017malwareanalysis} ... \LaTeX:\\ \url{https://www.vice.com/en_us/article/zmeyg8/ukraine-power-grid-malware-crashoverride-industroyer}
\bibitem{icsRussianHackingCyberWeapon} ... \LaTeX:\\ \url{http://blog.wallix.com/ics-security-russian-hacking}
\bibitem{usgovC2M2} ... \LaTeX:\\ \url{https://www.energy.gov/ceser/activities/cybersecurity-critical-energy-infrastructure/energy-sector-cybersecurity-0}

bediening werking schutsluizen pdf
\bibitem{varendoejesamenVeiligSluisvaren} ... \LaTeX:\\ \url{https://www.varendoejesamen.nl/storage/app/media/downloads/vlot-en-veilig-door-brug-en-sluis-.pdf}
\bibitem{vlaamsewaterwegen012014} ... \LaTeX:\\ \url{http://www.scarphout.be/assets/bedieningstijden2014.pdf}
\bibitem{bardetsluizenAmsterdam} ... \LaTeX:\\ \url{https://www.theobakker.net/pdf/sluizen.pdf}
\bibitem{dvsbedieningsluizenenbruggen} ... \LaTeX:\\ \url{http://www.watersportalmanak.nl/files/File/Brugbediengstijden_watersport.pdf}
\bibitem{crowstappenplanmachinerichtlijnen} ... \LaTeX:\\ \url{https://www.crow.nl/downloads/pdf/verkeer-en-vervoer/verkeersmanagement/verkeersregelinstallaties/stappenplan-machinerichtlijnen_web.aspx}
\bibitem{rijksoverheidrwsonderzoeksrapporten} ... \LaTeX:\\ \url{https://puc.overheid.nl/rijkswaterstaat/doc/PUC_95170_31/}
\bibitem{bedieningstijdensluizenenbruggen} ... \LaTeX:\\ \url{http://wsv.wsvdegors.nl/wp-content/uploads/2017/05/Bedieningstijden_201701.pdf}
\bibitem{sluiscomplexkornwerderzand} ... \LaTeX:\\ \url{https://www.commissiemer.nl/projectdocumenten/00004717.pdf}
\bibitem{varengroningendrenthe} ... \LaTeX:\\ \url{https://tasmanroutes.nl/wp-content/uploads/docs/1900-bedieningstijden-groningen-drenthe.pdf}
\bibitem{wallemuldergetijhoogteverschillen} ... \LaTeX:\\ \url{http://www.vliz.be/docs/groterede/GR21_Zeesluis.pdf}
\bibitem{biemandssluizenLith} ... \LaTeX:\\ \url{https://www.bhic.nl/media/document/file/rien-biemans-sluis-en-stuw-bij-lith.pdf}
\bibitem{weilerburgers11062018zoutindringingschutsluizen} ... \LaTeX:\\ \url{https://www.nattekunstwerkenvandetoekomst.nl/upload/documents/tinymce/KpNK-2017-SKW-01c001-v1-Zoutindringing-door-schutsluizen-overzicht-projecten-en-aanzet-formulering-tbv-netwerkmodellen.pdf}
\bibitem{rwsrichtlijnvaarwegen2011} ... \LaTeX:\\ \url{https://www.arnhemspeil.nl/nap/dok/2011-12-00-rijkswaterstaat-richtlijnen-vaarwegen.pdf}
\bibitem{slotterwisscha02062016swimwaywaddenzee} ... \LaTeX:\\ \url{https://rijkewaddenzee.nl/wp-content/uploads/2016/08/Inventarisatie-toestand-vispasseerbaarheid-zoet-zout-overgangen-Waddenzee-2-6-2016-PRW-rapportage-Definitief.pdf}
\bibitem{visvriendelijksluisbeheer06012014} ... \LaTeX:\\ \url{http://www.nevepaling.nl/files/Image//nederlands/informatiecentrum/2014-definitieve-voorkeursvariantennotitie-visvriendelijk-sluisbeheer-afsluitdijk-en-houtribdijk//2014_definitieve_voorkeursvariantennotitie_visvriendelijk_sluisbeheer_afsluitdijk_en_houtribdijk.pdf}
\bibitem{ID} ... \LaTeX:\\ \url{https://www.ifv.nl/kennisplein/Documents/20120614-BwNL-Handboek-brandbeveiligingsinstallaties.pdf}
\bibitem{aubel112016sluiscomplexHeumen} ... \LaTeX:\\ \url{https://ienc-kennisportaal.nl/wp-content/uploads/2017/01/Objectbeschrijving-Heumen.pdf}
\bibitem{arends052005schutsluisstolwijkschevliet} ... \LaTeX:\\ \url{https://library.wur.nl/edepot/websites/stolwijkersluis/presentatie-data/data/pdf/TUDelft-bouwhistorisch-onderzoek.pdf}
\bibitem{multidomeinbediening} ... \LaTeX:\\ \url{https://www.icentrale.nl/wp-content/uploads/bsk-pdf-manager/2019/01/20170929_Project-2.02-Deliverable-Gehele-werkpakket-2.02.pdf}
\bibitem{historischesluizen} ... \LaTeX:\\ \url{https://www.stowa.nl/sites/default/files/assets/PUBLICATIES/Publicaties%202000-2010/Publicaties%202000-2004/STOWA%202004-XX%20boekenreeks%2020.pdf}
\bibitem{nmmag2017centralebediening} ... \LaTeX:\\ \url{https://www.nm-magazine.nl/pdf/NM_Magazine_2017-3.pdf}
\bibitem{knooppuntvaarwegenFryslanGroningenDrenthe} ... \LaTeX:\\ \url{https://www.varendoejesamen.nl/storage/app/media/knooppunten/knooppuntenboekje_03_Friesland_Groningen_Drenthe.pdf}
\bibitem{afsluitdijk052015} ... \LaTeX:\\ \url{https://deafsluitdijk.nl/wp-content/uploads/2014/05/Plan-project-MER-Afsluitdijk.pdf}
\bibitem{ID} ... \LaTeX:\\ \url{file:///C:/Users/gally/Downloads/vaarroutekaart_provincie_drenthe.pdf}
\bibitem{ID} ... \LaTeX:\\ \url{file:///C:/Users/gally/Downloads/6227_watermanagement_nl_dv.pdf}
\bibitem{ID} ... \LaTeX:\\ \url{file:///C:/Users/gally/Downloads/applsci-11-00092-v3.pdf}
\bibitem{ID} ... \LaTeX:\\ \url{file:///C:/Users/gally/Downloads/Bedieningstijden_sluizen_en_bruggen_2004.pdf.pdf}
\bibitem{ID} ... \LaTeX:\\ \url{file:///C:/Users/gally/Downloads/Bedieningstijden.pdf}
\bibitem{ID} ... \LaTeX:\\ \url{file:///C:/Users/gally/Downloads/bijlagerapport_c_-_analyse_geavanceerd_definitief_v1_0.pdf}
\bibitem{ID} ... \LaTeX:\\ \url{file:///C:/Users/gally/Downloads/BIT-advies+Bediening+op+Afstand,+sluizen+en+bruggen+in+Friesland.pdf}
\bibitem{ID} ... \LaTeX:\\ \url{file:///C:/Users/gally/Downloads/conceptverordeningnautischbeheerzuid-holland.pdf}
\bibitem{ID} ... \LaTeX:\\ \url{file:///C:/Users/gally/Downloads/De_Deltawerken_Cultuurhistorie_ontwerpgeschiedenis_web-A.pdf}
\bibitem{ID} ... \LaTeX:\\ \url{file:///C:/Users/gally/Downloads/duurzaamheid_bij_de_ontwikkeling_van_reevesluis.pdf}
\bibitem{ID} ... \LaTeX:\\ \url{file:///C:/Users/gally/Downloads/gebruikershandleiding-databank-vismigratie.pdf}
\bibitem{ID} ... \LaTeX:\\ \url{file:///C:/Users/gally/Downloads/onderzoek_vispasseerbaarheid_sluizen_zuid_holland_2016_definitief_16-5-20171.pdf}
\bibitem{ID} ... \LaTeX:\\ \url{file:///C:/Users/gally/Downloads/rapport-veiligheid-van-op-afstand-bediende-bruggen.pdf}
\bibitem{ID} ... \LaTeX:\\ \url{file:///C:/Users/gally/Downloads/richtlijnen-vaarwegen-2017_tcm21-127359%20(1).pdf}
\bibitem{ID} ... \LaTeX:\\ \url{file:///C:/Users/gally/Downloads/richtlijnen-vaarwegen-2017_tcm21-127359.pdf}
\bibitem{ID} ... \LaTeX:\\ \url{file:///C:/Users/gally/Downloads/TPE_2342144_20210511223747_PEU_52457169.pdf}


Wat hebben alle bovenstaande rampen/ongelukken gemeen? Veiligheid.
Bij de therac waren er diverse problemen: communicatie, doorontwikkeling, controle en toetsing
Was het makkelijk te onderzoeken? Waarom?
Bij de boeing 737 crashes was het probleem van controle en communicatie naar medewerkers
Was het makkelijk te onderzoeken? Waarom?

Uit de evaluatie van de china explosion 2015 tianjin komt naar voren dat communicatie, transparantie en veiligheid niet altijd prioriteit hadden bij de lokale autoriteiten
Was het makkelijk te onderzoeken? Waarom?

Bij de tesla autopilot crashes komen soms onvoldoende onderbouwde ontwerpkeuzes naar voren die niet goed zij  afgewogen tegenover het gedrag van de bestuurder
vlucht 1951
Was het makkelijk te onderzoeken? Waarom?

De ramp in Tsjernobyl toont aan hoe autoriteiten een ramp in de doofpot proberen te stoppen
Was het makkelijk te onderzoeken? Waarom?



Wat heb ik geleerd
Ik heb erg veel geleerd van het veilig opzetten van VPN’s. Een VPN opzettenhad ik namelijk nog nooit gedaan. Het opzetten van SSH en het aanmaken vanVM’s was al bekend. Ook had ik nog nooit met UDP sockets geprogrammeerd.Verder heb ik geleerd hoe ik in de praktijk een VM in een VLAN kan zetten enhoe VLAN’s netwerken van elkaar kunnen scheiden.Het leukste onderdeel van het project, was dat wonderbaarlijk mijn gekozenoplossing elegant werkte. UDP Servers en clients zijn gerealiseerd met minderdan enkele regels logisch scipt. Ik had aan genomen dat het werken met socketsin shell absoluut rampzalig zou uitpakken. Ik ben blij dat het opdracht zo vrijwas, zodat ik experimenteel kon zijn met mijn implementatie.




\bibitem{ID} ... \LaTeX:\\ \url{https://www.uni-saarland.de/fileadmin/user_upload/Professoren/FreyG/DS_KT_GF_INCOM_May_2012.pdf}
vanaf 2.1 tot en met 5

\bibitem{ID} ... \LaTeX:\\ \url{http://www.lasid.ufba.br/publicacoes/artigos/Integrating+UML+and+UPPAAL+for+Designing,+Specifying+and+Verifying+Component-Based+Real-Time+Systems.pdf}

hf7
Reachability: i.e. some condition an posssibly be satisfied
Safety: i.e. some condition will never occur
Liveness: i.e. some condition wille eventually become true [] eventually or leadsto
hf 8
Het systeem is deadlockvrij
De wachttijd is altijd gelijk aan de invaarttijd _2x de nivlleertijd en de invaartijd van de overkant

\bibitem{ID} ... \LaTeX:\\ \url{https://www.diva-portal.org/smash/get/diva2:495691/FULLTEXT01.pdf}

blz 6 tot en met 10
\bibitem{ID} ... \LaTeX:\\ \url{https://www.cister-labs.pt/docs/formal_verification_of_aadl_models_using_uppaal/1331/view.pdf}

hf 3 geeft een voorbeeld van een template met guard en acies
De volgende automata worden gebruikt met hun lokale variabelen

De volgende globale variabelen

Een lijst met relevante einschappen van een schutsluis:

\bibitem{ID} ... \LaTeX:\\ \url{https://iopscience.iop.org/article/10.1088/1742-6596/1821/1/012031/pdf}

hf 5
deadlock

\bibitem{ID} ... \LaTeX:\\ \url{http://www.es.mdh.se/pdf_publications/2934.pdf}

hf 3 tool support
Modelling in UML
Code generation
Domain Model
Behaviour model
State Hierarchy
Transitions
Trigger methods
Time events
Effects
Requirements
Environment model

hf 4
\bibitem{ID} ... \LaTeX:\\ \url{https://files.ifi.uzh.ch/stiller/CLOSER%202014/WEBIST/WEBIST/Internet%20Technology/Full%20Papers/WEBIST_2014_130_CR.pdf}


\bibitem{ID} ... \LaTeX:\\ \url{https://files.ifi.uzh.ch/stiller/CLOSER%202014/WEBIST/WEBIST/Internet%20Technology/Full%20Papers/WEBIST_2014_130_CR.pdf}

4.2 5 en 6
Het Sluisbeheeerder model wordt getoond in fuguur[]. Het model is een uitbreiding van een schutsluis met alle condities en effecten. De kleuren in de automation verwijizen naar de kleuren in de staat van de automata . De template begint met een initiele lokatie start. De sluisbeheerder initieert het proces door een aangekomen schip te registreren metbehulp van een sychronizate met het channel... over de edge richtng de lokatie "aanmelden." Dit symboliseert een opstartprocedure, ook wordt een functie enqueeu_aanmeldLijst() gebruikt om de juiste waarden te geven aan lokale en globale avariabelen. De lokatie aanmelden regisseert het opstellen van schepen boven of beneden van de sluiskolk. De template Schip synchronizeerd met de template Sluisbeheerder met het channel move_down[id] of move_up[id] en bereikt daarmee de volgende lokatie afhankelijk af de sluis boven of beneden is worden de schepen die in de opstellijst voorkomen, max 2, klaargemaakt voor invaren.. De templates Stoplicht en sluisdeur synchroniseren met de channels ... call_Deur en call_stoplicht.
Het Sluisbeheerder model gebruikt de variabelen clock x, wachttijd_beneden, wachttijd_boven als invariant tussen de lokaties. Om op de hoogete te zijn van de invaar-/uitvaart van de verschillende schepen worden lijsten bijgehouden: list_wachtrij_beneden, list_pos_invaren_beneden, list_schepenInSluis, list_wachtrij_boven en list_pos_invaren_boven.

Het model voltooit de volgende transitie op basis van de waarde van de boolean sluis_bove en sluis_beneden. en de lokale klok variabele x.
Vanaf de locatie invaarverbod_gecontroleerd  wordt gecontroleerd of er nog invarende schepen zijn die in de sluiskolk passen.
Op de lokatie sluiskolk gereed zijn er 1 of meer schepen in de sluis. Als er nog plek is in de sluiskolk n er is nog een schip klaar om in te varen dan wordt dit gecontroleerd, de functie enqueu() voegt het schip toe aan de queue van de sluiskolk. De functie deque() verwijdert de schip van de lijst met invarende schepen. De variabele sluis_boven of sluis_beneden is waar, bij de switch voor het sluiten van deuren en het aanroepe van het stoplicht nr gelang de positie van  de laate binnenvarende schip (boven of beneden). Hierna bereikt de automation sluiskolk_afgesloten.



De lokatie start_nivelleren kiest op basis van de variabelen sluis_boven en de variabelen sluis_beneden het nivellereingsprogramma.
Heet nivellereingsprogramma is Aof B. De keuze voor het programma wordt bepaald door de variabelen van het schip dat in de sluis zit.

De lokatie klaarmaken_voor_openen wordt bereikt als de   hoogte van de sluis  door het nivellereingsprogramma is bereikt.
De positie van de kluis is bepaald door de schepen in de sluis. Vanuit deze lokatie wordt gekeken off de stoplichten gereed moeten worden gemaakt en of de sluisdeuren open mogen.
Hierna volgt een transiie waarin de stoplichte op groen worden gezet en de sluisdeuren worden geopend voor de uitvaart van de schepen in de sluis.
Als alle schepen zijn uitgevaren die uit moeten varen, worden de stoplichten op groen gezet en de deuren gesloten.


De lokatie uitvaren_toegestaan heeft een verbinding(edge) met de lokatie sluis_afsluiten.
Er is een select statement, e:id_t gebruikt als onderdeel van het prototocol om alle uitvarende schepen uit de queue van de sluiskolk te halen, en wordt dan ook gebruikt door de synchronisatie met de channel leave om de schepen uit de sluiskolk te begeleiden. De edge hieraan gekoppeld bevat de functie deque() om de variabelen  van de sluiskolk te resetten.

Vanuit de positie van de sluis worden de schepen gesignaleerd op een invaarverbod en worden de deuren van de sluis gesloten.
De lokatie sluiskolk_afgesloten is bereikt.

Ship [guards, invariants, assignents, synchronizations, properties,aannames]
De template Schip begint bij de Init lokatie. De lokatie is verbonden met de lokatie aangekomen met een edge waarbij een synchronizatie wordt aangeroepen met de template sluisbeheerder. De clock wordt op nul gezet. De lokatie aangekomen is verbonden met de lokatie aangemeld. De edge bevat een synchronizatie waarmee de edge een synchronizatie uitvoert met de template Sluisbheheerder.
De volgende lokatie is  controleren. De edge waarmee de lokatie aangemeld in verbinding staat met de lokatie cnotroleren heeft een synchronisatie voor de template Sluisbeheerder. De lokatie controleren heeft ook een edge met de lokatie wachten. Een schip max maximaal 30 seconden wachten op de lokatie wachten voordat er een mogelijkheid is om opniew in aanmerking te komen voor een controle. Als een schip langer dan 30 tijdseenheden moet wachten de is er een mogelijkheid voor het schip te vertrekken. Hierbij eindigt het schip het invaarproces. Een schip kan dus na aanvaren maximaal 20 seconden wachten om toestemming te krijgen voor een positie invaren anders wordt deze verwezen naar een wachtrij.
Hierna volgdde lokate invarene. De lokatie invarene implieert dat een schip in een invaarproces is dat eindigt in de lokatie gestopt.
Hierop volgd de lokatie nivelleer_start. Hierop wordt een nivelleer_proces gestart. Daarbij is ee synchronisatie met de template Sluisbeheerder.
De lokatie nivelleer_stop is een lokate waarin het nivelleerproces al is gestopt. Van hieruit is er een edge met de lokatie klaar voor vertrek. De edge synchroniseert hiermee met de template Sluisbeheerder.
De lokatie klaar_voor_vertrek is verbonden met de lokatie Init. Met een guard x>=3 tijdseenheden mag een schip vertrekken.


Deur
De deur bevat de volgende lokaties: dicht, openend, open en sluitende.
Een deur sluit niet in een enkele actie. Het proces die een deur dooploopt zijn de processen openend en sluitende. De finale lokaties zijn open en dicht.

Nivelleermachine
De nivelleermachine begint bij de lokatie uit. Met een synchronisatie wordt een nivelleermachine aangezet. De automatie kiest een programma en werkt deze uit in de lokatie bezig. Als ht programma is afgerond volgt de lokatie klaar. Na elk nivelleerproces wordt de machine uitgezet

Stoplicht
Een stoplicht heeft twee lokaties: rood en groen.



%%%%%%%%%%%%%%%%%%%%%%%%%%%%%%%%%%%%%%%%%%%%%%%%%%%%%%%%%%%%%%%%%

Bijlage A performance
\bibitem{kumarUppaalDMAMACProtocol} ... \LaTeX:\\ \url{https://home.hvl.no/ansatte/aaks/articles/2015IKT617.pdf}

test specification
\bibitem{larsenRealtimeUppaalTesting} ... \LaTeX:\\ \url{https://d-nb.info/987511998/34}

sheet 24 tot 65
\bibitem{proenza102008UppaalModelChecker} ... \LaTeX:\\ \url{http://ppedreiras.av.it.pt/resources/empse0809/slides/TheUppaalModelChecker-Julian.pdf}


2.3.4.2
4.7

coffie apparaat

\bibitem{uppaalCoffeeMachine} ... \LaTeX:\\ \url{https://www.comp.nus.edu.sg/~cs5270/Notes/chapt6a.pdf}



what is a good software specification
\bibitem{fvaandrager2322010Goodmodel} ... \LaTeX:\\ \url{http://www.cs.ru.nl/~fvaan/PV/what_is_a_good_model.html#:~:text=A%20good%20model%20has%20a%20clearly%20specified%20purpose%20and%20(ideally,code%20generation%2C%20and%20test%20generation.}

\bibitem{onix01102022devopmodel} ... \LaTeX:\\ \url{https://onix-systems.com/blog/7-basic-software-development-models-which-one-to-choose}
\bibitem{sulemani04012021softwareprocesmodel} ... \LaTeX:\\ \url{https://www.educative.io/blog/software-process-model-types}

\bibitem{globalluxsoft18102017softdev} ... \LaTeX:\\ \url{https://medium.com/globalluxsoft/5-popular-software-development-models-with-their-pros-and-cons-12a486b569dc}
\bibitem{wiegers30052022SRS} ... \LaTeX:\\ \url{https://www.jamasoftware.com/blog/characteristics-of-excellent-requirements/}
\bibitem{muller06092020goodspecification} ... \LaTeX:\\ \url{https://www.gaudisite.nl/ValidationOfRequirementsSlides.pdf}
\bibitem{informit30062008reqmanagement} ... \LaTeX:\\ \url{https://www.informit.com/articles/article.aspx?p=1152528&seqNum=4}
\bibitem{altexsoft15092020writingSRS} ... \LaTeX:\\ \url{https://www.altexsoft.com/blog/software-requirements-specification/}
\bibitem{ID} ... \LaTeX:\\ \url{E:\Backup Mijn Documenten\Hogeschool vakken\TINLab advnced algorithms\tinlab_advanced_algoriths\achtergrondinfo research}
sheet 28 transitorische relaties vertalen van ctl naar ltl
\bibitem{ID} ... \LaTeX:\\ \url{file:///E:/Backup%20Mijn%20Documenten/Hogeschool%20vakken/TINLab%20advnced%20algorithms/tinlab_advanced_algoriths/achtergrondinfo%20research/buchi/lec16_Buchi+LTL.pdf}
\bibitem{ID} ... \LaTeX:\\ \url{file:///E:/Backup%20Mijn%20Documenten/Hogeschool%20vakken/TINLab%20advnced%20algorithms/tinlab_advanced_algoriths/achtergrondinfo%20research/buchi/lect4.pdf}
\bibitem{ID} ... \LaTeX:\\ \url{file:///E:/Backup%20Mijn%20Documenten/Hogeschool%20vakken/TINLab%20advnced%20algorithms/tinlab_advanced_algoriths/achtergrondinfo%20research/buchi/lecture8.pdf}
transitie relaties in LTL sheet 8
\bibitem{ID} ... \LaTeX:\\ \url{file:///E:/Backup%20Mijn%20Documenten/Hogeschool%20vakken/TINLab%20advnced%20algorithms/tinlab_advanced_algoriths/achtergrondinfo%20research/CS%20267%20Automated%20Verification/l2.pdf}
\bibitem{ID} ... \LaTeX:\\ \url{file:///E:/Backup%20Mijn%20Documenten/Hogeschool%20vakken/TINLab%20advnced%20algorithms/tinlab_advanced_algoriths/achtergrondinfo%20research/FORMAL%20METHODS/slide3.pdf}
\bibitem{ID} ... \LaTeX:\\ \url{file:///E:/Backup%20Mijn%20Documenten/Hogeschool%20vakken/TINLab%20advnced%20algorithms/tinlab_advanced_algoriths/achtergrondinfo%20research/FORMAL%20METHODS/slide4.pdf}
hf 4.2
\bibitem{ID} ... \LaTeX:\\ \url{file:///E:/Backup%20Mijn%20Documenten/Hogeschool%20vakken/TINLab%20advnced%20algorithms/tinlab_advanced_algoriths/achtergrondinfo%20research/properties%20ctl/Chapter-4-Formal-Methods-LTL-CTL-TRAFFIC-LIGHT-EXAMPLE-pages-18-24.pdf}


\bibitem{ID} ... \LaTeX:\\ \url{E:\Backup Mijn Documenten\Hogeschool vakken\TINLab advnced algorithms\tinlab_advanced_algoriths\lesmateriaal\modelchecking.pdf}





\bibitem{ID} ... \LaTeX:\\ \url{E:\Backup Mijn Documenten\Hogeschool vakken\TINLab advnced algorithms\tinlab_advanced_algoriths\lesmateriaal}

parallelle compositie
\bibitem{ID} ... \LaTeX:\\ \url{file:///E:/Backup%20Mijn%20Documenten/Hogeschool%20vakken/TINLab%20advnced%20algorithms/tinlab_advanced_algoriths/achtergrondinfo%20research/properties%20ctl/model.pdf}

Urgent locations
Is hetzelfde als het toevoegen van een clock x, met een invariant x<=o op de locatie. Zolang een systeem in een urgente locatie zit mag er geen tijd verstrijken
Bjivoorbeeld als een sluis klaar is engeen schpeen in de sluis. Dan moet er een urgentie zijn dat alle schepen waar mogelijk worden opgesteld voor invaren. Als er geen schepen in de wachtrij en er staan geenschepen klaar om in te varen dn is er misschien urgentie om aan de andere kant schepen op te halen.
Commited locations
Als een of meerdere locaties ingesteld zijn als committed. Een committed state kan niet vertragen  en de volgende transitie moet een transitie zijn waarin de uitgaande edge komt van een committed edge


zeno gedrag: de mogelijkheid dat in een eindige hoeveelheid tijd een oneindig antal handelingen kan worden verricht.
Bijvoorbeeld tijdens het nivelleren
Bij het opstellen van schepen
Bij het laten wachten van schepen
Bij het invaren van schepen
\bibitem{ID} ... \LaTeX:\\ \url{file:///E:/Backup%20Mijn%20Documenten/Hogeschool%20vakken/TINLab%20advnced%20algorithms/tinlab_advanced_algoriths/achtergrondinfo%20research/properties%20ctl/std.pdf}



\bibitem{ID} ... \LaTeX:\\ \url{https://wayback.archive-it.org/9650/20200409062940/http:/p3-raw.greenpeace.org/international/Global/international/publications/nuclear/2016/Nuclear_Scars.pdf}
\bibitem{ID} ... \LaTeX:\\ \url{https://bdtechtalks.com/2020/07/29/self-driving-tesla-car-deep-learning/}



%%%%%%%%%%%%%%%%%%%%%%%%%%%%%%%%%%%%%%%%%%%%%%%%%%%%%%%%%%%%%%%%%
critical safety systems chemicals
\bibitem{esc16022021scsdeveloping} ... \LaTeX:\\ \url{https://esc.uk.net/safety-critical-systems}
\bibitem{oecd2008chemsafeperfindct} ... \LaTeX:\\ \url{https://www.oecd.org/chemicalsafety/chemical-accidents/41269710.pdf}
\bibitem{issa2003chemicalsID} ... \LaTeX:\\ \url{https://safety-work.org/fileadmin/safety-work/articles/Verwechslung_von_Chemikalien/Stoffverwechslung_e.pdf}
\bibitem{sommerville2008CriticalSystems} ... \LaTeX:\\ \url{https://ifs.host.cs.st-andrews.ac.uk/Books/SE9/Web/Dependability/CritSys.html}
\bibitem{identifyinglaboratoryHazards} ... \LaTeX:\\ \url{https://www.acs.org/content/dam/acsorg/about/governance/committees/chemicalsafety/publications/identifying-and-evaluating-hazards-in-research-laboratories.pdf}
\bibitem{ID} ... \LaTeX:\\ \url{https://www.computer.org/csdl/magazine/so/2017/04/mso2017040049/13rRUxCitHw}
\bibitem{ID} ... \LaTeX:\\ \url{https://msquair.files.wordpress.com/2012/06/assca-guiding-philosophic-principles-on-the-design-and-acquisition-of-safety-critical-systems-v1-6.pdf}
\bibitem{ID} ... \LaTeX:\\ \url{https://epsc.be/Documents/PS+Fundamentals/_/EPSC_Process%20Safety%20Fundamentals%20-%20Booklet_March2021.pdf}
\bibitem{winceckCriticalToSafety} ... \LaTeX:\\ \url{https://www.icheme.org/media/8976/xxiv-poster-11.pdf}
\bibitem{chambersHazardAnalysisSCS} ... \LaTeX:\\ \url{https://crpit.scem.westernsydney.edu.au/confpapers/CRPITV55Chambers.pdf}
\bibitem{rslater1998SCSAnalysis} ... \LaTeX:\\ \url{https://users.ece.cmu.edu/~koopman/des_s99/safety_critical/}




%%%%%%%%%%%%%%%%%%%%%%%%%%%%%%%%%%%%%%%%%%%%%%%%%%%%%%%%%%%%%%%%%
critical safety systems airplanes
\bibitem{ID} ... \LaTeX:\\ \url{file:///C:/Users/gally/Downloads/AGARDAG300.pdf}
\bibitem{brat2015verifysafetyflightcritical} ... \LaTeX:\\ \url{https://arxiv.org/abs/1502.02605}
\bibitem{knightchallengessafetyCritical} ... \LaTeX:\\ \url{https://users.encs.concordia.ca/~ymzhang/courses/reliability/ICSE02Knight.pdf}
\bibitem{ID} ... \LaTeX:\\ \url{https://www.jstor.org/stable/44682826}
\bibitem{johnson2006devsafetycritical} ... \LaTeX:\\ \url{http://www.dcs.gla.ac.uk/~johnson/teaching/safety/slides/pt2.pdf}
\bibitem{yeagerSafetyCritical} ... \LaTeX:\\ \url{https://sites.google.com/site/cis115textbook/safety-critical-systems}
\bibitem{ID} ... \LaTeX:\\ \url{https://www.dau.edu/tools/se-brainbook/Pages/Design%20Considerations/Critical-Safety-Item.aspx}
\bibitem{ID} ... \LaTeX:\\ \url{https://mcdpinc.com/safety-critical-systems}
\bibitem{fallsafedesign} ... \LaTeX:\\ \url{https://faculty.up.edu/lulay/MEStudentPage/failsafe.pdf}
\bibitem{2008manualflightsafetyParts} ... \LaTeX:\\ \url{https://www.enidine.com/CorporateSite/media/itt/Resources/Distributors/EndUserDocuments/Suppliers_Documents/QAM03_Rev_E.pdf}
\bibitem{arForce2015VerificationExpectations} ... \LaTeX:\\ \url{https://daytonaero.com/wp-content/uploads/AC-17-01.pdf}
\bibitem{harvardRiskResearchaircraft} ... \LaTeX:\\ \url{https://rmas.fad.harvard.edu/pages/chartered-private-aircraft-0}
\bibitem{ID} ... \LaTeX:\\ \url{https://pubmed.ncbi.nlm.nih.gov/7966484/}
\bibitem{ID} ... \LaTeX:\\ \url{https://nebula.esa.int/content/assessment-methodology-certification-safety-gnc-critical-space-systems}
\bibitem{ID} ... \LaTeX:\\ \url{https://www.aopa.org/training-and-safety/online-learning/safety-spotlights/aircraft-systems}
\bibitem{lalaArchitecturalPrinciples} ... \LaTeX:\\ \url{https://www.cs.unc.edu/~anderson/teach/comp790/papers/safety_critical_arch.pdf}
\bibitem{andersenromanski2011verificationsafetycritical} ... \LaTeX:\\ \url{https://queue.acm.org/detail.cfm?id=2024356}
\bibitem{aviatioLawGeneralDefinitions} ... \LaTeX:\\ \url{https://www.law.cornell.edu/cfr/text/14/1.1}
\bibitem{ntsb2006safetyreportTransportAirplanes} ... \LaTeX:\\ \url{http://libraryonline.erau.edu/online-full-text/ntsb/safety-reports/SR06-02.pdf}
\bibitem{mitNotesSafetyCritical} ... \LaTeX:\\ \url{https://www.cs.uct.ac.za/mit_notes/human_computer_interaction/htmls/ch02s10.html}
\bibitem{ID} ... \LaTeX:\\ \url{https://flightsafety.org/}
\bibitem{kochenfender2020aisafetycritical} ... \LaTeX:\\ \url{https://engineering.stanford.edu/magazine/article/mykel-kochenderfer-ai-and-safety-critical-systems}
\bibitem{humanfactorsAviationSafety} ... \LaTeX:\\ \url{https://www.faasafety.gov/files/gslac/courses/content/258/1097/AMT_Handbook_Addendum_Human_Factors.pdf}
\bibitem{eulegislator2000safetyregulation} ... \LaTeX:\\ \url{https://www.eurocontrol.int/sites/default/files/2019-06/src-doc-1-e1.0.pdf}
\bibitem{ID} ... \LaTeX:\\ \url{http://aerossurance.com/safety-management/critical-maintenance-tasks/}
\bibitem{gao2020aviationcybersecurity} ... \LaTeX:\\ \url{https://www.gao.gov/assets/gao-21-86.pdf}
\bibitem{roleofcodeinsoftware} ... \LaTeX:\\ \url{https://criticalsoftware.com/en/news/coding-the-skies}
\bibitem{airlie2018designparameters} ... \LaTeX:\\ \url{https://aviation.stackexchange.com/questions/46677/what-are-the-design-parameters-for-airliner-safety}
\bibitem{ID} ... \LaTeX:\\ \url{https://www.cantwell.senate.gov/news/press-releases/cantwells-comprehensive-bipartisan-bicameral-aircraft-safety-and-certification-reforms-signed-into-law}
\bibitem{ID} ... \LaTeX:\\ \url{https://www.forbes.com/advisor/travel-rewards/737-max-what-is-safety-anyway/}
\bibitem{humanfactors2010Aviation} ... \LaTeX:\\ \url{https://www.tandfonline.com/doi/full/10.1080/00140130903521587}
\bibitem{ID} ... \LaTeX:\\ \url{https://www.doi.gov/aviation/safety}
\bibitem{courtin2018safetyconsiderations} ... \LaTeX:\\ \url{https://dspace.mit.edu/bitstream/handle/1721.1/118438/ICAT_2018_07_Christoper%20Courtin_Report.pdf?sequence=1&isAllowed=y}
\bibitem{ID} ... \LaTeX:\\ \url{https://www.defence.gov.au/dasp/Docs/Manuals/7001053/eTAMMweb/1049.htm}
\bibitem{prentice2014Failedaviationprogram} ... \LaTeX:\\ \url{https://www.aviationpros.com/aircraft/commercial-airline/article/10239806/staying-legal-another-failed-faa-safety-program}
\bibitem{civilAviationConsulting} ... \LaTeX:\\ \url{https://www.iata.org/en/services/consulting/safety-operations/}
\bibitem{olivercalvardpotocnik2017Aviationautomation} ... \LaTeX:\\ \url{https://hbr.org/2017/09/the-tragic-crash-of-flight-af447-shows-the-unlikely-but-catastrophic-consequences-of-automation}
\bibitem{ID} ... \LaTeX:\\ \url{https://www.infosys.com/industries/communication-services/documents/landing-gear-design-and-development.pdf}
\bibitem{airforce2000SystemSafety} ... \LaTeX:\\ \url{https://www.acqnotes.com/Attachments/AF_System-Safety-HNDBK.pdf}
\bibitem{fed2019SafeSecureSUAS} ... \LaTeX:\\ \url{https://www.transportation.gov/testimony/state-airline-safety-federal-oversight-commercial-aviation}
\bibitem{AirportSafety} ... \LaTeX:\\ \url{https://www.federalregister.gov/documents/2019/02/13/2019-00758/safe-and-secure-operations-of-small-unmanned-aircraft-systems}
\bibitem{ID} ... \LaTeX:\\ \url{https://archive.etsc.eu/documents/safety%20in%20airports.pdf}
\bibitem{passengerSafety} ... \LaTeX:\\ \url{https://journals.sagepub.com/doi/pdf/10.1177/002029400403700202}
\bibitem{ID} ... \LaTeX:\\ \url{https://www.unmannedsystems.ca/wp-content/uploads/2019/01/DRAFT-AC-922-001-RPAS-SAFETY-ASSURANCE.pdf}
\bibitem{ID} ... \LaTeX:\\ \url{https://www.ccsdualsnap.com/pressure-switches-in-aerospace-applications/}
\bibitem{britishColumbia2020GuideSafetyCritical} ... \LaTeX:\\ \url{https://www.egbc.ca/getmedia/78073fda-5a83-4f0f-b12f-0a40dcbbc29d/EGBC-Safety-Critical-Software-V1-0.pdf.aspx}
\bibitem{uscongres2019aircraftresearch} ... \LaTeX:\\ \url{https://readwrite.com/2018/12/21/air-travel-is-far-safer-than-you-think-heres-why/}
\bibitem{ID} ... \LaTeX:\\ \url{https://fas.org/sgp/crs/misc/R45939.pdf}
\bibitem{ariAssociation2018} ... \LaTeX:\\ \url{https://cdn.ymaws.com/www.astna.org/resource/collection/4392B20B-D0DB-4E76-959C-6989214920E9/ASTNA_Safety_Position_Paper_2018_FINAL.pdf}
\bibitem{ID} ... \LaTeX:\\ \url{https://transportation.house.gov/imo/media/doc/2020.09.15%20FINAL%20737%20MAX%20Report%20for%20Public%20Release.pdf}
\bibitem{ID} ... \LaTeX:\\ \url{https://www.transportstyrelsen.se/globalassets/global/luftfart/seminarier_och_information/seminarier-2016/luftvardighet-camo-och-145-verkstader/11b-critical-task-fpl.pdf}
\bibitem{fox2005HelicopterSafety} ... \LaTeX:\\ \url{https://www.h-a-c.ca/IHSS_Helicopter_Safety_History_05.pdf}
\bibitem{fulvio1993safetycriticalsystems} ... \LaTeX:\\ \url{https://assembly.coe.int/nw/xml/XRef/X2H-Xref-ViewHTML.asp?FileID=7144&lang=EN}
\bibitem{rakas2018criticalsystemsFailures} ... \LaTeX:\\ \url{https://www.skybrary.aero/index.php/Cockpit_Automation_-_Advantages_and_Safety_Challenges}
\bibitem{ID} ... \LaTeX:\\ \url{https://ntrs.nasa.gov/citations/20120014507}
\bibitem{cockpitAutomation} ... \LaTeX:\\ \url{https://www.sciencedirect.com/science/article/abs/pii/S092575351730601X}
\bibitem{ID} ... \LaTeX:\\ \url{https://www.semanticscholar.org/paper/Safety-critical-avionics-for-the-777-primary-flight-Yeh/8facf90f4a9051c3ab8ce11e39d0893118268d90}
\bibitem{ID} ... \LaTeX:\\ \url{https://www.easa.europa.eu/faq/19013}
\bibitem{belcaatrovalidateSafetyCritical} ... \LaTeX:\\ \url{https://ntrs.nasa.gov/api/citations/20120014507/downloads/20120014507.pdf}
\bibitem{ID,	civilAviationsAuthority = {author},	ALTeditor = {editor},	title = {title},	date = {date},	url = {"https://publicapps.caa.co.uk/modalapplication.aspx?catid=1&pagetype=65&appid=11&mode=list&type=subcat&id=32}
	\bibitem{ID} ... \LaTeX:\\ \url{https://definitions.uslegal.com/f/flight-safety-critical-aircraft-part/}
	\bibitem{ID} ... \LaTeX:\\ \url{https://nbaa.org/nbaa-aviation-groups-ask-congress-to-prevent-5g-interference-to-critical-safety-systems/}
	\bibitem{ID} ... \LaTeX:\\ \url{https://www.dlr.de/ft/en/desktopdefault.aspx/tabid-1360/1856_read-36215/}
	\bibitem{ID} ... \LaTeX:\\ \url{https://www.fsd.lrg.tum.de/research/safety-critical/}
	\bibitem{knight2010SafetyCritical} ... \LaTeX:\\ \url{https://ieeexplore.ieee.org/document/1007998}
	\bibitem{AAA052005IdentifySCHitems} ... \LaTeX:\\ \url{https://www.faa.gov/about/office_org/headquarters_offices/ast/licenses_permits/media/RLVGuide_01-05_05.pdf}
	\bibitem{devTopics01032020} ... \LaTeX:\\ \url{https://smallbusinessprogramming.com/safety-critical-software-15-things-every-developer-should-know/}
	\bibitem{ID} ... \LaTeX:\\ \url{https://coreavi.com/the-future-of-safety-critical-systems-in-the-emerging-autonomous-world/}
	\bibitem{valdes2018SafetybyAutomation} ... \LaTeX:\\ \url{https://www.intechopen.com/chapters/59838}
	\bibitem{2015whensafetymanagementsystemsfail} ... \LaTeX:\\ \url{https://verticalmag.com/features/whensafetymanagementsystemsfail/}
	
	
	%%%%%%%%%%%%%%%%%%%%%%%%%%%%%%%%%%%%%%%%%%%%%%%%%%%%%%%%%%%%%%%%%
	critical safety systems fireworks
	\bibitem{ID} ... \LaTeX:\\ \url{https://www.hsdl.org/c/firework-safety/}
	\bibitem{ID} ... \LaTeX:\\ \url{https://www.cpsc.gov/Safety-Education/Safety-Education-Centers/Fireworks}
	\bibitem{ID} ... \LaTeX:\\ \url{https://www.seattletimes.com/subscribe/signup-offers/?pw=redirect&subsource=paywall&return=https://www.seattletimes.com/opinion/editorials/firework-safety-even-more-critical-after-heat-wave/}
	\bibitem{ID} ... \LaTeX:\\ \url{https://www.nrcan.gc.ca/sites/www.nrcan.gc.ca/files/mineralsmetals/pdf/mms-smm/expl-expl/20170828-G05-09E_ACC.pdf}
	\bibitem{fireworksinjuries2021duringcovid} ... \LaTeX:\\ \url{https://www.prnewswire.com/news-releases/fireworks-related-injuries-and-deaths-spiked-during-the-covid-19-pandemic-301322243.html}
	\bibitem{osha2017safetymanagementexplosives} ... \LaTeX:\\ \url{https://www.osha.gov/sites/default/files/publications/OSHA3912.pdf}
	\bibitem{ID} ... \LaTeX:\\ \url{https://www.firelinx.com/wp-content/uploads/2021/02/FLX-Issues-in-Firing-System-Safety.pdf}
	\bibitem{ID} ... \LaTeX:\\ \url{http://www.eig2.org.uk/wp-content/uploads/WTOFD-Blue-Guide.pdf}
	\bibitem{hse2014explosiveregulations} ... \LaTeX:\\ \url{https://www.hse.gov.uk/explosives/er2014-fireworks-retail-prem.pdf}
	\bibitem{ID} ... \LaTeX:\\ \url{https://www.firerescue1.com/firefighter-safety/articles/11-fireworks-safety-videos-from-the-serious-to-the-humorous-fHy0M4pT2gjcQ8jA/}
	\bibitem{pirone2016lessonslearnedfireworks} ... \LaTeX:\\ \url{https://www.aidic.it/cet/16/53/044.pdf}
	\bibitem{ID} ... \LaTeX:\\ \url{http://www.alarmascasas.com.mx/sites/default/files/85006-0061%20--%20FireWorks%20Brochure.pdf}
	\bibitem{ID} ... \LaTeX:\\ \url{https://www.ehs.ufl.edu/programs/fire/fireworks/}
	\bibitem{ID} ... \LaTeX:\\ \url{https://www.interlogix.com.au/documents/FireWorks%20Features%20and%20Operation%20(fire%20only).pdf}
	\bibitem{ID} ... \LaTeX:\\ \url{https://townhall.virginia.gov/l/GetFile.cfm?File=C:%5CTownHall%5Cdocroot%5CGuidanceDocs%5C960%5CGDoc_DFP_4448_v1.pdf}
	\bibitem{ID} ... \LaTeX:\\ \url{https://ec.europa.eu/growth/sectors/chemicals/specific-chemicals_en}
	\bibitem{safehandlepyrotechnics} ... \LaTeX:\\ \url{http://www.iiakm.org/ojakm/articles/2015/volume3_3/OJAKM_Volume3_3pp27-36.pdf}
	\bibitem{ID} ... \LaTeX:\\ \url{https://www.bristol.gov.uk/documents/20182/1175006/Fireworks+in+retail+premises/6aa6ee24-5b74-43b4-a1d9-747689b1dbc9}
	\bibitem{ID} ... \LaTeX:\\ \url{https://www.eversys.com.br/imagens/uploads/arqs/bra_arquivos/04-software-gerenciador-fireworks-brochura.pdf}
	\bibitem{ID} ... \LaTeX:\\ \url{http://www.doiserbia.nb.rs/img/doi/0354-9836/2016/0354-98361500050G.pdf}
	\bibitem{costinFireworksEmbedded} ... \LaTeX:\\ \url{http://s3.eurecom.fr/docs/wisec14_Costin.pdf}
	\bibitem{ID} ... \LaTeX:\\ \url{https://www.firetechsystems.com/assets/uploads/2018/09/FireWorks-Brochure.pdf}
	\bibitem{ritz2020firesafety} ... \LaTeX:\\ \url{https://blog.ritzsafety.com/fireworks-safety-tips}
	\bibitem{lawson2018criticalsystems} ... \LaTeX:\\ \url{https://www.engineerlive.com/content/fire-detection-and-protection-through-safety-critical-systems}
	%%%%%%%%%%%%%%%%%%%%%%%%%%%%%%%%%%%%%%%%%%%%%%%%%%%%%%%%%%%%%%%%%
%    
%    algemene vragen
%    oorzaken
%    @online{gates18112020boeingcrisis,	ALTauthor = {author},	ALTeditor = {editor},	title = {title},	date = {date},	url = {"https://www.seattletimes.com/business/boeing-aerospace/what-led-to-boeings-737-max-crisis-a-qa/"},}
%    @online{boeing737maxsoftwareprobles,	ALTauthor = {author},	ALTeditor = {editor},	title = {title},	date = {date},	url = {"https://www.schneier.com/blog/archives/2019/04/excellent_analy.html"},}
%    fout in de software
%    @online{avetisov19032019boeingmalwarestate,	ALTauthor = {author},	ALTeditor = {editor},	title = {title},	date = {date},	url = {"https://www.forbes.com/sites/georgeavetisov/2019/03/19/malware-at-30000-feet-what-the-737-max-says-about-the-state-of-airplane-software-security/?sh=4d26f7052a9e"},}
%    het nationaal veiligheidsbelang
%    @online{thompson23112020nationalsecurityboeing,	ALTauthor = {author},	ALTeditor = {editor},	title = {title},	date = {date},	url = {"https://www.forbes.com/sites/lorenthompson/2020/11/23/five-reasons-return-of-boeings-737-max-to-service-is-important-to-national-security/?sh=2128ea552018"},}
%    falend toezicht
%    @online{gates21032019FAAControlSystem,	ALTauthor = {author},	ALTeditor = {editor},	title = {title},	date = {date},	url = {"https://www.seattletimes.com/business/boeing-aerospace/failed-certification-faa-missed-safety-issues-in-the-737-max-system-implicated-in-the-lion-air-crash/"},}
%    onderzoeksrapport
%    @online{faa18112020boeingreview,	ALTauthor = {author},	ALTeditor = {editor},	title = {title},	date = {date},	url = {"https://www.faa.gov/foia/electronic_reading_room/boeing_reading_room/media/737_RTS_Summary.pdf"},}
%    @online{wiki737maxgroundings,	ALTauthor = {author},	ALTeditor = {editor},	title = {title},	date = {date},	url = {"https://en.wikipedia.org/wiki/Boeing_737_MAX_groundings"},}
%    veiligheidsrisico's
%    menselijke fouten
%    @online{campbell02052019boengcrashhumanerrors,	ALTauthor = {author},	ALTeditor = {editor},	title = {title},	date = {date},	url = {"https://www.theverge.com/2019/5/2/18518176/boeing-737-max-crash-problems-human-error-mcas-faa"},}
%    overzicht van crashes
%    @online{hawkins22032019737maxairplanes,	ALTauthor = {author},	ALTeditor = {editor},	title = {title},	date = {date},	url = {"https://www.theverge.com/2019/3/22/18275736/boeing-737-max-plane-crashes-grounded-problems-info-details-explained-reasons"},}
%    veiligheidsopmerking
%    @online{thomas30082020737safest,	ALTauthor = {author},	ALTeditor = {editor},	title = {title},	date = {date},	url = {"https://www.airlineratings.com/news/boeings-737-max-will-one-safest-aircraft-history/"},}
%    aanpassingen
%    @online{boeing737maxdisplay,	ALTauthor = {author},	ALTeditor = {editor},	title = {title},	date = {date},	url = {"https://www.boeing.com/commercial/737max/737-max-software-updates.page"},}
%    waarschuwingen//output signalen
%    @online{fehrm24112020737changes,	ALTauthor = {author},	ALTeditor = {editor},	title = {title},	date = {date},	url = {"https://leehamnews.com/2020/11/24/boeing-737-max-changes-beyond-mcas/"},}
%    software gerelateerde fouten
%    @online{travis18042019737maxsoftwaredevop,	ALTauthor = {author},	ALTeditor = {editor},	title = {title},	date = {date},	url = {"https://spectrum.ieee.org/aerospace/aviation/how-the-boeing-737-max-disaster-looks-to-a-software-developer"},}
%    onderzoeksrapport
%    de rol van de publieke opinie
%    @online{barnett05052019737maxcrisis,	ALTauthor = {author},	ALTeditor = {editor},	title = {title},	date = {date},	url = {"https://pubsonline.informs.org/do/10.1287/orms.2019.05.05/full/"},}
%    onderzoek van europese luchtvaart agentschap
%    @online{easa27012021737maxsafereturn,	ALTauthor = {author},	ALTeditor = {editor},	title = {title},	date = {date},	url = {"https://www.easa.europa.eu/newsroom-and-events/news/easa-declares-boeing-737-max-safe-return-service-europe"},}
%    veiligheidsvraagstuk
%    @online{touitou11032019737tragedies,	ALTauthor = {author},	ALTeditor = {editor},	title = {title},	date = {date},	url = {"https://phys.org/news/2019-03-boeing-max-safety-tragedies.html"},}
%    artikel over sensoren
%    @online{hemmerdinger02022021737maxdeliveries,	ALTauthor = {author},	ALTeditor = {editor},	title = {title},	date = {date},	url = {"https://www.flightglobal.com/airframers/boeing-delays-737-max-10-deliveries-two-years-to-2023/142245.article"},}
%    goedkeuring van europese luchtvaart autoriteiten
%    advies aan de faa
%    @online{bielby27022021faaimprovesafety,	ALTauthor = {author},	ALTeditor = {editor},	title = {title},	date = {date},	url = {"https://www.hstoday.us/subject-matter-areas/airport-aviation-security/oig-tells-faa-to-improve-safety-oversight-following-boeing-737-max-review/"},}
%    @online{boyle18112020737maxupgrade,	ALTauthor = {author},	ALTeditor = {editor},	title = {title},	date = {date},	url = {"https://www.geekwire.com/2020/faas-go-ahead-737-maxs-return-flight-kicks-off-massive-software-upgrade/"},}
%    @online{bergstraburgess122019737maxMcasAlgorithm,	ALTauthor = {author},	ALTeditor = {editor},	title = {title},	date = {date},	url = {"https://www.researchgate.net/publication/338420944_A_Promise_Theoretic_Account_of_the_Boeing_737_Max_MCAS_Algorithm_Affair"},}
%    achtergrond informatie
%    @online{737mcas,	ALTauthor = {author},	ALTeditor = {editor},	title = {title},	date = {date},	url = {"http://www.b737.org.uk/mcas.htm"},}
%    algemeen vertrouwen
%    @online{newburger17052019boeingcrisis,	ALTauthor = {author},	ALTeditor = {editor},	title = {title},	date = {date},	url = {"https://www.cnbc.com/2019/05/16/what-you-need-to-know-about-boeings-737-max-crisis.html"},}
%    toestemming europese autoriteiten
%    problemen
%    @online{arstechnica22012020737problems,	ALTauthor = {author},	ALTeditor = {editor},	title = {title},	date = {date},	url = {"https://arstechnica.com/information-technology/2020/01/737-max-fix-slips-to-summer-and-thats-just-one-of-boeings-problems/"},}
%    uitgebreid artikel over de onderzoeken en het vliegverbod
%    @online{german190620217372yaftergrounded,	ALTauthor = {author},	ALTeditor = {editor},	title = {title},	date = {date},	url = {"https://www.cnet.com/news/boeing-737-max-8-all-about-the-aircraft-flight-ban-and-investigations/"},}
%    computers als oorzaak
%    lessons learned
%    @online{beningo02052019boeinglessons,	ALTauthor = {author},	ALTeditor = {editor},	title = {title},	date = {date},	url = {"https://www.designnews.com/electronics-test/5-lessons-learn-boeing-737-max-fiasco"},}
%    @online{duran05042019boeingspof,	ALTauthor = {author},	ALTeditor = {editor},	title = {title},	date = {date},	url = {"https://www.eurocontrol.int/publication/effects-network-extra-standby-aircraft-and-boeing-737-max-grounding"},}
%    single point of failure
%    @online{ID,	ALTauthor = {author},	ALTeditor = {editor},	title = {title},	date = {date},	url = {"https://dmd.solutions/blog/2019/04/05/how-a-single-point-of-failure-spof-in-the-mcas-software-could-have-caused-the-boeing-737-max-crash-in-ethiopia/"},}
%    @online{makichuck24012021737fearflying,	ALTauthor = {author},	ALTeditor = {editor},	title = {title},	date = {date},	url = {"https://asiatimes.com/2021/01/boeings-737-max-and-the-fear-of-flying/"},}
%    lijst van tehnische aanpassingen
%    @online{caa737modifications,	ALTauthor = {author},	ALTeditor = {editor},	title = {title},	date = {date},	url = {"https://www.caa.co.uk/Consumers/Guide-to-aviation/Boeing-737-MAX/"},}
%    @online{oestergaard14122020boeingdeliveries,	ALTauthor = {author},	ALTeditor = {editor},	title = {title},	date = {date},	url = {"https://dsm.forecastinternational.com/wordpress/2020/12/14/airbus-and-boeing-report-november-2020-commercial-aircraft-orders-and-deliveries/"},}
%    code lek
%    @online{reenberg787flaws,	ALTauthor = {author},	ALTeditor = {editor},	title = {title},	date = {date},	url = {"https://www.wired.com/story/boeing-787-code-leak-security-flaws/"},}
%    @online{fitch16092020737backlogrisks,	ALTauthor = {author},	ALTeditor = {editor},	title = {title},	date = {date},	url = {"https://www.fitchratings.com/research/corporate-finance/boeing-737-max-return-backlog-risks-remain-16-09-2020"},}
%    Cultuurverandering, deregulatie, systeemwijziging of gewoon een kwestie van competentie
%    @online{willis27082020737maxfailures,	ALTauthor = {author},	ALTeditor = {editor},	title = {title},	date = {date},	url = {"https://www.aerospacetestinginternational.com/features/what-broke-the-737-max.html"},}
%    extra aanpassingen
%    @online{ostrower11062020more737changes,	ALTauthor = {author},	ALTeditor = {editor},	title = {title},	date = {date},	url = {"https://theaircurrent.com/aviation-safety/boeings-737-max-software-done-but-regulators-plot-more-changes-after-jets-return/"},}
%    wat ging er mis een analyse van een ex-iloot
%    De utoriteiten waren op de hoogte
%    @online{hruska13122019faaknown737crashrate,	ALTauthor = {author},	ALTeditor = {editor},	title = {title},	date = {date},	url = {"https://www.extremetech.com/extreme/303373-the-faa-knew-the-737-max-was-dangerous-and-kept-it-flying-anyway"},}
%    kwaliteiten van het alarmsysteem niet goed bekend
%    @online{bloomberg26092019failedpred,	ALTauthor = {author},	ALTeditor = {editor},	title = {title},	date = {date},	url = {"https://time.com/5687473/boeing-737-alarm-system/"},}
%    @online{whiteman09072020boengcancelstock,	ALTauthor = {author},	ALTeditor = {editor},	title = {title},	date = {date},	url = {"https://www.nasdaq.com/articles/boeing-gets-dealt-another-737-max-cancellation-blow.-what-it-means-for-boeing-stock-2020"},}
%    @online{leopold09192019boeingreliability,	ALTauthor = {author},	ALTeditor = {editor},	title = {title},	date = {date},	url = {"https://www.eetimes.com/boeing-crashes-highlight-a-worsening-reliability-crisis/"},}
%    veiligheidsvraagstuk
%    @online{koenig11122019737crashesnofix,	ALTauthor = {author},	ALTeditor = {editor},	title = {title},	date = {date},	url = {"https://www.latimes.com/business/story/2019-12-11/faa-boeing-737-max-crashes"},}
%    probleemanalyse, veiligheidsvraagstuk
%    @online{dohertylindeman15032019737problems,	ALTauthor = {author},	ALTeditor = {editor},	title = {title},	date = {date},	url = {"https://www.politico.com/story/2019/03/15/boeing-737-max-grounding-1223072"},}
%    falend toezicht
%    @online{stodder02102019corruptoversight,	ALTauthor = {author},	ALTeditor = {editor},	title = {title},	date = {date},	url = {"https://www.pogo.org/analysis/2019/10/corrupted-oversight-the-faa-boeing-and-the-737-max/"},}
%    @online{afacwaLostSafeguards,	ALTauthor = {author},	ALTeditor = {editor},	title = {title},	date = {date},	url = {"https://www.afacwa.org/the_inside_story_of_mcas_seattle_times"},}
%    doelstellingen en veiligheidsvraagstukken
%    @online{swayne18032019profitssafety,	ALTauthor = {author},	ALTeditor = {editor},	title = {title},	date = {date},	url = {"https://www.marxist.com/737-max-scandal-boeing-putting-profits-before-safety.htm"},}
%    @online{freed26022021liftaustraliaban,	ALTauthor = {author},	ALTeditor = {editor},	title = {title},	date = {date},	url = {"https://finance.yahoo.com/news/australia-lifts-ban-boeing-737-035817682.html?guccounter=1&guce_referrer=aHR0cHM6Ly93d3cuZ29vZ2xlLmNvbS8&guce_referrer_sig=AQAAAHZCJYy_0A5VS2WiPoCvH4xdrRNkmkdsv5EWJ2RLIz_AS-rxsTty6AF1_HlmJiRyWYqCXDi4p0Xs4isYkNkCq2Pfo-pQ60Xz_IfTNjm4FgoZiBMC4zpZlB6F0fwecrjE_ujAXZzG4xPJnWCd8-G3VLlPTY8h3H31eQ1i8hY9AIyy"},}
%    autoriteiten krijgen tik op de vingers
%    @online{reed15032019softwareattention,	ALTauthor = {author},	ALTeditor = {editor},	title = {title},	date = {date},	url = {"https://medium.com/@jpaulreed/the-737max-and-why-software-engineers-should-pay-attention-a041290994bd"},}
%    @online{news17032019softwareexplains,	ALTauthor = {author},	ALTeditor = {editor},	title = {title},	date = {date},	url = {"https://news.ycombinator.com/item?id=19414775"},}
%    @online{legget21122020eu737maxsafe,	ALTauthor = {author},	ALTeditor = {editor},	title = {title},	date = {date},	url = {"https://www.bbc.com/news/55366320"},}
%    @online{marketscreener0103221737chinarecertification,	ALTauthor = {author},	ALTeditor = {editor},	title = {title},	date = {date},	url = {"https://www.marketscreener.com/news/latest/China-studies-Boeing-737-MAX-recertification-wants-safety-concerns-fully-addressed--32569125/"},}
%    motor in brand
%    @online{euractiv22022021737firegrounds,	ALTauthor = {author},	ALTeditor = {editor},	title = {title},	date = {date},	url = {"https://www.euractiv.com/section/aviation/news/boeing-grounds-777s-after-engine-fire/"},}
%    @online{benny18022019737returnUAE,	ALTauthor = {author},	ALTeditor = {editor},	title = {title},	date = {date},	url = {"https://gulfnews.com/business/aviation/uae-airspace-to-see-return-of-boeing-737-max-1.1613627548923"},}
%    motor in brand gevlogen
%    @online{biersmichel22022021777grounds,	ALTauthor = {author},	ALTeditor = {editor},	title = {title},	date = {date},	url = {"https://techxplore.com/news/2021-02-boeing-urges-grounding-777s.html"},}
%    @online{ID,	ALTauthor = {author},	ALTeditor = {editor},	title = {title},	date = {date},	url = {"https://www.politico.eu/article/uk-temporarily-bans-some-boeing-aircraft-after-pratt-whitney-engine-incidents/"},}
%    @online{reuters23022021777metalfatigue,	ALTauthor = {author},	ALTeditor = {editor},	title = {title},	date = {date},	url = {"https://www.timeslive.co.za/news/world/2021-02-23-damage-to-united-boeing-777-engine-consistent-with-metal-fatigue--ntsb/"},}
%    faa was niet kritisch genoeg
%    @online{ID,	ALTauthor = {author},	ALTeditor = {editor},	title = {title},	date = {date},	url = {"https://federalnewsnetwork.com/government-news/2021/02/federal-watchdog-blasts-faa-over-certification-of-boeing-jet/"},}
%    
%    
%    
%    
%    
%    
%    
%    \cite{bnnvara13062018malirapport}
%    \cite{eucal11012021malimissieverlengd}
%    \cite{nos21052014zorgenmalimissie}
%    \cite{meijnders}
%    \cite{bnrwebredactie}
%    \cite{keultjes01062016malimissiecoalitie}
%    \cite{veenhof18012019}
%    
%    \cite{isitman06012016militair}
%    \cite{nporadio11072016filmdemissie}
%    \cite{parlementairmonitor15122013mortierongeluk}
%    
%    sollicitatie
%    de bureaucratie
%    aankomst
%    interview van de burgerbevolking
%    steun van de bevolking minuut 15:00
%    de organisatie minuut 23:00
%    De militaire briefing minuut 34:00
%    prioriteit minuut 39:00
%    briefing minuut 40:00
%    de communicatie met ministerie over inlichten minuut 44:00
%    @online{DemissieFilm,	ALTauthor = {author},	ALTeditor = {editor},	title = {title},	date = {date},	url = {"https://www.2doc.nl/documentaires/series/2doc/2016/juli/de-missie.html"},}
%    \cite{DemissieFilm}
%    
%    
%    
%    
%    
%    
%    verhaal van brandweermannen
%    
%    \cite{staff31082015tanjinblastunrevealed}
%    artikel
%    
%    \cite{chinafile18082015tanjinexplosion}
%    invloed van social media
%    \bibitem{ID} ... \LaTeX:\\ \url{https://www.economist.com/asia/2015/08/18/a-blast-in-tianjin-sets-off-an-explosion-online}
%    \cite{}
%    \bibitem{ID} ... \LaTeX:\\ \url{https://america.cgtn.com/2015/08/12/explosion-reported-in-tianjin-china}
%    \cite{}
%    \bibitem{ID} ... \LaTeX:\\ \url{https://factcheck.afp.com/no-photo-was-taken-chinese-city-tianjin-august-2015}
%    \cite{}
%    vergelijking van twee rampen
%    \bibitem{ID} ... \LaTeX:\\ \url{https://airshare.air-inc.com/how-does-the-beirut-explosion-compare-to-tianjin}
%    \cite{}
%    overheid en media
%    \bibitem{ID} ... \LaTeX:\\ \url{https://newbloommag.net/2015/08/17/tianjin-explosion/}
%    \cite{}
%    chemische industrie ondeer de loep
%    \bibitem{ID} ... \LaTeX:\\ \url{https://www.voanews.com/east-asia-pacific/tianjin-blast-puts-spotlight-chemical-industry}
%    \cite{}
%    \bibitem{ID} ... \LaTeX:\\ \url{https://abcnews.go.com/International/apocalyptic-aftermath-devastating-images-tianjin-china-explosions/story?id=33057017}
%    \cite{}
%    \bibitem{ID} ... \LaTeX:\\ \url{https://www.reachingoutacrossdurham.co.uk/osk/tianjin-explosion-2021}
%    \cite{}
%    \bibitem{pinghuang2410201TanjinFactreport} ... \LaTeX:\\ \url{https://aiche.onlinelibrary.wiley.com/doi/abs/10.1002/prs.11789}
%    \cite{pinghuang2410201TanjinFactreport}
%    \bibitem{ID} ... \LaTeX:\\ \url{https://www.automotivelogistics.media/thousands-of-cars-destroyed-in-tianjin-port-explosions/13570.article}
%    \cite{}
%    \bibitem{ID} ... \LaTeX:\\ \url{https://www.joc.com/port-news/asian-ports/port-tianjin/tianjin-port-explosions-could-be-most-expensive-maritime-disaster_20150826.html}
%    \cite{}
%    \bibitem{ID} ... \LaTeX:\\ \url{https://www.bloomberg.com/news/articles/2015-08-12/explosion-in-northern-china-shatters-windows-causes-injuries}
%    \cite{}
%    \bibitem{ID} ... \LaTeX:\\ \url{https://unece.org/fileadmin/DAM/env/documents/2016/TEIA/OECD_WGCA_24-27_OCT_2016/Session_3_Zhao_-__Introduction_of_Tianjin_Accident_-_Jinsong_Zhao.pdf}
%    \cite{}
%    gemaakte fouten
%    \bibitem{portoTanjinExplosionSight} ... \LaTeX:\\ \url{https://porteconomicsmanagement.org/pemp/contents/part6/port-resilience/site-2015-tianjin-port-explosions/}
%    \cite{portoTanjinExplosionSight}
%    \bibitem{ID} ... \LaTeX:\\ \url{https://www.alamy.com/stock-image-tianjin-china-17th-aug-2015-tianjin-explosion-aftermath-blast-site-165334778.html}
%    \cite{}
%    \bibitem{ID} ... \LaTeX:\\ \url{https://www.popularmechanics.com/technology/news/a16871/massive-explosions-china-city-of-tianjin/}
%    \cite{}
%    \bibitem{imago17082015TanjinApartmentImages} ... \LaTeX:\\ \url{https://www.imago-images.com/st/0080815934}
%    \cite{imago17082015TanjinApartmentImages}
%    \bibitem{trager14082015Chemicalblast} ... \LaTeX:\\ \url{https://www.chemistryworld.com/news/deadly-chemical-blast-at-chinese-port/8857.article}
%    \cite{trager14082015Chemicalblast}
%    \bibitem{pangeramo27082015TanjinExplosion} ... \LaTeX:\\ \url{https://www.process-worldwide.com/tianjin-explosion-from-chemical-perspective-insights-and-backgrounds-a-502381/}
%    \cite{pangeramo27082015TanjinExplosion}
%    vergelijking met andere explosies
%    \bibitem{ap06082020ammaniumnitrate} ... \LaTeX:\\ \url{https://apnews.com/article/lebanon-fires-us-news-explosions-middle-east-53f4206a7f1db0812262a15d22e1e58f}
%    \cite{ap06082020ammaniumnitrate}
%    invloed van de ramp op de industrie
%    \bibitem{morris14082015TanjinIndustryImpact} ... \LaTeX:\\ \url{https://fortune.com/2015/08/14/tianjin-port-explosion-shipping-delays/}
%    \cite{morris14082015TanjinIndustryImpact}
%    is er sprake van een doofpot
%    \bibitem{milesyu20082015exposingtoxicgovlines} ... \LaTeX:\\ \url{https://www.washingtontimes.com/news/2015/aug/20/inside-china-tianjin-explosions-cover-up-exposes-b/}
%    \cite{milesyu20082015exposingtoxicgovlines}
%    eigendomsverzekering
%    \bibitem{artemis30032016tanjininsurance} ... \LaTeX:\\ \url{https://www.artemis.bm/news/tianjin-explosions-property-insurance-loss-could-reach-3-5bn-swiss-re/}
%    \cite{artemis30032016tanjininsurance}
%    \bibitem{aidenxiatanjinblast} ... \LaTeX:\\ \url{https://www.thechinastory.org/yearbooks/yearbook-2015/forum-the-abyss-%E5%9D%8E/tianjin-explosions/}
%    \cite{aidenxiatanjinblast}
%    effecten op de lange termijn
%    \bibitem{danwangTanjinflexreport} ... \LaTeX:\\ \url{https://www.flexport.com/blog/tianjin-explosion-effect-on-supply-chains/}
%    \cite{danwangTanjinflexreport}
%    \bibitem{keyHighlightsTanjin} ... \LaTeX:\\ \url{https://www.cicm.org.my/images/articles/CICM-Article-on-Tianjin-Blast-Oct2015.pdf}
%    \cite{keyHighlightsTanjin}
%    lessons learned
%    \bibitem{ID} ... \LaTeX:\\ \url{https://www.genre.com/knowledge/blog/lessons-from-the-tianjin-explosion-en.html}
%    \cite{}
%    \bibitem{ID} ... \LaTeX:\\ \url{https://www.ft.com/content/ad62904c-44ce-11e5-b3b2-1672f710807b}
%    \cite{}
%    \bibitem{hartley13082015videofootage} ... \LaTeX:\\ \url{https://www.huffingtonpost.co.uk/2015/08/13/tianjin-explosion-china-shocking-footage-caught-on-camera_n_7980888.html}
%    \cite{hartley13082015videofootage}
%    \bibitem{odonnel01062017firetanjinblast2015} ... \LaTeX:\\ \url{https://www.thatsmags.com/china/post/19189/massive-fire-rocks-tianjin-port}
%    \cite{odonnel01062017firetanjinblast2015}
%    gevolgen voor de industrie
%    \bibitem{ID} ... \LaTeX:\\ \url{https://www.everstream.ai/risk-center/special-reports/the-jiangsu-yancheng-explosion/}
%    \cite{}
%    \bibitem{fan15082015newyorkermistrustchina} ... \LaTeX:\\ \url{https://www.newyorker.com/news/news-desk/after-tianjin-an-outbreak-of-mistrust-in-china}
%    \cite{fan15082015newyorkermistrustchina}
%    framing vanuit de chinese media
%    \bibitem{yanlidongchinamediaframingTanjin} ... \LaTeX:\\ \url{https://www.neliti.com/publications/101997/the-chinese-media-framing-of-the-2015s-tianjin-explosion}
%    \cite{yanlidongchinamediaframingTanjin}
%    \bibitem{evans27092017TnjinInsurance} ... \LaTeX:\\ \url{https://www.reinsurancene.ws/chinese-insurers-settle-1-5-billion-tianjin-blast-claims/}
%    \cite{evans27092017TnjinInsurance}
%    niewsartikel
%    \bibitem{jasi26032019chineschemplant} ... \LaTeX:\\ \url{https://www.thechemicalengineer.com/news/update-78-confirmed-dead-after-chinese-chemicals-plant-explosion/}
%    \cite{jasi26032019chineschemplant}
%    \bibitem{shiqingTanjinExecutiveSentence} ... \LaTeX:\\ \url{https://www.caixinglobal.com/2016-11-10/chinese-executive-receives-suspended-death-sentence-over-2015-tianjin-warehouse-blast-101006325.html}
%    \cite{shiqingTanjinExecutiveSentence}
%    toegang tot de ramplplek vanuit de okale journalistiek
%    \bibitem{sophiebeach15082015} ... \LaTeX:\\ \url{https://chinadigitaltimes.net/2015/08/he-xiaoxin-how-far-can-i-go-and-how-much-can-i-do/}
%    \cite{sophiebeach15082015}
%    artikel
%    \bibitem{ID} ... \LaTeX:\\ \url{https://www.wnpr.org/post/china-examines-aftermath-immense-twin-explosions-killed-dozens}
%    \cite{}
%    \bibitem{hamzeh05082020BeirutBlast} ... \LaTeX:\\ \url{https://theconversation.com/what-is-ammonium-nitrate-the-chemical-that-exploded-in-beirut-143979}
%    \cite{hamzeh05082020BeirutBlast}
%    \bibitem{chemwatch18082015TanjiinExplosion} ... \LaTeX:\\ \url{https://chemicalwatch.com/36730/nationwide-inspections-in-china-follow-tianjin-explosion}
%    \cite{chemwatch18082015TanjiinExplosion}
%    \bibitem{thehindu15062019chinaExplosion} ... \LaTeX:\\ \url{https://www.thehindu.com/news/international/investigation-begun-into-china-gas-explosion-as-toll-rises/article34818324.ece}
%    \cite{thehindu15062019chinaExplosion}
%    \bibitem{santagotimes24032019chinablast} ... \LaTeX:\\ \url{https://santiagotimes.cl/2019/03/24/64-killed-600-injured-in-china-chemical-plant-blast/}
%    \cite{santagotimes24032019chinablast}
%    oorzaken
%    \bibitem{klingecorp28042020causedTanjin} ... \LaTeX:\\ \url{https://klingecorp.com/blog/what-caused-the-tianjin-explosions/}
%    \cite{klingecorp28042020causedTanjin}
%    case study
%    \bibitem{mcgarryExplosions2017 ... \LaTeX:\\ \url{https://www.preventionweb.net/educational/view/57235}
%    	\cite{mcgarryExplosions2017}
%    	niewsartikel
%    	\bibitem{roswnfeld13082015TanjinReports} ... \LaTeX:\\ \url{https://www.cnbc.com/2015/08/12/explosion-in-tianjin-china.html}
%    	\cite{roswnfeld13082015TanjinReports}
%    	chronologische uiteenzetting
%    	\bibitem{aria12082015explosionaTanjin} ... \LaTeX:\\ \url{https://www.aria.developpement-durable.gouv.fr/wp-content/files_mf/A46803_a46803_fiche_impel_006.pdf}
%    	\cite{aria12082015explosionaTanjin}
%    	corruptie
%    	\bibitem{ID} ... \LaTeX:\\ \url{https://www.nytimes.com/2015/08/31/world/asia/behind-tianjin-tragedy-a-company-that-flouted-regulations-and-reaped-profits.html}
%    	\cite{}
%    	mismanagement als oorzaak
%    	\bibitem{ID} ... \LaTeX:\\ \url{https://www.nytimes.com/2016/02/06/world/asia/tianjin-explosions-were-result-of-mismanagement-china-finds.html}
%    	\cite{}
%    	\bibitem{ID} ... \LaTeX:\\ \url{https://cen.acs.org/articles/94/web/2016/02/Chinese-Investigators-Identify-Cause-Tianjin.html}
%    	\cite{}
%    	autoriteiten publiceren onderoeksrapport
%    	\bibitem{tremblay11022016chineseInvestigatorsTanjin} ... \LaTeX:\\ \url{https://cen.acs.org/articles/94/i7/Chinese-Investigators-Identify-Cause-Tianjin.html}
%    	\cite{tremblay11022016chineseInvestigatorsTanjin}
%    	fotos van de rampplek
%    	\bibitem{taylor13082015TanjinExplosianAftermath} ... \LaTeX:\\ \url{https://www.theatlantic.com/photo/2015/08/photos-of-the-aftermath-of-the-massive-explosions-in-tianjin-china/401228/}
%    	\cite{taylor13082015TanjinExplosianAftermath}
%    	\bibitem{ID} ... \LaTeX:\\ \url{https://edition.cnn.com/2015/08/13/asia/china-tianjin-explosions/index.html}
%    	\cite{}
%    	niuwesartiekel}
%    \bibitem{associatedPresss13082013} ... \LaTeX:\\ \url{https://www.cbc.ca/news/world/china-explosion-tianjin-1.3189455}
%    \cite{associatedPresss13082013}
%    verantwoordelijke
%    \bibitem{ID} ... \LaTeX:\\ \url{https://www.thestar.com/news/world/2016/11/09/chinese-executive-gets-death-sentence-over-tianjin-explosion-in-2015.html}
%    \cite{}
%    risicobeperking/controle
%    \bibitem{ID} ... \LaTeX:\\ \url{https://www.swissre.com/en/china/news-insights/articles/analysis-of-tianjin-port-explosion-china.html}
%    \cite{}
%    censuur
%    \bibitem{ID} ... \LaTeX:\\ \url{https://foreignpolicy.com/2015/09/10/censored-china-young-survivor-tianjin-explosion-viral-post/}
%    \cite{}
%    censuur
%    \bibitem{ID} ... \LaTeX:\\ \url{https://qz.com/756872/a-year-after-the-tianjin-blast-public-mourning-and-discussion-about-it-are-still-censored-in-china/}
%    \cite{}
%    verschillende artikelen
%    \bibitem{ID} ... \LaTeX:\\ \url{https://www.scmp.com/topics/tianjin-warehouse-explosion-2015}
%    \cite{}
%    \bibitem{ID} ... \LaTeX:\\ \url{https://www.wsj.com/articles/BL-CJB-27664}
%    \cite{}
%    \bibitem{ID} ... \LaTeX:\\ \url{https://www.nbcnews.com/news/world/tianjin-explosions-californian-witness-filmed-dramatic-china-blasts-n409701}
%    \cite{}
%    \bibitem{ID} ... \LaTeX:\\ \url{https://ui.adsabs.harvard.edu/abs/2016AGUFM.S13D..06P/abstract}
%    afwikkeling van de ramp
%    \cite{}
%    \bibitem{ID} ... \LaTeX:\\ \url{https://chinadialogue.net/en/pollution/9188-back-to-the-blast-zone-one-year-after-the-tianjin-explosion/}
%    \cite{}
%    \bibitem{ID} ... \LaTeX:\\ \url{https://www.wired.com/2015/08/chinas-huge-tianjin-explosion-looked-like-space/}
%    \cite{}
%    \bibitem{ID} ... \LaTeX:\\ \url{https://www.abc.net.au/news/2015-08-13/explosion-rocks-north-chinese-city-of-tianjin/6693336?nw=0}
%    \cite{}
%    ambtenaren onderzocht
%    
%    risico-inschatting
%    \bibitem{ID} ... \LaTeX:\\ \url{https://www.mdpi.com/2071-1050/12/3/1169/htm}
%    \cite{}
%    \bibitem{ID} ... \LaTeX:\\ \url{https://www.mdpi.com/2071-1050/12/3/1169/htm}
%    \cite{}
%    \bibitem{ID} ... \LaTeX:\\ \url{https://www.cbsnews.com/news/tianjin-port-china-massive-explosion-hundreds-injured/}
%    \cite{}
%    \bibitem{ID} ... \LaTeX:\\ \url{https://www.hkjcdpri.org.hk/download/casestudies/Tianjin_CASE.pdf}
%    \cite{}
%    \bibitem{ID} ... \LaTeX:\\ \url{https://time.com/3996168/tianjin-explosion-china-pictures/}
%    \cite{}
%    onderzoeksrapport
%    \bibitem{ID} ... \LaTeX:\\ \url{https://www.hfw.com/Tianjin-Port-explosion-August-2015}
%    \cite{}
%    \bibitem{un20082015InvestigationTanjin} ... \LaTeX:\\ \url{https://news.un.org/en/story/2015/08/506912-following-tianjin-explosion-un-expert-calls-china-ensure-transparent}
%    \cite{un20082015InvestigationTanjin}
%    \bibitem{france2412082015TnjinExplosion} ... \LaTeX:\\ \url{https://www.france24.com/en/20150812-huge-explosions-rock-chinese-city-tianjin}
%    \cite{france2412082015TnjinExplosion}
%    \bibitem{npr14082015TanjinCause} ... \LaTeX:\\ \url{https://choice.npr.org/index.html?origin=https://www.npr.org/2015/08/14/432280627/what-caused-the-warehouse-explosions-in-tianjin-china}
%    \cite{npr14082015TanjinCause}
%    123 verantwoordelijken
%    \bibitem{bbc05022016TanjinResponsibles} ... \LaTeX:\\ \url{https://www.bbc.com/news/world-asia-china-35506311}
%    \cite{bbc05022016TanjinResponsibles}
%    \bibitem{ID} ... \LaTeX:\\ \url{https://www.washingtonpost.com/gdpr-consent/?next_url=https%3a%2f%2fwww.washingtonpost.com%2fnews%2fworldviews%2fwp%2f2015%2f08%2f12%2fvideos-show-chinese-city-of-tianjin-rocked-by-enormous-explosion%2f}
%    \cite{}
%    lang artiekel
%    \bibitem{CBodeen15082015TanjinExplosion} ... \LaTeX:\\ \url{https://www.businessinsider.com/the-chemical-explosion-in-china-killed-more-than-100-people-and-the-devastation-is-unreal-2015-8?international=true&r=US&IR=T}
%    \cite{CBodeen15082015TanjinExplosion}
%    \bibitem{ID} ... \LaTeX:\\ \url{https://pubmed.ncbi.nlm.nih.gov/27311537/}
%    \cite{}
%    \bibitem{reutersTanjinInsurance} ... \LaTeX:\\ \url{https://www.reuters.com/article/us-china-blast-insurance-idUSKCN0QM0N220150817}
%    \cite{reutersTanjinInsurance}
%    \bibitem{yu082016evaluationTanjin2015} ... \LaTeX:\\ \url{https://www.sciencedirect.com/science/article/abs/pii/S0305417916300079}
%    \cite{yu082016evaluationTanjin2015}
%    \bibitem{wiki2015TanjinExplosions} ... \LaTeX:\\ \url{https://en.wikipedia.org/wiki/2015_Tianjin_explosions}
%    \cite{wiki2015TanjinExplosions}
%    \bibitem{bbc17082015whathappenedTanjin} ... \LaTeX:\\ \url{https://www.bbc.com/news/world-asia-china-33844084}
%    \cite{bbc17082015whathappenedTanjin}
%    \bibitem{mortimer19082016taijinexplosioncrater} ... \LaTeX:\\ \url{https://www.independent.co.uk/news/world/asia/tianjin-explosion-photos-china-chemical-factory-accident-crater-revealed-a7199591.html}
%    \cite{mortimer19082016taijinexplosioncrater}
%    veiigheidshandhaving
%    \bibitem{internationallabourofficeChmControlTooliit} ... \LaTeX:\\ \url{https://www.ilo.org/legacy/english/protection/safework/ctrl_banding/toolkit/main_guide.pdf}
%    \cite{internationallabourofficeChmControlTooliit}
%    \bibitem{ID} ... \LaTeX:\\ \url{https://echa.europa.eu/documents/10162/21332507/guide_chemical_safety_sme_en.pdf}
%    \cite{}
%    \bibitem{euTaxationCustomsICSC} ... \LaTeX:\\ \url{https://ec.europa.eu/taxation_customs/dds2/SAMANCTA/EN/Safety/AppendixD_EN.htm}
%    \cite{euTaxationCustomsICSC}
%    \bibitem{iloWHOChemSafetyCards} ... \LaTeX:\\ \url{https://www.ilo.org/safework/info/publications/WCMS_113134/lang--en/index.htm}
%    \cite{iloWHOChemSafetyCards}
%    
%    
%    
%    Wat is er gebeurd?
%    @online{wikiSchipholbrand,	ALTauthor = {author},	ALTeditor = {editor},	title = {title},	date = {date},	url = {"https://nl.wikipedia.org/wiki/Schipholbrand"},}
%    artikel
%    @online{schipholbrand27102005video,	ALTauthor = {author},	ALTeditor = {editor},	title = {title},	date = {date},	url = {"https://www.youtube.com/watch?v=1i-hfEzxFfk"},}
%    psychologische gevolgen
%    rapport
%    @online{onderzoeksraad2610schipholoost,	ALTauthor = {author},	ALTeditor = {editor},	title = {title},	date = {date},	url = {"https://www.onderzoeksraad.nl/nl/page/392/brand-cellencomplex-schiphol-oost-nacht-van-26-op-27-oktober"},}
%    artikel met video
%    herdenking
%    impact op de persoon
%    herdenking
%    @online{schipholbrandvideoargos,	ALTauthor = {author},	ALTeditor = {editor},	title = {title},	date = {date},	url = {"https://www.vpro.nl/argos/speel~POMS_VPRO_461907~schadevergoeding-voor-ex-verdachte-schipholbrand~.html"},}
%    chronologie
%    @online{nunl30052023feitenoverzicht,	ALTauthor = {author},	ALTeditor = {editor},	title = {title},	date = {date},	url = {"https://www.nu.nl/binnenland/3355935/feitenoverzicht-schipholbrand-en-rechtszaken.html"},}
%    tijdlijn
%    @online{ID,	ALTauthor = {author},	ALTeditor = {editor},	title = {title},	date = {date},	url = {"https://www.singeluitgeverijen.nl/isbn/de-schipholbrand/"},}
%    vervolgens van ministers
%    beeldanalyse en reconstructie
%    @online{ID,	ALTauthor = {author},	ALTeditor = {editor},	title = {title},	date = {date},	url = {"https://eenvandaag.avrotros.nl/item/schipholbrand-niet-ontstaan-in-cel-11/"},}
%    herdenking
%    korte samenvatting
%    rapport
%    artikel
%    verwijzing naar het rapport vanuit de politieke oppositie
%    beeld vanuit de gevangenisbewaarder
%    nationaliteit slachtoffers schipholbrand
%    verblijfsvergunning voor de slachtoffers
%    gen schadevergoeding voor de verdachte
%    verdachte voor de rechter
%    geen schadevergoeding voor verdachte
%    artikel wat ging er mis bji de schipholbrand
%    brand veroorzaakt door een peuk
%    smaadschrift
%    bewakers worden niet vervolgd
%    proces schipholbrand moet over en de brandveilgheid moet worden verbeterd
%    de rol van het parlement in de evaluatie
%    @online{parlementairemonitorschipholbrand,	ALTauthor = {author},	ALTeditor = {editor},	title = {title},	date = {date},	url = {"https://www.parlementairemonitor.nl/9353000/1/j9vvij5epmj1ey0/vi3aof7awcxg"},}
%    onderzoeksmemo
%    herdenking
%    @online{ID,	ALTauthor = {author},	ALTeditor = {editor},	title = {title},	date = {date},	url = {"https://archief.ntr.nl/nova/page/detail/uitzendingen/3847/Den%20Haag%20Vandaag_%20herdenking%20Schipholbrand.html"},}
%    herdenking
%    invloed van de ramp op samenleving
%    @online{videonpoNOVA13112008,	ALTauthor = {author},	ALTeditor = {editor},	title = {title},	date = {date},	url = {"https://www.npostart.nl/heropen-onderzoek-schipholbrand/13-11-2008/POMS_NTR_103332"},}
%    opmerkelijk rapport gestolen in de nasleep
%    @online{rizoomes01052014schipholbrand,	ALTauthor = {author},	ALTeditor = {editor},	title = {title},	date = {date},	url = {"https://www.rizoomes.nl/brandweer/brand-cellencomplex-schiphol/"},}
%    
%    
%    
%    publicaties
%    @online{heuvelkroesschipholbrandcamerabeelden,	ALTauthor = {author},	ALTeditor = {editor},	title = {title},	date = {date},	url = {"http://www.msnp.nl/downloads/Onderzoeksmemo%20beeldanalyse%20Schipholbrand%20prot.pdf"},}
%    Wat waren de regels destijds?
%    Waren de autoriteiten in staat om op tijd in te grijpen of om erger te voorkomen?
%    Wat is er gedaan om de veiligheid van illegalen en gevangenissbewaarders te verbeteren
%    
%    
%    
%    algemene vragen
%    oorzaken
%    \bibitem{gates18112020boeingcrisis} ... \LaTeX:\\ \url{https://www.seattletimes.com/business/boeing-aerospace/what-led-to-boeings-737-max-crisis-a-qa/}
%    \cite{gates18112020boeingcrisis}
%    \bibitem{boeing737maxsoftwareprobles} ... \LaTeX:\\ \url{https://www.schneier.com/blog/archives/2019/04/excellent_analy.html}
%    \cite{boeing737maxsoftwareprobles}
%    fout in de software
%    \bibitem{avetisov19032019boeingmalwarestate} ... \LaTeX:\\ \url{https://www.forbes.com/sites/georgeavetisov/2019/03/19/malware-at-30000-feet-what-the-737-max-says-about-the-state-of-airplane-software-security/?sh=4d26f7052a9e}
%    \cite{avetisov19032019boeingmalwarestate}
%    het nationaal veiligheidsbelang
%    \bibitem{thompson23112020nationalsecurityboeing} ... \LaTeX:\\ \url{https://www.forbes.com/sites/lorenthompson/2020/11/23/five-reasons-return-of-boeings-737-max-to-service-is-important-to-national-security/?sh=2128ea552018}
%    \cite{thompson23112020nationalsecurityboeing}
%    falend toezicht
%    \bibitem{gates21032019FAAControlSystem} ... \LaTeX:\\ \url{https://www.seattletimes.com/business/boeing-aerospace/failed-certification-faa-missed-safety-issues-in-the-737-max-system-implicated-in-the-lion-air-crash/}
%    \cite{gates21032019FAAControlSystem}
%    onderzoeksrapport
%    \bibitem{faa18112020boeingreview} ... \LaTeX:\\ \url{https://www.faa.gov/foia/electronic_reading_room/boeing_reading_room/media/737_RTS_Summary.pdf}
%    \cite{faa18112020boeingreview}
%    \bibitem{wiki737maxgroundings} ... \LaTeX:\\ \url{https://en.wikipedia.org/wiki/Boeing_737_MAX_groundings}
%    \cite{wiki737maxgroundings}
%    veiligheidsrisico's
%    menselijke fouten
%    \bibitem{campbell02052019boengcrashhumanerrors} ... \LaTeX:\\ \url{https://www.theverge.com/2019/5/2/18518176/boeing-737-max-crash-problems-human-error-mcas-faa}
%    \cite{campbell02052019boengcrashhumanerrors}
%    overzicht van crashes
%    \bibitem{hawkins22032019737maxairplanes} ... \LaTeX:\\ \url{https://www.theverge.com/2019/3/22/18275736/boeing-737-max-plane-crashes-grounded-problems-info-details-explained-reasons}
%    \cite{hawkins22032019737maxairplanes}
%    veiligheidsopmerking
%    \bibitem{thomas30082020737safest} ... \LaTeX:\\ \url{https://www.airlineratings.com/news/boeings-737-max-will-one-safest-aircraft-history/}
%    \cite{thomas30082020737safest}
%    aanpassingen
%    \bibitem{boeing737maxdisplay} ... \LaTeX:\\ \url{https://www.boeing.com/commercial/737max/737-max-software-updates.page}
%    \cite{boeing737maxdisplay}
%    waarschuwingen//output signalen
%    \bibitem{fehrm24112020737changes} ... \LaTeX:\\ \url{https://leehamnews.com/2020/11/24/boeing-737-max-changes-beyond-mcas/}
%    \cite{fehrm24112020737changes}
%    software gerelateerde fouten
%    \bibitem{travis18042019737maxsoftwaredevop} ... \LaTeX:\\ \url{https://spectrum.ieee.org/aerospace/aviation/how-the-boeing-737-max-disaster-looks-to-a-software-developer}
%    \cite{travis18042019737maxsoftwaredevop}
%    onderzoeksrapport
%    de rol van de publieke opinie
%    \bibitem{barnett05052019737maxcrisis} ... \LaTeX:\\ \url{https://pubsonline.informs.org/do/10.1287/orms.2019.05.05/full/}
%    \cite{barnett05052019737maxcrisis}
%    onderzoek van europese luchtvaart agentschap
%    \bibitem{easa27012021737maxsafereturn} ... \LaTeX:\\ \url{https://www.easa.europa.eu/newsroom-and-events/news/easa-declares-boeing-737-max-safe-return-service-europe}
%    \cite{easa27012021737maxsafereturn}
%    \veiligheidsvraagstuk
%    \bibitem{touitou11032019737tragedies} ... \LaTeX:\\ \url{https://phys.org/news/2019-03-boeing-max-safety-tragedies.html}
%    \cite{touitou11032019737tragedies}
%    artikel over sensoren
%    \bibitem{hemmerdinger02022021737maxdeliveries} ... \LaTeX:\\ \url{https://www.flightglobal.com/airframers/boeing-delays-737-max-10-deliveries-two-years-to-2023/142245.article}
%    \cite{hemmerdinger02022021737maxdeliveries}
%    goedkeuring van europese luchtvaart autoriteiten
%    advies aan de faa
%    \bibitem{bielby27022021faaimprovesafety} ... \LaTeX:\\ \url{https://www.hstoday.us/subject-matter-areas/airport-aviation-security/oig-tells-faa-to-improve-safety-oversight-following-boeing-737-max-review/}
%    \cite{bielby27022021faaimprovesafety}
%    \bibitem{boyle18112020737maxupgrade} ... \LaTeX:\\ \url{https://www.geekwire.com/2020/faas-go-ahead-737-maxs-return-flight-kicks-off-massive-software-upgrade/}
%    \cite{boyle18112020737maxupgrade}
%    \bibitem{bergstraburgess122019737maxMcasAlgorithm} ... \LaTeX:\\ \url{https://www.researchgate.net/publication/338420944_A_Promise_Theoretic_Account_of_the_Boeing_737_Max_MCAS_Algorithm_Affair}
%    \cite{bergstraburgess122019737maxMcasAlgorithm}
%    achtergrond informatie
%    \bibitem{737mcas} ... \LaTeX:\\ \url{http://www.b737.org.uk/mcas.htm}
%    \cite{737mcas}
%    algemeen vertrouwen
%    \bibitem{newburger17052019boeingcrisis} ... \LaTeX:\\ \url{https://www.cnbc.com/2019/05/16/what-you-need-to-know-about-boeings-737-max-crisis.html}
%    \cite{newburger17052019boeingcrisis}
%    toestemming europese autoriteiten
%    problemen
%    \bibitem{arstechnica22012020737problems} ... \LaTeX:\\ \url{https://arstechnica.com/information-technology/2020/01/737-max-fix-slips-to-summer-and-thats-just-one-of-boeings-problems/}
%    \cite{arstechnica22012020737problems}
%    uitgebreid artikel over de onderzoeken en het vliegverbod
%    \bibitem{german190620217372yaftergrounded} ... \LaTeX:\\ \url{https://www.cnet.com/news/boeing-737-max-8-all-about-the-aircraft-flight-ban-and-investigations/}
%    \cite{german190620217372yaftergrounded}
%    computers als oorzaak
%    lessons learned
%    \bibitem{beningo02052019boeinglessons} ... \LaTeX:\\ \url{https://www.designnews.com/electronics-test/5-lessons-learn-boeing-737-max-fiasco}
%    \cite{beningo02052019boeinglessons}
%    \bibitem{duran05042019boeingspof} ... \LaTeX:\\ \url{https://www.eurocontrol.int/publication/effects-network-extra-standby-aircraft-and-boeing-737-max-grounding}
%    \cite{duran05042019boeingspof}
%    single point of failure
%    \bibitem{ID} ... \LaTeX:\\ \url{https://dmd.solutions/blog/2019/04/05/how-a-single-point-of-failure-spof-in-the-mcas-software-could-have-caused-the-boeing-737-max-crash-in-ethiopia/}
%    \cite{}
%    \bibitem{makichuck24012021737fearflying} ... \LaTeX:\\ \url{https://asiatimes.com/2021/01/boeings-737-max-and-the-fear-of-flying/}
%    \cite{makichuck24012021737fearflying}
%    lijst van tehnische aanpassingen
%    \bibitem{caa737modifications} ... \LaTeX:\\ \url{https://www.caa.co.uk/Consumers/Guide-to-aviation/Boeing-737-MAX/}
%    \cite{caa737modifications}
%    \bibitem{oestergaard14122020boeingdeliveries} ... \LaTeX:\\ \url{https://dsm.forecastinternational.com/wordpress/2020/12/14/airbus-and-boeing-report-november-2020-commercial-aircraft-orders-and-deliveries/}
%    \cite{oestergaard14122020boeingdeliveries}
%    code lek
%    \bibitem{reenberg787flaws} ... \LaTeX:\\ \url{https://www.wired.com/story/boeing-787-code-leak-security-flaws/}
%    \cite{reenberg787flaws}
%    \bibitem{fitch16092020737backlogrisks} ... \LaTeX:\\ \url{https://www.fitchratings.com/research/corporate-finance/boeing-737-max-return-backlog-risks-remain-16-09-2020}
%    \cite{fitch16092020737backlogrisks}
%    Cultuurverandering, deregulatie, systeemwijziging of gewoon een kwestie van competentie
%    \bibitem{willis27082020737maxfailures} ... \LaTeX:\\ \url{https://www.aerospacetestinginternational.com/features/what-broke-the-737-max.html}
%    \cite{willis27082020737maxfailures}
%    extra aanpassingen
%    \bibitem{ostrower11062020more737changes} ... \LaTeX:\\ \url{https://theaircurrent.com/aviation-safety/boeings-737-max-software-done-but-regulators-plot-more-changes-after-jets-return/}
%    \cite{ostrower11062020more737changes}
%    wat ging er mis een analyse van een ex-iloot
%    De utoriteiten waren op de hoogte
%    \bibitem{hruska13122019faaknown737crashrate} ... \LaTeX:\\ \url{https://www.extremetech.com/extreme/303373-the-faa-knew-the-737-max-was-dangerous-and-kept-it-flying-anyway}
%    \cite{hruska13122019faaknown737crashrate}
%    kwaliteiten van het alarmsysteem niet goed bekend
%    \bibitem{bloomberg26092019failedpred} ... \LaTeX:\\ \url{https://time.com/5687473/boeing-737-alarm-system/}
%    \cite{bloomberg26092019failedpred}
%    \bibitem{whiteman09072020boengcancelstock} ... \LaTeX:\\ \url{https://www.nasdaq.com/articles/boeing-gets-dealt-another-737-max-cancellation-blow.-what-it-means-for-boeing-stock-2020}
%    \cite{whiteman09072020boengcancelstock}
%    \bibitem{leopold09192019boeingreliability} ... \LaTeX:\\ \url{https://www.eetimes.com/boeing-crashes-highlight-a-worsening-reliability-crisis/}
%    \cite{leopold09192019boeingreliability}
%    veiligheidsvraagstuk
%    \bibitem{koenig11122019737crashesnofix} ... \LaTeX:\\ \url{https://www.latimes.com/business/story/2019-12-11/faa-boeing-737-max-crashes}
%    \cite{koenig11122019737crashesnofix}
%    probleemanalyse, veiligheidsvraagstuk
%    \bibitem{dohertylindeman15032019737problems} ... \LaTeX:\\ \url{https://www.politico.com/story/2019/03/15/boeing-737-max-grounding-1223072}
%    \cite{dohertylindeman15032019737problems}
%    falend toezicht
%    \bibitem{stodder02102019corruptoversight} ... \LaTeX:\\ \url{https://www.pogo.org/analysis/2019/10/corrupted-oversight-the-faa-boeing-and-the-737-max/}
%    \cite{stodder02102019corruptoversight}
%    \bibitem{afacwaLostSafeguards} ... \LaTeX:\\ \url{https://www.afacwa.org/the_inside_story_of_mcas_seattle_times}
%    \cite{afacwaLostSafeguards}
%    doelstellingen en veiligheidsvraagstukken
%    \bibitem{swayne18032019profitssafety} ... \LaTeX:\\ \url{https://www.marxist.com/737-max-scandal-boeing-putting-profits-before-safety.htm}
%    \cite{swayne18032019profitssafety}
%    \bibitem{freed26022021liftaustraliaban} ... \LaTeX:\\ \url{https://finance.yahoo.com/news/australia-lifts-ban-boeing-737-035817682.html?guccounter=1&guce_referrer=aHR0cHM6Ly93d3cuZ29vZ2xlLmNvbS8&guce_referrer_sig=AQAAAHZCJYy_0A5VS2WiPoCvH4xdrRNkmkdsv5EWJ2RLIz_AS-rxsTty6AF1_HlmJiRyWYqCXDi4p0Xs4isYkNkCq2Pfo-pQ60Xz_IfTNjm4FgoZiBMC4zpZlB6F0fwecrjE_ujAXZzG4xPJnWCd8-G3VLlPTY8h3H31eQ1i8hY9AIyy}
%    \cite{freed26022021liftaustraliaban}
%    autoriteiten krijgen tik op de vingers
%    \bibitem{reed15032019softwareattention} ... \LaTeX:\\ \url{https://medium.com/@jpaulreed/the-737max-and-why-software-engineers-should-pay-attention-a041290994bd}
%    \cite{reed15032019softwareattention}
%    \bibitem{news17032019softwareexplains} ... \LaTeX:\\ \url{https://news.ycombinator.com/item?id=19414775}
%    \cite{news17032019softwareexplains}
%    \bibitem{legget21122020eu737maxsafe} ... \LaTeX:\\ \url{https://www.bbc.com/news/55366320}
%    \cite{legget21122020eu737maxsafe}
%    \bibitem{marketscreener0103221737chinarecertification} ... \LaTeX:\\ \url{https://www.marketscreener.com/news/latest/China-studies-Boeing-737-MAX-recertification-wants-safety-concerns-fully-addressed--32569125/}
%    \cite{marketscreener0103221737chinarecertification}
%    motor in brand
%    \bibitem{euractiv22022021737firegrounds} ... \LaTeX:\\ \url{https://www.euractiv.com/section/aviation/news/boeing-grounds-777s-after-engine-fire/}
%    \cite{euractiv22022021737firegrounds}
%    \bibitem{benny18022019737returnUAE} ... \LaTeX:\\ \url{https://gulfnews.com/business/aviation/uae-airspace-to-see-return-of-boeing-737-max-1.1613627548923}
%    \cite{benny18022019737returnUAE}
%    motor in brand gevlogen
%    \bibitem{biersmichel22022021777grounds} ... \LaTeX:\\ \url{https://techxplore.com/news/2021-02-boeing-urges-grounding-777s.html}
%    \cite{biersmichel22022021777grounds}
%    \bibitem{ID} ... \LaTeX:\\ \url{https://www.politico.eu/article/uk-temporarily-bans-some-boeing-aircraft-after-pratt-whitney-engine-incidents/}
%    \cite{}
%    \bibitem{reuters23022021777metalfatigue} ... \LaTeX:\\ \url{https://www.timeslive.co.za/news/world/2021-02-23-damage-to-united-boeing-777-engine-consistent-with-metal-fatigue--ntsb/}
%    \cite{reuters23022021777metalfatigue}
%    faa was niet kritisch genoeg
%    \bibitem{ID} ... \LaTeX:\\ \url{https://federalnewsnetwork.com/government-news/2021/02/federal-watchdog-blasts-faa-over-certification-of-boeing-jet/}
%    \cite{}
%    
%    
%    
%    algemene vragen
%    oorzaken
%    \bibitem{{gates18112020boeingcrisis} ... \LaTeX:\\ \url{https://www.seattletimes.com/business/boeing-aerospace/what-led-to-boeings-737-max-crisis-a-qa/}
%    	
%    	\bibitem{{boeing737maxsoftwareprobles} ... \LaTeX:\\ \url{https://www.schneier.com/blog/archives/2019/04/excellent_analy.html}
%    		fout in de software
%    		\bibitem{{avetisov19032019boeingmalwarestate} ... \LaTeX:\\ \url{https://www.forbes.com/sites/georgeavetisov/2019/03/19/malware-at-30000-feet-what-the-737-max-says-about-the-state-of-airplane-software-security/?sh=4d26f7052a9e}
%    			het nationaal veiligheidsbelang
%    			\bibitem{{thompson23112020nationalsecurityboeing} ... \LaTeX:\\ \url{https://www.forbes.com/sites/lorenthompson/2020/11/23/five-reasons-return-of-boeings-737-max-to-service-is-important-to-national-security/?sh=2128ea552018}
%    				falend toezicht
%    				\bibitem{{gates21032019FAAControlSystem} ... \LaTeX:\\ \url{https://www.seattletimes.com/business/boeing-aerospace/failed-certification-faa-missed-safety-issues-in-the-737-max-system-implicated-in-the-lion-air-crash/}
%    					onderzoeksrapport
%    					\bibitem{{faa18112020boeingreview} ... \LaTeX:\\ \url{https://www.faa.gov/foia/electronic_reading_room/boeing_reading_room/media/737_RTS_Summary.pdf}
%    						\bibitem{{wiki737maxgroundings} ... \LaTeX:\\ \url{https://en.wikipedia.org/wiki/Boeing_737_MAX_groundings}
%    							veiligheidsrisico's
%    							menselijke fouten
%    							\bibitem{{campbell02052019boengcrashhumanerrors} ... \LaTeX:\\ \url{https://www.theverge.com/2019/5/2/18518176/boeing-737-max-crash-problems-human-error-mcas-faa}
%    								overzicht van crashes
%    								\bibitem{{hawkins22032019737maxairplanes} ... \LaTeX:\\ \url{https://www.theverge.com/2019/3/22/18275736/boeing-737-max-plane-crashes-grounded-problems-info-details-explained-reasons}
%    									veiligheidsopmerking
%    									\bibitem{{thomas30082020737safest} ... \LaTeX:\\ \url{https://www.airlineratings.com/news/boeings-737-max-will-one-safest-aircraft-history/}
%    										aanpassingen
%    										\bibitem{{boeing737maxdisplay} ... \LaTeX:\\ \url{https://www.boeing.com/commercial/737max/737-max-software-updates.page}
%    											waarschuwingen//output signalen
%    											\bibitem{{fehrm24112020737changes} ... \LaTeX:\\ \url{https://leehamnews.com/2020/11/24/boeing-737-max-changes-beyond-mcas/}
%    												software gerelateerde fouten
%    												\bibitem{{travis18042019737maxsoftwaredevop} ... \LaTeX:\\ \url{https://spectrum.ieee.org/aerospace/aviation/how-the-boeing-737-max-disaster-looks-to-a-software-developer}
%    													onderzoeksrapport
%    													de rol van de publieke opinie
%    													\bibitem{{barnett05052019737maxcrisis} ... \LaTeX:\\ \url{https://pubsonline.informs.org/do/10.1287/orms.2019.05.05/full/}
%    														onderzoek van europese luchtvaart agentschap
%    														\bibitem{{easa27012021737maxsafereturn} ... \LaTeX:\\ \url{https://www.easa.europa.eu/newsroom-and-events/news/easa-declares-boeing-737-max-safe-return-service-europe}
%    															veiligheidsvraagstuk
%    															\bibitem{{touitou11032019737tragedies} ... \LaTeX:\\ \url{https://phys.org/news/2019-03-boeing-max-safety-tragedies.html}
%    																artikel over sensoren
%    																\bibitem{{hemmerdinger02022021737maxdeliveries} ... \LaTeX:\\ \url{https://www.flightglobal.com/airframers/boeing-delays-737-max-10-deliveries-two-years-to-2023/142245.article}
%    																	goedkeuring van europese luchtvaart autoriteiten
%    																	advies aan de faa
%    																	\bibitem{{bielby27022021faaimprovesafety} ... \LaTeX:\\ \url{https://www.hstoday.us/subject-matter-areas/airport-aviation-security/oig-tells-faa-to-improve-safety-oversight-following-boeing-737-max-review/}
%    																		\bibitem{{boyle18112020737maxupgrade} ... \LaTeX:\\ \url{https://www.geekwire.com/2020/faas-go-ahead-737-maxs-return-flight-kicks-off-massive-software-upgrade/}
%    																			\bibitem{{bergstraburgess122019737maxMcasAlgorithm} ... \LaTeX:\\ \url{https://www.researchgate.net/publication/338420944_A_Promise_Theoretic_Account_of_the_Boeing_737_Max_MCAS_Algorithm_Affair}
%    																				achtergrond informatie
%    																				\bibitem{{737mcas} ... \LaTeX:\\ \url{http://www.b737.org.uk/mcas.htm}
%    																					algemeen vertrouwen
%    																					\bibitem{{newburger17052019boeingcrisis} ... \LaTeX:\\ \url{https://www.cnbc.com/2019/05/16/what-you-need-to-know-about-boeings-737-max-crisis.html}
%    																						toestemming europese autoriteiten
%    																						problemen
%    																						\bibitem{{arstechnica22012020737problems} ... \LaTeX:\\ \url{https://arstechnica.com/information-technology/2020/01/737-max-fix-slips-to-summer-and-thats-just-one-of-boeings-problems/}
%    																							uitgebreid artikel over de onderzoeken en het vliegverbod
%    																							\bibitem{{german190620217372yaftergrounded} ... \LaTeX:\\ \url{https://www.cnet.com/news/boeing-737-max-8-all-about-the-aircraft-flight-ban-and-investigations/}
%    																								computers als oorzaak
%    																								lessons learned
%    																								\bibitem{{beningo02052019boeinglessons} ... \LaTeX:\\ \url{https://www.designnews.com/electronics-test/5-lessons-learn-boeing-737-max-fiasco}
%    																									\bibitem{{duran05042019boeingspof} ... \LaTeX:\\ \url{https://www.eurocontrol.int/publication/effects-network-extra-standby-aircraft-and-boeing-737-max-grounding}
%    																										single point of failure
%    																										\bibitem{{ID} ... \LaTeX:\\ \url{https://dmd.solutions/blog/2019/04/05/how-a-single-point-of-failure-spof-in-the-mcas-software-could-have-caused-the-boeing-737-max-crash-in-ethiopia/}
%    																											\bibitem{{makichuck24012021737fearflying} ... \LaTeX:\\ \url{https://asiatimes.com/2021/01/boeings-737-max-and-the-fear-of-flying/}
%    																												lijst van tehnische aanpassingen
%    																												\bibitem{{caa737modifications} ... \LaTeX:\\ \url{https://www.caa.co.uk/Consumers/Guide-to-aviation/Boeing-737-MAX/}
%    																													\bibitem{{oestergaard14122020boeingdeliveries} ... \LaTeX:\\ \url{https://dsm.forecastinternational.com/wordpress/2020/12/14/airbus-and-boeing-report-november-2020-commercial-aircraft-orders-and-deliveries/}
%    																														code lek
%    																														\bibitem{{reenberg787flaws} ... \LaTeX:\\ \url{https://www.wired.com/story/boeing-787-code-leak-security-flaws/}
%    																															\bibitem{{fitch16092020737backlogrisks} ... \LaTeX:\\ \url{https://www.fitchratings.com/research/corporate-finance/boeing-737-max-return-backlog-risks-remain-16-09-2020}
%    																																Cultuurverandering, deregulatie, systeemwijziging of gewoon een kwestie van competentie
%    																																\bibitem{{willis27082020737maxfailures} ... \LaTeX:\\ \url{https://www.aerospacetestinginternational.com/features/what-broke-the-737-max.html}
%    																																	extra aanpassingen
%    																																	\bibitem{{ostrower11062020more737changes} ... \LaTeX:\\ \url{https://theaircurrent.com/aviation-safety/boeings-737-max-software-done-but-regulators-plot-more-changes-after-jets-return/}
%    																																		wat ging er mis een analyse van een ex-iloot
%    																																		De utoriteiten waren op de hoogte
%    																																		\bibitem{{hruska13122019faaknown737crashrate} ... \LaTeX:\\ \url{https://www.extremetech.com/extreme/303373-the-faa-knew-the-737-max-was-dangerous-and-kept-it-flying-anyway}
%    																																			kwaliteiten van het alarmsysteem niet goed bekend
%    																																			\bibitem{{bloomberg26092019failedpred} ... \LaTeX:\\ \url{https://time.com/5687473/boeing-737-alarm-system/}
%    																																				\bibitem{{whiteman09072020boengcancelstock} ... \LaTeX:\\ \url{https://www.nasdaq.com/articles/boeing-gets-dealt-another-737-max-cancellation-blow.-what-it-means-for-boeing-stock-2020}
%    																																					\bibitem{{leopold09192019boeingreliability} ... \LaTeX:\\ \url{https://www.eetimes.com/boeing-crashes-highlight-a-worsening-reliability-crisis/}
%    																																						veiligheidsvraagstuk
%    																																						\bibitem{{koenig11122019737crashesnofix} ... \LaTeX:\\ \url{https://www.latimes.com/business/story/2019-12-11/faa-boeing-737-max-crashes}
%    																																							probleemanalyse, veiligheidsvraagstuk
%    																																							\bibitem{{dohertylindeman15032019737problems} ... \LaTeX:\\ \url{https://www.politico.com/story/2019/03/15/boeing-737-max-grounding-1223072}
%    																																								falend toezicht
%    																																								\bibitem{{stodder02102019corruptoversight} ... \LaTeX:\\ \url{https://www.pogo.org/analysis/2019/10/corrupted-oversight-the-faa-boeing-and-the-737-max/}
%    																																									\bibitem{{afacwaLostSafeguards} ... \LaTeX:\\ \url{https://www.afacwa.org/the_inside_story_of_mcas_seattle_times}
%    																																										doelstellingen en veiligheidsvraagstukken
%    																																										\bibitem{{swayne18032019profitssafety} ... \LaTeX:\\ \url{https://www.marxist.com/737-max-scandal-boeing-putting-profits-before-safety.htm}
%    																																											\bibitem{{freed26022021liftaustraliaban} ... \LaTeX:\\ \url{https://finance.yahoo.com/news/australia-lifts-ban-boeing-737-035817682.html?guccounter=1&guce_referrer=aHR0cHM6Ly93d3cuZ29vZ2xlLmNvbS8&guce_referrer_sig=AQAAAHZCJYy_0A5VS2WiPoCvH4xdrRNkmkdsv5EWJ2RLIz_AS-rxsTty6AF1_HlmJiRyWYqCXDi4p0Xs4isYkNkCq2Pfo-pQ60Xz_IfTNjm4FgoZiBMC4zpZlB6F0fwecrjE_ujAXZzG4xPJnWCd8-G3VLlPTY8h3H31eQ1i8hY9AIyy}
%    																																												autoriteiten krijgen tik op de vingers
%    																																												\bibitem{{reed15032019softwareattention} ... \LaTeX:\\ \url{https://medium.com/@jpaulreed/the-737max-and-why-software-engineers-should-pay-attention-a041290994bd}
%    																																													\bibitem{{news17032019softwareexplains} ... \LaTeX:\\ \url{https://news.ycombinator.com/item?id=19414775}
%    																																														\bibitem{{legget21122020eu737maxsafe} ... \LaTeX:\\ \url{https://www.bbc.com/news/55366320}
%    																																															\bibitem{{marketscreener0103221737chinarecertification} ... \LaTeX:\\ \url{https://www.marketscreener.com/news/latest/China-studies-Boeing-737-MAX-recertification-wants-safety-concerns-fully-addressed--32569125/}
%    																																																motor in brand
%    																																																\bibitem{{euractiv22022021737firegrounds} ... \LaTeX:\\ \url{https://www.euractiv.com/section/aviation/news/boeing-grounds-777s-after-engine-fire/}
%    																																																	\bibitem{{benny18022019737returnUAE} ... \LaTeX:\\ \url{https://gulfnews.com/business/aviation/uae-airspace-to-see-return-of-boeing-737-max-1.1613627548923}
%    																																																		motor in brand gevlogen
%    																																																		\bibitem{{biersmichel22022021777grounds} ... \LaTeX:\\ \url{https://techxplore.com/news/2021-02-boeing-urges-grounding-777s.html}
%    																																																			\bibitem{{ID} ... \LaTeX:\\ \url{https://www.politico.eu/article/uk-temporarily-bans-some-boeing-aircraft-after-pratt-whitney-engine-incidents/}
%    																																																				\bibitem{{reuters23022021777metalfatigue} ... \LaTeX:\\ \url{https://www.timeslive.co.za/news/world/2021-02-23-damage-to-united-boeing-777-engine-consistent-with-metal-fatigue--ntsb/}
%    																																																					faa was niet kritisch genoeg
%    																																																					\bibitem{{ID} ... \LaTeX:\\ \url{https://federalnewsnetwork.com/government-news/2021/02/federal-watchdog-blasts-faa-over-certification-of-boeing-jet/}
%    																																																						\cite{}
    																																																						
  
\end{thebibliography}

%
%\newpage
%\bibliography{references}
%

\vfill
    %(verplicht) bronvermeldingen
  \addcontentsline{toc}{chapter}{\bibname}
  \hoofdstuk{Persoonlijk verslag}

In de evaluatie reflecteer je over je eigen afstudeerproces. Daarbij
moet je vooral letten op de leereffecten. Welke competenties had je
nodig? Welke competenties kwam je tekort en moest je zelf verwerven?
Waren dit algemene of specifieke competenties?  Voldeden de
beroepscompetenties aan de standaard van het \emph{HBO-I} (analyseren,
adviseren, ontwerpen, realiseren en beheren)?  Vielen de algemene
competenties in de vijf categorieën van de \emph{Dublin
Descriptoren}\footnote{Dublin Descriptoren zijn eisen aan de
competenties voor de bachelor en master studies aan universiteiten en
hogescholen in Europa.} zoals het verkrijgen van kennis en inzicht,
het toepassen van kennis en inzicht, het maken van onderbouwde keuzen
(oordeelsvorming), het communiceren (schriftelijk en mondeling) en het
verkrijgen van leervaardigheden?

  %(verplicht) reflectie op het leerproces
  \appendix
  \bijlage{Achtergrond materiaal}

In de bijlagen komen alle gegevens die nodig zijn voor de
onderbouwing, maar die de leesbaarheid van het hoofdverslag verlagen.


%\bijlage{Onderzoeksgegevens}

%\bijlage{UML diagrammen}

%\bijlage{Gebruikshandleiding}



   %(optioneel) brainstormdocument


%  \include{}   %(optioneel) bijlagen
%  \include{}   %(optioneel) bijlagen
%  \include{}   %(optioneel) bijlagen
\fi
\end{document}       %Einde van het document
